
\documentstyle[amstex,amssymb,preprint,tighten,aps]{revtex}
%%%%%%%%%%%%%%%%%%%%%%%%%%%%%%%%%%%%%%%%%%%%%%%%%%%%%%%%%%%%%%%%%%%%%%%%%%%%%%%%%%%%%%%%%%%%%%%%%%%%%%%%%%%%%%%%%%%%%%%%%%%%
%TCIDATA{OutputFilter=LATEX.DLL}
%TCIDATA{Created=Wednesday, May 16, 2001 14:30:46}
%TCIDATA{LastRevised=Friday, December 28, 2001 15:24:49}
%TCIDATA{<META NAME="GraphicsSave" CONTENT="32">}
%TCIDATA{<META NAME="DocumentShell" CONTENT="Articles\SW\REVTeX - APS and AIP Article">}
%TCIDATA{Language=American English}
%TCIDATA{CSTFile=revtxtci.cst}

\newtheorem{theorem}{Theorem}
\newtheorem{acknowledgement}[theorem]{Acknowledgement}
\newtheorem{algorithm}[theorem]{Algorithm}
\newtheorem{axiom}[theorem]{Axiom}
\newtheorem{claim}[theorem]{Claim}
\newtheorem{conclusion}[theorem]{Conclusion}
\newtheorem{condition}[theorem]{Condition}
\newtheorem{conjecture}[theorem]{Conjecture}
\newtheorem{corollary}[theorem]{Corollary}
\newtheorem{criterion}[theorem]{Criterion}
\newtheorem{definition}[theorem]{Definition}
\newtheorem{example}[theorem]{Example}
\newtheorem{exercise}[theorem]{Exercise}
\newtheorem{lemma}[theorem]{Lemma}
\newtheorem{notation}[theorem]{Notation}
\newtheorem{problem}[theorem]{Problem}
\newtheorem{proposition}[theorem]{Proposition}
\newtheorem{remark}[theorem]{Remark}
\newtheorem{solution}[theorem]{Solution}
\newtheorem{summary}[theorem]{Summary}

\begin{document}
\title{Equations of Motion}
\author{Brian Weiner}
\date{11/13/01}
\maketitle

\section{Parameterization}

\subsection{Equivalence Classes}

Operators belonging to ${\cal B}_{1}\left( {\cal H}^{N}\right) $ can be
classified into equivalence classes $\left[ f\right] _{N\text{ }}$by the
action of the contraction map $\Xi :{\cal B}_{1}\left( {\cal H}^{N}\right)
\rightarrow L_{1}\left( {\rm Z}\right) $. The equivalence classes are affine
spaces and this class structure of ${\cal B}_{1}\left( {\cal H}^{N}\right) $
can be viewed in terms of an affine bundle ${\frak B}_{aff}$ with fibers $%
\left[ f\right] _{N\text{ }}$ and base $L_{1}\left( {\rm Z}\right) $,
elements of this fiber bundle can be denoted by $\left( Y,f\right) $ where $%
Y\in \left[ f\right] _{N\text{ }}$ and $f\in L_{1}\left( {\rm Z}\right) $.

The space of trace class operators ${\cal B}_{1}\left( {\cal H}^{N}\right) $
has the natural decomposition wrt to the linear map $\Xi $ as 
\begin{equation}
{\cal B}_{1}\left( {\cal H}^{N}\right) =\ker \Xi \oplus \ker \Xi ^{c}.
\end{equation}%
$\lceil $Note that the compliment of a subspace is {\bf not }unique (see %
\cite{Geroch}) and leads to different representations of the base, as we
discuss in the following. The non uniqueness is expressed in 
\begin{eqnarray}
\left[ \ker \Xi ^{c}\right] _{\mu } &=&\left\{ u+\varphi _{\mu }\left(
u\right) ;u\in \ker \Xi ^{c}\right\}  \nonumber \\
\varphi _{\mu } &:&\ker \Xi ^{c}\longrightarrow \ker \Xi
\end{eqnarray}%
with $\left[ \ker \Xi ^{c}\right] _{\mu }$ also being a compliment when $%
\ker \Xi ^{c}$ is a given compliment. We take $\ker \Xi ^{c}$ to be the {\em %
canonical }compliment made up of vectors that have {\bf no }components in $%
\ker \Xi $ and in the following the term compliment will refer to the
canonical compliment unless otherwise stated$\rfloor $. As $\Xi $ is a
linear map the compliment of its kernel is isomorphic to its range i.e. $%
\ker \Xi ^{c}$ $\cong L_{1}\left( {\rm Z}\right) $ $\lceil $assuming that $%
\Xi $ is onto $L_{1}\left( {\rm Z}\right) \rfloor $ and we will describe
this isomorphism by 
\begin{equation}
\Gamma :L_{1}\left( {\rm Z}\right) \longrightarrow \ker \Xi ^{c}.
\end{equation}%
Using $\Gamma $ the equivalence class $\left[ f\right] _{N\text{ }}$ can be
expressed as 
\begin{equation}
\left[ f\right] _{N\text{ }}=\left\{ Y:Y=\Gamma \left( f\right) +k,\quad
k\in \ker \Xi \right\}
\end{equation}%
and we call $\Gamma \left( f\right) $ the irreducible part of the operator $%
Y $ (obviously $\Gamma \left( f\right) \in \left[ f\right] _{N\text{ }}$ and
all operators belonging to a given equivalence class have the same
irreducible part). The class structure of ${\cal B}_{1}\left( {\cal H}%
^{N}\right) $ can thus also be viewed in terms of a vector bundle ${\frak B}%
_{vec}$ with fibers $\ker \Xi $ and base $\ker \Xi ^{c}$ and hence the
vector bundle ${\frak B}_{vec}$ is equivalent to the affine bundle ${\frak B}%
_{aff}$ in the preceding manner and can be thought of as a
''representation'' of it. The elements of this vector bundle can be denoted
by $\left( k,\Gamma \left( f\right) \right) $ where $k\in \ker \Xi $ and $%
\Gamma \left( f\right) \in \ker \Xi ^{c}$.

We can more generally introduce an injective map 
\begin{equation}
Y_{ref}:L_{1}\left( {\rm Z}\right) \longrightarrow {\cal B}_{1}\left( {\cal H%
}^{N}\right)
\end{equation}%
where $Y_{ref}\left( f\right) \in \left[ f\right] _{N}$, (this implies that $%
Y_{ref}\left( f_{1}\right) $ and $Y_{ref}\left( f_{2}\right) $ belong to
different equivalence classes when $f_{1}\neq f_{2}$ even if they are just
scalar multiples of each other), then the equivalence class $\left[ f\right]
_{N\text{ }}$ can be expressed as 
\begin{equation}
\left[ f\right] _{N\text{ }}=\left\{ Y:Y=Y_{ref}\left( f\right) +k,\quad
k\in \ker \Xi \right\}
\end{equation}%
and the elements of ${\frak B}_{vec}$ can also be denoted by $\left(
k,Y_{ref}\left( f\right) \right) $ where $k\in \ker \Xi $ and $Y_{ref}\left(
f\right) \in \left[ f\right] _{N}$ (clearly $Y_{ref}\left( f\right) -\Gamma
\left( f\right) \in \ker \Xi $). We can think of $\left\{ \left( k,\Gamma
\left( f\right) \right) \right\} $ and $\left\{ \left( k,Y_{ref}\left(
f\right) \right) \right\} $ as equivalent representations of ${\frak B}%
_{vec} $, with $\left\{ \left( k,\Gamma \left( f\right) \right) \right\} $
being the defining representation. Note it is the base which has different
representations, and these representations are related by the super operator 
$\hat{T}_{Y_{ref}}$ 
\begin{equation}
\left( k,Y_{ref}\left( f\right) \right) =\left( k,\hat{T}_{Y_{ref}}\left(
\Gamma \left( f\right) \right) \right)
\end{equation}%
where the composition $\hat{T}_{Y_{ref}}\circ \Gamma $ is an injection from $%
L_{1}\left( {\rm Z}\right) $ to ${\cal B}_{1}\left( {\cal H}^{N}\right) $
and 
\begin{eqnarray}
Y_{ref}\left( f\right) &=&\hat{T}_{Y_{ref}}\left( \Gamma \left( f\right)
\right) =\Gamma \left( f\right) +\varphi _{Y_{ref}}\left( \Gamma \left(
f\right) \right)  \label{newbase} \\
\varphi _{Y_{ref}} &:&\ker \Xi ^{c}\longrightarrow \ker \Xi  \nonumber
\end{eqnarray}%
$\lceil $Note that $\varphi _{Y_{ref}}$ does not have to be injective$%
\rfloor $.

\subsection{Cosets}

The pure states of ${\cal H}^{N}$ can be parameterized non redundantly by
the action of representations, (which we will not explicitly signify and
correspond to the different representations of ${\frak B}_{vec}$ described
in the preceding), of ${\frak G}_{den}$ and ${\frak G}_{fib}$ on a pure
reference state in the following manner 
\begin{equation}
|\Psi ({\cal V}_{fib}({\cal V}_{den}),\omega ,{\cal V}_{den}\rangle =e^{%
{\cal \breve{V}}_{fib}\left[ {\cal V}_{den}\right] }e^{i\breve{\omega}}e^{%
{\cal \breve{\beta}}\left( {\cal V}_{den}\right) }e^{{\cal \breve{V}}%
_{den}}|\Psi _{ref}\left( \gamma _{0}\right) \rangle  \label{param}
\end{equation}%
where ${\cal \breve{V}}_{fib}\left[ {\cal V}_{den}\right] $ denotes a
coordinate diagonal operator belonging to ${\cal B}\left( {\cal H}%
^{N}\right) $ that maps $\left[ \gamma \right] _{N}$ $\equiv \left[ {\cal V}%
_{den}\right] $ to itself i.e. 
\begin{eqnarray}
&&\Xi \left( e^{{\cal \breve{V}}_{fib}\left[ {\cal V}_{den}\right]
}D^{N}\left( \gamma \right) e^{{\cal \breve{V}}_{fib}\left[ {\cal V}_{den}%
\right] }\right) \equiv \Xi \left( \widehat{e^{{\cal \breve{V}}_{fib}\left[ 
{\cal V}_{den}\right] }}D^{N}\left( \gamma \right) \right) =\gamma  \nonumber
\\
&&D^{N}\left( \gamma \right) =e^{{\cal \breve{V}}_{den}}|\Psi _{ref}\left(
\gamma _{0}\right) \rangle \left\langle \Psi _{ref}\left( \gamma _{0}\right)
\right| e^{{\cal \breve{V}}_{den}}\equiv \widehat{e^{{\cal \breve{V}}_{den}}}%
\left( |\Psi _{ref}\left( \gamma _{0}\right) \rangle \left\langle \Psi
_{ref}\left( \gamma _{0}\right) \right| \right)
\end{eqnarray}%
and ${\cal \breve{\beta}}\left( {\cal V}_{den}\right) =$ ${\cal \breve{V}}%
_{fib}({\cal V}_{den})+i{\cal \breve{\omega}}({\cal V}_{den})$, where ${\cal 
\breve{V}}_{fib}({\cal V}_{den})$ and ${\cal \breve{\omega}}({\cal V}_{den})$
are fixed paths of intra equivalence class maps and $\widehat{e^{\breve{X}}}$
denotes the super operator induced by $e^{\breve{X}}$. ({\em Note} the
difference in meaning between ${\cal \breve{V}}_{fib}({\cal V}_{den})$ and $%
{\cal \breve{V}}_{fib}\left[ {\cal V}_{den}\right] $). It should be observed
that the map $\widehat{e^{{\cal \breve{V}}_{den}}}$ does not change the $%
\ker \Xi $ component of $D^{N}\left( \gamma _{0}\right) =|\Psi _{ref}\left(
\gamma _{0}\right) \rangle \left\langle \Psi _{ref}\left( \gamma _{0}\right)
\right| $, while $\widehat{e^{{\cal \breve{\beta}}\left( {\cal V}%
_{den}\right) }}$ does and corresponds to the map $\hat{T}_{Y_{ref}}\left(
\gamma \right) $ described for the general case in Eq. $\left( \ref{newbase}%
\right) $.

$\lceil ${\em Note}: If a reference state with density $\gamma _{0}$ is
given then the constructions described below and in Appendix A show that $%
\left[ \gamma \right] _{N}\equiv \left[ {\cal V}_{den}\right] _{N}$ in a $%
1-1 $ unique fashion$\rfloor $.

$\lceil ${\em Note}: In the following that $\left| \breve{Y}\right) $
denotes the operator $\breve{Y}$ when viewed as an element of the Banach
space ${\frak B}_{1}\left( {\cal H}^{N}\right) \rfloor $.

The factor $e^{{\cal \breve{V}}_{den}}$ in the transformation Eq. $\left( %
\ref{param}\right) $ belongs to the coset decomposition of ${\frak G}_{den}$%
, (in the representation by super operators: ${\frak B}_{1}\left( {\cal H}%
^{N}\right) \rightarrow {\frak B}_{1}\left( {\cal H}^{N}\right) )$, given by 
\begin{eqnarray}
&&{\frak G}_{\gamma _{0}}=\left\{ \hat{T}\in {\frak G}_{den}{\frak ;\,}\hat{T%
}\left| D_{ref}^{N}\left( \gamma _{0}\right) \right) =\left|
D_{ref}^{N}\left( \gamma _{0}\right) \right) \right\}  \nonumber \\
&&{\frak Co}_{den}={\frak G}_{den}/{\frak G}_{\gamma _{0}};\quad \widehat{e^{%
{\cal \breve{V}}_{den}}}\in {\frak Co}_{den}
\end{eqnarray}%
where 
\begin{equation}
\left| D_{ref}^{N}\left( \gamma _{0}\right) \right) \equiv \left| \left|
\Psi _{ref}\left( \gamma _{0}\right) \right\rangle \left\langle \Psi
_{ref}\left( \gamma _{0}\right) \right| \right) ,
\end{equation}%
and produces the following orbit of pure states 
\begin{eqnarray}
\hat{T}\left| D_{ref}^{N}\left( \gamma _{0}\right) \right) &=&\left|
D_{0}^{N}\left( \gamma \right) \right) ;\,\hat{T}\in {\frak Co}_{den} 
\nonumber \\
\left| D_{0}^{N}\left( \gamma \right) \right) &\equiv &\left| \left| \Psi
_{0}\left( \gamma \right) \right\rangle \left\langle \Psi _{0}\left( \gamma
\right) \right| \right) .
\end{eqnarray}%
The factor $e^{{\cal \breve{\beta}}\left( {\cal V}_{den}\right) }$ produces
a representation of the base 
\begin{equation}
\left| D_{ref}^{N}\left( \gamma \right) \right) =\widehat{e^{{\cal \breve{%
\beta}}\left( {\cal V}_{den}\right) }}\left| D_{0}^{N}\left( \gamma \right)
\right) \equiv e^{{\cal \breve{\beta}}\left( {\cal V}_{den}\right) }\left|
\Psi _{0}\left( \gamma \right) \right\rangle \left\langle \Psi _{0}\left(
\gamma \right) \right| e^{{\cal \breve{\beta}}\left( {\cal V}_{den}\right)
^{\dagger }}
\end{equation}%
The factor $e^{{\cal \breve{V}}_{fib}\left[ {\cal V}_{den}\right] }e^{i%
\breve{\omega}}$ belongs to the coset decomposition of ${\frak G}_{fib}$
given by 
\begin{eqnarray}
&&{\frak G}_{\gamma }=\left\{ \hat{T}\in {\frak G}_{fib}{\frak ;\,}\hat{T}%
\left| D_{ref}^{N}\left( \gamma \right) \right) =\left| D_{ref}^{N}\left(
\gamma \right) \right) \right\}  \nonumber \\
&&{\frak Co}_{fib}={\frak G}_{fib}/{\frak G}_{\gamma };\quad \widehat{e^{%
{\cal \breve{V}}_{fib}({\cal V}_{den})+i\breve{\omega}}}\in {\frak Co}_{fib}
\end{eqnarray}%
where 
\begin{equation}
\left| D_{ref}^{N}\left( \gamma \right) \right) \equiv \left| \left| \Psi
_{ref}\left( \gamma \right) \right\rangle \left\langle \Psi _{ref}\left(
\gamma \right) \right| \right) ,
\end{equation}%
and produces the following orbit of pure states%
\begin{eqnarray}
\hat{T}\left| D_{ref}^{N}\left( \gamma \right) \right) &=&\left| D^{N}\left(
\gamma \right) \right) ;\,\hat{T}\in {\frak Co}_{fib}  \nonumber \\
\left| D^{N}\left( \gamma \right) \right) &\equiv &\left| \left| \Psi \left(
\gamma \right) \right\rangle \left\langle \Psi \left( \gamma \right) \right|
\right) .
\end{eqnarray}%
The transformations appearing in Eq. $\left( \ref{param}\right) $ are
representations of the cosets, different $\gamma _{0}$ give equivalent
representations. Most often we start from a reference $\left| \Psi
_{0}\left( \gamma \right) \right\rangle \left\langle \Psi _{0}\left( \gamma
\right) \right| $ and act on it with the cosets ${\frak Co}_{den}$ and $%
{\frak Co}_{fib}$, as represented above, and set ${\cal \breve{\beta}}\left( 
{\cal V}_{den}\right) =0$ and thus use the representation of the base
induced by $\left| \Psi _{0}\left( \gamma \right) \right\rangle \left\langle
\Psi _{0}\left( \gamma \right) \right| $.

It is possible to combine the factors $e^{{\cal \breve{\beta}}\left( {\cal V}%
_{den}\right) }$ and $e^{{\cal \breve{V}}_{den}}$ together as coset
representatives of another but equivalent representation of ${\frak Co}%
_{den}\subset {\frak G}_{den}$ as we do in the section after next.

\subsection{Complex Analytic Parameterization}

In this section we set ${\cal \breve{\beta}}\left( {\cal V}_{den}\right) =0$%
, but all of the proceeding is valid when ${\cal \breve{\beta}}\left( {\cal V%
}_{den}\right) \neq 0$. We can introduce a complex analytic parameterization
by partitioning the unitary part of the transformation into parts conjugate
to $e^{{\cal \breve{V}}_{fib}({\cal V}_{den})}$ and $e^{{\cal \breve{V}}%
_{den}}$ respectively in the following way 
\begin{eqnarray}
|\Psi ({\cal V}_{fib}({\cal V}_{den}),\omega ,{\cal V}_{den}\rangle &\equiv
&|\Psi ({\cal V}_{fib}({\cal V}_{den}),\omega _{fib}({\cal V}_{den}),{\cal V}%
_{den},\omega \left( {\cal V}_{den}\right) \rangle  \nonumber \\
&\equiv &e^{{\cal \breve{T}}_{fib}({\cal V}_{den})}e^{{\cal \breve{T}}_{den}(%
{\cal V}_{den})}|\Psi _{ref}\left( \gamma _{0}\right) \rangle  \nonumber \\
{\cal \breve{T}}_{fib}({\cal V}_{den}) &=&{\cal \breve{V}}_{fib}({\cal V}%
_{den})+i\breve{\omega}_{fib}({\cal V}_{den})  \nonumber \\
{\cal \breve{T}}_{den}({\cal V}_{den}) &=&{\cal \breve{V}}_{den}+i\breve{%
\omega}_{den}({\cal V}_{den})
\end{eqnarray}%
We can also introduce an overall complex analytic parameterization by
setting 
\begin{eqnarray}
{\cal \breve{T}} &{\cal =}&{\cal \breve{T}}_{fib}({\cal V}_{den})+{\cal 
\breve{T}}_{den}({\cal V}_{den})  \nonumber \\
{\cal \breve{V}} &{\cal =}&{\cal \breve{V}}_{fib}({\cal V}_{den})+{\cal 
\breve{V}}_{den}  \nonumber \\
\breve{\omega} &=&\breve{\omega}_{fib}({\cal V}_{den})+\breve{\omega}_{den}(%
{\cal V}_{den})
\end{eqnarray}%
so that 
\begin{equation}
|\Psi ({\cal V}_{fib}({\cal V}_{den}),\omega ,{\cal V}_{den}\rangle =e^{%
{\cal \breve{T}}}|\Psi _{ref}\left( \gamma _{0}\right) \rangle
\end{equation}%
In order to obtain a complex analytic parameterization that explicitly
refers to the vector bundle structure we must have equal degrees of freedom
for ${\cal \breve{V}}_{fib}({\cal V}_{den})$ and $\breve{\omega}_{fib}({\cal %
V}_{den})$ i.e. ${\cal V}_{fib}({\cal V}_{den})$ and $\omega _{fib}({\cal V}%
_{den})$ must belong to vector spaces of the same dimension, otherwise it is
not possible to define ${\cal \breve{T}}_{fib}({\cal V}_{den})$. The same
holds for ${\cal V}_{den}$ and $\omega _{den}({\cal V}_{den})$ in order that 
${\cal \breve{T}}_{den}({\cal V}_{den})$ can be defined. It should be noted
that $\breve{\omega}_{den}({\cal V}_{den})$ produces an intra fiber
transformation not an inter fiber one.

\subsection{Representations of ${\frak G}_{den}$ and ${\frak G}_{fib}$}

We can incorporate the effect of ${\cal \breve{\beta}}\left( {\cal V}%
_{den}\right) $ in the action of ${\frak G}_{den}$ then each ${\cal \breve{%
\beta}}\left( {\cal V}_{den}\right) ={\cal \breve{V}}_{fib}({\cal V}_{den})+i%
\breve{\omega}\left( {\cal V}_{den}\right) $ (i.e. a specific fixed path of $%
{\cal \breve{V}}_{fib}$ and $\breve{\omega}$'s) determines different, but
equivalent, representations of ${\frak G}_{den}$ as super operators mapping $%
{\cal B}_{1}\left( {\cal H}^{N}\right) \rightarrow {\cal B}_{1}\left( {\cal H%
}^{N}\right) $ 
\begin{equation}
{\frak R}\left( {\frak G}_{den}\right) =\left( 
\begin{array}{ccc}
\hat{t}_{K_{OD}}\left( \hat{T}_{S_{D}}\right) & 0 & 0 \\ 
0 & \hat{t}_{K_{D}}\left( \hat{T}_{S_{D}}\right) & 0 \\ 
0 & 0 & \hat{T}_{S_{D}}%
\end{array}%
\right)
\end{equation}%
where%
\begin{equation}
\hat{T}_{S_{D}}\equiv \hat{T}\left( {\cal V}_{den}\right) =e^{{\cal \hat{V}}%
_{den}};\qquad \hat{t}_{K_{D}}\left( \hat{T}_{S_{D}}\right) \equiv \hat{t}%
_{K_{D}}\left( {\cal V}_{den}\right) =\widehat{e^{{\cal \breve{V}}_{fib}(%
{\cal V}_{den})+i\check{\omega}\left( {\cal V}_{den}\right) }}
\end{equation}%
On the other hand the representation of the group ${\frak G}_{fib}$ involves
all possible $\hat{\omega}$'s. 
\begin{equation}
{\frak R}\left( {\frak G}_{fib}\right) =\left( 
\begin{array}{ccc}
\hat{t}_{K_{OD}}^{\prime }\left( \hat{T}_{K_{D}}\right) & 0 & 0 \\ 
0 & \hat{T}_{K_{D}} & 0 \\ 
0 & 0 & \hat{P}_{K_{S}}%
\end{array}%
\right)
\end{equation}%
where%
\begin{equation}
\hat{T}_{K_{D}}\equiv \hat{T}\left( {\cal V}_{fib}({\cal V}_{den}),\omega
\right) =\widehat{e^{{\cal \breve{V}}_{fib}\left[ {\cal V}_{den}\right] +i%
\check{\omega}}}
\end{equation}

\section{Equation of Motion}

\subsection{Single Functional variable}

The Lagrangian 
\begin{eqnarray}
{\cal L}\left( \Psi \left( t\right) ,\dot{\Psi}\left( t\right) \right) &=&%
\frac{i}{2}\frac{\left( \left\langle \Psi \left( t\right) |\dot{\Psi}\left(
t\right) \right\rangle -\left\langle \dot{\Psi}\left( t\right) |\Psi \left(
t\right) \right\rangle \right) }{\left\langle \Psi \left( t\right) |\Psi
\left( t\right) \right\rangle }-\frac{\left\langle \Psi \left( t\right)
|H\Psi \left( t\right) \right\rangle }{\left\langle \Psi \left( t\right)
|\Psi \left( t\right) \right\rangle }  \nonumber \\
&\equiv &\frac{1}{{\cal N}\left( \Psi \left( t\right) \right) }\left\{ {\cal %
T}\left( \Psi \left( t\right) ,\dot{\Psi}\left( t\right) \right) {\cal -H}%
\left( \Psi \left( t\right) \right) \right\} ={\cal K}\left( \Psi \left(
t\right) ,\dot{\Psi}\left( t\right) \right) -{\cal E}\left( \Psi \left(
t\right) \right)  \nonumber \\
&\equiv &{\cal K}\left( \Psi \left( t\right) ,\dot{\Psi}\left( t\right)
\right) -{\cal E}\left( \Psi \left( t\right) \right)  \label{La}
\end{eqnarray}%
where 
\begin{eqnarray}
{\cal N}\left( \Psi \left( t\right) \right) &=&\left\langle \Psi \left(
t\right) |\Psi \left( t\right) \right\rangle  \nonumber \\
{\cal T}\left( \Psi \left( t\right) ,\dot{\Psi}\left( t\right) \right)
&=&\left\langle \Psi \left( t\right) |\dot{\Psi}\left( t\right)
\right\rangle -\left\langle \dot{\Psi}\left( t\right) |\Psi \left( t\right)
\right\rangle  \nonumber \\
{\cal H}\left( \Psi \left( t\right) \right) &=&\left\langle \Psi \left(
t\right) |H\Psi \left( t\right) \right\rangle  \nonumber \\
{\cal K}\left( \Psi \left( t\right) ,\dot{\Psi}\left( t\right) \right) &=&%
\frac{{\cal T}\left( \Psi \left( t\right) ,\dot{\Psi}\left( t\right) \right) 
}{{\cal N}\left( \Psi \left( t\right) \right) }  \nonumber \\
{\cal E}\left( \Psi \left( t\right) \right) &=&\frac{{\cal H}\left( \Psi
\left( t\right) \right) }{{\cal N}\left( \Psi \left( t\right) \right) }
\end{eqnarray}%
can be expressed in terms of the preceding parameters and their time
derivative by using the chain rule on eq. $\left( \ref{La}\right) $ as 
\begin{equation}
{\cal L}\left( {\cal \varkappa }\left( t\right) ,{\cal \dot{\varkappa}}%
\left( t\right) \right) =\int {\cal \dot{\varkappa}}\left( {\sf z},t\right) 
{\cal Z}\left( {\sf z},t\right) d{\sf z}-{\cal E}\left( {\cal \varkappa }%
\left( t\right) \right)  \label{ladef}
\end{equation}%
where ${\cal \varkappa }\equiv \left( {\cal V},\omega \right) \equiv \left( 
{\cal V}_{den},{\cal V}_{fib}({\cal V}_{den}),\omega \right) ,$ ${\sf z}%
=\kappa _{1},\ldots ,\kappa _{N}={\bf x}_{1},\ldots ,{\bf x}_{N_{e}},{\bf X}%
_{1},\ldots ,{\bf X}_{N_{n}}$ and 
\begin{equation}
{\cal Z}\left( {\sf z},t\right) =\frac{i}{2}\left\langle \Psi \left(
t\right) |\Psi \left( t\right) \right\rangle ^{-1}\left\{ \left\langle \Psi
\left( t\right) \left| \frac{\delta \Psi \left( t\right) }{\delta {\cal %
\varkappa }\left( {\sf z}\right) }\right. \right\rangle -\left\langle \left. 
\frac{\delta \Psi \left( t\right) }{\delta {\cal \varkappa }\left( {\sf z}%
\right) }\right| \Psi \left( t\right) \right\rangle \right\}
\end{equation}%
The Euler-Lagrange equation for the trajectory ${\cal \varkappa }\left(
t\right) $ that produces a stationary action, $\delta ^{\left( 1\right) }%
{\cal A}\left( {\cal \varkappa },\tau \right) =0$, is then 
\begin{equation}
\frac{\delta {\cal L}\left( {\cal \varkappa },{\cal \dot{\varkappa}}\right) 
}{\delta {\cal \varkappa }\left( {\sf z}\right) }-\frac{d}{dt}\left\{ \frac{%
\delta {\cal L}\left( {\cal \varkappa },{\cal \dot{\varkappa}}\right) }{%
\delta {\cal \dot{\varkappa}}\left( {\sf z}\right) }\right\} \stackrel{a.e.}{%
=}0  \label{EL}
\end{equation}%
($a.e.$ with respect to ${\sf z}$ not $t$, this is implied in all following
equations${\sf )}$. This is a second ODE\ which is made clearer if one
evaluates the time derrivative term to obtain 
\begin{equation}
\int \frac{\delta ^{2}{\cal L}\left( {\cal \varkappa },{\cal \dot{\varkappa}}%
\right) }{\delta {\cal \dot{\varkappa}}\left( {\sf z}\right) \delta {\cal 
\dot{\varkappa}}\left( {\sf z}^{\prime }\right) }\frac{d^{2}{\cal \varkappa }%
\left( {\sf z}^{\prime }\right) }{dt^{2}\quad }d{\sf z}^{\prime }+\int \frac{%
\delta ^{2}{\cal L}\left( {\cal \varkappa },{\cal \dot{\varkappa}}\right) }{%
\delta {\cal \dot{\varkappa}}\left( {\sf z}\right) \delta {\cal \varkappa }%
\left( {\sf z}^{\prime }\right) }\frac{d{\cal \varkappa }\left( {\sf z}%
^{\prime }\right) }{dt\quad }d{\sf z}^{\prime }-\frac{\delta {\cal L}\left( 
{\cal \varkappa },{\cal \dot{\varkappa}}\right) }{\delta {\cal \varkappa }%
\left( {\sf z}\right) }=0
\end{equation}
$\lceil $The Hamiltonian form of this equation (Hamiltons Equations) is 
\begin{equation}
\left. 
\begin{array}{c}
\frac{d{\cal \varkappa }\left( {\sf z}\right) }{dt}=\frac{\delta {\cal H}%
\left( {\cal \varkappa },{\cal Z}\right) }{\delta {\cal Z}\left( {\sf z}%
\right) } \\ 
\frac{d{\cal Z}\left( {\sf z}\right) }{dt}=-\frac{\delta {\cal H}\left( 
{\cal \varkappa },{\cal Z}\right) }{\delta {\cal \varkappa }\left( {\sf z}%
\right) }%
\end{array}%
\right\}  \label{HE}
\end{equation}%
where the Hamiltonian ${\cal H}\left( {\cal \varkappa },{\cal Z}\right) $ is
given by the Legendre transformation of the Lagrangian 
\begin{equation}
{\cal H}\left( {\cal \varkappa },{\cal Z}\right) ={\frak L}\left( {\cal L}%
\left( {\cal \varkappa },{\cal \dot{\varkappa}}\left( {\cal Z}\right)
\right) \right) =\int {\cal \dot{\varkappa}}\left( {\cal Z}\right) \left( 
{\sf z}\right) {\cal Z}\left( {\sf z}\right) d{\sf z-}{\cal L}\left( {\cal %
\varkappa },{\cal \dot{\varkappa}}\left( {\cal Z}\right) \right) .
\label{Legend}
\end{equation}%
where the value ${\cal \dot{\varkappa}}\left( {\cal Z}\right) \left( {\sf z}%
\right) $ of ${\cal \dot{\varkappa}}\left( {\sf z}\right) $ is determined
through the following relationship that defines the conjugate momenta ${\cal %
Z}$ 
\begin{equation}
{\cal Z\left( {\sf z}\right) }=\frac{\delta {\cal L}\left( {\cal \varkappa },%
{\cal \dot{\varkappa}}\right) }{\delta {\cal \dot{\varkappa}}\left( {\sf z}%
\right) }  \label{conjmom}
\end{equation}%
$\lceil $The correspondence between Hamiltons and Lagranges EOM\ is given by
considering 
\begin{equation}
d{\cal H=}\int \left\{ \frac{\delta {\cal H}}{\delta {\cal \varkappa }\left( 
{\sf z}\right) }\delta {\cal \varkappa }\left( {\sf z}\right) +\frac{\delta 
{\cal H}}{\delta {\cal Z}\left( {\sf z}\right) }\delta {\cal Z}\left( {\sf z}%
\right) \right\} d{\sf z}
\end{equation}%
and the differential of the expression eq. $\left( \ref{Legend}\right) $ 
\begin{equation}
d{\cal H=}\int \left\{ -\frac{\delta {\cal L}}{\delta {\cal \varkappa }%
\left( {\sf z}\right) }\delta {\cal \varkappa }\left( {\sf z}\right) +{\cal 
\dot{\varkappa}}\left( {\sf z}\right) \delta {\cal Z}\left( {\sf z}\right)
\right\} d{\sf z}
\end{equation}%
then equating coefficients of $\delta {\cal \varkappa }\left( {\sf z}\right) 
$ and $\delta {\cal Z}\left( {\sf z}\right) $ to give 
\begin{eqnarray}
\frac{\delta {\cal H}}{\delta {\cal \varkappa }\left( {\sf z}\right) } &=&-%
\frac{\delta {\cal L}}{\delta {\cal \varkappa }\left( {\sf z}\right) }
\label{HE1} \\
\frac{\delta {\cal H}}{\delta {\cal Z}\left( {\sf z}\right) } &=&{\cal \dot{%
\varkappa}}\left( {\sf z}\right)  \label{HE2}
\end{eqnarray}%
then by using Lagranges equation eq. $\left( \ref{EL}\right) $for a
stationary point and the definition of the conjugate momenta eq. $\left( \ref%
{conjmom}\right) $one can see that eq.$\left( \ref{HE1}\right) $ becomes 
\begin{equation}
\frac{\delta {\cal H}}{\delta {\cal \varkappa }\left( {\sf z}\right) }=-%
\frac{\delta {\cal L}}{\delta {\cal \varkappa }\left( {\sf z}\right) }=\frac{%
d{\cal Z}\left( {\sf z}\right) }{dt}
\end{equation}%
thus leading to Hamiltons equations 
\begin{equation}
\frac{\delta {\cal H}}{\delta {\cal Z}\left( {\sf z}\right) }={\cal \dot{%
\varkappa}}\left( {\sf z}\right) ;\qquad \frac{\delta {\cal H}}{\delta {\cal %
\varkappa }\left( {\sf z}\right) }=\,-\stackrel{\cdot }{{\cal Z}}\left( {\sf %
z}\right)
\end{equation}%
which can be written in matrix form as 
\begin{equation}
\left( 
\begin{array}{cc}
{\Bbb O} & -{\Bbb I} \\ 
{\Bbb I} & {\Bbb O}%
\end{array}%
\right) \left( 
\begin{array}{c}
{\cal \dot{\varkappa}} \\ 
\stackrel{\cdot }{{\cal Z}}%
\end{array}%
\right) =\left( 
\begin{array}{c}
\frac{\delta {\cal H}}{\delta {\cal Z}\left( {\sf z}\right) } \\ 
\frac{\delta {\cal H}}{\delta {\cal \varkappa }}%
\end{array}%
\right)
\end{equation}%
In our case the Lagrangian defined by eq. $\left( \ref{La}\right) $ produces
the Hamiltonian 
\begin{equation}
{\cal H}\left( {\cal \varkappa },{\cal Z}\right) ={\cal E}\left( {\cal %
\varkappa }\right)
\end{equation}%
which is conjugate momentum independent. Hamiltons equations eq. $\left( \ref%
{HE}\right) $ are equivalent to Lagranges equation eq. $\left( \ref{EL}%
\right) $ {\em when} the conjugate momenta is well defined and corresponds
in a 1-1 fashion with ${\cal \dot{\varkappa}}$. This occurs {\em only} when
the Legendre transformation ${\frak L}:$ ${\cal L\longrightarrow H}$ is non
singular. If ${\frak L}$ is non singular then Lagranges equations can be
formulated as Hamiltons equations in terms of ${\cal \varkappa }$ and ${\cal %
Z}$ and the Euler-Lagrange equation eq. $\left( \ref{EL}\right) $ can be
written is terms of the conjugate momenta (if it is well defined) as 
\begin{equation}
\frac{\delta {\cal L}\left( {\cal \varkappa },{\cal \dot{\varkappa}}\right) 
}{\delta {\cal \varkappa }\left( {\sf z}\right) }-\frac{d{\cal Z\left( {\sf z%
}\right) }}{dt}\stackrel{a.e.}{=}0\rfloor .
\end{equation}%
The Euler-Lagrange equation eq. $\left( \ref{EL}\right) $ gives 
\begin{eqnarray}
&&\frac{\delta {\cal L}\left( {\cal \varkappa },{\cal \dot{\varkappa}}%
\right) }{\delta {\cal \varkappa }\left( {\sf z}\right) }=\int {\cal \dot{%
\varkappa}}\left( {\sf z}^{\prime }\right) \frac{\delta {\cal Z}\left( {\sf z%
}^{\prime }\right) }{\delta {\cal \varkappa }\left( {\sf z}\right) }d{\sf z}%
^{\prime }-\frac{\delta {\cal E}\left( {\cal \varkappa }\right) }{\delta 
{\cal \varkappa }\left( {\sf z}\right) } \\
&&\frac{d}{dt}\left\{ \frac{\delta {\cal L}\left( {\cal \varkappa },{\cal 
\dot{\varkappa}}\right) }{\delta {\cal \dot{\varkappa}}\left( {\sf z}\right) 
}\right\} =\frac{d{\cal Z}\left( {\sf z}\right) }{dt}=\int \frac{\delta 
{\cal Z}\left( {\sf z}\right) }{\delta {\cal \varkappa }\left( {\sf z}%
^{\prime }\right) }{\cal \dot{\varkappa}}\left( {\sf z}^{\prime }\right) d%
{\sf z}^{\prime }
\end{eqnarray}%
the Euler-Lagrange equation eq. $\left( \ref{EL}\right) ${\em \ for} ${\cal %
\varkappa }\left( {\sf z}\right) $ becomes 
\begin{eqnarray}
&&\int {\cal \dot{\varkappa}}\left( {\sf z}^{\prime }\right) \frac{\delta 
{\cal Z}\left( {\sf z}^{\prime }\right) }{\delta {\cal \varkappa }\left( 
{\sf z}\right) }d{\sf z}^{\prime }-\frac{\delta {\cal E}\left( {\cal %
\varkappa }\right) }{\delta {\cal \varkappa }\left( {\sf z}\right) }-\int 
\frac{\delta {\cal Z}\left( {\sf z}\right) }{\delta {\cal \varkappa }\left( 
{\sf z}^{\prime }\right) }{\cal \dot{\varkappa}}\left( {\sf z}^{\prime
}\right) d{\sf z}^{\prime }=0 \\
&\Rightarrow &\int \left\{ \left( \frac{\delta {\cal Z}\left( {\sf z}%
^{\prime }\right) }{\delta {\cal \varkappa }\left( {\sf z}\right) }-\frac{%
\delta {\cal Z}\left( {\sf z}\right) }{\delta {\cal \varkappa }\left( {\sf z}%
^{\prime }\right) }\right) {\cal \dot{\varkappa}}\left( {\sf z}^{\prime
}\right) \right\} d{\sf z}^{\prime }-\frac{\delta {\cal E}\left( {\cal %
\varkappa }\right) }{\delta {\cal \varkappa }\left( {\sf z}\right) }=0
\end{eqnarray}%
An antisymmetric metric tensor $\eta $ can be defined in terms of ${\cal Z}$
as%
\begin{equation}
\eta \left( {\sf z,z}^{\prime }\right) =\frac{\delta {\cal Z}\left( {\sf z}%
^{\prime }\right) }{\delta {\cal \varkappa }\left( {\sf z}\right) }-\frac{%
\delta {\cal Z}\left( {\sf z}\right) }{\delta {\cal \varkappa }\left( {\sf z}%
^{\prime }\right) }  \label{metric}
\end{equation}%
leading to the equation of motion for ${\cal \varkappa }\left( {\sf z}%
\right) $ 
\begin{equation}
\int \eta \left( {\sf z,z}^{\prime }\right) {\cal \dot{\varkappa}}\left( 
{\sf z}^{\prime }\right) d{\sf z}^{\prime }=\frac{\delta {\cal E}\left( 
{\cal \varkappa }\right) }{\delta {\cal \varkappa }\left( {\sf z}\right) }
\label{EOM}
\end{equation}

The metric term $\eta \left( {\sf z,z}^{\prime }\right) $ is antisymmetric
wrt to interchange of ${\sf z}$ and ${\sf z}^{\prime }$, i.e. $\eta \left( 
{\sf z,z}^{\prime }\right) =-\eta \left( {\sf z}^{\prime }{\sf ,z}\right) $
and thus $\eta =-\eta ^{t}$ and can be further developed as 
\begin{eqnarray}
&&\frac{\delta {\cal Z}\left( {\sf z}^{\prime }\right) }{\delta {\cal %
\varkappa }\left( {\sf z}\right) }=\frac{i}{2}\left\langle \Psi |\Psi
\right\rangle ^{-1}\left\{ \left\langle \frac{\delta \Psi }{\delta {\cal %
\varkappa }\left( {\sf z}\right) }\left| \frac{\delta \Psi }{\delta {\cal %
\varkappa }\left( {\sf z}^{\prime }\right) }\right. \right\rangle
+\left\langle \Psi \left| \frac{\delta ^{2}\Psi }{\delta {\cal \varkappa }%
\left( {\sf z}\right) \delta {\cal \varkappa }\left( {\sf z}^{\prime
}\right) }\right. \right\rangle \right.  \nonumber \\
&&\left. -\left\langle \frac{\delta \Psi }{\delta {\cal \varkappa }\left( 
{\sf z}^{\prime }\right) }\left| \frac{\delta \Psi }{\delta {\cal \varkappa }%
\left( {\sf z}\right) }\right. \right\rangle -\left\langle \frac{\delta
^{2}\Psi }{\delta {\cal \varkappa }\left( {\sf z}\right) \delta {\cal %
\varkappa }\left( {\sf z}^{\prime }\right) }\left| \Psi \right.
\right\rangle -\frac{\left\{ \left\langle \Psi \left| \frac{\delta \Psi }{%
\delta {\cal \varkappa }\left( {\sf z}\right) }\right. \right\rangle \right.
+\left. \left\langle \left. \frac{\delta \Psi }{\delta {\cal \varkappa }%
\left( {\sf z}\right) }\right| \Psi \right\rangle \right\} }{\left\langle
\Psi |\Psi \right\rangle }\right\} \\
&&  \nonumber \\
&&\frac{\delta {\cal Z}\left( {\sf z}\right) }{\delta {\cal \varkappa }%
\left( {\sf z}^{\prime }\right) }=\frac{i}{2}\left\langle \Psi |\Psi
\right\rangle ^{-1}\left\{ \left\langle \frac{\delta \Psi }{\delta {\cal %
\varkappa }\left( {\sf z}^{\prime }\right) }\left| \frac{\delta \Psi }{%
\delta {\cal \varkappa }\left( {\sf z}\right) }\right. \right\rangle
+\left\langle \Psi \left| \frac{\delta ^{2}\Psi }{\delta {\cal \varkappa }%
\left( {\sf z}^{\prime }\right) \delta {\cal \varkappa }\left( {\sf z}%
\right) }\right. \right\rangle \right.  \nonumber \\
&&\left. -\left\langle \frac{\delta \Psi }{\delta {\cal \varkappa }\left( 
{\sf z}\right) }\left| \frac{\delta \Psi }{\delta {\cal \varkappa }\left( 
{\sf z}^{\prime }\right) }\right. \right\rangle -\left\langle \frac{\delta
^{2}\Psi }{\delta {\cal \varkappa }\left( {\sf z}^{\prime }\right) \delta 
{\cal \varkappa }\left( {\sf z}\right) }\left| \Psi \right. \right\rangle -%
\frac{\left\{ \left\langle \Psi \left| \frac{\delta \Psi }{\delta {\cal %
\varkappa }\left( {\sf z}^{\prime }\right) }\right. \right\rangle \right.
+\left. \left\langle \left. \frac{\delta \Psi }{\delta {\cal \varkappa }%
\left( {\sf z}^{\prime }\right) }\right| \Psi \right\rangle \right\} }{%
\left\langle \Psi |\Psi \right\rangle }\right\} \\
&&  \nonumber \\
&&\frac{\delta {\cal Z}\left( {\sf z}^{\prime }\right) }{\delta {\cal %
\varkappa }\left( {\sf z}\right) }-\frac{\delta {\cal Z}\left( {\sf z}%
\right) }{\delta {\cal \varkappa }\left( {\sf z}^{\prime }\right) }%
=i\left\langle \Psi |\Psi \right\rangle ^{-1}\left\{ \left\langle \frac{%
\delta \Psi }{\delta {\cal \varkappa }\left( {\sf z}\right) }\left| \frac{%
\delta \Psi }{\delta {\cal \varkappa }\left( {\sf z}^{\prime }\right) }%
\right. \right\rangle -\left\langle \frac{\delta \Psi }{\delta {\cal %
\varkappa }\left( {\sf z}^{\prime }\right) }\left| \frac{\delta \Psi }{%
\delta {\cal \varkappa }\left( {\sf z}\right) }\right. \right\rangle \right.
\nonumber \\
&&+\frac{1}{2}\left. \frac{\left\{ \left\langle \Psi \left| \frac{\delta
\Psi }{\delta {\cal \varkappa }\left( {\sf z}^{\prime }\right) }\right.
\right\rangle \right. +\left. \left\langle \left. \frac{\delta \Psi }{\delta 
{\cal \varkappa }\left( {\sf z}^{\prime }\right) }\right| \Psi \right\rangle
-\left\langle \Psi \left| \frac{\delta \Psi }{\delta {\cal \varkappa }\left( 
{\sf z}\right) }\right. \right\rangle -\left\langle \left. \frac{\delta \Psi 
}{\delta {\cal \varkappa }\left( {\sf z}\right) }\right| \Psi \right\rangle
\right\} }{\left\langle \Psi |\Psi \right\rangle }\right\} =  \nonumber \\
&&\qquad \qquad \qquad \qquad \qquad \qquad i\left\{ \frac{\delta }{\delta 
{\cal \tilde{\varkappa}}\left( {\sf z}^{\prime }\right) }\frac{\delta }{%
\delta {\cal \varkappa }\left( {\sf z}\right) }-\frac{\delta }{\delta {\cal 
\tilde{\varkappa}}\left( {\sf z}\right) }\frac{\delta }{\delta {\cal %
\varkappa }\left( {\sf z}^{\prime }\right) }\right\} \left. \ln \left\langle
\Psi \left( {\cal \tilde{\varkappa}}\right) |\Psi \left( {\cal \varkappa }%
\right) \right\rangle \right| _{{\cal \tilde{\varkappa}=\varkappa }}
\end{eqnarray}

\subsection{Many Functional Variables}

The Lagrangian eq. $\left( \ref{La}\right) $ can be explicitly expressed in
terms of more than one functional parameter $\left( {\cal V},\omega \right) $
as 
\begin{equation}
{\cal L}\left( {\cal V},\omega ,{\cal \dot{V}},\dot{\omega}\right) =\int
\left\{ {\cal \dot{V}}\left( {\sf z}\right) {\cal Z}\left( {\cal V}\right)
\left( {\sf z}\right) +{\cal \dot{\omega}}\left( {\sf z}\right) {\cal Z}%
\left( {\cal \omega }\right) \left( {\sf z}\right) \right\} d{\sf z}-{\cal E}%
\left( {\cal V},\omega \right)
\end{equation}%
where 
\begin{equation}
{\cal Z}\left( f\right) \left( {\sf z}\right) =\frac{i}{2}\left\langle \Psi
|\Psi \right\rangle ^{-1}\left\{ \left\langle \Psi \left| \frac{\delta \Psi 
}{\delta f\left( {\sf z}\right) }\right. \right\rangle \right. -\left.
\left\langle \left. \frac{\delta \Psi }{\delta f\left( {\sf z}\right) }%
\right| \Psi \right\rangle \right\} ;\quad f={\cal V}\text{ or }\omega
\end{equation}%
The term ${\cal Z}\left( f\right) $ is the {\em conjugate momentum}
associated with the variable $f$ as 
\begin{equation}
{\cal Z}\left( f\right) =\frac{\delta {\cal L}}{\delta \dot{f}}
\end{equation}%
The Euler-Lagrange equation for stationary action for each variable $f\left( 
{\sf z}\right) $ is then 
\begin{equation}
\frac{\delta {\cal L}\left( {\cal \varkappa },{\cal \dot{\varkappa}}\right) 
}{\delta f\left( {\sf z}\right) }-\frac{d}{dt}\left\{ \frac{\delta {\cal L}%
\left( {\cal \varkappa },{\cal \dot{\varkappa}}\right) }{\delta \dot{f}%
\left( {\sf z}\right) }\right\} =0
\end{equation}%
noting that 
\begin{equation}
\frac{\delta {\cal L}\left( {\cal \varkappa },{\cal \dot{\varkappa}}\right) 
}{\delta {\cal V}\left( {\sf z}\right) }=\int \left\{ {\cal \dot{V}}\left( 
{\sf z}^{\prime }\right) \frac{\delta {\cal Z}\left( {\cal V}\right) \left( 
{\sf z}^{\prime }\right) }{\delta {\cal V}\left( {\sf z}\right) }+{\cal \dot{%
\omega}}\left( {\sf z}^{\prime }\right) \frac{\delta {\cal Z}\left( {\cal %
\omega }\right) \left( {\sf z}^{\prime }\right) }{\delta {\cal V}\left( {\sf %
z}\right) }\right\} d{\sf z}^{\prime }-\frac{\delta {\cal E}\left( {\cal V}%
,\omega \right) }{\delta {\cal V}\left( {\sf z}\right) }
\end{equation}%
and%
\begin{equation}
\frac{d}{dt}\left\{ \frac{\delta {\cal L}\left( {\cal \varkappa },{\cal \dot{%
\varkappa}}\right) }{\delta {\cal \dot{V}}\left( {\sf z}\right) }\right\} =%
\frac{d{\cal Z}\left( {\cal V}\right) \left( {\sf z}\right) }{dt}=\int
\left\{ \frac{\delta {\cal Z}\left( {\cal V}\right) \left( {\sf z}\right) }{%
\delta {\cal V}\left( {\sf z}^{\prime }\right) }{\cal \dot{V}}\left( {\sf z}%
^{\prime }\right) +\frac{\delta {\cal Z}\left( {\cal V}\right) \left( {\sf z}%
\right) }{\delta {\cal \omega }\left( {\sf z}^{\prime }\right) }{\cal \dot{%
\omega}}\left( {\sf z}^{\prime }\right) \right\} d{\sf z}^{\prime }
\end{equation}%
the Euler-Lagrange equation eq.$\left( \ref{EL}\right) ${\em \ for} ${\cal V}%
\left( {\sf z}\right) $ becomes 
\begin{eqnarray}
&&\int \left\{ {\cal \dot{V}}\left( {\sf z}^{\prime }\right) \frac{\delta 
{\cal Z}\left( {\cal V}\right) \left( {\sf z}^{\prime }\right) }{\delta 
{\cal V}\left( {\sf z}\right) }+{\cal \dot{\omega}}\left( {\sf z}^{\prime
}\right) \frac{\delta {\cal Z}\left( {\cal \omega }\right) \left( {\sf z}%
^{\prime }\right) }{\delta {\cal V}\left( {\sf z}\right) }\right\} d{\sf z}%
^{\prime }-\frac{\delta {\cal E}\left( {\cal V},\omega \right) }{\delta 
{\cal V}\left( {\sf z}\right) }-  \nonumber \\
&&\qquad \qquad \qquad \qquad \qquad \qquad \int \left\{ \frac{\delta {\cal Z%
}\left( {\cal V}\right) \left( {\sf z}\right) }{\delta {\cal V}\left( {\sf z}%
^{\prime }\right) }{\cal \dot{V}}\left( {\sf z}^{\prime }\right) +\frac{%
\delta {\cal Z}\left( {\cal V}\right) \left( {\sf z}\right) }{\delta {\cal %
\omega }\left( {\sf z}^{\prime }\right) }{\cal \dot{\omega}}\left( {\sf z}%
^{\prime }\right) \right\} d{\sf z}^{\prime }=0\Rightarrow  \nonumber \\
&\quad &\int \left\{ \frac{\delta {\cal Z}\left( {\cal V}\right) \left( {\sf %
z}^{\prime }\right) }{\delta {\cal V}\left( {\sf z}\right) }-\frac{\delta 
{\cal Z}\left( {\cal V}\right) \left( {\sf z}\right) }{\delta {\cal V}\left( 
{\sf z}^{\prime }\right) }\right\} {\cal \dot{V}}\left( {\sf z}^{\prime
}\right) d{\sf z}^{\prime }+  \nonumber \\
&&\qquad \qquad \qquad \quad \qquad \quad \int \left\{ \frac{\delta {\cal Z}%
\left( {\cal \omega }\right) \left( {\sf z}^{\prime }\right) }{\delta {\cal V%
}\left( {\sf z}\right) }-\frac{\delta {\cal Z}\left( {\cal V}\right) \left( 
{\sf z}\right) }{\delta {\cal \omega }\left( {\sf z}^{\prime }\right) }%
\right\} {\cal \dot{\omega}}\left( {\sf z}^{\prime }\right) d{\sf z}^{\prime
}-\frac{\delta {\cal E}\left( {\cal \varkappa }\right) }{\delta {\cal V}%
\left( {\sf z}\right) }=0
\end{eqnarray}%
defining metric tensor components 
\begin{eqnarray}
\eta _{{\cal VV}}\left( {\sf z,z}^{\prime }\right) &=&\frac{\delta {\cal Z}%
\left( {\cal V}\right) \left( {\sf z}^{\prime }\right) }{\delta {\cal V}%
\left( {\sf z}\right) }-\frac{\delta {\cal Z}\left( {\cal V}\right) \left( 
{\sf z}\right) }{\delta {\cal V}\left( {\sf z}^{\prime }\right) }  \nonumber
\\
\eta _{{\cal V\omega }}\left( {\sf z,z}^{\prime }\right) &=&\frac{\delta 
{\cal Z}\left( {\cal \omega }\right) \left( {\sf z}^{\prime }\right) }{%
\delta {\cal V}\left( {\sf z}\right) }-\frac{\delta {\cal Z}\left( {\cal V}%
\right) \left( {\sf z}\right) }{\delta {\cal \omega }\left( {\sf z}^{\prime
}\right) }  \label{mt}
\end{eqnarray}%
this becomes 
\begin{equation}
\int \left\{ \eta _{{\cal VV}}\left( {\sf z,z}^{\prime }\right) {\cal \dot{V}%
}\left( {\sf z}^{\prime }\right) +\eta _{{\cal V\omega }}\left( {\sf z,z}%
^{\prime }\right) {\cal \dot{\omega}}\left( {\sf z}^{\prime }\right)
\right\} d{\sf z}^{\prime }=\frac{\delta {\cal E}\left( {\cal \varkappa }%
\right) }{\delta {\cal V}\left( {\sf z}\right) }
\end{equation}%
In an analogous fashion the Euler-Lagrange equation eq.$\left( \ref{EL}%
\right) ${\em \ for} ${\cal \omega }\left( {\sf z}\right) $ becomes 
\begin{equation}
\int \left\{ \eta _{{\cal \omega \omega }}\left( {\sf z,z}^{\prime }\right) 
{\cal \dot{\omega}}\left( {\sf z}^{\prime }\right) +\eta _{\omega {\cal V}%
}\left( {\sf z,z}^{\prime }\right) {\cal \dot{V}}\left( {\sf z}^{\prime
}\right) \right\} d{\sf z}^{\prime }=\frac{\delta {\cal E}\left( {\cal %
\varkappa }\right) }{\delta {\cal \omega }\left( {\sf z}\right) }%
\Leftrightarrow \frac{\delta {\cal L}\left( {\cal \varkappa },{\cal \dot{%
\varkappa}}\right) }{\delta \omega \left( {\sf z}\right) }-\frac{d}{dt}%
\left\{ \frac{\delta {\cal L}\left( {\cal \varkappa },{\cal \dot{\varkappa}}%
\right) }{\delta \dot{\omega}\left( {\sf z}\right) }\right\} =0
\end{equation}%
leading to the equation of motion 
\begin{equation}
\left( 
\begin{array}{cc}
\eta _{{\cal VV}} & \eta _{{\cal V\omega }} \\ 
\eta _{\omega {\cal V}} & \eta _{\omega \omega }%
\end{array}%
\right) \left( 
\begin{array}{c}
{\cal \dot{V}} \\ 
{\cal \dot{\omega}}%
\end{array}%
\right) =\left( 
\begin{array}{c}
\nabla _{{\cal V}}{\cal E} \\ 
\nabla _{\omega }{\cal E}%
\end{array}%
\right)
\end{equation}%
The metric tensor components have the property $\eta _{{\cal V\omega }%
}\left( {\sf z,z}^{\prime }\right) =-\eta _{{\cal \omega V}}\left( {\sf z}%
^{\prime },{\sf z}\right) $, $\eta _{{\cal VV}}\left( {\sf z,z}^{\prime
}\right) =-\eta _{{\cal VV}}\left( {\sf z}^{\prime },{\sf z}\right) $ and $%
\eta _{{\cal \omega \omega }}\left( {\sf z,z}^{\prime }\right) =-\eta _{%
{\cal \omega \omega }}\left( {\sf z}^{\prime },{\sf z}\right) $ thus $\eta _{%
{\cal V\omega }}=-\eta _{{\cal \omega V}}^{t}$, $\eta _{{\cal VV}}=-\eta _{%
{\cal VV}}^{t}$ and $\eta _{{\cal \omega \omega }}=-\eta _{{\cal \omega
\omega }}^{t}$. In our case $\eta _{{\cal VV}}=\eta _{{\cal \omega \omega }%
}=0$ and even ${\cal Z}\left( {\cal V}\right) =0$ i.e. the generalized
momentum conjugate to ${\cal V}$ is {\em identically zero} for all $t$. The
metric tensor expressions in our case become 
\begin{eqnarray}
\eta _{{\cal VV}}\left( {\sf z,z}^{\prime }\right) &=&\eta _{{\cal \omega
\omega }}=0  \nonumber \\
\eta _{{\cal V\omega }}\left( {\sf z,z}^{\prime }\right) &=&\frac{\delta 
{\cal Z}\left( \omega \right) \left( {\sf z}^{\prime }\right) }{\delta {\cal %
V}\left( {\sf z}\right) }  \nonumber \\
\eta _{\omega {\cal V}}\left( {\sf z,z}^{\prime }\right) &=&-\frac{\delta 
{\cal Z}\left( \omega \right) \left( {\sf z}\right) }{\delta {\cal V}\left( 
{\sf z}^{\prime }\right) }
\end{eqnarray}%
If $\xi =\eta ^{-1}$ exists then the Lagrangian is said to be {\em non
degenerate }if it does not exist it is {\em degenerate}.

The condition 
\begin{equation}
\delta _{\omega }^{\left( 1\right) }{\cal A}\left( {\cal V},\omega \right) =0
\label{eomw}
\end{equation}
determines equations 
\begin{equation}
\eta _{\omega {\cal V}}{\cal \dot{V}+}\eta _{\omega \omega }{\cal \dot{\omega%
}}=\nabla _{\omega }{\cal E}
\end{equation}%
which in our case become 
\begin{equation}
\eta _{\omega {\cal V}}{\cal \dot{V}}=\nabla _{\omega }{\cal E}  \label{elw}
\end{equation}%
for values of ${\cal V}$ and $\omega $ that produce stationary action with
respect to small variations {\em of }$\omega .$ These can be solved in many
ways and are in general combinations of algebraic and differential
equations. At one extreme purely {\em algebraic} equations for $\omega $ 
\begin{equation}
\delta _{\omega }^{\left( 1\right) }{\cal A}\left( {\cal V}_{0},\omega
\right) =0
\end{equation}%
are obtained by fixing ${\cal V}={\cal V}_{0}$. Solving this algebraic
equation gives a solution $\omega _{\#}\left( {\cal V}_{0}\right) $ which
satisfies the Euler-Lagrange equation for $\omega $ for the fixed trajectory 
${\cal V}_{0}$. At the other extreme the equation eq. $\left( \ref{eomw}%
\right) $ can be viewed as an ordinary differential equation for ${\cal V}$
for a fixed value of $\omega =\omega _{0}$ i.e. viewed as 
\begin{equation}
\delta _{\omega }^{\left( 1\right) }{\cal A}\left( {\cal V},\omega
_{0}\right) =0
\end{equation}%
If one finds a solution ${\cal V}_{\#}\left( \omega _{0}\right) $ for this
ODE for ${\cal V}$ so that ${\cal \dot{V}}_{\#}=\eta _{\omega {\cal V}%
}^{-1}\nabla _{\omega }{\cal E}$ then 
\begin{equation}
\left( 
\begin{array}{cc}
0 & \eta _{{\cal V\omega }} \\ 
\eta _{\omega {\cal V}} & 0%
\end{array}%
\right) \left( 
\begin{array}{c}
\eta _{\omega {\cal V}}^{-1}\nabla _{\omega }{\cal E} \\ 
{\cal \dot{\omega}}%
\end{array}%
\right) -\left( 
\begin{array}{c}
\nabla _{{\cal V}}{\cal E} \\ 
\nabla _{\omega }{\cal E}%
\end{array}%
\right) =\left( 
\begin{array}{c}
\eta _{{\cal V\omega }}{\cal \dot{\omega}-}\nabla _{{\cal V}}{\cal E} \\ 
0%
\end{array}%
\right)
\end{equation}%
which shows that for the trajectory ${\cal V}_{\#}\left( \omega _{0}\right) $
the action is stationary with respect to small variations $\delta \omega $
about the fixed trajectory $\omega _{0}$, i.e. $\delta _{\omega }{\cal A}%
\left( {\cal V}_{\#}\left( \omega _{0}\right) ,\omega _{0}\right) =0.$ This
is a surprising result since this means that $\omega _{0}$ is a solution for
the Euler-Lagrange equations for $\omega $ for the fixed value ${\cal V}%
_{\#}\left( \omega _{0}\right) $, i.e. $\omega _{0}=\omega _{\#}\left( {\cal %
V}_{\#}\left( \omega _{0}\right) \right) $ and we have found a trajectory $%
{\cal V}_{\#}\left( \omega _{0}\right) $ that makes $\omega _{0}$ a solution
by solving an ODE for ${\cal V}$! Eq. $\left( \ref{elw}\right) $ can also be
solved by simultaneously varying ${\cal V}$ and $\omega $, resulting in
mixed algebraic and ODE's.

In a similar fashion the condition that gives equations for values of ${\cal %
V}$ and $\omega $ that produce stationary action with respect to small
variations {\em of } ${\cal V}$%
\begin{equation}
\delta _{{\cal V}}^{\left( 1\right) }{\cal A}\left( {\cal V},\omega \right)
=0
\end{equation}%
which in our case are 
\begin{equation}
\eta _{{\cal V\omega }}{\cal \dot{\omega}}=\nabla _{{\cal V}}{\cal E}
\end{equation}%
can also be solved as purely algebraic equations 
\begin{equation}
\delta _{{\cal V}}^{\left( 1\right) }{\cal A}\left( {\cal V},\omega
_{0}\right) =0
\end{equation}%
that determine a stationary trajectory ${\cal V}_{\#}\left( \omega
_{0}\right) $ for any fixed trajectory $\omega _{0}$ or can be considered as
an ordinary differential equations for ${\cal \omega }$, i.e. 
\begin{equation}
\delta _{{\cal V}}^{\left( 1\right) }{\cal A}\left( {\cal V}_{0},\omega
\right) =0
\end{equation}%
with a solution $\omega _{\#}\left( {\cal V}_{0}\right) $, that makes ${\cal %
V}_{0}$ a solution of the Euler-Lagrange equation for ${\cal V}$ for a fixed
a trajectory $\omega _{\#}\left( {\cal V}_{0}\right) $ i.e. ${\cal V}_{0}=%
{\cal V}_{\#}\left( \omega _{\#}\left( {\cal V}_{0}\right) \right) $ or by
simultaneously varying ${\cal V}$ and $\omega $ in mixed algebraic and ODE's.

The above analysis shows a type of duality between the two extreme
strategies of solving the stationarity problem, one either finds a
trajectory ${\cal V}_{\#}\left( \omega _{0}\right) $ that makes the action
stationary wrt to small variations of $\omega $ about $\left( {\cal V}%
_{\#}\left( \omega _{0}\right) ,\omega _{0}\right) $ or one finds a
trajectory $\omega _{\#}\left( {\cal V}_{0}\right) $ that makes the action
stationary wrt to small variations of $\omega $ about $\left( {\cal V}%
_{0},\omega _{\#}\left( {\cal V}_{0}\right) \right) $. The corresponding
situation holds for small ${\cal V}$ variations. Of course when one iterates
between the equations for ${\cal V}$ and $\omega $ and finds a self
consistent solution by either solving the algebraic equations or the ODE's
one arrives at a stationary point of the action wrt to small variations of $%
{\cal V}$ and $\omega $.

\subsection{EOM for ${\cal V}_{den},{\cal V}_{fib},\protect\omega _{den}$
and $\protect\omega _{fib}$}

The complete EOM eq. $\left( \ref{EOM}\right) $ in matrix integral form for
all the variables is 
\begin{equation}
\left( 
\begin{array}{cccc}
0 & 0 & \eta _{{\cal V}_{den}\omega _{den}} & \eta _{{\cal V}_{den}\omega
_{fib}} \\ 
0 & 0 & \eta _{{\cal V}_{fib}\omega _{den}} & \eta _{{\cal V}_{fib}\omega
_{fib}} \\ 
\eta _{\omega _{den}{\cal V}_{den}} & \eta _{\omega _{den}{\cal V}_{fib}} & 0
& 0 \\ 
\eta _{\omega _{fib}{\cal V}_{den}} & \eta _{\omega _{fib}{\cal V}_{fib}} & 0
& 0%
\end{array}%
\right) \left( 
\begin{array}{c}
{\cal \dot{V}}_{den} \\ 
{\cal \dot{V}}_{fib} \\ 
\dot{\omega}_{den} \\ 
\dot{\omega}_{fib}%
\end{array}%
\right) -\left( 
\begin{array}{c}
\nabla _{{\cal V}_{den}}{\cal E} \\ 
\nabla _{{\cal V}_{fib}}{\cal E} \\ 
\nabla _{\omega _{den}}{\cal E} \\ 
\nabla _{\omega _{fib}}{\cal E}%
\end{array}%
\right) =\left( 
\begin{array}{c}
0 \\ 
0 \\ 
0 \\ 
0%
\end{array}%
\right)
\end{equation}%
This equation can be written in terms of $\left\{ {\cal V}_{den},{\cal V}%
_{fib}({\cal V}_{den}),\omega \right\} $ as%
\begin{equation}
\left( 
\begin{array}{ccc}
0 & 0 & \eta _{{\cal V}_{den}\omega } \\ 
0 & 0 & \eta _{{\cal V}_{fib}\omega } \\ 
\eta _{\omega {\cal V}_{den}} & \eta _{\omega {\cal V}_{fib}} & 0%
\end{array}%
\right) \left( 
\begin{array}{c}
{\cal \dot{V}}_{den} \\ 
{\cal \dot{V}}_{fib} \\ 
\dot{\omega}%
\end{array}%
\right) -\left( 
\begin{array}{c}
\nabla _{{\cal V}_{den}}{\cal E} \\ 
\nabla _{{\cal V}_{fib}}{\cal E} \\ 
\nabla _{\omega }{\cal E}%
\end{array}%
\right) =\left( 
\begin{array}{c}
\nabla _{{\cal V}_{den}}{\cal \tilde{L}} \\ 
\nabla _{{\cal V}_{fib}}{\cal \tilde{L}} \\ 
\nabla _{\omega }{\cal \tilde{L}}%
\end{array}%
\right) =\left( 
\begin{array}{c}
0 \\ 
0 \\ 
0%
\end{array}%
\right)  \label{FullEOM}
\end{equation}%
where $\omega =\left( \omega _{den},\omega _{fib}\right) $ and ${\cal \tilde{%
L}}$ is defined in Appendix F. Using the relationships between the second
derivatives the metric matrix can be expressed in fewer terms as 
\begin{equation}
\left( 
\begin{array}{ccc}
0 & 0 & \eta _{{\cal V}_{den}\omega } \\ 
0 & 0 & \eta _{{\cal V}_{fib}\omega } \\ 
\eta _{\omega {\cal V}_{den}} & \eta _{\omega {\cal V}_{fib}} & 0%
\end{array}%
\right) =\left( 
\begin{array}{ccc}
0 & 0 & \eta _{{\cal V}_{den}\omega } \\ 
0 & 0 & \eta _{{\cal V}_{fib}\omega } \\ 
-\eta _{{\cal V}_{den}\omega }^{t} & -\eta _{{\cal V}_{fib}\omega }^{t} & 0%
\end{array}%
\right)  \label{EOMVW}
\end{equation}%
If one finds a vector for eq. $\left( \ref{FullEOM}\right) $ such that 
\begin{eqnarray}
&&\left( 
\begin{array}{c}
\nabla _{{\cal V}_{den}}{\cal \tilde{L}} \\ 
\nabla _{{\cal V}_{fib}}{\cal \tilde{L}} \\ 
\nabla _{\omega }{\cal \tilde{L}}%
\end{array}%
\right) =  \nonumber \\
&&\left( 
\begin{array}{ccc}
0 & 0 & \eta _{{\cal V}_{den}\omega } \\ 
0 & 0 & \eta _{{\cal V}_{fib}\omega } \\ 
\eta _{\omega {\cal V}_{den}} & \eta _{\omega {\cal V}_{fib}} & 0%
\end{array}%
\right) \left( 
\begin{array}{c}
{\cal \dot{V}}_{den} \\ 
{\cal \dot{V}}_{fib}^{\#} \\ 
\dot{\omega}^{\#}%
\end{array}%
\right) -\left( 
\begin{array}{c}
\nabla _{{\cal V}_{den}}{\cal E} \\ 
\nabla _{{\cal V}_{fib}}{\cal E} \\ 
\nabla _{\omega }{\cal E}%
\end{array}%
\right) =\left( 
\begin{array}{c}
\eta _{{\cal V}_{den}\omega }\dot{\omega}^{\#}-\nabla _{{\cal V}_{den}}{\cal %
E} \\ 
0 \\ 
0%
\end{array}%
\right)  \label{res}
\end{eqnarray}%
for a given ${\cal V}_{den}$ (we can think of the RHS\ of eq. $\left( \ref%
{res}\right) $ as a residue vector). Then ${\cal A}$ is stationary wrt
variations of ${\cal V}_{fib}$ and $\omega $ around the trajectories ${\cal V%
}_{fib}^{\#}$ and $\omega ^{\#}$ for the fixed trajectory ${\cal V}_{den}$
and we have found the values of the time dependent density functional along
the trajectory $\gamma \left( t\right) \equiv {\cal V}_{den}\left( t\right) $%
. This is equivalent to solving the algebraic equation 
\begin{equation}
\eta _{{\cal V}_{fib}\omega }\dot{\omega}-\nabla _{{\cal V}_{fib}}{\cal E}=0
\label{vfibeq}
\end{equation}%
self consistently for ${\cal V}_{fib}$, with fixed $\omega $, and 
\begin{equation}
\eta _{\omega {\cal V}_{den}}{\cal \dot{V}}_{den}+\eta _{\omega {\cal V}%
_{fib}}{\cal \dot{V}}_{fib}-\nabla _{\omega }{\cal E}=0  \label{omegaeq}
\end{equation}%
self consistently for $\omega $, with fixed ${\cal V}_{fib}$, or treating
them as coupled ODEs for $\omega $ and ${\cal V}_{fib}$%
\begin{equation}
\left( 
\begin{array}{cc}
0 & \eta _{{\cal V}_{fib}\omega } \\ 
\eta _{\omega {\cal V}_{fib}} & 0%
\end{array}%
\right) \left( 
\begin{array}{c}
{\cal \dot{V}}_{fib} \\ 
\dot{\omega}%
\end{array}%
\right) -\left( 
\begin{array}{c}
\nabla _{{\cal V}_{fib}}{\cal E} \\ 
\nabla _{\omega }{\cal E-}\eta _{\omega {\cal V}_{den}}{\cal \dot{V}}_{den}%
\end{array}%
\right) =0.
\end{equation}%
If one has found solutions $\left( {\cal V}_{fib}^{\#}\left( {\cal V}%
_{den}\right) ,\omega ^{\#}\left( {\cal V}_{den}\right) \right) $ for all $%
{\cal V}_{den}$ the equation 
\begin{equation}
\eta _{{\cal V}_{den}\omega }\dot{\omega}^{\#}-\nabla _{{\cal V}_{den}}{\cal %
E}=0  \label{gammaeq}
\end{equation}%
can be used to obtain the solution trajectory ${\cal V}_{den}^{\#}$. This is
all discussed in more detail later.

\section{Partitioning of the EOM}

\subsection{Holding the trajectory ${\cal V}_{den}$ constant}

First we consider the case when the trajectory ${\cal V}_{den}\left(
t\right) $ is constant (i.e. so that ${\cal V}_{den}\left( t\right) $ is no
longer a variational parameter) then the stationary action condition 
\begin{eqnarray}
&&\delta _{{\cal V}_{fib}({\cal V}_{den}),\omega }^{\left( 1\right) }{\cal A}%
\left( {\cal V}_{den},{\cal V}_{fib}({\cal V}_{den}),\omega \right) =0 
\nonumber \\
&&{\cal V}_{den}=fixed\,\,trajectory
\end{eqnarray}%
leads to the following Euler-Lagrange equations associated with a {\em %
degenerate }Lagrangian 
\begin{eqnarray}
\left( 
\begin{array}{cc}
0 & \eta _{{\cal V}_{fib}\omega } \\ 
\eta _{\omega {\cal V}_{fib}} & 0%
\end{array}%
\right) \left( 
\begin{array}{c}
{\cal \dot{V}}_{fib} \\ 
\dot{\omega}%
\end{array}%
\right) &=&\left( 
\begin{array}{c}
\nabla _{{\cal V}_{fib}}{\cal E} \\ 
\nabla _{\omega }{\cal E}%
\end{array}%
\right) \equiv  \nonumber \\
\left( 
\begin{array}{ccc}
0 & \eta _{{\cal V}_{fib}\omega _{den}} & \eta _{{\cal V}_{fib}\omega _{fib}}
\\ 
\eta _{\omega _{den}{\cal V}_{fib}} & 0 & 0 \\ 
\eta _{\omega _{fib}{\cal V}_{fib}} & 0 & 0%
\end{array}%
\right) \left( 
\begin{array}{c}
{\cal \dot{V}}_{fib} \\ 
\dot{\omega}_{den} \\ 
\dot{\omega}_{fib}%
\end{array}%
\right) &=&\left( 
\begin{array}{c}
\nabla _{{\cal V}_{fib}}{\cal E} \\ 
\nabla _{\omega _{den}}{\cal E} \\ 
\nabla _{\omega _{fib}}{\cal E}%
\end{array}%
\right)
\end{eqnarray}%
The same equation is obtained if the metric tensor is constant with respect
to variations in ${\cal V}_{den}$ which implies that the parameterization of 
$\Psi $ is linear in ${\cal V}_{den}.$

In order to solve this non linear ODE one must find a special solution and
solutions to the homogeneous equation 
\begin{equation}
\left( 
\begin{array}{ccc}
0 & \eta _{{\cal V}_{fib}\omega _{den}} & \eta _{{\cal V}_{fib}\omega _{fib}}
\\ 
\eta _{\omega _{den}{\cal V}_{fib}} & 0 & 0 \\ 
\eta _{\omega _{fib}{\cal V}_{fib}} & 0 & 0%
\end{array}%
\right) \left( 
\begin{array}{c}
{\cal \dot{V}}_{fib} \\ 
\dot{\omega}_{den} \\ 
\dot{\omega}_{fib}%
\end{array}%
\right) =0
\end{equation}%
This is the same as finding the generalized inverse of metric matrix on the
compliment of its kernel and then adding to the solution vector the subspace
of vectors belonging to the kernel. The ''driving vector'' on the RHS must
be in the range of the metric matrix for a solution to exist (a subsidiary
conditions must be invoked as shown below). The kernel of the fiber metric
matrix is spanned by eigen vectors of this matrix of the form%
\begin{equation}
\left( 
\begin{array}{c}
0 \\ 
k_{den} \\ 
k_{fib}%
\end{array}%
\right) =\left( 
\begin{array}{c}
0 \\ 
k_{den} \\ 
-\eta _{{\cal V}_{fib}\omega _{den}}^{-1}\eta _{{\cal V}_{fib}\omega
_{den}}k_{den}%
\end{array}%
\right)
\end{equation}%
as 
\begin{equation}
\left( 
\begin{array}{ccc}
0 & \eta _{{\cal V}_{fib}\omega _{den}} & \eta _{{\cal V}_{fib}\omega _{fib}}
\\ 
\eta _{\omega _{den}{\cal V}_{fib}} & 0 & 0 \\ 
\eta _{\omega _{fib}{\cal V}_{fib}} & 0 & 0%
\end{array}%
\right) \left( 
\begin{array}{c}
0 \\ 
k_{den} \\ 
k_{fib}%
\end{array}%
\right) =0
\end{equation}%
implies that 
\begin{equation}
\eta _{{\cal V}_{fib}\omega _{den}}k_{den}+\eta _{{\cal V}_{fib}\omega
_{fib}}k_{fib}=0\Rightarrow k_{fib}=-\eta _{{\cal V}_{fib}\omega
_{fib}}^{-1}\eta _{{\cal V}_{fib}\omega _{den}}k_{den}
\end{equation}%
The dimension of the kernel of the metric matrix is equal to $\dim \left\{
\ker \eta \right\} $, while the dimension of the domain and range of the
metric matrix is $\dim \left\{ \ker \eta ^{c}\right\} =\dim \left\{ {\bf W}%
_{den}\right\} $ where ${\bf W}_{den}$ is the space of $\omega _{den}$
vectors.

\subsection{Partitioning of the EOM wrt ${\cal V}_{den}$ and ${\cal V}_{fib}(%
{\cal V}_{den}),$ $\protect\omega _{den}({\cal V}_{den}),$ $\protect\omega %
_{fib}({\cal V}_{den})$}

The condition for the stationarity of the Action is 
\begin{equation}
\left( 
\begin{array}{cc}
0 & 
\begin{array}{ccc}
\qquad 0\qquad & \eta _{{\cal V}_{den}\omega _{den}} & \eta _{{\cal V}%
_{den}\omega _{fib}}%
\end{array}
\\ 
\begin{array}{c}
0 \\ 
-\eta _{{\cal V}_{den}\omega _{den}}^{t} \\ 
-\eta _{{\cal V}_{den}\omega _{fib}}^{t}%
\end{array}
& 
\begin{array}{ccc}
0 & \eta _{{\cal V}_{fib}\omega _{den}} & \eta _{{\cal V}_{fib}\omega _{fib}}
\\ 
-\eta _{{\cal V}_{fib}\omega _{den}}^{t} & 0 & 0 \\ 
-\eta _{{\cal V}_{fib}\omega _{fib}}^{t} & 0 & 0%
\end{array}%
\end{array}%
\right) \left( 
\begin{array}{c}
{\cal \dot{V}}_{den} \\ 
{\cal \dot{V}}_{fib} \\ 
\dot{\omega}_{den} \\ 
\dot{\omega}_{fib}%
\end{array}%
\right) =\left( 
\begin{array}{c}
\nabla _{{\cal V}_{den}}{\cal E} \\ 
\nabla _{{\cal V}_{fib}}{\cal E} \\ 
\nabla _{\omega _{den}}{\cal E} \\ 
\nabla _{\omega _{fib}}{\cal E}%
\end{array}%
\right)
\end{equation}%
which can also be written as 
\begin{eqnarray}
&&\left( 
\begin{array}{ccc}
0 & 0 & \eta _{{\cal V}_{den}\omega } \\ 
0 & 0 & \eta _{{\cal V}_{fib}\omega } \\ 
\eta _{\omega {\cal V}_{den}} & \eta _{\omega {\cal V}_{fib}} & 0%
\end{array}%
\right) \left( 
\begin{array}{c}
{\cal \dot{V}}_{den} \\ 
{\cal \dot{V}}_{fib} \\ 
\dot{\omega}%
\end{array}%
\right) =\left( 
\begin{array}{c}
\nabla _{{\cal V}_{den}}{\cal E} \\ 
\nabla _{{\cal V}_{fib}}{\cal E} \\ 
\nabla _{\omega }{\cal E}%
\end{array}%
\right)  \label{reom} \\
&&  \nonumber \\
&\equiv &\qquad \qquad \qquad \left( 
\begin{array}{cc}
0 & {\Bbb A}_{12} \\ 
{\Bbb A}_{21} & {\Bbb A}_{22}%
\end{array}%
\right) \left( 
\begin{array}{c}
{\bf x}_{1} \\ 
{\bf x}_{2}%
\end{array}%
\right) =\left( 
\begin{array}{c}
{\bf y}_{1} \\ 
{\bf y}_{2}%
\end{array}%
\right)
\end{eqnarray}%
where we introduce the shorthand notation ${\Bbb A},{\bf x}$, and ${\bf y}$
that was utilized in the preceding section. Eq. $\left( \ref{reom}\right) $.
This equation can be solved by first obtaining the partial solution 
\begin{equation}
\frame{$%
%TCIMACRO{
%\underset{}{\stackrel{}{{\Bbb A}_{22}\left( {\bf x}_{1},{\bf x}_{2}\right) {\bf x}_{2}={\bf y}_{2}\left( {\bf x}_{1},{\bf x}_{2}\right) -{\Bbb A}_{21}\left( {\bf x}_{1},{\bf x}_{2}\right) {\bf x}}_{1}}}%
%BeginExpansion
\mathrel{\mathop{\stackrel{}{{\Bbb A}_{22}\left( {\bf x}_{1},{\bf x}_{2}\right) {\bf x}_{2}={\bf y}_{2}\left( {\bf x}_{1},{\bf x}_{2}\right) -{\Bbb A}_{21}\left( {\bf x}_{1},{\bf x}_{2}\right) {\bf x}}_{1}}\limits_{}}%
%EndExpansion
$}
\end{equation}%
as described previously, which leads to the following equation for ${\cal V}%
_{fib}({\cal V}_{den}),$ $\omega _{den}({\cal V}_{den}),$ and $\omega _{fib}(%
{\cal V}_{den})$ for a fixed ${\cal V}_{den}$ 
\begin{eqnarray}
&&\left( 
\begin{array}{ccc}
0 & \eta _{{\cal V}_{fib}\omega _{den}} & \eta _{{\cal V}_{fib}\omega _{fib}}
\\ 
-\eta _{{\cal V}_{fib}\omega _{den}}^{t} & 0 & 0 \\ 
-\eta _{{\cal V}_{fib}\omega _{fib}}^{t} & 0 & 0%
\end{array}%
\right) \left( 
\begin{array}{c}
{\cal \dot{V}}_{fib} \\ 
\dot{\omega}_{den} \\ 
\dot{\omega}_{fib}%
\end{array}%
\right) =  \nonumber \\
&&\qquad \quad \left( 
\begin{array}{c}
\nabla _{{\cal V}_{fib}}{\cal E} \\ 
\nabla _{\omega _{den}}{\cal E} \\ 
\nabla _{\omega _{fib}}{\cal E}%
\end{array}%
\right) -\left( 
\begin{array}{c}
0 \\ 
-\eta _{{\cal V}_{den}\omega _{den}}^{t} \\ 
-\eta _{{\cal V}_{den}\omega _{fib}}^{t}%
\end{array}%
\right) \left( {\cal \dot{V}}_{den}\right) =\left( 
\begin{array}{c}
\nabla _{{\cal V}_{fib}}{\cal E} \\ 
\nabla _{\omega _{den}}{\cal E}+\eta _{{\cal V}_{den}\omega _{den}}^{t}{\cal 
\dot{V}}_{den} \\ 
\nabla _{\omega _{fib}}{\cal E}+\eta _{{\cal V}_{den}\omega _{fib}}^{t}{\cal 
\dot{V}}_{den}%
\end{array}%
\right)  \label{feom}
\end{eqnarray}%
i.e. 
\begin{equation}
\left( 
\begin{array}{ccc}
0 & \eta _{{\cal V}_{fib}\omega _{den}} & \eta _{{\cal V}_{fib}\omega _{fib}}
\\ 
-\eta _{{\cal V}_{fib}\omega _{den}}^{t} & 0 & 0 \\ 
-\eta _{{\cal V}_{fib}\omega _{fib}}^{t} & 0 & 0%
\end{array}%
\right) \left( 
\begin{array}{c}
{\cal \dot{V}}_{fib} \\ 
\dot{\omega}_{den} \\ 
\dot{\omega}_{fib}%
\end{array}%
\right) =\left( 
\begin{array}{c}
\nabla _{{\cal V}_{fib}}{\cal E} \\ 
\nabla _{\omega _{den}}{\cal E}+\eta _{{\cal V}_{den}\omega _{den}}^{t}{\cal 
\dot{V}}_{den} \\ 
\nabla _{\omega _{fib}}{\cal E}+\eta _{{\cal V}_{den}\omega _{fib}}^{t}{\cal 
\dot{V}}_{den}%
\end{array}%
\right)
\end{equation}%
This equation is singular and shows that the Lagrangian associated with the
constrained Action corresponding to a fixed trajectory ${\cal V}_{den}$ is
degenerate. The kernel of the metric tensor is spanned by vectors of the
form 
\begin{equation}
\left( 
\begin{array}{c}
0 \\ 
k_{di}\left( {\cal V}_{den}\right) \\ 
k_{fi}\left( {\cal V}_{den}\right)%
\end{array}%
\right) =\left( 
\begin{array}{c}
0 \\ 
k_{den\,i} \\ 
-\eta _{{\cal V}_{fib}\omega _{den}}^{-1}\eta _{{\cal V}_{fib}\omega
_{den}}k_{den\,i}%
\end{array}%
\right) \equiv k_{i}\left( {\cal V}_{den}\right) ;1\leq i\leq \dim \left\{
\ker \eta \left( {\cal V}_{den}\right) \right\}
\end{equation}%
The dimension of $\ker \eta \left( {\cal V}_{den}\right) $ can be seen to be
at least equal to $d_{den}$ where 
\begin{eqnarray}
\dim \left\{ {\bf V}_{den}\right\} &=&\dim \left\{ {\bf W}_{den}\right\}
=d_{den}  \nonumber \\
\dim \left\{ {\bf V}_{fib}\right\} &=&\dim \left\{ {\bf W}_{fib}\right\}
=d_{fib}
\end{eqnarray}%
and thus at most the dimension of $\ker \eta \left( {\cal V}_{den}\right)
^{c}$ is 
\begin{equation}
\dim \left\{ \ker \eta \left( {\cal V}_{den}\right) ^{c}\right\} =2d_{fib}
\end{equation}%
The generalized inverse of $\eta \left( {\cal V}_{den}\right) $ is then 
\begin{equation}
\eta \left( {\cal V}_{den}\right) ^{-1}=\sum_{1\leq i\leq d_{fib}}\beta
_{i}^{-1}\left( {\cal V}_{den}\right) \left\{ \left| s_{i}\left( {\cal V}%
_{den}\right) \right\rangle \left\langle s_{i+d_{fib}}^{d}\left( {\cal V}%
_{den}\right) \right| -\left| s_{i+d_{fib}}\left( {\cal V}_{den}\right)
\right\rangle \left\langle s_{i}^{d}\left( {\cal V}_{den}\right) \right|
\right\}
\end{equation}%
where the bases$\left\{ s_{i}\left( {\cal V}_{den}\right) \right\} $ spans
the complement of the kernel of $\eta \left( {\cal V}_{den}\right) $, (it
can be shown that $s_{i+d_{fib}}\left( {\cal V}_{den}\right) =s_{i}\left( 
{\cal V}_{den}\right) ^{\ast }$ since the tensor is real). The solution of
eq. $\left( \ref{feom}\right) $ is then the self consistent solution of 
\begin{eqnarray}
&&\left( 
\begin{array}{c}
{\cal \dot{V}}_{fib} \\ 
\dot{\omega}_{den} \\ 
\dot{\omega}_{fib}%
\end{array}%
\right) =\eta \left( {\cal V}_{den}\right) ^{-1}\left( 
\begin{array}{c}
\nabla _{{\cal V}_{fib}}{\cal E} \\ 
\nabla _{\omega _{den}}{\cal E}+\eta _{{\cal V}_{den}\omega _{den}}^{t}{\cal 
\dot{V}}_{den} \\ 
\nabla _{\omega _{fib}}{\cal E}+\eta _{{\cal V}_{den}\omega _{fib}}^{t}{\cal 
\dot{V}}_{den}%
\end{array}%
\right) +\sum_{1\leq i\leq m_{1}}\lambda _{i}\left( {\cal V}_{den}\right)
\left| k_{i}\left( {\cal V}_{den}\right) \right\rangle  \nonumber \\
&&  \nonumber \\
&=&\sum \beta _{i}^{-1}\left( {\cal V}_{den}\right) \left\{ \left|
s_{i}\left( {\cal V}_{den}\right) \right\rangle \left\langle s_{i}^{\ast
d}\left( {\cal V}_{den}\right) \right| \left. \left( 
\begin{array}{c}
\nabla _{{\cal V}_{fib}}{\cal E} \\ 
\nabla _{\omega _{den}}{\cal E}+\eta _{{\cal V}_{den}\omega _{den}}^{t}{\cal 
\dot{V}}_{den} \\ 
\nabla _{\omega _{fib}}{\cal E}+\eta _{{\cal V}_{den}\omega _{fib}}^{t}{\cal 
\dot{V}}_{den}%
\end{array}%
\right) \right\rangle -\right.  \nonumber \\
&&\left. \left| s_{i}^{\ast }\left( {\cal V}_{den}\right) \right\rangle
\left\langle s_{i}^{d}\left( {\cal V}_{den}\right) \right| \left. \left( 
\begin{array}{c}
\nabla _{{\cal V}_{fib}}{\cal E} \\ 
\nabla _{\omega _{den}}{\cal E}+\eta _{{\cal V}_{den}\omega _{den}}^{t}{\cal 
\dot{V}}_{den} \\ 
\nabla _{\omega _{fib}}{\cal E}+\eta _{{\cal V}_{den}\omega _{fib}}^{t}{\cal 
\dot{V}}_{den}%
\end{array}%
\right) \right\rangle \right\} +\sum_{1\leq i\leq m_{1}}\lambda _{i}\left( 
{\cal V}_{den}\right) \left| k_{i}\left( {\cal V}_{den}\right) \right\rangle
\nonumber \\
&\Longleftrightarrow &  \nonumber \\
&&\left( 
\begin{array}{c}
{\cal \dot{V}}_{fib} \\ 
\dot{\omega}%
\end{array}%
\right) =\sum \beta _{i}^{-1}\left( {\cal V}_{den}\right) \left\{ \left|
s_{i}\left( {\cal V}_{den}\right) \right\rangle \left\langle s_{i}^{\ast
d}\left( {\cal V}_{den}\right) \right| \left. \left( 
\begin{array}{c}
\nabla _{{\cal V}_{fib}}{\cal E} \\ 
\nabla _{\omega }{\cal E}+\eta _{{\cal V}_{den}\omega }^{t}{\cal \dot{V}}%
_{den}%
\end{array}%
\right) \right\rangle -\right.  \nonumber \\
&&\qquad \left. \left| s_{i}^{\ast }\left( {\cal V}_{den}\right)
\right\rangle \left\langle s_{i}^{d}\left( {\cal V}_{den}\right) \right|
\left. \left( 
\begin{array}{c}
\nabla _{{\cal V}_{fib}}{\cal E} \\ 
\nabla _{\omega }{\cal E}+\eta _{{\cal V}_{den}\omega }^{t}{\cal \dot{V}}%
_{den}%
\end{array}%
\right) \right\rangle \right\} +\sum_{1\leq i\leq m_{1}}\lambda _{i}\left( 
{\cal V}_{den}\right) \left| k_{i}\left( {\cal V}_{den}\right) \right\rangle
\end{eqnarray}%
with the constraint%
\begin{eqnarray}
&&\sum_{1\leq i\leq m_{1}}\left\{ \left| k_{di}\left( {\cal V}_{den}\right)
\right\rangle \left\langle k_{i}^{d}\left( {\cal V}_{den}\right) \right|
\left. \nabla _{\omega _{den}}{\cal E}+\eta _{{\cal V}_{den}\omega
_{den}}^{t}{\cal \dot{V}}_{den}\right\rangle \right. +  \nonumber \\
&&\quad \left. \left| k_{di}\left( {\cal V}_{den}\right) \right\rangle
\left\langle k_{i}^{d}\left( {\cal V}_{den}\right) \right| \left. \nabla
_{\omega _{fib}}{\cal E}+\eta _{{\cal V}_{den}\omega _{fib}}^{t}{\cal \dot{V}%
}_{den}\right\rangle \right\} =0  \nonumber \\
&\equiv &\sum_{1\leq i\leq m_{1}}\left\{ \left| k_{di}\left( {\cal V}%
_{den}\right) \right\rangle \left\langle k_{i}^{d}\left( {\cal V}%
_{den}\right) \right| \left. \nabla _{\omega }{\cal E}+\eta _{{\cal V}%
_{den}\omega }^{t}{\cal \dot{V}}_{den}\right\rangle \right\} =0  \nonumber \\
&\Rightarrow &\frame{$%
%TCIMACRO{
%\underset{}{\stackrel{}{\left\langle k_{i}^{d}\left( {\cal V}_{den}\right) |\nabla _{\omega }{\cal E}\left( {\cal V}_{den}\right) \right\rangle =-\left\langle k_{i}^{d}\left( {\cal V}_{den}\right) |\eta _{{\cal V}_{den}\omega }^{t}{\cal \dot{V}}_{den}\right\rangle \quad \forall i}}}%
%BeginExpansion
\mathrel{\mathop{\stackrel{}{\left\langle k_{i}^{d}\left( {\cal V}_{den}\right) |\nabla _{\omega }{\cal E}\left( {\cal V}_{den}\right) \right\rangle =-\left\langle k_{i}^{d}\left( {\cal V}_{den}\right) |\eta _{{\cal V}_{den}\omega }^{t}{\cal \dot{V}}_{den}\right\rangle \quad \forall i}}\limits_{}}%
%EndExpansion
$}
\end{eqnarray}%
where $\left\{ \lambda _{i}\left( {\cal V}_{den}\right) ;1\leq i\leq
m_{1}\right\} $ are arbitrary complex coefficients. In order to obtain $%
\left( 
\begin{array}{c}
{\cal V}_{fib}\left( {\cal V}_{den},\lambda \right) \\ 
\omega \left( {\cal V}_{den},\lambda \right)%
\end{array}%
\right) $ one must integrate this solution, which we discuss in depth later.

The self consistent solution for ${\cal \dot{V}}_{den}$ is obtained by using
eq. $\left( \ref{scx1}\right) $ 
\begin{equation}
\frame{$%
%TCIMACRO{
%\underset{}{\stackrel{}{{\Bbb A}_{12}{\Bbb A}_{22}^{-1}{\Bbb A}_{21}{\bf x}_{1\#\left( k+1\right) }={\Bbb A}_{12}{\Bbb A}_{22}^{-1}{\bf y}_{2}-{\bf y}_{1}+{\Bbb A}_{12}{\bf k}}}}%
%BeginExpansion
\mathrel{\mathop{\stackrel{}{{\Bbb A}_{12}{\Bbb A}_{22}^{-1}{\Bbb A}_{21}{\bf x}_{1\#\left( k+1\right) }={\Bbb A}_{12}{\Bbb A}_{22}^{-1}{\bf y}_{2}-{\bf y}_{1}+{\Bbb A}_{12}{\bf k}}}\limits_{}}%
%EndExpansion
$}
\end{equation}%
(which is equivalent to eq. $\left( \ref{gammaeq}\right) $ that implicitly
depends on ${\cal \dot{V}}_{den})$, and the notation 
\begin{equation}
{\cal V}_{d}={\cal V}_{den};\,{\cal V}_{f}={\cal V}_{fib};\,{\cal \omega }%
_{d}={\cal \omega }_{den};\,{\cal \omega }_{f}={\cal \omega }_{fib}
\end{equation}%
to arrive at 
\begin{eqnarray}
&&\left( 
\begin{array}{ccc}
0 & \eta _{{\cal V}_{d}\omega _{d}} & \eta _{{\cal V}_{d}\omega _{f}}%
\end{array}%
\right) \left( 
\begin{array}{ccc}
\left( \eta ^{-1}\right) _{{\cal V}_{f}{\cal V}_{f}} & \left( \eta
^{-1}\right) _{{\cal V}_{f}{\cal \omega }d} & \left( \eta ^{-1}\right) _{%
{\cal V}_{f}{\cal \omega }_{f}} \\ 
\left( \eta ^{-1}\right) _{{\cal \omega }_{d}{\cal V}_{f}} & \left( \eta
^{-1}\right) _{{\cal \omega }_{d}{\cal \omega }_{d}} & \left( \eta
^{-1}\right) _{{\cal \omega }_{d}{\cal \omega }_{f}} \\ 
\left( \eta ^{-1}\right) _{{\cal \omega }_{f}{\cal V}_{f}} & \left( \eta
^{-1}\right) _{{\cal \omega }_{f}{\cal \omega }_{d}} & \left( \eta
^{-1}\right) _{{\cal \omega }_{f}{\cal \omega }_{f}}%
\end{array}%
\right) \left( 
\begin{array}{c}
0 \\ 
-\eta _{{\cal V}_{d}\omega _{d}}^{t} \\ 
-\eta _{{\cal V}_{d}\omega _{f}}^{t}%
\end{array}%
\right) \left| {\cal \dot{V}}_{d}\right\rangle =  \nonumber \\
&&  \nonumber \\
&&\left( 
\begin{array}{ccc}
0 & \eta _{{\cal V}_{d}\omega _{d}} & \eta _{{\cal V}_{d}\omega _{f}}%
\end{array}%
\right) \left( 
\begin{array}{ccc}
\left( \eta ^{-1}\right) _{{\cal V}_{f}{\cal V}_{f}} & \left( \eta
^{-1}\right) _{{\cal V}_{f}{\cal \omega }d} & \left( \eta ^{-1}\right) _{%
{\cal V}_{f}{\cal \omega }_{f}} \\ 
\left( \eta ^{-1}\right) _{{\cal \omega }_{d}{\cal V}_{f}} & \left( \eta
^{-1}\right) _{{\cal \omega }_{d}{\cal \omega }_{d}} & \left( \eta
^{-1}\right) _{{\cal \omega }_{d}{\cal \omega }_{f}} \\ 
\left( \eta ^{-1}\right) _{{\cal \omega }_{f}{\cal V}_{f}} & \left( \eta
^{-1}\right) _{{\cal \omega }_{f}{\cal \omega }_{d}} & \left( \eta
^{-1}\right) _{{\cal \omega }_{f}{\cal \omega }_{f}}%
\end{array}%
\right) \left| \left( 
\begin{array}{c}
\nabla _{{\cal V}_{f}}{\cal E} \\ 
\nabla _{\omega _{f}}{\cal E} \\ 
\nabla _{\omega _{d}}{\cal E}%
\end{array}%
\right) \right\rangle  \nonumber \\
&&\qquad \qquad \qquad \qquad \qquad \qquad -\nabla _{{\cal V}_{d}}{\cal E+\,%
}\sum_{1\leq i\leq m_{1}}\lambda _{i}\left( {\cal V}_{d}\right) \left( 
\begin{array}{ccc}
0 & \eta _{{\cal V}_{d}\omega _{d}} & \eta _{{\cal V}_{d}\omega _{f}}%
\end{array}%
\right) \left| \left( 
\begin{array}{c}
0 \\ 
k_{d\,i} \\ 
k_{f\,i}%
\end{array}%
\right) \right\rangle
\end{eqnarray}%
which can be developed to give 
\begin{eqnarray}
&&\left\{ \eta _{{\cal V}_{d}\omega _{d}}\left( \eta ^{-1}\right) _{{\cal %
\omega }_{d}{\cal \omega }_{d}}\eta _{{\cal V}_{d}\omega _{d}}^{t}+\eta _{%
{\cal V}_{d}\omega _{d}}\left( \eta ^{-1}\right) _{{\cal \omega }_{d}{\cal %
\omega }_{f}}\eta _{{\cal V}_{d}\omega _{f}}^{t}\right. +\eta _{{\cal V}%
_{d}\omega _{f}}\left( \eta ^{-1}\right) _{{\cal \omega }_{f}{\cal \omega }%
_{d}}\eta _{{\cal V}_{d}\omega _{d}}^{t}+  \nonumber \\
&&\qquad \qquad \qquad \qquad \qquad \qquad \qquad \qquad \qquad \qquad
\qquad \left. +\eta _{{\cal V}_{d}\omega _{f}}\left( \eta ^{-1}\right) _{%
{\cal \omega }_{f}{\cal \omega }_{f}}\eta _{{\cal V}_{d}\omega
_{f}}^{t}\right\} \left| {\cal \dot{V}}_{d}\right\rangle =  \nonumber \\
&&  \nonumber \\
&&-\left\{ \eta _{{\cal V}_{d}\omega _{d}}\left( \eta ^{-1}\right) _{{\cal %
\omega }_{d}{\cal V}_{f}}+\eta _{{\cal V}_{d}\omega _{f}}\left( \eta
^{-1}\right) _{{\cal \omega }_{f}{\cal V}_{f}}\right\} \nabla _{{\cal V}%
_{fib}}{\cal E+}\left\{ \eta _{{\cal V}_{d}\omega _{d}}\left( \eta
^{-1}\right) _{{\cal \omega }_{d}{\cal \omega }_{d}}+\eta _{{\cal V}%
_{d}\omega _{f}}\left( \eta ^{-1}\right) _{{\cal \omega }_{f}{\cal \omega }%
_{d}}\right\} \nabla _{\omega _{f}}{\cal E}  \nonumber \\
&&\quad \qquad \quad \qquad \quad \qquad \quad \qquad \quad \qquad \quad
\quad +\left\{ \eta _{{\cal V}_{d}\omega _{d}}\left( \eta ^{-1}\right) _{%
{\cal \omega }_{f}{\cal \omega }_{d}}+\eta _{{\cal V}_{d}\omega _{f}}\left(
\eta ^{-1}\right) _{{\cal V}_{f}{\cal \omega }_{f}}\right\} \nabla _{\omega
_{d}}{\cal E}-\nabla _{{\cal V}_{d}}{\cal E}  \nonumber \\
&&\quad \quad \qquad \quad \qquad \quad \qquad \quad \qquad \quad \qquad
\quad \,{\cal +}\sum_{1\leq i\leq m_{1}}\lambda _{i}\left( {\cal V}%
_{d}\right) \left\{ \eta _{{\cal V}_{d}\omega _{d}}\left|
k_{d\,i}\right\rangle +\eta _{{\cal V}_{d}\omega _{f}}\left|
k_{f\,i}\right\rangle \right\}
\end{eqnarray}%
which can be written in a very concise form as 
\begin{equation}
{\frak M}\left( {\cal V}_{d},t\right) \left| {\cal \dot{V}}_{d}\right\rangle
={\frak R}\left( {\cal V}_{d},t\right)
\end{equation}%
where 
\begin{eqnarray}
{\frak M}_{1}\left( {\cal V}_{d},t\right) &=&\eta _{{\cal V}_{d}\omega
_{d}}\left( \eta ^{-1}\right) _{{\cal \omega }_{d}{\cal \omega }_{d}}\eta _{%
{\cal V}_{d}\omega _{d}}^{t}  \nonumber \\
{\frak M}_{2}\left( {\cal V}_{d},t\right) &=&\eta _{{\cal V}_{d}\omega
_{d}}\left( \eta ^{-1}\right) _{{\cal \omega }_{d}{\cal \omega }_{f}}\eta _{%
{\cal V}_{d}\omega _{f}}^{t}  \nonumber \\
{\frak M}_{3}\left( {\cal V}_{d},t\right) &=&\eta _{{\cal V}_{d}\omega
_{f}}\left( \eta ^{-1}\right) _{{\cal \omega }_{f}{\cal \omega }_{d}}\eta _{%
{\cal V}_{d}\omega _{d}}^{t}  \nonumber \\
{\frak M}_{4}\left( {\cal V}_{d},t\right) &=&\eta _{{\cal V}_{d}\omega
_{f}}\left( \eta ^{-1}\right) _{{\cal \omega }_{f}{\cal \omega }_{f}}\eta _{%
{\cal V}_{d}\omega _{f}}^{t}  \nonumber \\
{\frak M}\left( {\cal V}_{d},t\right) &=&{\frak M}_{1}\left( {\cal V}%
_{d},t\right) +{\frak M}_{2}\left( {\cal V}_{d},t\right) +{\frak M}%
_{3}\left( {\cal V}_{d},t\right) +{\frak M}_{4}\left( {\cal V}_{d},t\right)
\end{eqnarray}%
and the matrices $\left\{ {\frak M}_{j}\left( {\cal V}_{d},t\right) ;1\leq
j\leq 4\right\} $ all are $\dim \left\{ {\bf V}_{d}\right\} \times \dim
\left\{ {\bf V}_{d}\right\} $. In a similar fashion we decompose ${\frak R}%
\left( {\cal V}_{d},t\right) $ as 
\begin{eqnarray}
{\frak R}_{1}\left( {\cal V}_{d},t\right) &=&\left\{ \eta _{{\cal V}%
_{d}\omega _{d}}\left( \eta ^{-1}\right) _{{\cal \omega }_{d}{\cal V}%
_{f}}+\eta _{{\cal V}_{d}\omega _{f}}\left( \eta ^{-1}\right) _{{\cal \omega 
}_{f}{\cal V}_{f}}\right\} \nabla _{{\cal V}_{fib}}{\cal E}  \nonumber \\
{\frak R}_{2}\left( {\cal V}_{d},t\right) &=&\left\{ \eta _{{\cal V}%
_{d}\omega _{d}}\left( \eta ^{-1}\right) _{{\cal \omega }_{d}{\cal \omega }%
_{f}}+\eta _{{\cal V}_{d}\omega _{f}}\left( \eta ^{-1}\right) _{{\cal V}_{f}%
{\cal \omega }_{f}}\right\} \nabla _{\omega _{f}}{\cal E}  \nonumber \\
{\frak R}_{3}\left( {\cal V}_{d},t\right) &=&\left\{ \eta _{{\cal V}%
_{d}\omega _{f}}\left( \eta ^{-1}\right) _{{\cal \omega }_{f}{\cal \omega }%
_{d}}+\eta _{{\cal V}_{d}\omega _{f}}\left( \eta ^{-1}\right) _{{\cal V}_{f}%
{\cal \omega }_{d}}\right\} \nabla _{\omega _{d}}{\cal E}  \nonumber \\
{\frak R}_{4}\left( {\cal V}_{d},t\right) &=&\sum_{1\leq i\leq m_{1}}\lambda
_{i}\left( {\cal V}_{d}\right) \left\{ \eta _{{\cal V}_{d}\omega _{d}}\left|
k_{d\,i}\right\rangle +\eta _{{\cal V}_{d}\omega _{f}}\left|
k_{f\,i}\right\rangle \right\}  \nonumber \\
{\frak R}\left( {\cal V}_{d},t\right) &=&{\frak R}_{1}\left( {\cal V}%
_{d},t\right) +{\frak R}_{2}\left( {\cal V}_{d},t\right) +{\frak R}%
_{3}\left( {\cal V}_{d},t\right) +{\frak R}_{4}\left( {\cal V}_{d},t\right)
-\nabla _{{\cal V}_{d}}{\cal E}
\end{eqnarray}

\section{Solutions of Nonlinear Ordinary Differential Equations}

If eq. $\left( \ref{nl}\right) $ is a nonlinear ODE of the form 
\begin{equation}
{\cal A}\left( x\left( t\right) \right) \left| \dot{x}\left( t\right)
\right\rangle =\left| \nabla {\cal E}\left( x\left( t\right) \right)
\right\rangle ;\quad \ t\in \left[ t_{i},t_{f}\right]
\end{equation}%
we can first use the Taylor series expansion of the force about an arbitrary
trajectory $x_{0}$ to obtain a first order approximation of $x$ about $x_{0}$
\begin{equation}
\nabla {\cal E}\left( x\left( t\right) \right) =\nabla {\cal E}\left(
x_{0}\left( t\right) +y\left( t\right) \right) \simeq \nabla {\cal E}\left(
x_{0}\left( t\right) \right) +D_{x}^{2}{\cal E}\left( x_{0}\left( t\right)
\right) \left| y\left( t\right) \right\rangle  \label{gradexpan}
\end{equation}%
where 
\begin{equation}
D_{x}^{2}{\cal E}\left( x_{0}\left( t\right) \right) =D_{y}\nabla {\cal E}%
\left( x_{0}\left( t\right) \right)
\end{equation}%
which can be seen from%
\[
\nabla {\cal E}\left( x_{0}\left( t\right) \right) \left( \delta x\left(
t\right) \right) =\nabla {\cal E}\left( x_{0}\left( t\right) \right) \left(
y\left( t\right) \right) =\int \frac{\delta {\cal E}\left( x_{0}\left(
t\right) \right) }{\delta x\left( {\sf z}\right) }y\left( t,{\sf z}\right) d%
{\sf z} 
\]%
\[
\frac{\delta }{\delta x\left( {\sf z}\right) }=\frac{\delta }{\delta y\left( 
{\sf z}\right) }\text{ }\qquad \text{when }x\left( {\sf z}\right) =y\left( 
{\sf z}\right) +x_{0}\left( {\sf z}\right) \Rightarrow \qquad \qquad 
\]%
\begin{eqnarray}
D_{y}\nabla {\cal E}\left( x_{0}\left( t\right) \right) \left( y\left(
t\right) ,y\left( t\right) \right) &=&\int \frac{\delta ^{2}{\cal E}\left(
x_{0}\left( t\right) \right) }{\delta x\left( {\sf z}^{\prime }\right)
\delta x\left( {\sf z}\right) }y\left( t,{\sf z}\right) y\left( t,{\sf z}%
^{\prime }\right) d{\sf z}d{\sf z}^{\prime }  \nonumber \\
&=&D_{x}^{2}{\cal E}\left( x_{0}\left( t\right) \right) \left( y\left(
t\right) \otimes y\left( t\right) \right)
\end{eqnarray}%
Eq. $\left( \ref{gradexpan}\right) $gives us the linearized equation 
\begin{equation}
{\cal A}\left( x_{0}\left( t\right) \right) \left| \dot{y}\left( t\right)
\right\rangle =\nabla {\cal E}\left( x_{0}\left( t\right) \right) +D_{x}^{2}%
{\cal E}\left( x_{0}\left( t\right) \right) \left| y\left( t\right)
\right\rangle
\end{equation}%
which gives using the generalized inverse of ${\cal A}\left( x_{0}\left(
t\right) \right) $ 
\begin{eqnarray}
\left| \dot{y}\left( t\right) \right\rangle &=&{\cal A}\left( x_{0}\left(
t\right) \right) ^{-1}D_{x}^{2}{\cal E}\left( x_{0}\left( t\right) \right)
\left| y\left( t\right) \right\rangle +{\cal A}\left( x_{0}\left( t\right)
\right) ^{-1}\nabla {\cal E}\left( x_{0}\left( t\right) \right) +\left|
\lambda \left( t\right) \right\rangle  \nonumber \\
&\equiv &\left| \dot{y}\left( t\right) \right\rangle ={\frak H}\left(
t\right) \left| y\left( t\right) \right\rangle +\left| {\frak f}\left(
t\right) \right\rangle  \label{ODE}
\end{eqnarray}%
where 
\begin{eqnarray}
{\frak H}\left( t\right) &=&{\cal A}\left( x_{0}\left( t\right) \right)
^{-1}D_{x}^{2}{\cal E}\left( x_{0}\left( t\right) \right)  \nonumber \\
\left| {\frak f}\left( t\right) \right\rangle &=&{\cal A}\left( x_{0}\left(
t\right) \right) ^{-1}\nabla {\cal E}\left( x_{0}\left( t\right) \right)
+\left| \lambda \left( t\right) \right\rangle
\end{eqnarray}%
where 
\begin{equation}
\left| \lambda \left( t\right) \right\rangle =\sum_{1\leq i\leq
m_{1}}\lambda _{i}\left| k_{i}\left( x_{0}\left( t\right) \right)
\right\rangle \quad \in \ker \left\{ {\cal A}\left( x_{0}\left( t\right)
\right) \right\}
\end{equation}%
After finding a solution $y_{\#}\left( t\right) $ we start at a new
trajectory $x_{0}\left( t\right) +y_{\#}\left( t\right) $ and continue until
we have self consistency. One way to obtain a solution to the ODE\ eq. $%
\left( \ref{ODE}\right) $ is use the short time expansion 
\begin{equation}
\left| y\left( t_{k+1}\right) \right\rangle =\left| y\left( t_{k}\right)
\right\rangle +\left\{ {\frak H}\left( t_{k}\right) \left| y\left(
t_{k}\right) \right\rangle +{\frak f}\left( t_{k}\right) \right\} \Delta t
\end{equation}%
then iterate 
\begin{eqnarray}
\left| y\left( t_{k+2}\right) \right\rangle &=&\left| y\left( t_{k+1}\right)
\right\rangle +\left\{ {\frak H}\left( t_{k+1}\right) \left| y\left(
t_{k+1}\right) \right\rangle +{\frak f}\left( t_{k+1}\right) \right\} \Delta
t\qquad \qquad \qquad \qquad \qquad \qquad \qquad \,  \nonumber \\
&=&\left| y\left( t_{k}\right) \right\rangle +\left\{ {\frak H}\left(
t_{k}\right) \left| y\left( t_{k}\right) \right\rangle +{\frak f}\left(
t_{k}\right) \right\} \Delta t+  \nonumber \\
&&\left\{ {\frak H}\left( t_{k+1}\right) \left\{ \left| y\left( t_{k}\right)
\right\rangle +\left\{ {\frak H}\left( t_{k}\right) \left| y\left(
t_{k}\right) \right\rangle +{\frak f}\left( t_{k}\right) \right\} \Delta
t\right\} +{\frak f}\left( t_{k+1}\right) \right\} \Delta t  \nonumber \\
&=&\left| y\left( t_{k}\right) \right\rangle +\left\{ {\frak H}\left(
t_{k}\right) \left| y\left( t_{k}\right) \right\rangle +{\frak f}\left(
t_{k}\right) \right\} \Delta t  \nonumber \\
&&+{\frak H}\left( t_{k+1}\right) \left| y\left( t_{k}\right) \right\rangle
\Delta t+{\frak H}\left( t_{k+1}\right) {\frak H}\left( t_{k}\right) \left|
y\left( t_{k}\right) \right\rangle \Delta t\Delta t  \nonumber \\
&&\qquad \qquad \qquad \qquad \qquad \qquad +{\frak H}\left( t_{k+1}\right) 
{\frak f}\left( t_{k}\right) \Delta t\Delta t+{\frak f}\left( t_{k+1}\right)
\Delta t  \nonumber \\
&=&\left| y\left( t_{k}\right) \right\rangle +\left( {\frak H}\left(
t_{k}\right) +{\frak H}\left( t_{k+1}\right) \right) \left| y\left(
t_{k}\right) \right\rangle \Delta t+\left( {\frak f}\left( t_{k}\right) +%
{\frak f}\left( t_{k+1}\right) \right) \Delta t  \nonumber \\
&&+{\frak H}\left( t_{k+1}\right) {\frak H}\left( t_{k}\right) \left|
y\left( t_{k}\right) \right\rangle \Delta t\Delta t+{\frak H}\left(
t_{k+1}\right) {\frak f}\left( t_{k}\right) \Delta t\Delta t
\end{eqnarray}%
\begin{eqnarray}
\left| y\left( t_{k+3}\right) \right\rangle &=&\left| y\left( t_{k+2}\right)
\right\rangle +\left\{ {\frak H}\left( t_{k+2}\right) \left| y\left(
t_{k+2}\right) \right\rangle +{\frak f}\left( t_{k+2}\right) \right\} \Delta
t  \nonumber \\
&=&\left| y\left( t_{k}\right) \right\rangle +\left( {\frak H}\left(
t_{k}\right) +{\frak H}\left( t_{k+1}\right) \right) \left| y\left(
t_{k}\right) \right\rangle \Delta t+\left( {\frak f}\left( t_{k}\right) +%
{\frak f}\left( t_{k+1}\right) \right) \Delta t  \nonumber \\
&&+{\frak H}\left( t_{k+1}\right) {\frak H}\left( t_{k}\right) \left|
y\left( t_{k}\right) \right\rangle \Delta t\Delta t+{\frak H}\left(
t_{k+1}\right) {\frak f}\left( t_{k}\right) \Delta t\Delta t  \nonumber \\
&&+\left\{ {\frak H}\left( t_{k+2}\right) \left\{ \left| y\left(
t_{k}\right) \right\rangle +\left( {\frak H}\left( t_{k}\right) +{\frak H}%
\left( t_{k+1}\right) \right) \left| y\left( t_{k}\right) \right\rangle
\Delta t+\left( {\frak f}\left( t_{k}\right) +{\frak f}\left( t_{k+1}\right)
\right) \Delta t\right. \right.  \nonumber \\
&&\left. \left. +{\frak H}\left( t_{k+1}\right) {\frak H}\left( t_{k}\right)
\left| y\left( t_{k}\right) \right\rangle \Delta t\Delta t+{\frak H}\left(
t_{k+1}\right) {\frak f}\left( t_{k}\right) \Delta t\Delta t\right\} +{\frak %
f}\left( t_{k+2}\right) \right\} \Delta t  \nonumber \\
&=&\left| y\left( t_{k}\right) \right\rangle +\left( {\frak H}\left(
t_{k}\right) +{\frak H}\left( t_{k+1}\right) +{\frak H}\left( t_{k+2}\right)
\right) \left| y\left( t_{k}\right) \right\rangle \Delta t  \nonumber \\
&&+\left( {\frak H}\left( t_{k+1}\right) {\frak H}\left( t_{k}\right) +%
{\frak H}\left( t_{k+2}\right) {\frak H}\left( t_{k+1}\right) +{\frak H}%
\left( t_{k+2}\right) {\frak H}\left( t_{k}\right) \right) \left| y\left(
t_{k}\right) \right\rangle \Delta t\Delta t  \nonumber \\
&&+{\frak H}\left( t_{k+2}\right) {\frak H}\left( t_{k+1}\right) {\frak H}%
\left( t_{k}\right) \left| y\left( t_{k}\right) \right\rangle \Delta t\Delta
t\Delta t  \nonumber \\
&&+\left( {\frak f}\left( t_{k}\right) +{\frak f}\left( t_{k+1}\right) +%
{\frak f}\left( t_{k+2}\right) \right) \Delta t+{\frak H}\left(
t_{k+2}\right) \left( {\frak f}\left( t_{k}\right) +{\frak f}\left(
t_{k+1}\right) \right) \Delta t\Delta t  \nonumber \\
&&+{\frak H}\left( t_{k+1}\right) {\frak f}\left( t_{k}\right) \Delta
t\Delta t+{\frak H}\left( t_{k+2}\right) {\frak H}\left( t_{k+1}\right) 
{\frak f}\left( t_{k}\right) \Delta t\Delta t\Delta t
\end{eqnarray}%
Thus 
\begin{eqnarray}
&&\left| y\left( t_{k+r}\right) \right\rangle =\left| y\left( t_{k}\right)
\right\rangle +\sum_{1\leq i\leq r}{\frak f}\left( t_{k+i}\right) \Delta
+\sum_{1\leq i\leq r}{\frak H}\left( t_{k+i}\right) \Delta t\left| y\left(
t_{k}\right) \right\rangle  \nonumber \\
&&+\sum_{1\leq i_{1}<i_{2}\leq r}{\frak H}\left( t_{k+i_{2}}\right) {\frak H}%
\left( t_{k+i_{1}}\right) \Delta t\Delta t\left| y\left( t_{k}\right)
\right\rangle +\cdots +{\frak H}\left( t_{k+i_{r}}\right) \cdots {\frak H}%
\left( t_{k+i_{1}}\right) \Delta t\cdots \Delta t\left| y\left( t_{k}\right)
\right\rangle t  \nonumber \\
&&+\sum_{1\leq i_{1}<i_{2}\leq r}{\frak H}\left( t_{k+i_{2}}\right) {\frak f}%
\left( t_{k+i_{1}}\right) \Delta t\Delta t+\cdots +{\frak H}\left(
t_{k+i_{r}}\right) \cdots {\frak H}\left( t_{k+i_{2}}\right) {\frak f}\left(
t_{k+i_{1}}\right) \Delta t\cdots \Delta t
\end{eqnarray}%
which gives in the limit 
\begin{eqnarray}
&&\left| y_{\#}\left( t\right) \right\rangle ={\LARGE T}e^{\int_{t_{i}}^{t}%
{\frak H}\left( t^{\prime }\right) dt^{\prime }}\left| y\left( t_{i}\right)
\right\rangle +\int_{t_{i}}^{t}{\LARGE T}e^{\int_{t_{i}}^{t}{\frak H}\left(
t^{\prime }-s\right) dt^{\prime }}{\frak f}\left( s\right) ds  \nonumber \\
&\equiv &{\LARGE T}e^{\int_{t_{i}}^{t}{\cal A}\left( x_{0}\left( t^{\prime
}\right) \right) ^{-1}D_{y}\nabla {\cal E}\left( x_{0}\left( t^{\prime
}\right) \right) dt^{\prime }}\left| y\left( t_{i}\right) \right\rangle 
\nonumber \\
&&+\int_{t_{i}}^{t}{\LARGE T}e^{\int_{t_{i}}^{t}{\cal A}\left( x_{0}\left(
t^{\prime }-s\right) \right) ^{-1}D_{y}\nabla {\cal E}\left( x_{0}\left(
t^{\prime }-s\right) \right) dt^{\prime }}{\cal A}\left( x_{0}\left(
s\right) \right) ^{-1}\nabla {\cal E}\left( x_{0}\left( s\right) \right)
ds+\int_{t_{i}}^{t}\left| \lambda \left( s\right) \right\rangle ds
\end{eqnarray}%
where we use the generalized inverse ${\cal A}\left( x\left( t\right)
\right) ^{-1}$ and invoked ''tractability conditions'' 
\begin{equation}
{\cal A}\left( x\left( t\right) \right) ^{-1}{\cal A}\left( x\left( t\right)
\right) \left| \nabla {\cal E}\left( x\left( t\right) \right) \right\rangle
=P_{%
%TCIMACRO{\func{range}}%
%BeginExpansion
\mathop{\rm range}%
%EndExpansion
\left\{ {\cal A}\left( x\left( t\right) \right) \right\} ^{c}}\left| \nabla 
{\cal E}\left( x\left( t\right) \right) \right\rangle =0\quad \forall t
\end{equation}%
which in a linearized form are 
\begin{equation}
P_{%
%TCIMACRO{\func{range}}%
%BeginExpansion
\mathop{\rm range}%
%EndExpansion
\left\{ {\cal A}\left( x\left( t\right) \right) \right\} ^{c}}\left( \nabla 
{\cal E}\left( x_{0}\left( t\right) \right) +D_{x}^{2}{\cal E}\left(
x_{0}\left( t\right) \right) \left| y\left( t\right) \right\rangle \right)
=0\quad \forall t
\end{equation}

\appendix

\QSubDoc{Include Part}{%TCIDATA{LaTeXparent=0,0,EOMN.tex}
                      

\subsection{General Form of Partitioning}

Consider the matrix equation 
\begin{equation}
\left( 
\begin{array}{cc}
{\Bbb A}_{11} & {\Bbb A}_{12} \\ 
{\Bbb A}_{21} & {\Bbb A}_{22}%
\end{array}%
\right) \left( 
\begin{array}{c}
{\bf x}_{1} \\ 
{\bf x}_{2}%
\end{array}%
\right) =\left( 
\begin{array}{c}
{\bf y}_{1} \\ 
{\bf y}_{2}%
\end{array}%
\right)  \label{ME1}
\end{equation}%
that is equivalent to the coupled equations 
\begin{eqnarray}
{\Bbb A}_{11}{\bf x}_{1}+{\Bbb A}_{12}{\bf x}_{2} &=&{\bf y}_{1}  \nonumber
\\
{\Bbb A}_{21}{\bf x}_{1}+{\Bbb A}_{22}{\bf x}_{2} &=&{\bf y}_{2}
\end{eqnarray}%
we can solve for ${\bf x}_{1}$ from the top equation to obtain a ''fixed'' $%
{\bf x}_{1}$ 
\begin{equation}
{\bf x}_{1}={\Bbb A}_{11}^{-1}\left( {\bf y}_{1}-{\Bbb A}_{12}{\bf x}%
_{2}\right)  \label{x1eq}
\end{equation}%
and then substitute in the bottom equation to obtain a solution for ${\bf x}%
_{2}$ 
\begin{equation}
{\bf x}_{2}=\left( {\Bbb A}_{22}+{\Bbb A}_{21}{\Bbb A}_{11}^{-1}{\Bbb A}%
_{12}\right) ^{-1}\left( {\bf y}_{2}-{\Bbb A}_{21}{\Bbb A}_{11}^{-1}{\bf y}%
_{1}\right)  \label{x2sol}
\end{equation}%
which then determines ${\bf x}_{1}$ by eq. $\left( \ref{x1eq}\right) $. The
choice of which variables to call ${\bf x}_{1}$ and ${\bf x}_{2}$ in any
problem is of course quite arbitrary. One should also observe that the ${\bf %
x}_{2}$ obtained from eq. $\left( \ref{x2sol}\right) $ can also be expressed
as 
\begin{equation}
{\bf x}_{2}={\Bbb A}_{22}^{-1}\left( {\bf y}_{2}-{\Bbb A}_{21}{\bf x}%
_{1}\right)  \label{x2leq}
\end{equation}%
where ${\bf x}_{1}$ is given by eq. $\left( \ref{x1eq}\right) $. If neither $%
{\Bbb A}$ nor ${\bf y}$ depend on ${\bf x}$ and all the inverses exist then
a fixed solution for ${\bf x}_{1}$ and ${\bf x}_{2}$ is obtained in one
iteration.

The equations we deal with are, however, in general nonlinear i.e. 
\begin{eqnarray}
{\Bbb A}_{ij} &=&{\Bbb A}_{ij}\left( {\bf x}_{1},{\bf x}_{2}\right) ;\quad
1\leq i,j\leq 2  \nonumber \\
{\bf y}_{i} &=&{\bf y}_{i}\left( {\bf x}_{1},{\bf x}_{2}\right) ;\quad 1\leq
i\leq 2
\end{eqnarray}%
and we have to iterate to obtain a solution by fixing a value for ${\bf x}%
_{1}$ and solving for ${\bf x}_{2\text{ }}$to obtain 
\begin{equation}
{\bf x}_{2}=\left( {\Bbb A}_{22}+{\Bbb A}_{21}{\Bbb A}_{11}^{-1}{\Bbb A}%
_{12}\right) ^{-1}\left( {\bf y}_{2}-{\Bbb A}_{21}{\Bbb A}_{11}^{-1}{\bf y}%
_{1}\right) ={\Bbb A}_{22}^{-1}\left( {\bf y}_{2}-{\Bbb A}_{21}{\bf x}%
_{1}\right)  \label{x2parteq}
\end{equation}%
until one reaches self consistency or fixing a value for ${\bf x}_{2}$ and
solving for ${\bf x}_{1\text{ }}$to obtain%
\begin{equation}
{\bf x}_{1}={\Bbb A}_{11}^{-1}\left( {\bf y}_{1}-{\Bbb A}_{12}{\bf x}%
_{2}\right)  \label{x1parteq}
\end{equation}%
until one reaches self consistency. When we have a partial self consistent
solution for ${\bf x}_{1}$ and use a fixed ${\bf x}_{2}$ we can substitute
back into the equations to obtain 
\begin{eqnarray}
&&\left( 
\begin{array}{cc}
{\Bbb A}_{11} & {\Bbb A}_{12} \\ 
{\Bbb A}_{21} & {\Bbb A}_{22}%
\end{array}%
\right) \left( 
\begin{array}{c}
{\Bbb A}_{11}^{-1}\left( {\bf y}_{1}-{\Bbb A}_{12}{\bf x}_{2}\right) \\ 
{\bf x}_{2}%
\end{array}%
\right) -\left( 
\begin{array}{c}
{\bf y}_{1} \\ 
{\bf y}_{2}%
\end{array}%
\right)  \nonumber \\
&=&\left( 
\begin{array}{c}
{\Bbb A}_{11}{\Bbb A}_{11}^{-1}\left( {\bf y}_{1}-{\Bbb A}_{12}{\bf x}%
_{2}\right) +{\Bbb A}_{12}{\bf x}_{2} \\ 
{\Bbb A}_{21}{\Bbb A}_{11}^{-1}\left( {\bf y}_{1}-{\Bbb A}_{12}{\bf x}%
_{2}\right) +{\Bbb A}_{22}{\bf x}_{2}%
\end{array}%
\right) -\left( 
\begin{array}{c}
{\bf y}_{1} \\ 
{\bf y}_{2}%
\end{array}%
\right)  \nonumber \\
&=&\left( 
\begin{array}{c}
\left( {\bf y}_{1}-{\Bbb A}_{12}{\bf x}_{2}\right) +{\Bbb A}_{12}{\bf x}_{2}
\\ 
{\Bbb A}_{21}{\Bbb A}_{11}^{-1}\left( {\bf y}_{1}-{\Bbb A}_{12}{\bf x}%
_{2}\right) +{\Bbb A}_{22}{\bf x}_{2}%
\end{array}%
\right) -\left( 
\begin{array}{c}
{\bf y}_{1} \\ 
{\bf y}_{2}%
\end{array}%
\right)  \nonumber \\
&=&\left( 
\begin{array}{c}
{\bf 0} \\ 
{\Bbb A}_{21}{\Bbb A}_{11}^{-1}\left( {\bf y}_{1}-{\Bbb A}_{12}{\bf x}%
_{2}\right) +{\Bbb A}_{22}{\bf x}_{2}-{\bf y}_{2}%
\end{array}%
\right)  \label{baksub1}
\end{eqnarray}%
While when we have a partial self consistent solution for ${\bf x}_{2}$ and
use a fixed ${\bf x}_{1}$ we can substitute back into the equations to
obtain 
\begin{eqnarray}
&&\left( 
\begin{array}{cc}
{\Bbb A}_{11} & {\Bbb A}_{12} \\ 
{\Bbb A}_{21} & {\Bbb A}_{22}%
\end{array}%
\right) \left( 
\begin{array}{c}
{\bf x}_{1} \\ 
\left( {\Bbb A}_{22}+{\Bbb A}_{21}{\Bbb A}_{11}^{-1}{\Bbb A}_{12}\right)
^{-1}\left( {\bf y}_{2}-{\Bbb A}_{21}{\Bbb A}_{11}^{-1}{\bf y}_{1}\right)%
\end{array}%
\right) -\left( 
\begin{array}{c}
{\bf y}_{1} \\ 
{\bf y}_{2}%
\end{array}%
\right)  \nonumber \\
&=&\left( 
\begin{array}{c}
{\Bbb A}_{11}{\bf x}_{1}+{\Bbb A}_{12}\left( {\Bbb A}_{22}+{\Bbb A}_{21}%
{\Bbb A}_{11}^{-1}{\Bbb A}_{12}\right) ^{-1}\left( {\bf y}_{2}-{\Bbb A}_{21}%
{\Bbb A}_{11}^{-1}{\bf y}_{1}\right) \\ 
{\Bbb A}_{21}{\bf x}_{1}+{\Bbb A}_{22}\left( {\Bbb A}_{22}+{\Bbb A}_{21}%
{\Bbb A}_{11}^{-1}{\Bbb A}_{12}\right) ^{-1}\left( {\bf y}_{2}-{\Bbb A}_{21}%
{\Bbb A}_{11}^{-1}{\bf y}_{1}\right)%
\end{array}%
\right) -\left( 
\begin{array}{c}
{\bf y}_{1} \\ 
{\bf y}_{2}%
\end{array}%
\right)  \nonumber \\
&=&\left( 
\begin{array}{c}
{\Bbb A}_{11}{\bf x}_{1}+{\Bbb A}_{12}{\Bbb A}_{22}^{-1}\left\{ {\bf y}_{2}-%
{\Bbb A}_{21}{\bf x}_{1}\right\} \\ 
{\bf y}_{2}%
\end{array}%
\right) -\left( 
\begin{array}{c}
{\bf y}_{1} \\ 
{\bf y}_{2}%
\end{array}%
\right)  \nonumber \\
&=&\left( 
\begin{array}{c}
{\Bbb A}_{11}{\bf x}_{1}+{\Bbb A}_{12}{\Bbb A}_{22}^{-1}\left\{ {\bf y}_{2}-%
{\Bbb A}_{21}{\bf x}_{1}\right\} -{\bf y}_{1} \\ 
{\bf 0}%
\end{array}%
\right)  \label{baksub2}
\end{eqnarray}%
The partitioning displayed above obviously only works when the various
inverses exist.

We can obtain equivalent equations by multiplying both sides of eq. $\left( %
\ref{ME1}\right) $ by an invertible matrix e.g. 
\begin{equation}
\left( 
\begin{array}{cc}
{\Bbb B}_{11} & {\Bbb B}_{12} \\ 
{\Bbb B}_{21} & {\Bbb B}_{22}%
\end{array}%
\right) \left( 
\begin{array}{cc}
{\Bbb A}_{11} & {\Bbb A}_{12} \\ 
{\Bbb A}_{21} & {\Bbb A}_{22}%
\end{array}%
\right) \left( 
\begin{array}{c}
{\bf x}_{1} \\ 
{\bf x}_{2}%
\end{array}%
\right) =\left( 
\begin{array}{cc}
{\Bbb B}_{11} & {\Bbb B}_{12} \\ 
{\Bbb B}_{21} & {\Bbb B}_{22}%
\end{array}%
\right) \left( 
\begin{array}{c}
{\bf y}_{1} \\ 
{\bf y}_{2}%
\end{array}%
\right) =\left( 
\begin{array}{c}
{\bf z}_{1} \\ 
{\bf z}_{2}%
\end{array}%
\right)  \label{pre}
\end{equation}%
and then partitioning in the same manner. If ${\Bbb A}_{12}={\Bbb A}_{21}=0$
then the equations are uncoupled if they are linear, but are still coupled
if they are nonlinear. If ${\Bbb A}_{11}={\Bbb A}_{22}=0$ then one can
premultiply eq. $\left( \ref{pre}\right) $ by $\left( 
\begin{array}{cc}
{\Bbb O} & {\Bbb I} \\ 
{\Bbb I} & {\Bbb O}%
\end{array}%
\right) $ and then solve, but then the partial solutions do not give the
same back substituted properties displayed in eqs. $\left( \ref{baksub1}%
\right) $ and $\left( \ref{baksub2}\right) $ i.e. the residual vectors of
the original problem generally have zero components in different positions
or no zero components at all.

\subsection{Zero ${\Bbb A}_{22}$ Block}

The {\em nonlinear }equation eq. $\left( \ref{ME1}\right) $ for ${\bf x}_{1}$
and ${\bf x}_{2}$ 
\begin{equation}
\left( 
\begin{array}{cc}
{\Bbb A}_{11} & {\Bbb A}_{12} \\ 
{\Bbb A}_{21} & {\Bbb A}_{22}%
\end{array}%
\right) \left( 
\begin{array}{c}
{\bf x}_{1} \\ 
{\bf x}_{2}%
\end{array}%
\right) -\left( 
\begin{array}{c}
{\bf y}_{1} \\ 
{\bf y}_{2}%
\end{array}%
\right) =0
\end{equation}%
does not have be solved in the preceding manner, in fact when one of the
diagonal blocks is zero it may become impossible. In particular consider 
\begin{equation}
\left( 
\begin{array}{cc}
0 & {\Bbb A}_{12} \\ 
{\Bbb A}_{21} & {\Bbb A}_{22}%
\end{array}%
\right) \left( 
\begin{array}{c}
{\bf x}_{1} \\ 
{\bf x}_{2}%
\end{array}%
\right) -\left( 
\begin{array}{c}
{\bf y}_{1} \\ 
{\bf y}_{2}%
\end{array}%
\right) =0\Longleftrightarrow \left. 
\begin{array}{c}
\qquad \qquad {\Bbb A}_{12}{\bf x}_{2}-{\bf y}_{1}=0 \\ 
{\Bbb A}_{21}{\bf x}_{1}+{\Bbb A}_{22}{\bf x}_{2}-{\bf y}_{2}=0%
\end{array}%
\right\}
\end{equation}%
which gives for fixed ${\bf x}_{1}$ the equation 
\begin{equation}
{\Bbb A}_{22}{\bf x}_{2}={\bf y}_{2}-{\Bbb A}_{21}{\bf x}_{1}  \label{x2le}
\end{equation}%
for ${\bf x}_{2}$. A self consistent solution, ${\bf x}_{2\#}\left( {\bf x}%
_{1\left( k\right) }\right) $, is obtained by iterating ${\bf x}_{2}$ in eq. 
$\left( \ref{x2le}\right) $ at the fixed value ${\bf x}_{1}={\bf x}_{1\left(
k\right) }$ 
\begin{equation}
{\bf x}_{2\#}\left( {\bf x}_{1\left( k\right) }\right) ={\Bbb A}%
_{22}^{-1}\left( {\bf y}_{2}-{\Bbb A}_{21}{\bf x}_{1\left( k\right) }\right)
\label{x2}
\end{equation}%
If one substitutes back this partial solution, noting that all quantities, $%
{\Bbb A}_{12},{\Bbb A}_{21},{\Bbb A}_{22},{\bf y}_{1}$ and ${\bf y}_{2}$ are
evaluated at $\left( {\bf x}_{1\left( k\right) },{\bf x}_{2\#}\left( {\bf x}%
_{1\left( k\right) }\right) \right) $, we obtain 
\begin{eqnarray}
&&\left( 
\begin{array}{cc}
0 & {\Bbb A}_{12} \\ 
{\Bbb A}_{21} & {\Bbb A}_{22}%
\end{array}%
\right) \left( 
\begin{array}{c}
{\bf x}_{1\left( k\right) } \\ 
{\Bbb A}_{22}^{-1}\left( {\bf y}_{2}-{\Bbb A}_{21}{\bf x}_{1\left( k\right)
}\right)%
\end{array}%
\right) -\left( 
\begin{array}{c}
{\bf y}_{1} \\ 
{\bf y}_{2}%
\end{array}%
\right) =  \nonumber \\
&&\qquad \left( 
\begin{array}{c}
{\Bbb A}_{12}{\Bbb A}_{22}^{-1}\left( {\bf y}_{2}-{\Bbb A}_{21}{\bf x}%
_{1\left( k\right) }\right) -{\bf y}_{1} \\ 
{\Bbb A}_{21}{\bf x}_{1\left( k\right) }+{\bf y}_{2}-{\Bbb A}_{21}{\bf x}%
_{1\left( k\right) }-{\bf y}_{2}%
\end{array}%
\right) =\left( 
\begin{array}{c}
{\Bbb A}_{12}{\Bbb A}_{22}^{-1}\left( {\bf y}_{2}-{\Bbb A}_{21}{\bf x}%
_{1\left( k\right) }\right) -{\bf y}_{1} \\ 
{\bf 0}%
\end{array}%
\right)  \label{baksub}
\end{eqnarray}%
The self consistent solution, ${\bf x}_{1\#\left( k+1\right) }$, for ${\bf x}%
_{1}$ at ${\bf x}_{2\#}\left( {\bf x}_{1\left( k\right) }\right) $ is then
obtained by iterating ${\bf x}_{1\left( k+1\right) }$ in%
\begin{equation}
{\Bbb A}_{12}{\Bbb A}_{22}^{-1}\left( {\bf y}_{2}-{\Bbb A}_{21}{\bf x}%
_{1\left( k+1\right) }\right) -{\bf y}_{1}=0
\end{equation}%
or in the formal relationship 
\begin{equation}
{\bf x}_{1\left( k+1\right) }={\Bbb A}_{12}^{-1}{\Bbb A}_{22}{\Bbb A}%
_{21}^{-1}\left( {\bf y}_{1}-{\Bbb A}_{12}{\Bbb A}_{22}^{-1}{\bf y}%
_{2}\right)  \label{x1kplus1}
\end{equation}%
where all quantities are evaluated at $\left( {\bf x}_{1\left( k+1\right) },%
{\bf x}_{2\#}\left( {\bf x}_{1\left( k\right) }\right) \right) $ and left
inverses are used for ${\Bbb A}_{12}$ and ${\Bbb A}_{21}$. These two coupled
iterative sequences can then be repeated until self consistency is reached
for both ${\bf x}_{1}$ and ${\bf x}_{2}$ to obtain a solution of the full
non linear problem.

If ${\Bbb A}_{22}$ is singular then eq. $\left( \ref{x2}\right) $ becomes%
\begin{equation}
{\bf x}_{2\#}\left( {\bf x}_{1\left( k\right) },{\bf k}\right) ={\Bbb A}%
_{22}^{-1}\left( {\bf y}_{2}-{\Bbb A}_{21}{\bf x}_{1\left( k\right) }\right)
+{\bf k}  \label{x2eq}
\end{equation}%
where we use the generalized inverse ${\Bbb A}_{22}^{-1}$ of ${\Bbb A}_{22}$%
, ${\bf k}\in \ker {\Bbb A}_{22}$ and the subsidiary condition 
\begin{equation}
P_{\ker {\Bbb A}_{22}}\left( {\bf y}_{2}-{\Bbb A}_{21}{\bf x}_{1\left(
k\right) }\right) =0  \label{x2cond}
\end{equation}%
must be satisfied, with all terms evaluated at $\left( {\bf x}_{1\left(
k\right) },{\bf x}_{2\#}\left( {\bf x}_{1\left( k\right) },{\bf k}\right)
\right) $. Then the back substitution eq. $\left( \ref{baksub}\right) $
becomes 
\begin{eqnarray}
&&\left( 
\begin{array}{cc}
0 & {\Bbb A}_{12} \\ 
{\Bbb A}_{21} & {\Bbb A}_{22}%
\end{array}%
\right) \left( 
\begin{array}{c}
{\bf x}_{1\left( k\right) } \\ 
{\Bbb A}_{22}^{-1}\left( {\bf y}_{2}-{\Bbb A}_{21}{\bf x}_{1\left( k\right)
}\right) +{\bf k}%
\end{array}%
\right) -\left( 
\begin{array}{c}
{\bf y}_{1} \\ 
{\bf y}_{2}%
\end{array}%
\right) =  \nonumber \\
&&\qquad \left( 
\begin{array}{c}
{\Bbb A}_{12}{\Bbb A}_{22}^{-1}\left( {\bf y}_{2}-{\Bbb A}_{21}{\bf x}%
_{1\left( k\right) }\right) +{\Bbb A}_{12}{\bf k}-{\bf y}_{1} \\ 
P_{\ker {\Bbb A}_{22}}^{c}\left( {\bf y}_{2}-{\Bbb A}_{21}{\bf x}_{1\left(
k\right) }\right) +{\Bbb A}_{21}{\bf x}_{1\left( k\right) }-{\bf y}_{2}%
\end{array}%
\right) =  \nonumber \\
&&\qquad \qquad \qquad \qquad \left( 
\begin{array}{c}
{\Bbb A}_{12}{\Bbb A}_{22}^{-1}\left( {\bf y}_{2}-{\Bbb A}_{21}{\bf x}%
_{1\left( k\right) }\right) +{\Bbb A}_{12}{\bf k}-{\bf y}_{1} \\ 
{\bf 0}%
\end{array}%
\right)  \label{subbak}
\end{eqnarray}%
for any ${\bf k}$. The equation for a self consistent value, ${\bf x}%
_{1\#\left( k+1\right) }\left( {\bf x}_{2\#}\left( {\bf x}_{1\left( k\right)
},{\bf k}\right) \right) $, for ${\bf x}_{1}$ at ${\bf x}_{2}={\bf x}%
_{2\#}\left( {\bf x}_{1\left( k\right) },{\bf k}\right) $ obtained from eq. $%
\left( \ref{subbak}\right) $ is%
\begin{equation}
{\Bbb A}_{12}{\Bbb A}_{22}^{-1}{\Bbb A}_{21}{\bf x}_{1\#\left( k+1\right) }=%
{\Bbb A}_{12}{\Bbb A}_{22}^{-1}{\bf y}_{2}-{\bf y}_{1}+{\Bbb A}_{12}{\bf k}
\label{scx1}
\end{equation}%
which can be formally solved as 
\begin{equation}
{\bf x}_{1\#\left( k+1\right) }\left( {\bf x}_{1\left( k\right) },{\bf k}%
\right) ={\Bbb A}_{21}^{-1}{\Bbb A}_{22}{\Bbb A}_{12}^{-1}\left( {\Bbb A}%
_{12}{\Bbb A}_{22}^{-1}{\bf y}_{2}-{\bf y}_{1}\right) ={\Bbb A}_{21}^{-1}%
{\bf y}_{2}-{\Bbb A}_{21}^{-1}{\Bbb A}_{22}{\Bbb A}_{12}^{-1}{\bf y}_{1}
\end{equation}%
where ${\Bbb A}_{12},{\Bbb A}_{21},{\Bbb A}_{22},{\bf y}_{1},$ and ${\bf y}%
_{2}$ are evaluated at $\left( {\bf x}_{1\#\left( k+1\right) },{\bf x}%
_{2\#}\left( {\bf x}_{1\left( k\right) },{\bf k}\right) \right) $. This
value can then be substituted back into the equation for ${\bf x}_{2}$,
which can then be iterated to self consistency once more, the new self
consistent value of ${\bf x}_{2}$ can then be used to obtain a new self
consistent value of ${\bf x}_{1}$ and so on.

Instead\ of iterating to self consistency for a new\ ${\bf x}_{1\#\left(
k+1\right) }$ and then substituting back into the equation for ${\bf x}_{2}$%
, which would be the most efficient way to solve the total coupled problem,
one could find paths of self consistent solutions, ${\bf x}_{2\#}\left( {\bf %
x}_{1},{\bf k}\right) $, for all possible ${\bf x}_{1}$'s producing paths, $%
{\bf x}_{2\#}\left( {\bf x}_{1},{\bf k}\right) $, of ${\bf x}_{2}$'s that
satisfy 
\begin{equation}
{\bf x}_{2\#}\left( {\bf x}_{1},{\bf k}\right) ={\Bbb A}_{22}^{-1}\left( 
{\bf y}_{2}-{\Bbb A}_{21}{\bf x}_{1}\right) +{\bf k}
\end{equation}%
where 
\begin{eqnarray}
{\Bbb A}_{22} &=&{\Bbb A}_{22}\left( {\bf x}_{1},{\bf x}_{2\#}\left( {\bf x}%
_{1},{\bf k}\right) \right)  \nonumber \\
{\bf y}_{2} &=&{\bf y}_{2}\left( {\bf x}_{1},{\bf x}_{2\#}\left( {\bf x}_{1},%
{\bf k}\right) \right)  \nonumber \\
{\Bbb A}_{21} &=&{\Bbb A}_{21}\left( {\bf x}_{1},{\bf x}_{2\#}\left( {\bf x}%
_{1},{\bf k}\right) \right)  \nonumber \\
{\bf k} &=&{\bf k}\left( {\bf x}_{1}\right)
\end{eqnarray}%
and 
\begin{eqnarray}
\left( 
\begin{array}{cc}
0 & {\Bbb A}_{12} \\ 
{\Bbb A}_{21} & {\Bbb A}_{22}%
\end{array}%
\right) \left( 
\begin{array}{c}
{\bf x}_{1} \\ 
{\bf x}_{2\#}\left( {\bf x}_{1},{\bf k}\right)%
\end{array}%
\right) -\left( 
\begin{array}{c}
{\bf y}_{1} \\ 
{\bf y}_{2}%
\end{array}%
\right) &=&\left( 
\begin{array}{c}
{\Bbb A}_{12}{\bf x}_{2\#}\left( {\bf x}_{1},{\bf k}\right) -{\bf y}_{1} \\ 
{\Bbb A}_{21}{\bf x}_{1}+{\Bbb A}_{22}{\bf x}_{2\#}\left( {\bf x}_{1},{\bf k}%
\right) -{\bf y}_{2}%
\end{array}%
\right)  \nonumber \\
&=&\left( 
\begin{array}{c}
{\Bbb A}_{12}{\bf x}_{2\#}\left( {\bf x}_{1},{\bf k}\right) -{\bf y}_{1} \\ 
0%
\end{array}%
\right)
\end{eqnarray}%
where 
\begin{eqnarray}
{\Bbb A}_{12} &=&{\Bbb A}_{12}\left( {\bf x}_{1},{\bf x}_{2\#}\left( {\bf x}%
_{1},{\bf k}\right) \right)  \nonumber \\
{\bf y}_{1} &=&{\bf y}_{1}\left( {\bf x}_{1},{\bf x}_{2\#}\left( {\bf x}_{1},%
{\bf k}\right) \right)
\end{eqnarray}%
Then solve the nonlinear equation for ${\bf x}_{1}$ 
\begin{eqnarray}
&&{\Bbb A}_{12}{\bf x}_{2\#}\left( {\bf x}_{1},{\bf k}\right) -{\bf y}_{1}=0
\nonumber \\
&\equiv &{\Bbb A}_{12}{\Bbb A}_{22}^{-1}\left( {\bf y}_{2}-{\Bbb A}_{21}{\bf %
x}_{1}\right) +{\Bbb A}_{12}{\bf k}-{\bf y}_{1}=0  \label{X1eq}
\end{eqnarray}%
with formal solution 
\[
{\bf x}_{1\#}={\Bbb A}_{21}^{-1}{\bf y}_{2}-{\Bbb A}_{21}^{-1}{\Bbb A}_{22}%
{\Bbb A}_{12}^{-1}{\bf y}_{1}={\Bbb A}_{21}^{-1}\left( {\bf y}_{2}-{\Bbb A}%
_{22}{\Bbb A}_{12}^{-1}{\bf y}_{1}\right) 
\]%
(which corresponds to the equation we identified for ${\cal V}_{den}$ in eq. 
$\left( \ref{gammaeq}\right) $) to determine the exact ${\bf x}_{1}$, where
everything in eq. $\left( \ref{X1eq}\right) $ is evaluated along the paths $%
\left( {\bf x}_{1},{\bf x}_{2\#}\left( {\bf x}_{1},{\bf k}\right) \right) $.
The different paths of solutions ${\bf x}_{2\#}\left( {\bf x}_{1},{\bf k}%
\right) $ are labelled by ${\bf k}$ i.e. 
\begin{equation}
{\bf x}_{2\#}\left( {\bf x}_{1},{\bf k}\right) \neq {\bf x}_{2\#}\left( {\bf %
x}_{1},{\bf k}^{\prime }\right) \text{ for }{\bf k}\left( {\bf x}_{1}\right)
\neq {\bf k}^{\prime }\left( {\bf x}_{1}\right)
\end{equation}%
and an extra criterion must be used to determine unique paths. As ${\bf x}%
_{2\#}\left( {\bf x}_{1},{\bf k}\right) $ is always a self consistent
solution for ${\bf x}_{2}$ in eq. $\left( \ref{x2eq}\right) $ one can
observe that for any ${\bf x}_{1}$ we have 
\begin{eqnarray}
&&\left( 
\begin{array}{cc}
0 & {\Bbb A}_{12} \\ 
{\Bbb A}_{21} & {\Bbb A}_{22}%
\end{array}%
\right) \left( 
\begin{array}{c}
{\bf x}_{1} \\ 
{\bf x}_{2\#}\left( {\bf x}_{1},{\bf k}\right)%
\end{array}%
\right) -\left( 
\begin{array}{c}
{\bf y}_{1} \\ 
{\bf y}_{2}%
\end{array}%
\right) =\left( 
\begin{array}{c}
{\Bbb A}_{12}{\bf x}_{2\#}\left( {\bf x}_{1}\right) -{\bf y}_{1} \\ 
{\Bbb A}_{21}{\bf x}_{1}+{\Bbb A}_{22}{\bf x}_{2\#}\left( {\bf x}_{1},{\bf k}%
\right) -{\bf y}_{2}%
\end{array}%
\right)  \nonumber \\
&=&\left( 
\begin{array}{c}
{\Bbb A}_{12}{\bf x}_{2\#}\left( {\bf x}_{1}\right) -{\bf y}_{1} \\ 
{\Bbb A}_{21}{\bf x}_{1}+{\Bbb A}_{22}\left( {\Bbb A}_{22}^{-1}\left( {\bf y}%
_{2}-{\Bbb A}_{21}{\bf x}_{1}\right) +{\bf k}\right) -{\bf y}_{2}%
\end{array}%
\right) =\left( 
\begin{array}{c}
{\Bbb A}_{12}{\bf x}_{2\#}\left( {\bf x}_{1}\right) -{\bf y}_{1} \\ 
{\bf 0}%
\end{array}%
\right)
\end{eqnarray}%
i.e.%
\begin{equation}
\left( 
\begin{array}{cc}
0 & {\Bbb A}_{12} \\ 
{\Bbb A}_{21} & {\Bbb A}_{22}%
\end{array}%
\right) \left( 
\begin{array}{c}
{\bf x}_{1} \\ 
{\bf x}_{2\#}\left( {\bf x}_{1}\right)%
\end{array}%
\right) -\left( 
\begin{array}{c}
{\bf y}_{1} \\ 
{\bf y}_{2}%
\end{array}%
\right) =\left( 
\begin{array}{c}
{\Bbb A}_{12}{\bf x}_{2\#}\left( {\bf x}_{1}\right) -{\bf y}_{1} \\ 
0%
\end{array}%
\right)
\end{equation}
}

\QSubDoc{Include VecBun}{\input{VecBun.tex}}

\QSubDoc{Include DualOp}{%TCIDATA{LaTeXparent=0,0,EOMN.tex}
                      

\section{Dual Operator Spaces}

\subsection{Operators mapping ${\cal B}_{1}\left( {\cal H}^{N}\right)
\longrightarrow {\cal B}_{1}\left( {\cal H}^{N}\right) $}

The dual space of the Banach space, ${\cal B}_{1}\left( {\cal H}^{N}\right) $%
, of operators formed by the set of trace class operators equipped with the
norm%
\begin{equation}
\left\| Y\right\| _{1}=Tr\left\{ \left( Y^{\dagger }Y\right) ^{\frac{1}{2}%
}\right\}
\end{equation}%
is the Banach space of bounded operators ${\cal B}\left( {\cal H}^{N}\right) 
$ equipped with the norm%
\begin{equation}
\left\| X\right\| =%
%TCIMACRO{
%\underset{\left\| \left| \Psi \right\rangle \right\| _{{\cal H}^{N}}=1}{\sup }}%
%BeginExpansion
\mathrel{\mathop{\sup }\limits_{\left\| \left| \Psi \right\rangle \right\| _{{\cal H}^{N}}=1}}%
%EndExpansion
\left\| X\left| \Psi \right\rangle \right\| _{{\cal H}^{N}};\quad \left|
\Psi \right\rangle \in {\cal H}^{N}
\end{equation}%
where $\left\| \left| \Psi \right\rangle \right\| _{{\cal H}^{N}}=\left(
\left\langle \Psi |\Psi \right\rangle \right) ^{\frac{1}{2}}$. The duality
is expressed by%
\begin{equation}
f_{X}\left( Y\right) =Tr\left\{ \left( X^{\dagger }Y\right) \right\} \equiv
\left( X|Y\right) \quad \forall Y\in {\cal B}_{1}\left( {\cal H}^{N}\right)
\end{equation}%
Bases for these two spaces can be constructed from bases of ${\cal H}^{N}$
so that 
\begin{equation}
{\cal B}_{1}\left( {\cal H}^{N}\right) ={\rm cls}\left\{ \sum\limits_{1\leq
i,j\leq \infty }\lambda _{ij}\left| \Psi _{i}\right\rangle \left\langle \Psi
_{j}\right| \right\}  \label{TCBasis}
\end{equation}%
\begin{equation}
{\cal B}\left( {\cal H}^{N}\right) ={\rm cls}\left\{ \sum_{1\leq i,j\leq
\infty }\lambda _{ij}\left| \Psi _{i}\right\rangle \left\langle \Psi
_{j}\right| \right\}  \label{BoundBasis}
\end{equation}%
where ${\rm cls\equiv closed\,linear\,span,}$ $\left\{ \Psi _{i}:1\leq i\leq
\infty \right\} $ is a complete orthonormal basis for ${\cal H}^{N}$ and the
closures are taken wrt the norms $\left\| {}\right\| _{1}$ and $\left\|
{}\right\| $ respectively. The space$\ {\cal B}_{1}\left( {\cal H}%
^{N}\right) $ forms a two-sided ideal in ${\cal B}\left( {\cal H}^{N}\right) %
\cite{ideal}$. We can also introduce the common overcomplete basis $\left\{
\left| {\sf z}\right\rangle \left\langle {\sf z}^{\prime }\right| \right\} $
formed by the generalized eigenvectors of the combined coordinate operator
(nuclear space-spin and electron space-spin). Any bounded operator can be
expanded in this basis as 
\begin{equation}
X=\int X\left( {\sf z}{\bf ,}{\sf z}^{\prime }\right) \left| {\sf z}%
\right\rangle \left\langle {\sf z}^{\prime }\right| d{\sf z}d{\sf z}^{\prime
}  \label{overcomplete}
\end{equation}%
where {\sf z}${\bf \equiv }\kappa _{1,}\ldots ,\kappa _{N_{e}},\kappa
_{N_{e}+1,}\ldots ,\kappa _{N}$, and $X\left( {\sf z}{\bf ,}{\sf z}^{\prime
}\right) $ is the integral kernel of $X$ which can be a distribution or even
more singular. As we are dealing with dual spaces the equality sign in Eq.(%
\ref{overcomplete}) needs some discussion, in this expansion it means that%
\begin{equation}
\lim_{\Delta \rightarrow {\Bbb K{\Large \times }K}}\left\| X-\int_{\Delta
}X\left( {\sf z}{\bf ,}{\sf z}^{\prime }\right) \left| {\sf z}\right\rangle
\left\langle {\sf z}^{\prime }\right| d{\sf z}d{\sf z}^{\prime }\right\| =0
\end{equation}%
where ${\Bbb K}$ $=\left\{ {\sf z}\right\} $, i.e. the set of all
coordinates. As $\left\| X\right\| \leq \left\| X\right\| _{1}$ convergence
in the trace class topology implies convergence in the operator topology,
but obviously not vice-versa, which is consistent with ${\cal B}_{1}\left( 
{\cal H}^{N}\right) \subseteq {\cal B}\left( {\cal H}^{N}\right) $. Hence if
a trace class operator is expanded in the basis $\left\{ \left| {\sf z}%
\right\rangle \left\langle {\sf z}^{\prime }\right| \right\} $ as an element
of ${\cal B}_{1}\left( {\cal H}^{N}\right) $ it has the same expansion
coefficients, i.e. integral kernel, when considered as an element of ${\cal B%
}\left( {\cal H}^{N}\right) $ i.e. as a bounded operator and its integral
kernel is absolute integrable i.e. $\int \left| X\left( {\sf z}{\bf ,}{\sf z}%
^{\prime }\right) \right| d$\underline{$\kappa $}$d$\underline{$\kappa $}$%
^{\prime }{\bf <\infty }$. Those elements of \ ${\cal B}\left( {\cal H}%
^{N}\right) $\ that do not belong to ${\cal B}_{1}\left( {\cal H}^{N}\right) 
$ have integral expansions in terms of integral kernels that are not
absolute integrable.

We can consider more general bases $\left\{ X_{\iota };1\leq \iota \leq
\infty \right\} $, $\left\{ Y_{\iota };1\leq \iota \leq \infty \right\} $ of 
${\cal B}\left( {\cal H}^{N}\right) $ and ${\cal B}_{1}\left( {\cal H}%
^{N}\right) $ than those of Eqs.(\ref{TCBasis}) and (\ref{BoundBasis}) and
expand each basis element in terms of the coordinate overcomplete set as%
\begin{eqnarray}
X_{\iota } &=&\int X_{\iota }\left( {\sf z}{\bf ,}{\sf z}^{\prime }\right)
\left| {\sf z}\right\rangle \left\langle {\sf z}^{\prime }\right| d{\sf z}d%
{\sf z}^{\prime }{\bf \equiv }\int X_{\iota }\left( {\sf z}{\bf ,}{\sf z}%
^{\prime }\right) \left( {\sf z}{\bf ,}{\sf z}^{\prime }\right| d{\sf z}d%
{\sf z}^{\prime }{\bf \equiv }\left( X_{\iota }\right|  \nonumber \\
Y_{\iota } &=&\int Y_{\iota }\left( {\sf z}{\bf ,}{\sf z}^{\prime }\right)
\left| {\sf z}\right\rangle \left\langle {\sf z}^{\prime }\right| d{\sf z}d%
{\sf z}^{\prime }{\bf \equiv }\int Y_{\iota }\left( {\sf z}{\bf ,}{\sf z}%
^{\prime }\right) \left| {\sf z}{\bf ,}{\sf z}^{\prime }\right) d{\sf z}d%
{\sf z}^{\prime }{\bf \equiv }\left| Y_{\iota }\right)  \label{Overcomplete}
\end{eqnarray}%
where we have introduced the notation $\left| {\sf z}{\bf ,}{\sf z}^{\prime
}\right) \equiv \left| {\sf z}\right\rangle \left\langle {\sf z}^{\prime
}\right| $ and the {\it super vector} notation $\left| Y\right) $ for an
operator $Y\in {\cal B}_{1}\left( {\cal H}^{N}\right) $ and $\left( {\sf z}%
{\bf ,}{\sf z}^{\prime }\right| $ and $\left( X\right| $ for operators
belonging to the dual space ${\cal B}\left( {\cal H}^{N}\right) .$ The
operators $\left\{ \left| {\sf z}\right\rangle \left\langle {\sf z}^{\prime
}\right| \right\} $ are common to both spaces and thus one can consider them
as elements of either, however the type of convergence for the continuous
linear combinations of $\left| {\sf z}\right\rangle \left\langle {\sf z}%
^{\prime }\right| $ is implied by the use of either $\left( {\sf z}{\bf ,}%
{\sf z}^{\prime }\right| $ or $\left| {\sf z}{\bf ,}{\sf z}^{\prime }\right) 
$.

We now are in a position to express maps $\hat{T}$: ${\cal B}_{1}\left( 
{\cal H}^{N}\right) \longrightarrow {\cal B}_{1}\left( {\cal H}^{N}\right) $
(sometimes called {\it super operators}) in terms of these dual bases as%
\begin{equation}
\hat{T}=\sum\Sb 1\leq \iota \leq \infty  \\ 1\leq \kappa \leq \infty  \endSb %
T_{\kappa \iota }\left| Y_{\kappa }\right) \left( X_{\iota }\right|
\end{equation}%
where the bases $\left\{ X_{\kappa }\right\} $ and $\left\{ Y_{\iota
}\right\} $ have been constructed to be biorthogonal i.e.%
\begin{equation}
\left( X_{\kappa }|Y_{\iota }\right) =\delta _{\kappa \iota }
\end{equation}%
and $T_{\kappa \iota }=\left( X_{\kappa }|\hat{T}Y_{\iota }\right) $. (A
remark on notation: we denote operators :${\cal B}_{1}\left( {\cal H}%
^{N}\right) \longrightarrow {\cal B}_{1}\left( {\cal H}^{N}\right) $ with
the decoration $\hat{T}$, while operators :${\cal H}^{N}\rightarrow {\cal H}%
^{N}$ will be denoted by the normal math font $X$ or with the decoration $%
\breve{X}$ if the domain of operation is not self evident.) An important
point to note is that unlike operators on Hilbert spaces, which are self
dual, the adjoint operation on $\hat{T}$ changes it from an operator ${\cal B%
}_{1}\left( {\cal H}^{N}\right) \longrightarrow {\cal B}_{1}\left( {\cal H}%
^{N}\right) $ to an operator ${\cal B}\left( {\cal H}^{N}\right)
\longrightarrow {\cal B}\left( {\cal H}^{N}\right) .$ Using the expansions
Eq.(\ref{Overcomplete}) one can express $\hat{T}$ as%
\begin{eqnarray}
\hat{T} &=&\sum\Sb 1\leq \iota \leq \infty  \\ 1\leq \kappa \leq \infty 
\endSb T_{\kappa \iota }\int Y_{\iota }\left( {\sf z}_{1}{\bf ,}{\sf z}%
_{2}\right) \left| {\sf z}_{1}{\bf ,}{\sf z}_{2}\right) d{\sf z}_{1}d{\sf z}%
_{2}\int X_{\kappa }^{\ast }\left( {\sf z}_{3}{\bf ,}{\sf z}_{4}\right)
\left( {\sf z}_{3}{\bf ,}{\sf z}_{4}\right| d{\sf z}_{3}d{\sf z}_{4} 
\nonumber \\
&=&\sum\Sb 1\leq \iota \leq \infty  \\ 1\leq \kappa \leq \infty  \endSb %
T_{\kappa \iota }\int Y_{\iota }\left( {\sf z}_{1}{\bf ,}{\sf z}_{2}\right)
X_{\kappa }^{\ast }\left( {\sf z}_{3}{\bf ,}{\sf z}_{4}\right) \left| {\sf z}%
_{1}{\bf ,}{\sf z}_{2}\right) \left( {\sf z}_{3}{\bf ,}{\sf z}_{4}\right| d%
{\sf z}_{1}d{\sf z}_{2}d{\sf z}_{3}d{\sf z}_{4}  \nonumber \\
&=&\int T\left( {\sf z}_{1}{\bf ,}{\sf z}_{2},{\sf z}_{3}{\bf ,}{\sf z}%
_{4}\right) \left| {\sf z}_{1}{\bf ,}{\sf z}_{2}\right) \left( {\sf z}_{3}%
{\bf ,}{\sf z}_{4}\right| d{\sf z}_{1}d{\sf z}_{2}d{\sf z}_{3}d{\sf z}_{4}
\label{SuperOp}
\end{eqnarray}%
where 
\begin{equation}
T\left( {\sf z}_{1}{\bf ,}{\sf z}_{2},{\sf z}_{3}{\bf ,}{\sf z}_{4}\right)
=\sum\Sb 1\leq \iota \leq \infty  \\ 1\leq \kappa \leq \infty  \endSb %
T_{\kappa \iota }Y_{\iota }\left( {\sf z}_{1}{\bf ,}{\sf z}_{2}\right)
X_{\kappa }^{\ast }\left( {\sf z}_{3}{\bf ,}{\sf z}_{4}\right)
\end{equation}%
and thus%
\begin{equation}
T_{\kappa \iota }=\int T\left( {\sf z}_{1}{\bf ,}{\sf z}_{2},{\sf z}_{3}{\bf %
,}{\sf z}_{4}\right) Y_{\kappa }\left( {\sf z}_{1}{\bf ,}{\sf z}_{2}\right)
X_{\iota }^{\ast }\left( {\sf z}_{3}{\bf ,}{\sf z}_{4}\right) d{\sf z}_{1}d%
{\sf z}_{2}d{\sf z}_{3}d{\sf z}_{4}
\end{equation}%
We can express Eq.(\ref{SuperOp}) in a more transparent form as%
\begin{equation}
\hat{T}=\int T\left( {\sf z}_{1}{\bf ,}{\sf z}_{2},{\sf z}_{3}{\bf ,}{\sf z}%
_{4}\right) |\left| {\sf z}_{1}\right\rangle \left\langle {\sf z}_{2}\right| 
{\bf )}\left( \left| {\sf z}_{3}\right\rangle \left\langle {\sf z}%
_{4}\right| \right| d{\sf z}_{1}d{\sf z}_{2}d{\sf z}_{3}d{\sf z}_{4}
\end{equation}%
Note that $T\left( {\sf z}_{1}{\bf ,}{\sf z}_{2},{\sf z}_{3}{\bf ,}{\sf z}%
_{4}\right) $ inherits properties from $\left\{ X_{\kappa }^{\ast }\left( 
{\sf z}_{3}{\bf ,}{\sf z}_{4}\right) \right\} $ for the $\left( {\sf z}_{3}%
{\bf ,}{\sf z}_{4}\right) $ arguments and properties from $\left\{ Y_{\iota
}\left( {\sf z}_{1}{\bf ,}{\sf z}_{2}\right) \right\} $ for the $\left( {\sf %
z}_{1}{\bf ,}{\sf z}_{2}\right) $ arguments {\it and these properties can be
quite different}. Thus these operators do not, in general, have symmetry
properties relating the $\left( {\sf z}_{1}{\bf ,}{\sf z}_{2}\right) $
arguments to the $\left( {\sf z}_{3}{\bf ,}{\sf z}_{4}\right) $ arguments.

\subsection{Operators preserving the direct sum structure ${\cal B}%
_{1}\left( {\cal H}^{N}\right) =K_{OD}\oplus K_{D}\oplus S_{D}$}

The space of trace class operators ${\cal B}_{1}\left( {\cal H}^{N}\right) $
can be decomposed into the direct sum of three subspaces defined as%
\begin{eqnarray}
K_{OD} &=&{\rm cls}\left\{ Y\in {\cal B}_{1}\left( {\cal H}^{N}\right)
:Y=\int Y\left( {\sf z}{\bf ,}{\sf z}^{\prime }\right) \left| {\sf z}%
\right\rangle \left\langle {\sf z}^{\prime }\right| d{\sf z}d{\sf z}^{\prime
}{\bf ;\quad }Y\left( {\sf z},{\sf z}^{\prime }\right) {\bf =0}\text{%
\thinspace\ if }\,{\sf z}{\bf =}{\sf z}^{\prime }\right\} \\
K_{D} &=&{\rm cls}\left\{ Y\in {\cal B}_{1}\left( {\cal H}^{N}\right)
:Y=\int Y\left( {\sf z}\right) \left| {\sf z}\right\rangle \left\langle {\sf %
z}\right| d{\sf z}\right. ;  \nonumber \\
&&\left. \qquad \qquad \quad \int Y\left( \kappa _{1,}..,\kappa
_{N_{e}},\kappa _{N_{e}+1},..\kappa _{N}\right) d\kappa _{2}\ldots d\kappa
_{N_{e}}d\kappa _{N_{e}+2}\ldots d\kappa _{N}{\bf =0}\text{ }{\bf a.e.}%
\right\} \\
S_{D} &=&{\rm cls}\left\{ Y\in {\cal B}_{1}\left( {\cal H}^{N}\right)
:Y=\int Y\left( {\sf z}\right) \left| {\sf z}\right\rangle \left\langle {\sf %
z}\right| d{\sf z}\right. \text{ };  \nonumber \\
&&\left. \qquad \qquad \quad \int Y\left( \kappa _{1,}..,\kappa
_{N_{e}},\kappa _{N_{e}+1},..\kappa _{N}\right) d\kappa _{2}\ldots d\kappa
_{N_{e}}d\kappa _{N_{e}+2}\ldots d\kappa _{N}{\bf \neq 0\,a.e.}\right\}
\end{eqnarray}%
where ${\bf a.e.\equiv }$ almost everywhere i.e. except on sets of measure
zero. We can form bases for these three subspaces of ${\cal B}_{1}\left( 
{\cal H}^{N}\right) $,%
\begin{eqnarray}
K_{OD} &=&{\rm cls}\left\{ \breve{h}_{\iota }:\breve{h}_{\iota }=\int
h_{\iota }\left( {\sf z}{\bf ,}{\sf z}^{\prime }\right) \left| {\sf z}%
\right\rangle \left\langle {\sf z}^{\prime }\right| d{\sf z}d{\sf z}^{\prime
};{\sf z}{\bf \neq }{\sf z}^{\prime }\right\}  \label{hbasis} \\
K_{D} &=&{\rm cls}\left\{ \breve{k}_{\iota }:\breve{k}_{\iota }=\int
k_{\iota }\left( {\sf z}\right) \left| {\sf z}\right\rangle \left\langle 
{\sf z}\right| d{\sf z}\right.  \nonumber \\
&&\quad \quad \quad \left. \int k_{\iota }\left( \kappa _{2,}..,\kappa
_{N_{e}},\kappa _{N_{e}+1},..\kappa _{N}\right) d\kappa _{2}\ldots d\kappa
_{N_{e}}d\kappa _{N_{e}+2}\ldots d\kappa _{N}{\bf =0}\text{ }{\bf a.e.}%
\right\}  \label{kbasis} \\
S_{D} &=&{\rm cls}\left\{ \breve{s}_{\iota }:\breve{s}_{\iota }=\int
s_{\iota }\left( {\sf z}\right) \left| {\sf z}\right\rangle \left\langle 
{\sf z}\right| d{\sf z}\right. ;  \nonumber \\
&&\quad \left. \qquad \int s_{\iota }\left( \kappa _{2,}..,\kappa
_{N_{e}},\kappa _{N_{e}+1},..\kappa _{N}\right) d\kappa _{2}\ldots d\kappa
_{N_{e}}d\kappa _{N_{e}+2}\ldots d\kappa _{N}{\bf \neq 0\,a.e.}\right\}
\label{sbasis}
\end{eqnarray}

In the following we will be considering representations in both the
overcomplete continuous basis $\left\{ \left| {\cal K}\right) \right\} $ and
the complete discrete basis $\left\{ \left| {\bf \breve{h}}\right) ,\left| 
{\bf \breve{k}}\right) ,\left| {\bf \breve{s}}\right) \right\} $ of ${\cal B}%
_{1}\left( {\cal H}^{N}\right) $, where arrays of bases vectors are denoted
by $\left| {\cal K}\right) \equiv \left| {\sf z}_{1},{\sf z}_{2}\right)
,..,\left| {\sf z}_{1}^{\prime },{\sf z}_{2}^{\prime }\right) ,..;$forming a
continuous array of uncountably many basis vectors and $\left| {\bf \breve{h}%
}\right) \equiv \left| \breve{h}_{1}\right) ,..,\left| \breve{h}_{r}\right)
;\left| {\bf \breve{k}}\right) \equiv \left| \breve{k}_{1}\right) ,..,\left| 
\breve{k}_{r}\right) $and $\left| {\bf \breve{s}}\right) \equiv \left| 
\breve{s}_{1}\right) ,..,\left| \breve{s}_{r}\right) $ forming a discrete
array of countably many basis vectors, where in general $r=\infty $. The
form of the transformation from the discrete basis to the continuous one can
be deduced by using the biorthogonal bases $\left\{ \left\{ \left( {\bf 
\breve{h}}^{d}\right| ,\left( {\bf \breve{k}}^{d}\right| ,\left( {\bf \breve{%
s}}^{d}\right| \right\} ,\left\{ \left| {\bf \breve{h}}\right) ,\left| {\bf 
\breve{k}}\right) ,\left| {\bf \breve{s}}\right) \right\} \right\} $ and $%
\left\{ \left\{ \left( {\cal K}\right| \right\} ,\left\{ \left| {\cal K}%
\right) \right\} \right\} $, (note that $\left\{ {\bf \breve{h}}^{d},{\bf 
\breve{k}}^{d},{\bf \breve{s}}^{d}\right\} \in {\cal B}\left( {\cal H}%
^{N}\right) $, while $\left\{ {\bf \breve{h}},{\bf \breve{k}},{\bf \breve{s}}%
\right\} \in {\cal B}_{1}\left( {\cal H}^{N}\right) $), and the resolutions
of the identity that they form in the following manner 
\begin{equation}
\left| Y\right) =\left| {\bf \breve{h}}\right) \left( {\bf \breve{h}}%
^{d}|Y\right) +\left| {\bf \breve{k}}\right) \left( {\bf \breve{k}}%
^{d}|Y\right) +\left| {\bf \breve{s}}\right) \left( {\bf \breve{s}}%
^{d}|Y\right) =\left| {\cal K}\right) \left( {\cal K}|Y\right)
\end{equation}%
\begin{eqnarray}
\Rightarrow \left| Y\right) &=&\left| {\cal K}\right) \left( {\cal K}|{\bf 
\breve{h}}\right) \left( {\bf \breve{h}}^{d}|Y\right) +\left| {\cal K}%
\right) \left( {\cal K}|{\bf \breve{k}}\right) \left( {\bf \breve{k}}%
^{d}|Y\right) +\left| {\cal K}\right) \left( {\cal K}|{\bf \breve{s}}\right)
\left( {\bf \breve{s}}^{d}|Y\right)  \nonumber \\
&=&\left( \left| {\cal K}\right) \right) \left( 
\begin{array}{ccc}
\left( {\cal K}|{\bf \breve{h}}\right) & \left( {\cal K}|{\bf \breve{k}}%
\right) & \left( {\cal K}|{\bf \breve{s}}\right)%
\end{array}%
\right) \left( 
\begin{array}{c}
\left( {\bf \breve{h}}^{d}|Y\right) \\ 
\left( {\bf \breve{k}}^{d}|Y\right) \\ 
\left( {\bf \breve{s}}^{d}|Y\right)%
\end{array}%
\right)  \label{KtoH}
\end{eqnarray}%
and 
\begin{eqnarray}
\Rightarrow \left| Y\right) &=&\left| {\bf \breve{h}}\right) \left( {\bf 
\breve{h}}^{d}|{\cal K}\right) \left( {\cal K}|Y\right) +\left| {\bf \breve{k%
}}\right) \left( {\bf \breve{k}}^{d}|{\cal K}\right) \left( {\cal K}%
|Y\right) +\left| {\bf \breve{s}}\right) \left( {\bf \breve{s}}^{d}|{\cal K}%
\right) \left( {\cal K}|Y\right)  \nonumber \\
&=&\left( \left| {\bf \breve{h}}\right) ,\left| {\bf \breve{k}}\right)
,\left| {\bf \breve{s}}\right) \right) \left( 
\begin{array}{c}
\left( {\bf \breve{h}}^{d}|{\cal K}\right) \\ 
\left( {\bf \breve{k}}^{d}|{\cal K}\right) \\ 
\left( {\bf \breve{s}}^{d}|{\cal K}\right)%
\end{array}%
\right) \left( {\cal K}|Y\right)  \label{HtoK}
\end{eqnarray}%
where 
\begin{eqnarray}
\left| {\bf \breve{h}}\right) \left( {\bf \breve{h}}^{d}\right|
&=&\sum_{1\leq \iota \leq \infty }\left| \breve{h}_{\iota }\right) \left( 
\breve{h}_{\iota }^{d}\right| =\hat{P}_{K_{OD}}  \nonumber \\
\left| {\bf \breve{k}}\right) \left( {\bf \breve{k}}^{d}\right|
&=&\sum_{1\leq \iota \leq \infty }\left| \breve{k}_{\iota }\right) \left( 
\breve{k}_{\iota }^{d}\right| =\hat{P}_{K_{D}}  \nonumber \\
\left| {\bf \breve{s}}\right) \left( {\bf \breve{s}}^{d}\right|
&=&\sum_{1\leq \iota \leq \infty }\left| \breve{s}_{\iota }\right) \left( 
\breve{s}_{\iota }^{d}\right| =\hat{P}_{S_{D}}  \nonumber \\
I &=&\left| {\bf \breve{h}}\right) \left( {\bf \breve{h}}^{d}\right| +\left| 
{\bf \breve{k}}\right) \left( {\bf \breve{k}}^{d}\right| +\left| {\bf \breve{%
s}}\right) \left( {\bf \breve{s}}^{d}\right|  \nonumber \\
I &=&\left| {\cal K}\right) \left( {\cal K}\right| =\int \left| {\sf z}_{1},%
{\sf z}_{2}\right) \left( {\sf z}_{1},{\sf z}_{2}\right| d{\sf z}_{1}d{\sf z}%
_{2}
\end{eqnarray}%
$\hat{P}_{K_{OD}}$, $\hat{P}_{K_{D}}$ and $\hat{P}_{S_{D}}$ are the
projectors onto $K_{OD}$, $K_{D}$ and $S_{D}$ respectively and their
integral kernels are given by 
\begin{eqnarray}
P_{K_{OD}}\left( {\sf z}_{1},{\sf z}_{2},{\sf z}_{3},{\sf z}_{4}\right)
&=&\sum_{\kappa }\left( {\sf z}_{1},{\sf z}_{2}|\breve{h}_{\kappa }\right)
\left( \breve{h}_{\kappa }^{d}|{\sf z}_{3},{\sf z}_{4}\right) \\
P_{K_{D}}\left( {\sf z}_{1},{\sf z}_{2},{\sf z}_{3},{\sf z}_{4}\right)
&=&\sum_{\kappa }\left( {\sf z}_{1},{\sf z}_{2}|\breve{k}_{\kappa }\right)
\left( \breve{k}_{\kappa }^{d}|{\sf z}_{3},{\sf z}_{4}\right)  \nonumber \\
&=&\sum_{\kappa }\left( {\sf z}_{1},{\sf z}_{1}|\breve{k}_{\kappa }\right)
\left( \breve{k}_{\kappa }^{d}|{\sf z}_{3},{\sf z}_{3}\right) \delta \left( 
{\sf z}_{1}-{\sf z}_{2}\right) \delta \left( {\sf z}_{3}-{\sf z}_{4}\right) 
\nonumber \\
&\equiv &P_{K_{D}}\left( {\sf z}_{1},{\sf z}_{3}\right) \delta \left( {\sf z}%
_{1}-{\sf z}_{2}\right) \delta \left( {\sf z}_{3}-{\sf z}_{4}\right)
\end{eqnarray}%
\begin{eqnarray}
P_{S_{D}}\left( {\sf z}_{1},{\sf z}_{2},{\sf z}_{3},{\sf z}_{4}\right)
&=&\sum_{\kappa }\left( {\sf z}_{1},{\sf z}_{2}|\breve{s}_{\kappa }\right)
\left( \breve{s}_{\kappa }^{d}|{\sf z}_{3},{\sf z}_{4}\right)  \nonumber \\
&=&\sum_{\kappa }\left( {\sf z}_{1},{\sf z}_{1}|\breve{s}_{\kappa }\right)
\left( \breve{s}_{\kappa }^{d}|{\sf z}_{3},{\sf z}_{3}\right) \delta \left( 
{\sf z}_{1}-{\sf z}_{2}\right) \delta \left( {\sf z}_{3}-{\sf z}_{4}\right) 
\nonumber \\
&\equiv &P_{S_{D}}\left( {\sf z}_{1},{\sf z}_{3}\right) \delta \left( {\sf z}%
_{1}-{\sf z}_{2}\right) \delta \left( {\sf z}_{3}-{\sf z}_{4}\right)
\end{eqnarray}%
Note that these projectors are {\bf not }orthogonal projectors and {\it in
fact }there is no property of orthogonality nor adjoint operation on
operators acting in ${\cal B}_{1}\left( {\cal H}^{N}\right) $, as it is{\it %
\ not }a Hilbert space, unless ${\cal H}^{N}$ is finite dimensional. They do
however have the properties 
\begin{eqnarray}
\hat{P}_{A}^{2} &=&\hat{P}_{A}  \nonumber \\
\hat{P}_{A}\hat{P}_{B} &=&0\quad A\neq B
\end{eqnarray}%
The matrix transforming from the $\left\{ \left| {\bf \breve{h}}\right)
,\left| {\bf \breve{k}}\right) ,\left| {\bf \breve{s}}\right) \right\} $
basis to the $\left\{ \left| {\cal K}\right) \right\} $ basis obtained from
Eq.(\ref{KtoH}) is, 
\begin{equation}
{\Bbb S=}\left( 
\begin{array}{c}
\left( {\bf \breve{h}}^{d}|{\cal K}\right) \\ 
\left( {\bf \breve{k}}^{d}|{\cal K}\right) \\ 
\left( {\bf \breve{s}}^{d}|{\cal K}\right)%
\end{array}%
\right)
\end{equation}%
i.e.%
\begin{equation}
\left| {\cal K}\right) =\left( \left| {\bf \breve{h}}\right) ,\left| {\bf 
\breve{k}}\right) ,\left| {\bf \breve{s}}\right) \right) \left( 
\begin{array}{c}
\left( {\bf \breve{h}}^{d}|{\cal K}\right) \\ 
\left( {\bf \breve{k}}^{d}|{\cal K}\right) \\ 
\left( {\bf \breve{s}}^{d}|{\cal K}\right)%
\end{array}%
\right)
\end{equation}%
and using Eq.(\ref{HtoK}) the transformation from the $\left\{ \left| {\cal K%
}\right) \right\} $ to the $\left\{ \left| {\bf \breve{h}}\right) ,\left| 
{\bf \breve{k}}\right) ,\left| {\bf \breve{s}}\right) \right\} $ basis is 
\begin{equation}
{\Bbb S}^{-1}{\Bbb =}\left( 
\begin{array}{ccc}
\left( {\cal K}|{\bf \breve{h}}\right) & \left( {\cal K}|{\bf \breve{k}}%
\right) & \left( {\cal K}|{\bf \breve{s}}\right)%
\end{array}%
\right)
\end{equation}%
i.e.%
\begin{equation}
\left( \left| {\bf \breve{h}}\right) ,\left| {\bf \breve{k}}\right) ,\left| 
{\bf \breve{s}}\right) \right) =\left| {\cal K}\right) \left( 
\begin{array}{ccc}
\left( {\cal K}|{\bf \breve{h}}\right) & \left( {\cal K}|{\bf \breve{k}}%
\right) & \left( {\cal K}|{\bf \breve{s}}\right)%
\end{array}%
\right)
\end{equation}%
{\bf Note }that ${\Bbb S}^{-1}$ {\it is not }the transpose or adjoint matrix
of ${\Bbb S}$ as ${\cal B}_{1}\left( {\cal H}^{N}\right) $ is not self dual,
one can for instance can see that ${\Bbb S}^{-1}$ involves integral kernels $%
\left( {\sf z}_{1},{\sf z}_{2}|\breve{h}_{\iota }\right) $ of trace class
operators, while ${\Bbb S}$ involves integral kernels $\left( \breve{h}%
_{\iota }^{d}|{\sf z}_{1},{\sf z}_{2}\right) $ of bounded operators that do
not have to be of trace class. We also have the biorthogonality condition 
\begin{eqnarray}
&&\left( 
\begin{array}{c}
\left( {\bf \breve{h}}^{d}\right| \\ 
\left( {\bf \breve{k}}^{d}\right| \\ 
\left( {\bf \breve{s}}^{d}\right|%
\end{array}%
\right) \left( \left| {\bf \breve{h}}\right) ,\left| {\bf \breve{k}}\right)
,\left| {\bf \breve{s}}\right) \right) =\left( 
\begin{array}{ccc}
\left( {\bf \breve{h}}^{d}|{\bf \breve{h}}\right) & \left( {\bf \breve{h}}%
^{d}|{\bf \breve{k}}\right) & \left( {\bf \breve{h}}^{d}|{\bf \breve{s}}%
\right) \\ 
\left( {\bf \breve{k}}^{d}|{\bf \breve{h}}\right) & \left( {\bf \breve{k}}%
^{d}|{\bf \breve{k}}\right) & \left( {\bf \breve{k}}^{d}|{\bf \breve{s}}%
\right) \\ 
\left( {\bf \breve{s}}^{d}|{\bf \breve{h}}\right) & \left( {\bf \breve{s}}%
^{d}|{\bf \breve{k}}\right) & \left( {\bf \breve{s}}^{d}|{\bf \breve{s}}%
\right)%
\end{array}%
\right) ={\Bbb I} \\
&\Leftrightarrow &{\Bbb S}^{d}\left( 
\begin{array}{c}
\left( {\bf \breve{h}}^{d}\right| \\ 
\left( {\bf \breve{k}}^{d}\right| \\ 
\left( {\bf \breve{s}}^{d}\right|%
\end{array}%
\right) \left( \left| {\bf \breve{h}}\right) ,\left| {\bf \breve{k}}\right)
,\left| {\bf \breve{s}}\right) \right) {\Bbb S}=\left( {\cal K}|{\cal K}%
\right) ={\cal I=}{\Bbb S}^{d}{\Bbb S}
\end{eqnarray}%
where ${\cal I}$ is the integral kernel $\delta \left( {\sf z}_{1}-{\sf z}%
_{3}\right) \delta \left( {\sf z}_{2}-{\sf z}_{4}\right) $ and ${\Bbb S}%
^{d}: $ ${\cal B}\left( {\cal H}^{N}\right) \rightarrow {\cal B}\left( {\cal %
H}^{N}\right) $ is the transformation between the dual bases $\left( 
\begin{array}{c}
\left( {\bf \breve{h}}^{d}\right| \\ 
\left( {\bf \breve{k}}^{d}\right| \\ 
\left( {\bf \breve{s}}^{d}\right|%
\end{array}%
\right) \rightarrow \left( {\cal K}\right| $, which shows that ${\Bbb S}^{d}=%
{\Bbb S}^{-1}$. Hence we obtain the identities 
\begin{eqnarray}
&&{\cal \hat{M}}\breve{Y}=\left( \left| {\bf \breve{h}}\right) ,\left| {\bf 
\breve{k}}\right) ,\left| {\bf \breve{s}}\right) \right) {\Bbb SS}%
^{-1}\left( 
\begin{array}{ccc}
{\Bbb M}_{K_{OD}} & {\Bbb M}_{K_{OD}K_{D}} & {\Bbb M}_{K_{OD}S_{D}} \\ 
{\Bbb M}_{K_{D}K_{OD}} & {\Bbb M}_{K_{D}} & {\Bbb M}_{K_{D}S_{D}} \\ 
{\Bbb M}_{S_{D}K_{OD}} & {\Bbb M}_{S_{D}K_{D}} & {\Bbb T}_{S_{D}}%
\end{array}%
\right) {\Bbb SS}^{-1}\left( 
\begin{array}{c}
{\Bbb Y}_{_{K_{OD}}} \\ 
{\Bbb Y}_{_{K_{D}}} \\ 
{\Bbb Y}_{_{S_{D}}}%
\end{array}%
\right) = \\
&&  \nonumber \\
&&\left( \left| {\cal K}\right) \right) {\Bbb S}^{-1}\left( 
\begin{array}{ccc}
{\Bbb M}_{K_{OD}} & {\Bbb M}_{K_{OD}K_{D}} & {\Bbb M}_{K_{OD}S_{D}} \\ 
{\Bbb M}_{K_{D}K_{OD}} & {\Bbb M}_{K_{D}} & {\Bbb M}_{K_{D}S_{D}} \\ 
{\Bbb M}_{S_{D}K_{OD}} & {\Bbb M}_{S_{D}K_{D}} & {\Bbb T}_{S_{D}}%
\end{array}%
\right) {\Bbb S}\left( {\cal K}|Y\right) =\left( \left| {\cal K}\right)
\right) \left( {\cal K}|{\cal \hat{M}K}\right) \left( {\cal K}|Y\right)
\end{eqnarray}%
where 
\begin{equation}
\left( {\cal K}|{\cal \hat{M}K}\right) =\left( 
\begin{array}{ccc}
\left( {\cal K}|{\bf \breve{h}}\right) & \left( {\cal K}|{\bf \breve{k}}%
\right) & \left( {\cal K}|{\bf \breve{s}}\right)%
\end{array}%
\right) \left( 
\begin{array}{ccc}
{\Bbb M}_{K_{OD}} & {\Bbb M}_{K_{OD}K_{D}} & {\Bbb M}_{K_{OD}S_{D}} \\ 
{\Bbb M}_{K_{D}K_{OD}} & {\Bbb M}_{K_{D}} & {\Bbb M}_{K_{D}S_{D}} \\ 
{\Bbb M}_{S_{D}K_{OD}} & {\Bbb M}_{S_{D}K_{D}} & {\Bbb M}_{S_{D}}%
\end{array}%
\right) \left( 
\begin{array}{c}
\left( {\bf \breve{h}}^{d}|{\cal K}\right) \\ 
\left( {\bf \breve{k}}^{d}|{\cal K}\right) \\ 
\left( {\bf \breve{s}}^{d}|{\cal K}\right)%
\end{array}%
\right)
\end{equation}%
We also have that 
\begin{eqnarray}
&&{\cal \hat{M}}\breve{Y}=\left( \left| {\cal K}\right) \right) {\Bbb S}%
^{-1}\left( 
\begin{array}{ccc}
{\Bbb M}_{K_{OD}} & {\Bbb M}_{K_{OD}K_{D}} & {\Bbb M}_{K_{OD}S_{D}} \\ 
{\Bbb M}_{K_{D}K_{OD}} & {\Bbb M}_{K_{D}} & {\Bbb M}_{K_{D}S_{D}} \\ 
{\Bbb M}_{S_{D}K_{OD}} & {\Bbb M}_{S_{D}K_{D}} & {\Bbb M}_{S_{D}}%
\end{array}%
\right) {\Bbb SS}^{-1}\left( 
\begin{array}{c}
{\Bbb Y}_{_{K_{OD}}} \\ 
{\Bbb Y}_{_{K_{D}}} \\ 
{\Bbb Y}_{_{S_{D}}}%
\end{array}%
\right) =  \nonumber \\
&&  \nonumber \\
&&\left( \left| {\cal K}\right) \right) \left\{ \left( {\cal K}|{\cal \hat{M}%
}_{K_{OD}}{\cal K}\right) +\left( {\cal K}|{\cal \hat{M}}_{K_{OD}K_{D}}{\cal %
K}\right) +\left( {\cal K}|{\cal \hat{M}}_{K_{OD}S_{D}}{\cal K}\right)
+\right.  \nonumber \\
&&\qquad \left( {\cal K}|{\cal \hat{M}}_{K_{D}K_{OD}}{\cal K}\right) +\left( 
{\cal K}|{\cal \hat{M}}_{K_{D}}{\cal K}\right) +\left( {\cal K}|{\cal \hat{M}%
}_{K_{D}S_{D}}{\cal K}\right) +\left( {\cal K}|{\cal \hat{M}}_{S_{D}K_{OD}}%
{\cal K}\right) +  \nonumber \\
&&\qquad \qquad \left. \left( {\cal K}|{\cal \hat{M}}_{S_{D}K_{D}}{\cal K}%
\right) +\left( {\cal K}|{\cal \hat{M}}_{S_{D}}{\cal K}\right) \right\}
\left\{ \left( {\cal K}|Y_{K_{OD}}\right) +\left( {\cal K}|Y_{K_{D}}\right)
+\left( {\cal K}|Y_{S_{D}}\right) \right\}  \label{func}
\end{eqnarray}%
where 
\begin{equation}
{\cal \hat{M}}_{AB}=\hat{P}_{A}{\cal \hat{M}}\hat{P}_{B}\quad \text{and\quad 
}Y_{A}=\hat{P}_{A}Y
\end{equation}%
for various subspaces $A$ and $B$ of ${\cal B}_{1}\left( {\cal H}^{N}\right) 
$. The expression in Eq.(\ref{func}) can be further developed to give 
\begin{equation}
{\cal \hat{M}}\breve{Y}=\breve{Y}^{\prime }=\left( \left| {\cal K}\right)
\right) \left\{ \left( {\cal K}|Y_{K_{OD}}^{\prime }\right) +\left( {\cal K}%
|Y_{K_{D}}^{\prime }\right) +\left( {\cal K}|Y_{S_{D}}^{\prime }\right)
\right\}
\end{equation}%
where 
\begin{eqnarray}
\left( 
\begin{array}{c}
\left( {\cal K}|Y_{K_{OD}}^{\prime }\right) \\ 
\left( {\cal K}|Y_{K_{D}}^{\prime }\right) \\ 
\left( {\cal K}|Y_{S_{D}}^{\prime }\right)%
\end{array}%
\right) &=&\left( 
\begin{array}{ccc}
\left( {\cal K}|{\cal \hat{M}}_{K_{OD}}{\cal K}\right) & \left( {\cal K}|%
{\cal \hat{M}}_{K_{OD}K_{D}}{\cal K}\right) & \left( {\cal K}|{\cal \hat{M}}%
_{K_{OD}S_{D}}{\cal K}\right) \\ 
\left( {\cal K}|{\cal \hat{M}}_{K_{D}K_{OD}}{\cal K}\right) & \left( {\cal K}%
|{\cal \hat{M}}_{K_{D}}{\cal K}\right) & \left( {\cal K}|{\cal \hat{M}}%
_{S_{D}K_{OD}}{\cal K}\right) \\ 
\left( {\cal K}|{\cal \hat{M}}_{S_{D}K_{D}}{\cal K}\right) & \left( {\cal K}|%
{\cal \hat{M}}_{S_{D}K_{D}}{\cal K}\right) & \left( {\cal K}|{\cal \hat{M}}%
_{S_{D}}{\cal K}\right)%
\end{array}%
\right) \left( 
\begin{array}{c}
\left( {\cal K}|Y_{K_{OD}}\right) \\ 
\left( {\cal K}|Y_{K_{D}}\right) \\ 
\left( {\cal K}|Y_{S_{D}}\right)%
\end{array}%
\right)  \nonumber \\
&&
\end{eqnarray}

Transformations of ${\cal B}_{1}\left( {\cal H}^{N}\right) $ that preserve
the vector bundle structure of ${\frak B}_{vect}$ can be expressed as a sum%
\begin{eqnarray}
{\frak \hat{T}} &{\frak =}&{\frak \hat{T}}_{K_{OD}}+{\frak \hat{T}}_{K_{D}}+%
{\frak \hat{T}}_{S_{D}}+{\frak \hat{T}}_{K_{D}S_{D}}+{\frak \hat{T}}%
_{K_{OD}S_{D}}+{\frak \hat{T}}_{K_{OD}K_{D}}+{\frak \hat{T}}_{K_{D}K_{OD}} 
\nonumber \\
{\frak \hat{T}} &:&{\frak B}_{vect}\longrightarrow {\frak B}_{vect}
\end{eqnarray}%
where ${\frak \hat{T}}_{AB}$ denotes maps between subspaces $A$ and $B$.
These can be represented as matrices wrt the bases introduced in Eqs. (\ref%
{hbasis}),(\ref{kbasis}) and (\ref{sbasis}) as 
\begin{equation}
\left( 
\begin{array}{ccc}
{\Bbb T}_{K_{OD}} & {\Bbb T}_{K_{OD}K_{D}} & {\Bbb T}_{K_{OD}S_{D}} \\ 
{\Bbb T}_{K_{D}K_{OD}} & {\Bbb T}_{K_{D}} & {\Bbb T}_{K_{D}S_{D}} \\ 
{\Bbb O} & {\Bbb O} & {\Bbb T}_{S_{D}}%
\end{array}%
\right)
\end{equation}%
The non zero components can be expanded as%
\begin{eqnarray}
{\frak \hat{T}}_{K_{OD}} &=&\sum_{\kappa ,\iota }{\frak T}_{K_{OD}\kappa
\iota }\left| \breve{h}_{\kappa }\right) \left( \breve{h}_{\iota }^{d}\right|
\nonumber \\
{\frak \hat{T}}_{K_{D}} &=&\sum_{\kappa ,\iota }{\frak T}_{K_{D}\kappa \iota
}\left| \breve{k}_{\kappa }\right) \left( \breve{k}_{\iota }^{d}\right| 
\nonumber \\
{\frak \hat{T}}_{S_{D}} &=&\sum_{\kappa ,\iota }{\frak T}_{S_{D}\kappa \iota
}\left| \breve{s}_{\kappa }\right) \left( \breve{s}_{\iota }^{d}\right| 
\nonumber \\
{\frak \hat{T}}_{K_{D}S_{D}} &=&\sum_{\kappa ,\iota }{\frak T}%
_{K_{D}S_{D}\kappa \iota }\left| \breve{k}_{\kappa }\right) \left( \breve{s}%
_{\iota }^{d}\right|  \nonumber \\
{\frak \hat{T}}_{K_{OD}K_{D}} &=&\sum_{\kappa ,\iota }{\frak T}%
_{K_{OD}K_{D}\kappa \iota }\left| \breve{h}_{\kappa }\right) \left( \breve{k}%
_{\iota }^{d}\right|  \nonumber \\
{\frak \hat{T}}_{K_{D}K_{OD}} &=&\sum_{\kappa ,\iota }{\frak T}%
_{K_{D}K_{OD}\kappa \iota }\left| \breve{k}_{\kappa }\right) \left( \breve{h}%
_{\iota }^{d}\right|  \nonumber \\
{\frak \hat{T}}_{K_{OD}S_{D}} &=&\sum_{\kappa ,\iota }{\frak T}%
_{K_{OD}S_{D}\kappa \iota }\left| \breve{h}_{\kappa }\right) \left( \breve{s}%
_{\iota }^{d}\right|
\end{eqnarray}

In order to analyze the properties of Lie groups formed from such operators
it is convenient to express them as ${\frak \hat{T}}=\hat{I}+\hat{\Lambda}%
=e^{\hat{\Upsilon}}$, then%
\begin{eqnarray}
{\frak \hat{T}}_{K_{OD}} &=&\hat{P}_{K_{OD}}+\hat{\Lambda}_{K_{OD}}=\hat{P}%
_{K_{OD}}{\frak \hat{T}}\hat{P}_{K_{OD}}  \nonumber \\
{\frak \hat{T}}_{K_{D}} &=&\hat{P}_{K_{D}}+\hat{\Lambda}_{K_{D}}=\hat{P}%
_{K_{D}}{\frak \hat{T}}\hat{P}_{K_{D}}  \nonumber \\
{\frak \hat{T}}_{S_{D}} &=&\hat{P}_{S_{D}}+\hat{\Lambda}_{S_{D}}=\hat{P}%
_{S_{D}}{\frak \hat{T}}\hat{P}_{S_{D}}  \nonumber \\
{\frak \hat{T}}_{K_{D}S_{D}} &=&\hat{\Lambda}_{K_{D}S_{D}}=\hat{P}_{K_{D}}%
{\frak \hat{T}}\hat{P}_{S_{D}}  \nonumber \\
{\frak \hat{T}}_{K_{D}K_{OD}} &=&\hat{\Lambda}_{K_{D}K_{OD}}=\hat{P}_{K_{D}}%
{\frak \hat{T}}\hat{P}_{K_{OD}}  \nonumber \\
{\frak \hat{T}}_{K_{OD}K_{D}} &=&\hat{\Lambda}_{K_{OD}K_{D}}=\hat{P}_{K_{OD}}%
{\frak \hat{T}}\hat{P}_{K_{D}}  \nonumber \\
{\frak \hat{T}}_{K_{OD}S_{D}} &=&\hat{\Lambda}_{K_{OD}S_{D}}=\hat{P}_{K_{OD}}%
{\frak \hat{T}}\hat{P}_{S_{D}}
\end{eqnarray}%
(A notational note: in kernels where arguments are repeated we only display
one argument if this does not detract from the clarity of the exposition).
The integral kernels of ${\frak \hat{T}}_{K_{OD}}$, ${\frak \hat{T}}_{K_{D}}$%
, and ${\frak \hat{T}}_{S_{D}}$, are given by 
\begin{eqnarray}
{\frak T}_{K_{OD}}\left( {\sf z}_{1},{\sf z}_{2},{\sf z}_{3},{\sf z}%
_{4}\right) &=&\left( {\sf z}_{1},{\sf z}_{2}|{\frak \hat{T}}_{K_{OD}}{\sf z}%
_{3},{\sf z}_{4}\right) =\left( \left| {\sf z}_{1}\right\rangle \left\langle 
{\sf z}_{2}\right| |{\frak \hat{T}}_{K_{OD}}\left| {\sf z}_{3}\right\rangle
\left\langle {\sf z}_{4}\right| \right)  \nonumber \\
&=&\sum_{\kappa ,\iota }{\frak T}_{K_{OD}\kappa \iota }\left( \left| {\sf z}%
_{1}\right\rangle \left\langle {\sf z}_{2}\right| |\breve{h}_{\kappa
}\right) \left( \breve{h}_{\iota }^{d}|\left| {\sf z}_{3}\right\rangle
\left\langle {\sf z}_{4}\right| \right)  \nonumber \\
&=&\sum_{\kappa ,\iota }{\frak T}_{K_{OD}\kappa \iota }\breve{h}_{\kappa
}\left( {\sf z}_{1},{\sf z}_{2}\right) \breve{h}_{\iota }^{d\ast }\left( 
{\sf z}_{3},{\sf z}_{4}\right) \\
&=&P_{K_{OD}}\left( {\sf z}_{1},{\sf z}_{2},{\sf z}_{3},{\sf z}_{4}\right)
+\sum_{\kappa ,\iota }\Lambda _{K_{OD}\kappa \iota }\breve{h}_{\kappa
}\left( {\sf z}_{1},{\sf z}_{2}\right) \breve{h}_{\iota }^{d\ast }\left( 
{\sf z}_{3},{\sf z}_{4}\right)  \nonumber \\
&=&P_{K_{OD}}\left( {\sf z}_{1},{\sf z}_{2},{\sf z}_{3},{\sf z}_{4}\right)
+\Lambda _{K_{OD}}\left( {\sf z}_{1},{\sf z}_{2},{\sf z}_{3},{\sf z}%
_{4}\right)  \label{KOD}
\end{eqnarray}%
\begin{equation}
\qquad \qquad {\frak T}_{K_{D}}\left( {\sf z}_{1},{\sf z}_{2}\right)
=\sum_{\kappa ,\iota }{\frak T}_{K_{D}\kappa \iota }k_{\kappa }\left( {\sf z}%
_{1}\right) k_{\iota }^{d\ast }\left( {\sf z}_{2}\right) =P_{K_{D}}\left( 
{\sf z}_{1},{\sf z}_{2}\right) +\Lambda _{K_{D}}\left( {\sf z}_{1},{\sf z}%
_{2}\right)  \label{KD}
\end{equation}%
\begin{equation}
\qquad \qquad {\frak T}_{S_{D}}\left( {\sf z}_{1},{\sf z}_{2}\right)
=\sum_{\kappa ,\iota }{\frak T}_{S_{D}\kappa \iota }s_{\kappa }\left( {\sf z}%
_{1}\right) s_{\iota }^{d\ast }\left( {\sf z}_{2}\right) =P_{S_{D}}\left( 
{\sf z}_{1},{\sf z}_{2}\right) +\Lambda _{S_{D}}\left( {\sf z}_{1},{\sf z}%
_{2}\right)  \label{SD}
\end{equation}%
and those of ${\frak \hat{T}}_{K_{OD}K_{D}}$, ${\frak \hat{T}}_{K_{D}K_{OD}}$%
, ${\frak \hat{T}}_{K_{OD}S_{D}}$ and ${\frak \hat{T}}_{K_{D}S_{D}}$ by%
\begin{eqnarray}
{\frak T}_{K_{D}S_{D}}\left( {\sf z}_{1},{\sf z}_{2}\right) &=&\sum_{\kappa
,\iota }{\frak T}_{K_{D}S_{D}\kappa \iota }k_{\kappa }\left( {\sf z}%
_{1}\right) s_{\iota }^{d\ast }\left( {\sf z}_{2}\right)  \nonumber \\
&=&\sum_{\kappa ,\iota }\Lambda _{K_{D}S_{D}\kappa \iota }k_{\kappa }\left( 
{\sf z}_{1}\right) s_{\iota }^{d\ast }\left( {\sf z}_{2}\right) =\Lambda
_{K_{D}S_{D}}\left( {\sf z}_{1},{\sf z}_{2}\right)  \nonumber \\
{\frak T}_{K_{OD}S_{D}}\left( {\sf z}_{1},{\sf z}_{2},{\sf z}_{3}\right)
&=&\sum_{\kappa ,\iota }{\frak T}_{K_{OD}S_{D}\kappa \iota }h_{\kappa
}\left( {\sf z}_{1},{\sf z}_{2}\right) s_{\iota }^{d\ast }\left( {\sf z}%
_{3}\right)  \nonumber \\
&=&\sum_{\kappa ,\iota }\Lambda _{K_{OD}S_{D}\kappa \iota }h_{\kappa }\left( 
{\sf z}_{1},{\sf z}_{2}\right) s_{\iota }^{d\ast }\left( {\sf z}_{3}\right)
=\Lambda _{K_{OD}S_{D}}\left( {\sf z}_{1},{\sf z}_{2},{\sf z}_{3}\right) 
\nonumber \\
{\frak T}_{K_{D}K_{OD}}\left( {\sf z}_{1},{\sf z}_{2},{\sf z}_{3}\right)
&=&\sum_{\kappa ,\iota }{\frak T}_{K_{D}K_{OD}\kappa \iota }k_{\kappa
}\left( {\sf z}_{1}\right) h_{\iota }^{d\ast }\left( {\sf z}_{1},{\sf z}%
_{2}\right)  \nonumber \\
&=&\sum_{\kappa ,\iota }\Lambda _{K_{D}K_{OD}\kappa \iota }k_{\kappa }\left( 
{\sf z}_{1}\right) h_{\iota }^{d\ast }\left( {\sf z}_{2},{\sf z}_{3}\right)
=\Lambda _{K_{D}K_{OD}}\left( {\sf z}_{1},{\sf z}_{2},{\sf z}_{3}\right) 
\nonumber \\
{\frak T}_{K_{OD}K_{D}}\left( {\sf z}_{1},{\sf z}_{2},{\sf z}_{3}\right)
&=&\sum_{\kappa ,\iota }{\frak T}_{K_{OD}K_{D}\kappa \iota }h_{\kappa
}\left( {\sf z}_{1},{\sf z}_{2}\right) k_{\iota }^{d\ast }\left( {\sf z}%
_{3}\right)  \nonumber \\
&=&\sum_{\kappa ,\iota }\Lambda _{K_{OD}K_{D}\kappa \iota }h_{\kappa }\left( 
{\sf z}_{1},{\sf z}_{2}\right) k_{\iota }^{d\ast }\left( {\sf z}_{3}\right)
=\Lambda _{K_{OD}K_{D}}\left( {\sf z}_{1},{\sf z}_{2},{\sf z}_{3}\right)
\label{mixed}
\end{eqnarray}%
The super operator transformations, ${\frak \hat{T}}$, can be associated
with matrices of integral kernels as 
\begin{eqnarray}
{\frak \hat{T}\equiv }\left( 
\begin{array}{ccc}
{\frak T}_{K_{OD}} & {\frak T}_{K_{OD,}K_{D}} & {\frak T}_{K_{D},S_{D}} \\ 
{\frak T}_{K_{D},K_{OD}} & {\frak T}_{K_{D}} & {\frak T}_{K_{D}S_{D}} \\ 
0 & 0 & {\frak T}_{S_{D}}%
\end{array}%
\right) = &&\left( 
\begin{array}{ccc}
P_{K_{OD}} & 0 & 0 \\ 
0 & P_{K_{D}} & 0 \\ 
0 & 0 & P_{S_{D}}%
\end{array}%
\right)  \nonumber \\
&&+\left( 
\begin{array}{ccc}
\Lambda _{K_{OD}} & \Lambda _{K_{D}K_{OD}} & \Lambda _{K_{D}S_{D}} \\ 
\Lambda _{K_{D}K_{OD}} & \Lambda _{K_{D}} & \Lambda _{K_{OD}S_{D}} \\ 
0 & 0 & \Lambda _{S_{D}}%
\end{array}%
\right)
\end{eqnarray}%
which produce the transformed integral kernels $Y^{\prime }\left( {\sf z}%
_{1},{\sf z}_{2}\right) ={\frak \hat{T}}Y\left( {\sf z}_{1},{\sf z}%
_{2}\right) $ of trace class operators in the following manner%
\begin{equation}
\left( {\sf z}_{1},{\sf z}_{2}|{\frak \hat{T}}Y\right) =\left( {\sf z}_{1},%
{\sf z}_{2}|\breve{Y}_{K_{OD}}^{\prime }\right) +\left( {\sf z}_{1},{\sf z}%
_{2}|\breve{Y}_{K_{D}}^{\prime }\right) +\left( {\sf z}_{1},{\sf z}_{2}|%
\breve{Y}_{S_{D}}^{\prime }\right)  \label{transkern}
\end{equation}%
where%
\begin{eqnarray}
&&\left( 
\begin{array}{c}
\left( {\sf z}_{1},{\sf z}_{2}|\breve{Y}_{K_{OD}}^{\prime }\right) \\ 
\left( {\sf z}_{1},{\sf z}_{2}|\breve{Y}_{K_{D}}^{\prime }\right) \\ 
\left( {\sf z}_{1},{\sf z}_{2}|\breve{Y}_{S_{D}}^{\prime }\right)%
\end{array}%
\right) =  \nonumber \\
&&\int \left( 
\begin{array}{ccc}
{\frak T}_{K_{OD}}\left( {\sf z}_{1},{\sf z}_{2},{\sf z}_{3},{\sf z}%
_{4}\right) & {\frak T}_{K_{OD,}K_{D}}\left( {\sf z}_{1},{\sf z}_{2},{\sf z}%
_{3},{\sf z}_{4}\right) & {\frak T}_{K_{D},S_{D}}\left( {\sf z}_{1},{\sf z}%
_{2},{\sf z}_{3},{\sf z}_{4}\right) \\ 
{\frak T}_{K_{D},K_{OD}}\left( {\sf z}_{1},{\sf z}_{2},{\sf z}_{3},{\sf z}%
_{4}\right) & {\frak T}_{K_{D}}\left( {\sf z}_{1},{\sf z}_{2},{\sf z}_{3},%
{\sf z}_{4}\right) & {\frak T}_{K_{D}S_{D}}\left( {\sf z}_{1},{\sf z}_{2},%
{\sf z}_{3},{\sf z}_{4}\right) \\ 
0 & 0 & {\frak T}_{S_{D}}\left( {\sf z}_{1},{\sf z}_{2},{\sf z}_{3},{\sf z}%
_{4}\right)%
\end{array}%
\right)  \nonumber \\
&&\qquad \qquad \qquad \qquad \qquad \qquad \qquad \qquad \qquad \qquad
\qquad \qquad \times \left( 
\begin{array}{c}
Y_{K_{OD}}\left( {\sf z}_{3},{\sf z}_{4}\right) \\ 
Y_{K_{D}}\left( {\sf z}_{3},{\sf z}_{4}\right) \\ 
S_{K_{D}}\left( {\sf z}_{3},{\sf z}_{4}\right)%
\end{array}%
\right) d{\sf z}_{3}d{\sf z}_{4}  \label{kerneq}
\end{eqnarray}%
where 
\begin{eqnarray}
Y_{K_{OD}}\left( {\sf z}_{1,}{\sf z}_{2}\right) &=&\sum_{1\leq \iota \leq
\infty }\left( {\sf z}_{1,}{\sf z}_{2}|\breve{h}_{\iota }\right) \left( 
\breve{h}_{\iota }^{d}|\breve{Y}\right) =\int P_{K_{OD}}\left( {\sf z}_{1},%
{\sf z}_{2},{\sf z}_{3},{\sf z}_{4}\right) Y\left( {\sf z}_{3,}{\sf z}%
_{4}\right) d{\sf z}_{3}d{\sf z}_{4}  \nonumber \\
Y_{K_{D}}\left( {\sf z}_{1,}{\sf z}_{2}\right) &=&\sum_{1\leq \iota \leq
\infty }\left( {\sf z}_{1,}{\sf z}_{2}|\breve{k}_{\iota }\right) \left( 
\breve{k}_{\iota }^{d}|\breve{Y}\right) =\int P_{K_{D}}\left( {\sf z}_{1},%
{\sf z}_{2},{\sf z}_{3},{\sf z}_{4}\right) Y\left( {\sf z}_{3,}{\sf z}%
_{4}\right) d{\sf z}_{3}d{\sf z}_{4}  \nonumber \\
Y_{S_{D}}\left( {\sf z}_{1,}{\sf z}_{2}\right) &=&\sum_{1\leq \iota \leq
\infty }\left( {\sf z}_{1,}{\sf z}_{2}|\breve{s}_{\iota }\right) \left( 
\breve{s}_{\iota }^{d}|\breve{Y}\right) =\int P_{S_{D}}\left( {\sf z}_{1},%
{\sf z}_{2},{\sf z}_{3},{\sf z}_{4}\right) Y\left( {\sf z}_{3,}{\sf z}%
_{4}\right) d{\sf z}_{3}d{\sf z}_{4}
\end{eqnarray}

In the main body of this article we have considered maps, $\hat{T}$, that
map pure states to pure states and have the property $\hat{T}_{K_{OD}K_{D}}=%
\hat{T}_{K_{D}K_{OD}}=\hat{T}_{K_{OD}S_{D}}=0$. These transformations $\hat{T%
}$ are congruence transformations of ${\cal B}_{1}\left( {\cal H}^{N}\right) 
$ formed from multiplication operators (in the coordinate representation)
acting on ${\cal H}^{N}$ and act as {\it super multiplication operators }(in
the coordinate representation){\it \ }on ${\cal B}_{1}\left( {\cal H}%
^{N}\right) $ in the following fashion%
\begin{eqnarray}
\hat{T}\breve{X} &=&\breve{T}\breve{X}\breve{T}^{\dagger
}\Longleftrightarrow \text{ the integral kernel relationships}  \nonumber \\
&&\left( \left| {\sf z}_{1}\right\rangle \left\langle {\sf z}_{2}\right| |%
\hat{T}\breve{X}\right) =\left( \left| {\sf z}_{1}\right\rangle \left\langle 
{\sf z}_{2}\right| |\breve{T}\breve{X}\breve{T}^{\dagger }\right) =Tr\left\{
\left| {\sf z}_{2}\right\rangle \left\langle {\sf z}_{1}\right| \breve{T}%
\breve{X}\breve{T}^{\dagger }\right\}  \nonumber \\
\qquad &=&\left\langle {\sf z}_{1}|\breve{T}\breve{X}\breve{T}^{\dagger }%
{\sf z}_{2}\right\rangle =T\left( {\sf z}_{1}\right) X\left( {\sf z}_{1},%
{\sf z}_{2}\right) T^{\ast }\left( {\sf z}_{2}\right)  \nonumber \\
&=&\left( 1+\phi \left( {\sf z}_{1}\right) \right) ^{\frac{1}{2}}e^{i\omega
\left( {\sf z}_{1}\right) }X\left( {\sf z}_{1},{\sf z}_{2}\right) \left(
1+\phi \left( {\sf z}_{2}\right) \right) ^{\frac{1}{2}}e^{-i\omega \left( 
{\sf z}_{2}\right) }  \nonumber \\
&=&\left( 1+\phi \left( {\sf z}_{1}\right) \right) ^{\frac{1}{2}}\left(
1+\phi \left( {\sf z}_{2}\right) \right) ^{\frac{1}{2}}e^{i\left( \omega
\left( {\sf z}_{1}\right) -\omega \left( {\sf z}_{2}\right) \right) }X\left( 
{\sf z}_{1},{\sf z}_{2}\right)  \nonumber \\
&\equiv &{\rm T}\left( {\sf z}_{1},{\sf z}_{2}\right) X\left( {\sf z}_{1},%
{\sf z}_{2}\right)  \nonumber \\
&\Rightarrow &{\rm T}\left( {\sf z}_{1},{\sf z}_{2}\right) =\left( 1+\phi
\left( {\sf z}_{1}\right) \right) ^{\frac{1}{2}}\left( 1+\phi \left( {\sf z}%
_{2}\right) \right) ^{\frac{1}{2}}e^{i\left( \omega \left( {\sf z}%
_{1}\right) -\omega \left( {\sf z}_{2}\right) \right) }
\end{eqnarray}%
where, as in the main text, $T\left( {\sf z}_{1}\right) =\left( 1+\phi
\left( {\sf z}_{1}\right) \right) ^{\frac{1}{2}}e^{i\omega \left( {\sf z}%
_{1}\right) }$ $=e^{{\cal V}\left( {\sf z}_{1}\right) +i\omega \left( {\sf z}%
_{1}\right) }$ and hence the operators $\hat{T}$ can be expanded as 
\begin{eqnarray}
\hat{T} &=&\int {\rm T}\left( {\sf z}_{1},{\sf z}_{2}\right) \delta \left( 
{\sf z}_{1}-{\sf z}_{3}\right) \delta \left( {\sf z}_{2}-{\sf z}_{4}\right) |%
{\sf z}_{1},{\sf z}_{2}{\bf )}\left( {\sf z}_{3},{\sf z}_{4}\right| d{\sf z}%
_{1}d{\sf z}_{2}d{\sf z}_{3}d{\sf z}_{4}  \nonumber \\
&=&\int \left( 1+\phi \left( {\sf z}_{1}\right) \right) ^{\frac{1}{2}}\left(
1+\phi \left( {\sf z}_{2}\right) \right) ^{\frac{1}{2}}e^{i\left( \omega
\left( {\sf z}_{1}\right) -\omega \left( {\sf z}_{2}\right) \right) }|{\sf z}%
_{1},{\sf z}_{2}{\bf )}\left( {\sf z}_{1},{\sf z}_{2}\right| d{\sf z}_{1}d%
{\sf z}_{2}
\end{eqnarray}%
The operators $\hat{T}$ have the decomposition 
\begin{eqnarray}
\hat{T} &=&\hat{P}_{K_{OD}}\hat{T}\hat{P}_{K_{OD}}+\hat{P}_{K_{D}}\hat{T}%
\hat{P}_{K_{D}}+\hat{P}_{K_{D}}\hat{T}\hat{P}_{K_{S}}+\hat{P}_{S_{D}}\hat{T}%
\hat{P}_{S_{D}}  \nonumber \\
&\equiv &\hat{T}_{K_{OD}}+\hat{T}_{K_{D}}+\hat{T}_{K_{D}S_{D}}+\hat{T}%
_{S_{D}}
\end{eqnarray}%
one can see that $\hat{T}_{d}\equiv \hat{T}_{K_{D}}+\hat{T}_{K_{D}S_{D}}+%
\hat{T}_{K_{S}S_{D}}$ has the expansion 
\begin{equation}
\hat{T}_{d}=\int \left( 1+\phi \left( {\sf z}_{1}\right) \right) |{\sf z}%
_{1},{\sf z}_{1}{\bf )}\left( {\sf z}_{1},{\sf z}_{1}\right| d{\sf z}_{1}
\end{equation}%
We now consider $\hat{T}=\hat{I}+\hat{\Phi}$ and note that 
\begin{equation}
\hat{\Phi}_{d}=\int \phi \left( {\sf z}_{1}\right) |{\sf z}_{1},{\sf z}_{1}%
{\bf )}\left( {\sf z}_{1},{\sf z}_{1}\right| d{\sf z}_{1}
\end{equation}%
and the various projections and integral kernels of projections of $\hat{\Phi%
}_{d}$ are given by 
\begin{eqnarray}
\hat{\Phi}_{K_{D}} &=&\int \Phi _{K_{D}}\left( {\sf z}_{1},{\sf z}%
_{3}\right) |{\sf z}_{1},{\sf z}_{1}{\bf )}\left( {\sf z}_{3},{\sf z}%
_{3}\right| d{\sf z}_{1}d{\sf z}_{3}  \nonumber \\
\Phi _{K_{D}}\left( {\sf z}_{1},{\sf z}_{3}\right) &=&\int P_{K_{D}}\left( 
{\sf z}_{1},{\sf z}_{2}\right) \phi \left( {\sf z}_{2}\right)
P_{K_{D}}\left( {\sf z}_{2},{\sf z}_{3}\right) d{\sf z}_{2}  \nonumber \\
\hat{\Phi}_{K_{D}S_{D}} &=&\int \Phi _{K_{D}S_{D}}\left( {\sf z}_{1},{\sf z}%
_{3}\right) |{\sf z}_{1},{\sf z}_{1}{\bf )}\left( {\sf z}_{3},{\sf z}%
_{3}\right| d{\sf z}_{1}d{\sf z}_{3}  \nonumber \\
\Phi _{K_{D}S_{D}}\left( {\sf z}_{1},{\sf z}_{3}\right) &=&\int
P_{K_{D}}\left( {\sf z}_{1},{\sf z}_{2}\right) \phi \left( {\sf z}%
_{2}\right) P_{S_{D}}\left( {\sf z}_{2},{\sf z}_{3}\right) d{\sf z}_{2} 
\nonumber \\
\hat{\Phi}_{S_{D}} &=&\int \Phi _{S_{D}}\left( {\sf z}_{1},{\sf z}%
_{3}\right) |{\sf z}_{1},{\sf z}_{1}{\bf )}\left( {\sf z}_{3},{\sf z}%
_{3}\right| d{\sf z}_{1}d{\sf z}_{3}  \nonumber \\
\Phi _{S_{D}}\left( {\sf z}_{1},{\sf z}_{3}\right) &=&\int P_{S_{D}}\left( 
{\sf z}_{1},{\sf z}_{2}\right) \phi \left( {\sf z}_{2}\right)
P_{S_{D}}\left( {\sf z}_{2},{\sf z}_{3}\right) d{\sf z}_{2}  \label{kerproj}
\end{eqnarray}%
The super operator $\hat{T}$ can be associated with the following matrices
of integral kernels 
\begin{eqnarray}
&&\hat{T}{\frak \equiv }\left( 
\begin{array}{ccc}
\left( {\cal K}|\hat{T}_{K_{OD}}{\cal K}\right) & 0 & 0 \\ 
0 & \left( {\cal K}|\hat{T}_{K_{D}}{\cal K}\right) & \left( {\cal K}|\hat{T}%
_{K_{D}S_{D}}{\cal K}\right) \\ 
0 & 0 & \left( {\cal K}|\hat{T}_{S_{D}}{\cal K}\right)%
\end{array}%
\right) =  \nonumber \\
&&\left( 
\begin{array}{ccc}
\left( {\cal K}|\hat{P}_{K_{OD}}{\cal K}\right) & 0 & 0 \\ 
0 & \left( {\cal K}|\hat{P}_{K_{D}}{\cal K}\right) & 0 \\ 
0 & 0 & \left( {\cal K}|\hat{P}_{S_{D}}{\cal K}\right)%
\end{array}%
\right) +\left( 
\begin{array}{ccc}
\left( {\cal K}|\hat{\Phi}_{K_{OD}}{\cal K}\right) & 0 & 0 \\ 
0 & \left( {\cal K}|\hat{\Phi}_{K_{D}}{\cal K}\right) & \left( {\cal K}|\hat{%
\Phi}_{K_{D}S_{D}}{\cal K}\right) \\ 
0 & 0 & \left( {\cal K}|\hat{\Phi}_{S_{D}}{\cal K}\right)%
\end{array}%
\right)  \nonumber \\
&&
\end{eqnarray}%
The various integral kernels composing the transformed kernel evaluated in
the manner of Eqs.(\ref{transkern}) and (\ref{kerneq}) are 
\begin{equation}
Y^{\prime }\left( {\sf z}_{1},{\sf z}_{2}\right) =Y_{K_{OD}}^{\prime }\left( 
{\sf z}_{1},{\sf z}_{2}\right) +Y_{K_{D}}^{\prime }\left( {\sf z}_{1},{\sf z}%
_{2}\right) +Y_{S_{D}}^{\prime }\left( {\sf z}_{1},{\sf z}_{2}\right)
\end{equation}%
and are given by 
\begin{eqnarray}
Y_{K_{OD}}^{\prime }\left( {\sf z}_{1},{\sf z}_{2}\right) &=&\left( {\sf z}%
_{1},{\sf z}_{2}|\hat{T}\breve{Y}_{K_{OD}}\right) =\left( {\sf z}_{1},{\sf z}%
_{2}|\hat{T}_{K_{OD}}\breve{Y}_{K_{OD}}\right)  \nonumber \\
&=&\int \left\{ P_{K_{OD}}\left( {\sf z}_{1},{\sf z}_{2},{\sf z}_{3},{\sf z}%
_{4}\right) +\Phi _{K_{OD}}\left( {\sf z}_{1},{\sf z}_{2},{\sf z}_{3},{\sf z}%
_{4}\right) \right\} Y_{K_{OD}}\left( {\sf z}_{3},{\sf z}_{4}\right) d{\sf z}%
_{3}d{\sf z}_{4}  \nonumber \\
&=&\left( 1+\phi \left( {\sf z}_{1}\right) \right) ^{\frac{1}{2}}\left(
1+\phi \left( {\sf z}_{2}\right) \right) ^{\frac{1}{2}}e^{i\left( \omega
\left( {\sf z}_{1}\right) -\omega \left( {\sf z}_{2}\right) \right) }\left(
1-\delta \left( {\sf z}_{1}-{\sf z}_{2}\right) \right) Y_{K_{OD}}\left( {\sf %
z}_{1},{\sf z}_{2}\right)
\end{eqnarray}%
and 
\begin{eqnarray}
Y_{d}^{\prime }\left( {\sf z}_{1},{\sf z}_{1}\right) &=&\left( {\sf z}_{1},%
{\sf z}_{1}|\hat{T}\breve{Y}_{d}\right) =\left\langle {\sf z}_{1}|\breve{T}%
\breve{Y}_{d}\breve{T}^{\dagger }{\sf z}_{1}\right\rangle  \nonumber \\
&=&\int \left\{ P_{d}\left( {\sf z}_{1},{\sf z}_{1},{\sf z}_{3},{\sf z}%
_{3}\right) +\Phi _{K_{d}}\left( {\sf z}_{1},{\sf z}_{1},{\sf z}_{3},{\sf z}%
_{3}\right) \right\} Y_{d}\left( {\sf z}_{3},{\sf z}_{3}\right) d{\sf z}_{3}
\nonumber \\
&=&T\left( {\sf z}_{1}\right) Y_{d}\left( {\sf z}_{1},{\sf z}_{1}\right)
T^{\ast }\left( {\sf z}_{1}\right) =\left| T\left( {\sf z}_{1}\right)
\right| ^{2}Y_{d}\left( {\sf z}_{1}\right) =\left( 1+\phi \left( {\sf z}%
_{1}\right) \right) Y_{d}\left( {\sf z}_{1}\right)
\end{eqnarray}%
which splits into two parts $\left( {\sf z}_{1},{\sf z}_{1}|\hat{T}\breve{Y}%
_{d}\right) =\left( {\sf z}_{1},{\sf z}_{1}|\breve{Y}_{K_{D}}^{\prime
}\right) +\left( {\sf z}_{1},{\sf z}_{1}|\breve{Y}_{S_{D}}^{\prime }\right) $
given by 
\begin{eqnarray}
\left( {\sf z}_{1},{\sf z}_{1}|\breve{Y}_{K_{D}}^{\prime }\right) &=&\int
\left\{ \left( \delta \left( {\sf z}_{1}-{\sf z}_{2}\right) +\Phi
_{K_{D}}\left( {\sf z}_{1},{\sf z}_{2}\right) \right) Y_{K_{D}}\left( {\sf z}%
_{2}\right) +\Phi _{K_{D}S_{D}}\left( {\sf z}_{1},{\sf z}_{2}\right)
Y_{S_{D}}\left( {\sf z}_{2}\right) \right\} d{\sf z}_{2}  \nonumber \\
\left( {\sf z}_{1},{\sf z}_{1}|\breve{Y}_{S_{D}}^{\prime }\right) &=&\int
\left( \delta \left( {\sf z}_{1}-{\sf z}_{2}\right) +\Phi _{S_{D}}\left( 
{\sf z}_{1},{\sf z}_{2}\right) \right) Y_{S_{D}}\left( {\sf z}_{2}\right) d%
{\sf z}_{2}
\end{eqnarray}%
Collecting together the diagonal terms from the last two equations one
obtains 
\begin{eqnarray}
\left( {\sf z}_{1},{\sf z}_{1}|\hat{T}\breve{Y}_{d}\right) &=&\left( 1+\phi
\left( {\sf z}_{1}\right) \right) Y_{d}\left( {\sf z}_{1}\right)  \nonumber
\\
&=&\left\{ Y_{K_{D}}\left( {\sf z}_{1}\right) +Y_{S_{D}}\left( {\sf z}%
_{1}\right) \right\} +\int \left\{ \Phi _{K_{D}}\left( {\sf z}_{1},{\sf z}%
_{2}\right) Y_{K_{D}}\left( {\sf z}_{2}\right) \right.  \nonumber \\
&&+\left. \Phi _{K_{D}S_{D}}\left( {\sf z}_{1},{\sf z}_{2}\right)
Y_{S_{D}}\left( {\sf z}_{2}\right) +\Phi _{K_{D}S_{D}}\left( {\sf z}_{1},%
{\sf z}_{2}\right) Y_{S_{D}}\left( {\sf z}_{2}\right) \right\} d{\sf z}_{2}
\end{eqnarray}%
which gives the important relationship 
\begin{eqnarray}
\phi \left( {\sf z}_{1}\right) Y_{d}\left( {\sf z}_{1}\right) &=&\int
\left\{ \Phi _{K_{D}}\left( {\sf z}_{1},{\sf z}_{2}\right) Y_{K_{D}}\left( 
{\sf z}_{2}\right) \right.  \nonumber \\
&&+\left. \Phi _{K_{D}S_{D}}\left( {\sf z}_{1},{\sf z}_{2}\right)
Y_{S_{D}}\left( {\sf z}_{2}\right) +\Phi _{S_{D}}\left( {\sf z}_{1},{\sf z}%
_{2}\right) Y_{S_{D}}\left( {\sf z}_{2}\right) \right\} d{\sf z}_{2}
\end{eqnarray}%
The projected integral kernels of $\hat{\Phi}$ in the preceding equations
are given in Eq.$\left( \ref{kerproj}\right) $. To summarize, the action of
the $\hat{T}$ {\it super multiplication operator }(in the coordinate
representation) on an operator $\breve{Y}=\breve{Y}_{K_{OD}}+\breve{Y}_{d}=%
\breve{Y}_{K_{OD}}+\breve{Y}_{K_{D}}+\breve{Y}_{S_{D}}$, is given by%
\begin{eqnarray}
&&\left( 
\begin{array}{c}
\left( {\cal K}|\breve{Y}_{K_{OD}}^{\prime }\right) \\ 
\left( {\cal K}|\breve{Y}_{K_{D}}^{\prime }\right) \\ 
\left( {\cal K}|\breve{Y}_{S_{D}}^{\prime }\right)%
\end{array}%
\right) =\left( 
\begin{array}{ccc}
\left( {\cal K}|\hat{T}_{K_{OD}}{\cal K}\right) & 0 & 0 \\ 
0 & \left( {\cal K}|\hat{T}_{K_{D}}{\cal K}\right) & \left( {\cal K}|\hat{T}%
_{S_{D}K_{OD}}{\cal K}\right) \\ 
0 & 0 & \left( {\cal K}|\hat{T}_{S_{D}}{\cal K}\right)%
\end{array}%
\right) \left( 
\begin{array}{c}
\left( {\cal K}|\breve{Y}_{K_{OD}}\right) \\ 
\left( {\cal K}|\breve{Y}_{K_{D}}\right) \\ 
\left( {\cal K}|\breve{Y}_{S_{D}}\right)%
\end{array}%
\right) =  \nonumber \\
&&\left\{ \left( 
\begin{array}{ccc}
\left( {\cal K}|\hat{P}_{K_{OD}}{\cal K}\right) & 0 & 0 \\ 
0 & \left( {\cal K}|\hat{P}_{K_{D}}{\cal K}\right) & 0 \\ 
0 & 0 & \left( {\cal K}|\hat{P}_{S_{D}}{\cal K}\right)%
\end{array}%
\right) +\right.  \nonumber \\
&&\qquad \qquad \qquad \left. \left( 
\begin{array}{ccc}
\left( {\cal K}|\hat{\Phi}_{K_{OD}}{\cal K}\right) & 0 & 0 \\ 
0 & \left( {\cal K}|\hat{\Phi}_{K_{D}}{\cal K}\right) & \left( {\cal K}|\hat{%
\Phi}_{S_{D}K_{OD}}{\cal K}\right) \\ 
0 & 0 & \left( {\cal K}|\hat{\Phi}_{S_{D}}{\cal K}\right)%
\end{array}%
\right) \right\} \left( 
\begin{array}{c}
\left( {\cal K}|\breve{Y}_{K_{OD}}\right) \\ 
\left( {\cal K}|\breve{Y}_{K_{D}}\right) \\ 
\left( {\cal K}|\breve{Y}_{S_{D}}\right)%
\end{array}%
\right)
\end{eqnarray}%
and 
\begin{equation}
\left( {\cal K}|\breve{Y}^{\prime }\right) =\left( {\cal K}|\breve{Y}%
_{K_{OD}}^{\prime }\right) +\left( {\cal K}|\breve{Y}_{K_{D}}^{\prime
}\right) +\left( {\cal K}|\breve{Y}_{S_{D}}^{\prime }\right)
\end{equation}%
where 
\begin{eqnarray}
&&\left( {\sf z}_{1},{\sf z}_{1}|\breve{Y}_{K_{D}}^{\prime }\right) +\left( 
{\sf z}_{1},{\sf z}_{1}|\breve{Y}_{S_{D}}^{\prime }\right) =\left( 1+\phi
\left( {\sf z}_{1}\right) \right) Y_{d}\left( {\sf z}_{1}\right)  \nonumber
\\
&&\qquad \qquad \qquad \qquad =Y_{K_{D}}^{\prime }\left( {\sf z}_{1}\right)
+Y_{S_{D}}^{\prime }\left( {\sf z}_{1}\right)  \nonumber \\
&&\qquad \qquad \qquad \qquad +\Delta _{K_{D}}Y_{K_{D}}\left( {\sf z}%
_{1}\right) +\Delta _{K_{D}S_{D}}Y_{S_{D}}\left( {\sf z}_{1}\right) +\Delta
_{S_{D}}Y_{S_{D}}\left( {\sf z}_{1}\right)  \label{diagker}
\end{eqnarray}%
where the incremental change functions are 
\begin{eqnarray}
\Delta _{K_{D}}Y_{K_{D}}\left( {\sf z}_{1}\right) &=&\int \Phi
_{K_{D}}\left( {\sf z}_{1},{\sf z}_{2}\right) Y_{K_{D}}\left( {\sf z}%
_{2}\right) d{\sf z}_{2}  \nonumber \\
\Delta _{K_{D}S_{D}}Y_{S_{D}}\left( {\sf z}_{1}\right) &=&\int \Phi
_{K_{D}S_{D}}\left( {\sf z}_{1},{\sf z}_{2}\right) Y_{S_{D}}\left( {\sf z}%
_{2}\right) d{\sf z}_{2}  \nonumber \\
\Delta _{S_{D}}Y_{S_{D}}\left( {\sf z}_{1}\right) &=&\int \Phi
_{S_{D}}\left( {\sf z}_{1},{\sf z}_{2}\right) Y_{S_{D}}\left( {\sf z}%
_{2}\right) d{\sf z}_{2}  \label{incrfun}
\end{eqnarray}%
Note that in order for $\hat{T}$ $=I+\hat{\Phi}$ to be a vector bundle map :$%
{\frak B}_{vec}\rightarrow {\frak B}_{vec}$ the function $\phi $ defining
the integral kernel of $\hat{\Phi}$ must have the following properties 
\begin{equation}
\left( \breve{s}_{i}^{d}|\hat{\Phi}\breve{k}_{\kappa }\right) =\int s_{\iota
}^{d\ast }\left( {\sf z}\right) \phi \left( {\sf z}\right) k_{\kappa }\left( 
{\sf z}\right) d{\sf z}=0\quad \forall \iota ,\kappa
\end{equation}

\subsection{Intra and Inter Equivalence Class Transformations}

\subsubsection{Intra Equivalence Class Transformations}

In this section we utilize the decomposition of trace class operators 
\begin{equation}
\breve{Y}\left( f\right) =\breve{Y}_{K_{OD}}+\breve{Y}_{K_{D}}+\breve{\Gamma}%
\left( f\right)
\end{equation}%
adapted to the vector bundle structure of ${\frak B}_{vec}$, where $\Xi
\left( \breve{Y}\left( f\right) \right) =\Xi \left( \breve{\Gamma}\left(
f\right) \right) $ and $\Xi $ is an isomorphism from $S_{D}$ to $L_{1}\left( 
{\tt Z}\right) $. The equivalence class $\left[ f\right] _{N}$ is then given
by 
\begin{equation}
\left[ f\right] _{N}=\left\{ \breve{Y}_{K_{OD}}+\breve{Y}_{K_{D}}+\breve{%
\Gamma}\left( f\right) ;\breve{Y}_{K_{OD}}\in K_{OD},\breve{Y}_{K_{D}}\in
K_{D}\right\}
\end{equation}%
Transformations, $\hat{T}_{\ker },$ that map $\ker \Xi \rightarrow \ker \Xi $%
, where $\ker \Xi =K_{OD}\oplus K_{D}$, and have no effect on $S_{D\text{ }}$
are of the form 
\begin{eqnarray}
\hat{T}_{\ker } &=&\sum_{1\leq \kappa ,\iota \leq \infty }\left\{ {\rm T}%
_{\ker K_{OD}\kappa \iota }\left| \breve{h}_{\kappa }\right) \left( \breve{h}%
_{\iota }^{d}\right| +{\rm T}_{\ker K_{OD}K_{D}\kappa \iota }\left| \breve{h}%
_{\kappa }\right) \left( \breve{k}_{\iota }^{d}\right| \right.  \nonumber \\
&&\qquad \qquad \quad \left. +{\rm T}_{\ker K_{D}K_{OD}\kappa \iota }\left| 
\breve{k}_{\kappa }\right) \left( \breve{h}_{\iota }^{d}\right| +{\rm T}%
_{\ker K_{D}\kappa \iota }\left| \breve{k}_{\kappa }\right) \left( \breve{k}%
_{\iota }^{d}\right| +\left| \breve{s}_{\kappa }\right) \left( \breve{s}%
_{\iota }^{d}\right| \delta _{\iota \kappa }\right\} \\
&=&\hat{P}_{S_{D}}+\sum_{1\leq \kappa ,\iota \leq \infty }\left\{ {\rm T}%
_{\ker K_{OD}\kappa \iota }\left| \breve{h}_{\kappa }\right) \left( \breve{h}%
_{\iota }^{d}\right| +{\rm T}_{\ker K_{OD}K_{D}\kappa \iota }\left| \breve{h}%
_{\kappa }\right) \left( \breve{k}_{\iota }^{d}\right| \right.  \nonumber \\
&&\qquad \qquad \qquad \qquad \qquad \left. \,+{\rm T}_{\ker
K_{D}K_{OD}\kappa \iota }\left| \breve{k}_{\kappa }\right) \left( \breve{h}%
_{\iota }^{d}\right| +{\rm T}_{\ker K_{D}\kappa \iota }\left| \breve{k}%
_{\kappa }\right) \left( \breve{k}_{\iota }^{d}\right| \right\} \\
&=&\int {\rm T}_{\ker }\left( {\sf z}_{1},{\sf z}_{2},{\sf z}_{3},{\sf z}%
_{4}\right) |{\sf z}_{1},{\sf z}_{2}{\bf )}\left( {\sf z}_{3},{\sf z}%
_{4}\right| d{\sf z}_{1}d{\sf z}_{2}d{\sf z}_{3}d{\sf z}_{4}
\end{eqnarray}%
and map any given element of $\left[ f\right] _{N}$ into all others in $%
\left[ f\right] _{N}$.

Invertible transformations have the added property that they can be
expressed as exponentials 
\begin{equation}
\hat{T}_{\ker }=e^{\sum_{1\leq \kappa ,\iota \leq \infty }\left\{ {\rm %
\Upsilon }_{\ker K_{OD}\kappa \iota }\left| \breve{h}_{\kappa }\right)
\left( \breve{h}_{\iota }^{d}\right| +{\rm \Upsilon }_{\ker
K_{OD}K_{D}\kappa \iota }\left| \breve{h}_{\kappa }\right) \left( \breve{k}%
_{\iota }^{d}\right| +{\rm \Upsilon }_{\ker K_{D}K_{OD}\kappa \iota }\left| 
\breve{k}_{\kappa }\right) \left( \breve{h}_{\iota }^{d}\right| +{\rm %
\Upsilon }_{\ker K_{D}\kappa \iota }\left| \breve{k}_{\kappa }\right) \left( 
\breve{k}_{\iota }^{d}\right| \right\} }  \label{tker}
\end{equation}%
and form a representation of the group $GL\left( \ker \Xi \right) ,$ which
acts homogeneously\cite{homogen} in $\left[ f\right] _{N}$. More generally
we can form a representation of $GL\left( \ker \Xi \right) $ by invertible
super operators of the form 
\begin{equation}
\hat{T}_{\ker }^{\prime }=e^{\sum_{1\leq \iota ,\kappa \leq \infty }{\rm %
\Upsilon }_{\ker S_{D}}\left( \hat{T}_{\ker }\right) _{\kappa \iota }\left| 
\breve{s}_{\kappa }\right) \left( \breve{s}_{\iota }^{d}\right| }\hat{T}%
_{\ker }=\hat{t}\left( \hat{T}_{\ker }\right) \hat{T}_{\ker }
\end{equation}%
for coefficients $\left\{ {\rm \Upsilon }_{\ker S_{D}}\left( \hat{T}_{\ker
}\right) _{\kappa \iota }\right\} $ that are determined by the $\hat{T}%
_{\ker }$ given in Eq.$\left( \ref{tker}\right) $ and that satisfy the group
property%
\begin{equation}
\hat{t}\left( \hat{T}_{\ker 1}\right) \hat{t}\left( \hat{T}_{\ker 2}\right) =%
\hat{t}\left( \hat{T}_{\ker 1}\hat{T}_{\ker 2}\right)
\end{equation}%
but such transformations would not be intra fiber transformations. The
matrix coefficients $\left\{ {\rm \Upsilon }_{\ker K_{OD}\kappa \iota },{\rm %
\Upsilon }_{\ker K_{OD}K_{D}\kappa \iota },{\rm \Upsilon }_{\ker
K_{D}K_{OD}\kappa \iota },{\rm \Upsilon }_{\ker K_{D}\kappa \iota }\right\} $
parameterize the equivalence class in a redundant manner (i.e. two or more
sets of parameters are associated with a single element), however if we
start from a reference element of $\left[ f\right] _{N}$, which has the
decomposition, 
\begin{equation}
\breve{Y}_{ref}\left( f\right) =\breve{Y}_{refK_{OD}}+\breve{Y}_{refK_{D}}+%
\breve{\Gamma}\left( f\right)
\end{equation}%
we can generate all of $\left[ f\right] _{N}$ by acting on $\left| \breve{Y}%
_{ref}\left( f\right) \right) $ with the transformations 
\begin{equation}
\hat{T}_{\ker \breve{Y}_{ref}}=e^{\sum_{1\leq \kappa \leq \infty }\left\{
\Upsilon _{\ker K_{OD}\kappa \iota }\left| \breve{h}_{\kappa }\right) \left( 
\breve{Y}_{refK_{OD}}^{d}\right| +\Upsilon _{\ker K_{D}\kappa \iota }\left| 
\breve{k}_{\kappa }\right) \left( \breve{Y}_{refK_{D}}^{d}\right| \right\} }
\label{parker}
\end{equation}%
that form a coset as described in the summary section, where $\breve{h}_{1}=%
\breve{Y}_{refK_{OD}}$ and $\breve{k}_{1\iota }=\breve{Y}_{refK_{D}}$. This
gives a non redundant parameterization for the action of the group $GL\left(
\ker \Xi \right) $ on the element $\breve{Y}_{ref}\left( f\right) $, which
produces all the members of $\left[ f\right] _{N}$.

We now restrict attention to the subgroup, ${\frak G}_{fib}{\frak \subset }%
GL\left( \ker \Xi \right) $, of super multiplication operators, (in the
coordinate representation){\it \ }that can be represented by super operators
of the form $\hat{T}_{fib}=\hat{T}_{fibK_{OD}{\rm \ }}+\hat{T}_{fibK_{D}{\rm %
\ }}+\hat{P}_{S_{D}}$ where 
\begin{eqnarray}
\hat{T}_{fibK_{OD}{\rm \ }} &=&\hat{P}_{K_{OD}}+\hat{\Phi}_{K_{OD}} 
\nonumber \\
\hat{T}_{fibK_{D}{\rm \ }} &=&\hat{P}_{K_{D}}+\hat{\Phi}_{K_{D}}
\end{eqnarray}%
that map {\it pure states to pure states} and are generated by the functions 
$\phi _{K_{D}}$ and a fixed $\omega $ for $\hat{T}_{fibK_{OD}{\rm \ }}$,
that have the properties 
\begin{eqnarray}
\left( \breve{s}_{i}^{d}|\hat{\Phi}_{K_{D}}\breve{s}_{\kappa }\right)
&=&\int s_{\iota }^{d\ast }\left( {\sf z}\right) \phi _{K_{D}}\left( {\sf z}%
\right) s_{\kappa }\left( {\sf z}\right) d{\sf z}=0\quad \forall \iota
,\kappa  \nonumber \\
\left( \breve{k}_{i}^{d}|\hat{\Phi}_{K_{D}}\breve{s}_{\kappa }\right)
&=&\int k_{\iota }^{d\ast }\left( {\sf z}\right) \phi _{K_{D}}\left( {\sf z}%
\right) s_{\kappa }\left( {\sf z}\right) d{\sf z}=0\quad \forall \iota
,\kappa  \nonumber \\
\left( \breve{s}_{i}^{d}|\hat{\Phi}_{K_{D}}\breve{k}_{\kappa }\right)
&=&\int s_{\iota }^{d\ast }\left( {\sf z}\right) \phi _{K_{D}}\left( {\sf z}%
\right) k_{\kappa }\left( {\sf z}\right) d{\sf z}=0\quad \forall \iota
,\kappa  \nonumber \\
\left( \breve{k}_{i}^{d}|\hat{\Phi}_{K_{D}}\breve{k}_{\kappa }\right)
&=&\int k_{\iota }^{d\ast }\left( {\sf z}\right) \phi _{K_{D}}\left( {\sf z}%
\right) k_{\kappa }\left( {\sf z}\right) d{\sf z}\neq 0\quad \text{for some }%
\iota ,\kappa  \label{phicond}
\end{eqnarray}%
Functions with these properties can be expanded as 
\begin{eqnarray}
\phi _{K_{D}}\left( {\sf z}\right) &=&\sum_{1\leq \kappa ,\iota \leq \infty
}\phi _{K_{D}\kappa \iota }\left( {\sf z},{\sf z}|\breve{k}_{\kappa }\right)
\left( \breve{k}_{\iota }^{d}|{\sf z},{\sf z}\right)  \nonumber \\
&=&\int P_{K_{D}}\left( {\sf z},{\sf z}^{\prime }\right) \phi _{K_{D}}\left( 
{\sf z}^{\prime }\right) P_{K_{D}}\left( {\sf z}^{\prime },{\sf z}\right) d%
{\sf z}^{\prime }  \nonumber \\
\phi _{K_{D}\kappa \iota } &=&\int k_{\kappa }^{d\ast }\left( {\sf z}\right)
\phi _{K_{D}}\left( {\sf z}\right) k_{\iota }\left( {\sf z}\right) d{\sf z}
\label{PhiKD}
\end{eqnarray}%
However note that it is {\bf not }sufficient for functions to have
expansions of the form Eq.$\left( \ref{PhiKD}\right) $ in order for them to
satisfy the conditions expressed in Eq.$\left( \ref{phicond}\right) $, it is
only necessary. This can be seen, for instance, by considering the following
expansions with arbitrary $\phi _{K_{D}\kappa \iota }$ 
\begin{eqnarray}
&&\int s_{r}^{d\ast }\left( {\sf z}\right) \sum_{1\leq \kappa ,\iota \leq
\infty }\phi _{K_{D}\kappa \iota }\left( {\sf z},{\sf z}|\breve{k}_{\kappa
}\right) \left( \breve{k}_{\iota }^{d}|{\sf z},{\sf z}\right) s_{s}\left( 
{\sf z}\right) d{\sf z}  \nonumber \\
&=&\int \sum_{1\leq \kappa ,\iota \leq \infty }\phi _{K_{D}\kappa \iota
}\left( \breve{s}_{r}^{d}|{\sf z},{\sf z}\right) \left( {\sf z},{\sf z}|%
\breve{k}_{\kappa }\right) \left( \breve{k}_{\iota }^{d}|{\sf z},{\sf z}%
\right) \left( {\sf z},{\sf z}|\breve{s}_{s}\right) d{\sf z}
\end{eqnarray}%
which {\it might or might not} equal zero. The super operator $\hat{\Phi}%
_{K_{D}}$ is defined in terms of functions satisfying Eq.$\left( \ref%
{phicond}\right) $by the integral kernel 
\begin{equation}
\Phi _{K_{D}}\left( {\sf z}_{1,}{\sf z}_{2}\right) =\phi _{K_{D}}\left( {\sf %
z}_{1}\right) \delta \left( {\sf z}_{1}-{\sf z}_{2}\right)
\end{equation}%
The diagonal part of the transformed integral kernel $\left( {\sf z}_{1},%
{\sf z}_{1}|\breve{Y}_{d}^{\prime }\right) =\left( {\sf z}_{1},{\sf z}_{1}|%
\hat{T}_{fib}\breve{Y}_{d}\right) $ can be developed as 
\begin{eqnarray}
&&\left( {\sf z}_{1},{\sf z}_{1}|\breve{Y}_{K_{D}}^{\prime }\right) +\left( 
{\sf z}_{1},{\sf z}_{1}|\breve{\Gamma}\left( f\right) \right) =  \nonumber \\
&&\qquad \qquad \left( 1+\phi _{K_{D}}\left( {\sf z}_{1}\right) \right)
Y_{d}\left( {\sf z}_{1}\right) =  \nonumber \\
&&Y_{K_{D}}\left( {\sf z}_{1}\right) +\Gamma \left( f\right) \left( {\sf z}%
_{1}\right) +\Delta _{K_{D}}Y_{K_{D}}\left( {\sf z}_{1}\right) +\Delta
_{K_{D}S_{D}}\Gamma \left( f\right) \left( {\sf z}_{1}\right) +\Delta
_{S_{D}}\Gamma \left( f\right) \left( {\sf z}_{1}\right)  \nonumber \\
&&\qquad \qquad \qquad \qquad \qquad \qquad \quad =Y_{K_{D}}\left( {\sf z}%
_{1}\right) +\Gamma \left( f\right) \left( {\sf z}_{1}\right) +\Delta
_{K_{D}}Y_{K_{D}}\left( {\sf z}_{1}\right)  \label{diagexp}
\end{eqnarray}%
the last line in the above equation results from 
\begin{eqnarray}
\Phi _{K_{D}S_{D}}\left( {\sf z}_{1},{\sf z}_{2}\right) &=&\int
P_{K_{D}}\left( {\sf z}_{1},{\sf z}_{1}^{\prime }\right) \phi _{K_{D}}\left( 
{\sf z}_{1}^{\prime }\right) P_{S_{D}}\left( {\sf z}_{1}^{\prime },{\sf z}%
_{2}\right) d{\sf z}_{1}^{\prime }=0  \label{KDSD} \\
\Phi _{S_{D}}\left( {\sf z}_{1},{\sf z}_{2}\right) &=&\int P_{S_{D}}\left( 
{\sf z}_{1},{\sf z}_{1}^{\prime }\right) \phi _{K_{D}}\left( {\sf z}%
_{1}^{\prime }\right) P_{S_{D}}\left( {\sf z}_{1}^{\prime },{\sf z}%
_{2}\right) d{\sf z}_{1}^{\prime }=0  \label{PhiSD}
\end{eqnarray}%
and the definitions of the incremental functions Eq.$\left( \ref{incrfun}%
\right) $. Thus Eq.$\left( \ref{diagexp}\right) $ provides the identity 
\begin{equation}
\phi _{K_{D}}\left( {\sf z}_{1}\right) Y_{d}\left( {\sf z}_{1}\right)
=\Delta _{K_{D}}Y_{K_{D}}\left( {\sf z}_{1}\right) =\phi _{K_{D}}\left( {\sf %
z}_{1}\right) Y_{K_{D}}\left( {\sf z}_{1}\right)
\end{equation}%
as a result of Eqs.$\left( \ref{KDSD}\right) $, $\left( \ref{PhiSD}\right) $
and 
\begin{eqnarray}
\Delta _{K_{D}}Y_{K_{D}}\left( {\sf z}_{1}\right) &=&\int \Phi
_{K_{D}}\left( {\sf z}_{1},{\sf z}_{2}\right) Y_{K_{D}}\left( {\sf z}%
_{2}\right) d{\sf z}_{2}  \nonumber \\
&=&\int \phi _{K_{D}}\left( {\sf z}_{1}\right) \delta \left( {\sf z}_{1}-%
{\sf z}_{2}\right) Y_{K_{D}}\left( {\sf z}_{2}\right) d{\sf z}_{2}=\phi
_{K_{D}}\left( {\sf z}_{1}\right) Y_{K_{D}}\left( {\sf z}_{1}\right)
\end{eqnarray}%
The integral kernel of the super operator $\hat{T}_{fibK_{D}{\rm \ }}$ thus
has the expansions 
\begin{eqnarray}
\hat{T}_{fibK_{D}{\rm \ }} &=&\hat{P}_{K_{D}}+\hat{\Phi}_{K_{D}} \\
&=&\int \left\{ P_{K_{D}}\left( {\sf z}_{1},{\sf z}_{2}\right) +\phi
_{K_{D}}\left( {\sf z}_{1}\right) \delta \left( {\sf z}_{1}-{\sf z}%
_{2}\right) \right\} \left| {\sf z}_{1},{\sf z}_{2}\right) \left( {\sf z}%
_{1},{\sf z}_{2}\right| d{\sf z}_{1}d{\sf z}_{2}
\end{eqnarray}%
The integral kernel of the super operator $\hat{T}_{fibK_{OD}{\rm \ }}=\hat{P%
}_{K_{OD}}+\hat{\Phi}_{K_{OD}}$ which is determined up to phase by $\phi
_{K_{D}}\left( {\sf z}_{1}\right) $ has the expansion 
\begin{equation}
\Phi _{K_{OD}}\left( {\sf z}_{1},{\sf z}_{2},{\sf z}_{3,}{\sf z}_{4}\right)
=\sum_{1\leq \kappa ,\iota \leq \infty }\Phi _{K_{OD}\kappa \iota }\left( 
{\sf z}_{1},{\sf z}_{2}|\breve{h}_{\kappa }\right) \left( \breve{h}_{\iota
}^{d}|{\sf z}_{3},{\sf z}_{4}\right)
\end{equation}%
where 
\begin{equation}
\Phi _{K_{OD}\kappa \iota }=\int k_{\kappa }^{d\ast }\left( {\sf z}_{1},{\sf %
z}_{2}\right) \Phi _{K_{OD}}\left( {\sf z}_{1,}{\sf z}_{2},{\sf z}_{3,}{\sf z%
}_{4}\right) k_{\iota }\left( {\sf z}_{3},{\sf z}_{4}\right) d{\sf z}_{1}d%
{\sf z}_{2}d{\sf z}_{3}d{\sf z}_{4}
\end{equation}%
and has the projection properties 
\begin{eqnarray}
&&\Phi _{K_{OD}}\left( {\sf z}_{1,}{\sf z}_{2},{\sf z}_{3,}{\sf z}_{4}\right)
\nonumber \\
&=&\int P_{K_{OD}}\left( {\sf z}_{1},{\sf z}_{2},{\sf z}_{1}^{\prime },{\sf z%
}_{2}^{\prime }\right) \Phi _{K_{OD}}\left( {\sf z}_{1}^{\prime },{\sf z}%
_{2}^{\prime },{\sf z}_{3}^{\prime },{\sf z}_{4}^{\prime }\right)
P_{K_{OD}}\left( {\sf z}_{3}^{\prime },{\sf z}_{4}^{\prime },{\sf z}_{3},%
{\sf z}_{4}\right) d{\sf z}_{1}^{\prime }d{\sf z}_{2}^{\prime }d{\sf z}%
_{3}^{\prime }d{\sf z}_{4}^{\prime }  \nonumber \\
&=&\int P_{K_{OD}}\left( {\sf z}_{1},{\sf z}_{2},{\sf z}_{1}^{\prime },{\sf z%
}_{2}^{\prime }\right) \Phi _{K_{OD}}\left( {\sf z}_{1}^{\prime },{\sf z}%
_{2}^{\prime },{\sf z}_{3}^{\prime },{\sf z}_{4}^{\prime }\right) d{\sf z}%
_{1}^{\prime }d{\sf z}_{2}^{\prime }  \nonumber \\
&=&\int \Phi _{K_{OD}}\left( {\sf z}_{1},{\sf z}_{2},{\sf z}_{3}^{\prime },%
{\sf z}_{4}^{\prime }\right) P_{K_{OD}}\left( {\sf z}_{3}^{\prime },{\sf z}%
_{4}^{\prime },{\sf z}_{3},{\sf z}_{4}\right) d{\sf z}_{3}^{\prime }d{\sf z}%
_{4}^{\prime }
\end{eqnarray}%
The action of $\hat{T}_{fibK_{OD}}$ on $\hat{Y}_{K_{OD}}$ is given by 
\begin{eqnarray}
&&Y_{K_{OD}}^{\prime }\left( {\sf z}_{1,}{\sf z}_{2}\right) =\int
T_{fibK_{OD}}\left( {\sf z}_{1,}{\sf z}_{2},{\sf z}_{3},{\sf z}_{4}\right)
Y_{K_{OD}}\left( {\sf z}_{3,}{\sf z}_{4}\right) d{\sf z}_{3}d{\sf z}_{4} 
\nonumber \\
&=&\int \left\{ P_{K_{OD}}\left( {\sf z}_{1,}{\sf z}_{2},{\sf z}_{3},{\sf z}%
_{4}\right) +\Phi _{K_{OD}}\left( {\sf z}_{1,}{\sf z}_{2},{\sf z}_{3},{\sf z}%
_{4}\right) \right\} Y_{K_{OD}}\left( {\sf z}_{3},{\sf z}_{4}\right) d{\sf z}%
_{3}d{\sf z}_{4}  \nonumber \\
&=&\int \left( 1+\phi _{K_{D}}\left( {\sf z}_{1}\right) \right) ^{\frac{1}{2}%
}\left( 1+\phi _{K_{D}}\left( {\sf z}_{3}\right) \right) ^{\frac{1}{2}%
}e^{i\left( \omega \left( {\sf z}_{1}\right) -\omega \left( {\sf z}%
_{3}\right) \right) }  \nonumber \\
&&\times \left( 1-\delta \left( {\sf z}_{1}-{\sf z}_{2}\right) \right)
\left( 1-\delta \left( {\sf z}_{3}-{\sf z}_{4}\right) \right) \delta \left( 
{\sf z}_{1}-{\sf z}_{3}\right) \delta \left( {\sf z}_{3}-{\sf z}_{4}\right)
Y_{K_{OD}}\left( {\sf z}_{3},{\sf z}_{4}\right) d{\sf z}_{3}d{\sf z}_{4} 
\nonumber \\
&=&\left( 1+\phi _{K_{D}}\left( {\sf z}_{1}\right) \right) ^{\frac{1}{2}%
}\left( 1+\phi _{K_{D}}\left( {\sf z}_{2}\right) \right) ^{\frac{1}{2}%
}\left( 1-\delta \left( {\sf z}_{1}-{\sf z}_{2}\right) \right) e^{i\left(
\omega \left( {\sf z}_{1}\right) -\omega \left( {\sf z}_{2}\right) \right)
}Y_{K_{OD}}\left( {\sf z}_{1},{\sf z}_{2}\right)  \label{Yprime}
\end{eqnarray}%
where the phase $\omega $ can be {\it arbitrarily} chosen and represents the
factor from the subgroup ${\cal G}_{\Xi }$. The factor multiplying $%
Y_{K_{OD}}\left( {\sf z}_{1,}{\sf z}_{2}\right) $ in the last line of Eq.$%
\left( \ref{Yprime}\right) $ can be expanded further to give%
\begin{equation}
Eq.\left( \ref{Yprime}\right) =\left( 1+\Lambda _{_{K_{D}}}\left( {\sf z}%
_{1,}{\sf z}_{2}\right) \right) e^{i\left( \omega \left( {\sf z}_{1}\right)
-\omega \left( {\sf z}_{2}\right) \right) }\left( 1-\delta \left( {\sf z}%
_{1}-{\sf z}_{2}\right) \right)  \label{xpan}
\end{equation}%
where $\Lambda _{_{K_{D}}}\left( {\sf z}_{1,}{\sf z}_{2}\right) $ is a power
series in $\left( \phi _{K_{D}}\left( {\sf z}_{1}\right) +\phi
_{K_{D}}\left( {\sf z}_{2}\right) +\phi _{K_{D}}\left( {\sf z}_{1}\right)
\phi _{K_{D}}\left( {\sf z}_{2}\right) \right) $. Hence one can make the
identification 
\begin{equation}
\Phi _{K_{OD}}\left( {\sf z}_{1,}{\sf z}_{2},{\sf z}_{3},{\sf z}_{4}\right)
=\Lambda _{_{K_{D}}}\left( {\sf z}_{1,}{\sf z}_{2}\right) \delta \left( {\sf %
z}_{1}-{\sf z}_{3}\right) \delta \left( {\sf z}_{2}-{\sf z}_{4}\right)
\left( 1-\delta \left( {\sf z}_{1}-{\sf z}_{2}\right) \right)
\end{equation}%
To recap the complete super operator $\hat{T}_{fib}$ is given by 
\begin{equation}
\hat{T}_{fib}=\hat{T}_{fibK_{OD}{\rm \ }}+\hat{T}_{fibK_{D}{\rm \ }}+\hat{P}%
_{S_{D}}
\end{equation}%
where the various components are defined by the integral kernels in the
preceding equations. One should note that in contrast to $GL\left( \ker \Xi
\right) $ the subgroup ${\frak G}_{fib}$ does not act homogeneously on $%
\left[ f\right] _{N}$.

A non redundant parameterization for the action of the subgroup ${\frak G}%
_{fib}$ on a reference operator $\left| \breve{k}_{ref}\right) $ is obtained
by replacing $\phi _{K_{D}}\left( {\sf z}_{1}\right) $ in Eq.$\left( \ref%
{PhiKD}\right) $ by functions that have the additional properties 
\begin{equation}
\left( \breve{k}_{i}^{d}|\hat{\Phi}_{K_{D}}\breve{k}_{\kappa }\right) =\int
k_{\iota }^{d\ast }\left( {\sf z}\right) \phi _{K_{D}}\left( {\sf z}\right)
k_{\kappa }\left( {\sf z}\right) d{\sf z}=0;\quad 2\leq \kappa \leq \infty
,1\leq \iota \leq \infty  \label{phinr}
\end{equation}%
where $\left| \breve{k}_{1}\right) \equiv \left| \breve{k}_{ref}\right) =%
\hat{P}_{K_{D}}\left| \breve{Y}_{ref}\left( f\right) \right) $ to giving the
transformation $\hat{T}_{fib}\left( \breve{k}_{ref}\right) $. The functions
in Eq.$\left( \ref{phinr}\right) $ have the expansion 
\begin{eqnarray}
\phi _{refK_{D}}\left( {\sf z}_{1}\right) &=&\sum_{2\leq \kappa \leq \infty
}\phi _{refK_{D\kappa }}\left( {\sf z}_{1},{\sf z}_{1}|\breve{k}_{\kappa
}\right) \left( \breve{k}_{ref}^{d}|{\sf z}_{1},{\sf z}_{1}\right)  \nonumber
\\
\phi _{refK_{D\kappa }} &=&\int k_{\kappa }^{d\ast }\left( {\sf z}%
_{1}\right) \phi _{K_{D}}\left( {\sf z}_{1}\right) k_{ref}\left( {\sf z}%
_{1}\right) d{\sf z}_{1};\quad 2\leq \kappa \leq \infty  \label{refkd}
\end{eqnarray}%
Unlike the non redundant parameterization of $GL\left( \ker \Xi \right) $,
which produced all of $\left[ f\right] _{N}$ from the reference operator $%
\breve{Y}_{ref}\left( f\right) ,$ ${\frak G}_{fib}$ {\it does not. }However
if $\breve{Y}_{ref}\left( \gamma \right) $ is a pure state this
parameterization of {\it \ }${\frak G}_{fib}$ {\it does }produce all pure
states belonging to $\left[ \gamma \right] _{N}$ when acting on $\left| 
\breve{Y}_{ref}\left( \gamma \right) \right) \forall \gamma \in L_{1}\left( 
{\tt Z}\right) _{+}$.

\subsubsection{Inter Equivalence Class Transformations}

We now wish to consider transformations that map between equivalence
classes, and in particular to find a non redundant parameterization of $%
S_{D} $ by such transformations acting on a reference element. The most
obvious source of such transformations is the subgroup $GL\left(
S_{D}\right) $ represented by super operators of the form 
\begin{eqnarray}
\hat{T}_{S}=e^{\sum_{1\leq \kappa ,\iota \leq \infty }\Upsilon _{S_{D}\kappa
\iota }\left| \breve{s}_{\kappa }\right) \left( \breve{s}_{\iota
}^{d}\right| } &=&\hat{P}_{K_{OD}}+\hat{P}_{K_{D}}+\hat{t}_{S}  \nonumber \\
\hat{t}_{S} &=&\hat{P}_{S_{D}}\hat{T}_{S}\hat{P}_{S_{D}}
\end{eqnarray}%
More generally, however, we can form a reducible representation of $GL\left(
S_{D}\right) $ by invertible super operators of the form 
\begin{eqnarray}
\hat{T}_{S} &=&\hat{T}_{K_{OD}}\left( \hat{t}_{S}\right) +\hat{P}_{K_{D}}+%
\hat{t}_{S}  \nonumber \\
\hat{T}_{K_{OD}}\left( \hat{t}_{S}\right) &=&\hat{P}_{K_{OD}}\hat{T}_{S}\hat{%
P}_{K_{OD}}
\end{eqnarray}%
which are both inter fiber and intra fiber, where $\hat{T}_{K_{OD}}$ is a
representation of $GL\left( S_{D}\right) $. Considering these more general
representations of $GL\left( S_{D}\right) $ allows us to construct non
redundant parameterizations of $S_{D}$ using the a coset of the subgroup $%
{\frak G}_{den}\subset $ $GL\left( S_{D}\right) $ of super multiplication
operators (in the coordinate representation),{\it \ }%
\begin{equation}
\hat{T}_{den}\left( \phi _{S_{D}}\right) =\hat{T}_{K_{OD}{\rm \ }}\left(
\phi _{S_{D}},\omega \right) +\hat{P}_{K_{D}}+\hat{T}_{S_{D}{\rm \ }}\left(
\phi _{S_{D}}\right)  \label{tden}
\end{equation}%
generated by a $\phi _{S_{D}}$, where $\omega $ can be chosen to be zero,
that have the properties 
\begin{eqnarray}
\left( \breve{s}_{i}^{d}|\hat{\Phi}_{S_{D}}\breve{s}_{ref}\right) &=&\int
s_{\iota }^{d\ast }\left( {\sf z}\right) \phi _{S_{D}}\left( {\sf z}\right)
s_{ref}\left( {\sf z}\right) d{\sf z}\neq 0;\quad \text{for some }\iota
;1\leq \iota \leq \infty  \nonumber \\
\left( \breve{s}_{i}^{d}|\hat{\Phi}_{S_{D}}\breve{s}_{\kappa }\right)
&=&\int s_{\iota }^{d\ast }\left( {\sf z}\right) \phi _{S_{D}}\left( {\sf z}%
\right) s_{\kappa }\left( {\sf z}\right) d{\sf z}=0;\quad 1\leq \iota \leq
\infty ,2\leq \kappa \leq \infty  \nonumber \\
\left( \breve{k}_{i}^{d}|\hat{\Phi}_{S_{D}}\breve{s}_{\kappa }\right)
&=&\int k_{\iota }^{d\ast }\left( {\sf z}\right) \phi _{S_{D}}\left( {\sf z}%
\right) s_{\kappa }\left( {\sf z}\right) d{\sf z}=0;\quad 1\leq \iota
,\kappa \leq \infty  \nonumber \\
\left( \breve{s}_{i}^{d}|\hat{\Phi}_{S_{D}}\breve{k}_{\kappa }\right)
&=&\int s_{\iota }^{d\ast }\left( {\sf z}\right) \phi _{S_{D}}\left( {\sf z}%
\right) k_{\kappa }\left( {\sf z}\right) d{\sf z}=0;\quad 1\leq \iota
,\kappa \leq \infty  \nonumber \\
\left( \breve{k}_{i}^{d}|\hat{\Phi}_{S_{D}}\breve{k}_{\kappa }\right)
&=&\int k_{\iota }^{d\ast }\left( {\sf z}\right) \phi _{S_{D}}\left( {\sf z}%
\right) k_{\kappa }\left( {\sf z}\right) d{\sf z}=0;\quad 1\leq \iota
,\kappa \leq \infty
\end{eqnarray}%
{\it together} with a fixed choice of phase $\omega $ for $\hat{T}_{K_{OD}%
{\rm \ }}\left( \phi _{S_{D}}\right) $, where $\breve{s}_{1}=\breve{\Gamma}%
\left( f_{ref}\right) =\hat{P}_{S_{D}}\breve{Y}_{ref}$ \ and $\Xi \left( 
\breve{Y}_{ref}\right) =f_{ref}$. This coset is described in the summary
section. The function $\phi _{S_{D}}$ has the expansion and properties 
\begin{eqnarray}
\phi _{S_{D}}\left( {\sf z}\right) &=&\sum_{1\leq \kappa \leq \infty }\phi
_{S_{D}\kappa ref}\left( {\sf z},{\sf z}|\breve{s}_{\kappa }\right) \left( 
\breve{\Gamma}\left( f_{ref}\right) ^{d}|{\sf z},{\sf z}\right)  \nonumber \\
&=&\int P_{S_{D}}\left( {\sf z},{\sf z}^{\prime }\right) \phi _{S_{D}}\left( 
{\sf z}^{\prime }\right) P_{S_{D}}\left( {\sf z}^{\prime },{\sf z}\right) d%
{\sf z}^{\prime }  \nonumber \\
\phi _{S_{D}\kappa ref} &=&\int s_{\kappa }^{d\ast }\left( {\sf z}\right)
\phi _{S_{D}}\left( {\sf z}\right) \breve{\Gamma}\left( f_{ref}\right)
\left( {\sf z}\right) d{\sf z}
\end{eqnarray}%
The super operators $\hat{T}_{den}\left( \phi _{S_{D}}\right) $ generate
reference elements for all equivalence classes in a unique fashion and form
a non redundant parameterization for the action of the subgroup ${\frak G}%
_{den}$ {\it for a fixed choice of phase function }$\omega $ i.e. 
\begin{equation}
\hat{T}_{den}\left( \phi _{S_{D}}\right) :\left[ f_{ref}\right]
_{N}\rightarrow \left[ f^{\prime }\right] _{N};\quad f\neq f^{\prime }
\end{equation}%
The integral kernels of the super operators appearing in Eq.$\left( \ref%
{tden}\right) $, with suitable replacements of $\phi _{K_{D}}$ by $\phi
_{S_{D}}$, are given by%
\begin{eqnarray}
T_{K_{OD}{\rm \ }}\left( \phi _{S_{D}}\right) \left( {\sf z}_{1,}{\sf z}_{2},%
{\sf z}_{3},{\sf z}_{4}\right) &=&P_{K_{OD}}\left( {\sf z}_{1,}{\sf z}_{2},%
{\sf z}_{3},{\sf z}_{4}\right) +\Phi _{K_{OD}}\left( {\sf z}_{1,}{\sf z}_{2},%
{\sf z}_{3},{\sf z}_{4}\right)  \nonumber \\
\Phi _{K_{OD}}\left( {\sf z}_{1,}{\sf z}_{2},{\sf z}_{3},{\sf z}_{4}\right)
&=&\Lambda _{_{K_{S}}}\left( {\sf z}_{1},{\sf z}_{2}\right) \delta \left( 
{\sf z}_{1}-{\sf z}_{3}\right) \delta \left( {\sf z}_{2}-{\sf z}_{4}\right)
\left( 1-\delta \left( {\sf z}_{1}-{\sf z}_{2}\right) \right)  \label{tden2}
\end{eqnarray}%
(notice the subscript $K_{S\text{ }}$ in the definition of $\Lambda
_{_{K_{S}}}$) and 
\begin{equation}
T_{S_{D}{\rm \ }}\left( \phi _{K_{S}}\right) \left( {\sf z}_{1,}{\sf z}%
_{2}\right) =P_{K_{S}}\left( {\sf z}_{1,}{\sf z}_{2}\right) +\phi
_{K_{S}}\left( {\sf z}_{1}\right) \delta \left( {\sf z}_{1}-{\sf z}%
_{2}\right)  \label{tden3}
\end{equation}%
If $\left| \breve{Y}_{ref}\right) $ is a pure state then $\hat{T}%
_{den}\left( \phi _{S_{D}}\right) \left| \breve{Y}_{ref}\right) $ is also a
pure state different from $\left| \breve{Y}_{ref}\right) $. The set $\left\{ 
\hat{T}_{den}\left( \phi _{S_{D}}\right) \breve{Y}_{ref};\quad \forall \phi
_{S_{D}},\text{ fixed }\omega \right\} $ does not produce all pure states 
{\it but it does contain a unique pure state representative for each
equivalence class. }{\bf In the vector bundle classification of }${\cal B}%
_{1}\left( {\cal H}^{N}\right) $ {\bf these different} {\bf paths of pure
state represenatives provide different representations of the base of the
vector bundle }${\frak B}_{vect}$.

\section{Summary}

We now summarize the salient features of this document and relate them to
ideas introduced in the main body of paper I. The group of transformation, $%
{\frak G}_{\Xi }$, that preserves the vector bundle structure, ${\frak B}%
_{vect}$, of ${\cal B}_{1}\left( {\cal H}^{N}\right) $ is composed of maps
that have the decomposition 
\begin{equation}
{\frak \hat{T}}{\frak =}{\frak \hat{T}}_{K_{OD}}+{\frak \hat{T}}_{K_{D}}+%
{\frak \hat{T}}_{S_{D}}+{\frak \hat{T}}_{K_{D}S_{D}}+{\frak \hat{T}}%
_{K_{OD}S_{D}}+{\frak \hat{T}}_{K_{OD}K_{D}}+{\frak \hat{T}}_{K_{D}K_{OD}}
\end{equation}%
and refer to the direct decomposition of ${\cal B}_{1}\left( {\cal H}%
^{N}\right) $%
\begin{equation}
{\cal B}_{1}\left( {\cal H}^{N}\right) =\ker \Xi \oplus S_{D}=K_{OD}\oplus
K_{D}\oplus S_{D}
\end{equation}%
The subgroup $GL\left( \ker \Xi \right) $ maps each equivalence class to
itself and are thus intra fiber transformations, while the subgroup $%
GL\left( S_{D}\right) $, {\it in general},{\it \ }maps between equivalence
classes, but contains the subgroup, ${\frak G}_{den}$, of inter fiber
transformations. These subgroups can be used to give non redundant
parameterizations of $\ker \Xi $ and $S_{D}\equiv $ $\ker \Xi ^{c}$ by
choosing a reference state $\breve{Y}_{ref}\in {\cal B}_{1}\left( {\cal H}%
^{N}\right) $, where $\Xi \left( \breve{Y}_{ref}\right) =f_{ref}$, which can
be decomposed as 
\begin{eqnarray}
Y_{ref} &=&Y_{ref\,K_{OD}}+Y_{ref\,K_{D}}+\breve{\Gamma}\left( f_{ref}\right)
\\
\left| \breve{\Gamma}\left( f_{ref}\right) \right) &=&\hat{P}_{S_{D}}\left| 
\breve{Y}_{ref}\right)  \label{ref}
\end{eqnarray}%
and then forming left cosets, ${\frak Co}_{\ker \Xi }$ and ${\frak Co}%
_{S_{D}}$ of $GL\left( \ker \Xi \right) $ and $GL\left( S_{D}\right) $ with
respect to the isotropy subgroups of $\left| \breve{Y}_{ref\,K_{D}}\right) $
and $\left| \breve{\Gamma}\left( f_{ref}\right) \right) $ respectively in
the following manner. 
\begin{eqnarray}
{\frak G}_{Y_{ref\,d}} &=&\left\{ {\frak \hat{T};\hat{T}}\left| \breve{Y}%
_{ref\,K_{D}}\right) =\left| \breve{Y}_{ref\,K_{D}}\right) \right\} 
\nonumber \\
{\frak Co}_{\ker \Xi } &=&GL\left( \ker \Xi \right) /{\frak G}_{\breve{Y}%
_{ref\,K_{D}}} \\
{\frak G}_{\Gamma \left( f_{ref}\right) } &=&\left\{ {\frak \hat{T};\hat{T}}%
\left| \breve{\Gamma}\left( f_{ref}\right) \right) =\left| \breve{\Gamma}%
\left( f_{ref}\right) \right) \right\}  \nonumber \\
{\frak Co}_{S_{D}} &=&GL\left( S_{D}\right) /{\frak G}_{\Gamma \left(
f_{ref}\right) }
\end{eqnarray}%
One can see that the parameterization of $\ker \Xi $ given by ${\frak C}%
_{\ker \Xi }$, which does not depend on $f$ , (but it does implicitly if a
reference $\breve{Y}_{ref}$ is used as in Eq.$\left( \ref{ref}\right) $
rather than a reference $\breve{Y}_{ref\,K_{D}}$ that is totally contained
in $\ker \Xi $), and is provided by Eq.$\left( \ref{parker}\right) $.

\begin{itemize}
\item If $\phi \ $defines a vector bundle map then $\int P_{S_{D}}\left( 
{\sf z}_{1},{\sf z}_{1}^{\prime }\right) \phi \left( {\sf z}_{1}^{\prime
}\right) P_{K_{D}}\left( {\sf z}_{1}^{\prime },{\sf z}_{2}\right) d{\sf z}%
_{1}^{\prime }=0\quad \forall {\sf z}_{1},{\sf z}_{2}$

If $\phi \bot _{s}Y\quad \forall Y\in {\cal B}_{1}\left( {\cal H}^{N}\right) 
$ then $\Phi _{S_{D}}\left( {\sf z}_{1},{\sf z}_{2}\right) =\int
P_{S_{D}}\left( {\sf z}_{1},{\sf z}_{1}^{\prime }\right) \phi \left( {\sf z}%
_{1}^{\prime }\right) P_{S_{D}}\left( {\sf z}_{1}^{\prime },{\sf z}%
_{2}\right) d{\sf z}_{1}^{\prime }=0\quad \forall {\sf z}_{1},{\sf z}_{2}$

If $\phi \bot _{s}Y\quad \forall Y\in K_{D}$ then $\int P_{S_{D}}\left( {\sf %
z}_{1},{\sf z}_{1}^{\prime }\right) \phi \left( {\sf z}_{1}^{\prime }\right)
P_{K_{D}}\left( {\sf z}_{1}^{\prime },{\sf z}_{2}\right) d{\sf z}%
_{1}^{\prime }=0\quad \forall {\sf z}_{1},{\sf z}_{2}\Longleftrightarrow $ $%
\phi \ $defines a vector bundle map

If $\phi \not{\bot}_{s}Y$ then $\Phi _{S_{D}}\left( {\sf z}_{1},{\sf z}%
_{2}\right) =\int P_{S_{D}}\left( {\sf z}_{1},{\sf z}_{1}^{\prime }\right)
\phi \left( {\sf z}_{1}^{\prime }\right) P_{S_{D}}\left( {\sf z}_{1}^{\prime
},{\sf z}_{2}\right) d{\sf z}_{1}^{\prime }\neq 0\quad \forall {\sf z}_{1},%
{\sf z}_{2}$

If $\phi $ defines an element of ${\frak G}_{fib}$ then $\Phi _{S_{D}}\left( 
{\sf z}_{1},{\sf z}_{2}\right) =\int P_{S_{D}}\left( {\sf z}_{1},{\sf z}%
_{1}^{\prime }\right) \phi \left( {\sf z}_{1}^{\prime }\right)
P_{S_{D}}\left( {\sf z}_{1}^{\prime },{\sf z}_{2}\right) d{\sf z}%
_{1}^{\prime }=0\quad \forall {\sf z}_{1},{\sf z}_{2}$ and $\int
P_{K_{D}}\left( {\sf z}_{1},{\sf z}_{1}^{\prime }\right) \phi \left( {\sf z}%
_{1}^{\prime }\right) P_{S_{D}}\left( {\sf z}_{1}^{\prime },{\sf z}%
_{2}\right) d{\sf z}_{1}^{\prime }=0\quad \forall {\sf z}_{1},{\sf z}_{2}$

If $\phi $ defines an element of ${\frak G}_{den}$ then $\int
P_{K_{D}}\left( {\sf z}_{1},{\sf z}_{1}^{\prime }\right) \phi \left( {\sf z}%
_{1}^{\prime }\right) P_{S_{D}}\left( {\sf z}_{1}^{\prime },{\sf z}%
_{2}\right) d{\sf z}_{1}^{\prime }=0\quad \forall {\sf z}_{1},{\sf z}_{2}$
and $\Phi _{K_{D}}\left( {\sf z}_{1},{\sf z}_{2}\right) =\int
P_{K_{D}}\left( {\sf z}_{1},{\sf z}_{1}^{\prime }\right) \phi \left( {\sf z}%
_{1}^{\prime }\right) P_{K_{D}}\left( {\sf z}_{1}^{\prime },{\sf z}%
_{2}\right) d{\sf z}_{1}^{\prime }=0\quad \forall {\sf z}_{1},{\sf z}_{2}$

\item The subgroup ${\frak G}_{\ker }\subset {\frak G}_{\Xi }$ introduced in
the main section of paper I is formed by super multiplication operators
generated by congruence transformations composed of operators $\in {\cal B}%
_{1}\left( {\cal H}^{N}\right) $ that are diagonal in the coordinate
representation. The operators $\hat{T}\left( \phi \right) $ belonging to the
representation of ${\frak G}_{\ker }$ describe in the preceding have the
properties 
\begin{eqnarray}
\hat{T}\left( \phi ,\omega \right) &:&K_{OD}\rightarrow K_{OD}  \nonumber \\
\hat{T}\left( \phi ,\omega \right) &:&K_{D}\rightarrow K_{D}  \nonumber \\
\hat{T}\left( \phi ,\omega \right) &:&S_{D}\rightarrow K_{D}  \nonumber \\
\hat{T}\left( \phi ,\omega \right) &:&S_{D}\rightarrow S_{D}
\end{eqnarray}%
The subgroup that acts homogeneously on pure states is ${\frak G}_{pure}=$ $%
{\frak G}_{fib}\times {\frak G}_{den}\subset {\frak G}_{\ker }$. The
subgroup ${\frak G}_{\ker }$ can be factored as ${\frak G}_{\ker }={\frak Co}%
_{kfd}\times {\frak G}_{pure}$ where the coset ${\frak Co}_{kfd}{\frak =G}%
_{\ker }/\left( {\frak G}_{fib}\times {\frak G}_{den}\right) $ and contains
maps from $S_{D}\rightarrow K_{D}$. The group ${\frak G}_{\ker }$ can also
be factorized into products of positive and unitary operators as 
\begin{equation}
{\frak G}_{\ker }={\cal G}_{\ker }\times {\cal G}_{\Xi }
\end{equation}

\item The subgroup ${\frak G}_{den}$ can be represented by operators of the
form 
\begin{equation}
{\frak G}_{den}=\left( 
\begin{array}{ccc}
\hat{t}_{K_{OD}}\left( \hat{T}_{S_{D}}\right) & 0 & 0 \\ 
0 & \hat{t}_{K_{D}}\left( \hat{T}_{S_{D}}\right) & 0 \\ 
0 & 0 & \hat{T}_{S_{D}}%
\end{array}%
\right)
\end{equation}%
where $\hat{t}_{K_{OD}}$ and $\hat{t}_{K_{D}}$ are representations of $%
{\frak G}_{den}$. The representation $\hat{t}_{K_{OD}}$ is given by%
\begin{equation}
\hat{t}_{K_{OD}}\left( \hat{T}_{S_{D}}\right) \left( \breve{Y}%
_{K_{OD}}\right) =e^{{\cal \breve{V}+}i\breve{\omega}}\breve{Y}_{K_{OD}}e^{%
{\cal \breve{V}-}i\breve{\omega}}
\end{equation}%
for ${\cal V=V}\left( \phi _{S_{D}},\phi _{K_{D}0}\right) $ where $\phi
_{K_{D}0}$ and $\omega $ are fixed. Note that $\hat{T}_{S_{D}}$ is not a
super multiplication operator unless $\hat{t}_{K_{D}}$ is the identity
representation i.e. $\phi _{K_{D}0}=0$. The subgroup ${\frak G}_{fib}$ can
be represented by 
\begin{equation}
{\frak G}_{fib}=\left( 
\begin{array}{ccc}
\hat{t}_{K_{OD}}^{\prime }\left( \hat{T}_{K_{D}}\right) & 0 & 0 \\ 
0 & \hat{T}_{K_{D}} & 0 \\ 
0 & 0 & \hat{P}_{K_{S}}%
\end{array}%
\right)
\end{equation}%
where $\hat{t}_{K_{OD}}^{\prime }$ and $\hat{t}_{S_{D}}^{\prime }$ are
representations of ${\frak G}_{fib}$. The representation $\hat{t}%
_{K_{OD}}^{\prime }$ is given by 
\begin{equation}
\hat{t}_{K_{OD}}^{\prime }\left( \hat{T}_{K_{D}}\right) \left(
Y_{K_{OD}}\right) =e^{{\cal \breve{V}+}i\breve{\omega}}Y_{K_{OD}}e^{{\cal 
\breve{V}-}i\breve{\omega}}
\end{equation}%
for ${\cal V=V}\left( \phi _{K_{D}}\right) $ where $\omega $ is fixed.

\item Restricting attention to the subgroup ${\frak G}_{fib}$ of super
multiplication operators that transform equivalence classes to themselves
i.e. 
\begin{equation}
\hat{T}_{fib}:\left( \left| \breve{k}\right) ,\left| \Gamma \left( f\right)
\right) \right) \rightarrow \left( \left| \breve{k}^{\prime }\right) ,\left|
\Gamma \left( f\right) \right) \right) ;\quad \hat{T}_{fib}\in {\frak G}%
_{fib}
\end{equation}%
contained in $GL\left( \ker \Xi \right) $, leads to the non redundant
parameterization of pure states belonging to an equivalence class as
described in preceding. The parameterization for this action of ${\frak G}%
_{fib}$ on a reference element is given in Eq.$\left( \ref{refkd}\right) $
and the discussion preceding this equation, where the form of function, $%
\phi _{refK_{D}}$, generating the super operator transformation $\hat{T}%
_{fib}\left( \phi _{refK_{D}}\right) $ is described. The unitary subgroup $%
{\cal G}_{\Xi }$ is contained in ${\frak G}_{fib}$ and we have used the
factorization ${\frak G}_{fib}={\cal G}_{fib}\times {\cal G}_{\Xi }$
extensively in section V of paper I. The function $\phi _{refK_{D}}$ is
strongly orthogonal to $Y$, the integral kernel of any pure state, as by
construction $\hat{P}_{S_{D}}\hat{\Phi}_{refK_{D}}\hat{P}_{S_{D}}=0$ (see
Eqs.$\left( \ref{diagker}\right) $,$\left( \ref{kerproj}\right) $and $\left( %
\ref{PhiKD}\right) )$ thus has no affect on $\hat{P}_{S_{D}}Y$ for any $Y$,
while $\hat{\Phi}_{refK_{D}}:K_{D}\rightarrow K_{D}$ and $\hat{\Phi}%
_{refK_{OD}}:K_{OD}\rightarrow K_{OD}$, which means that their integral
kernels are automatically strongly orthogonal to $P_{K_{OD}}Y$ and $%
P_{K_{D}}Y$. {\it Note }that it is possible for a more general $\phi $ (that
has non zero components, $P_{S_{D}}\phi )$, to be strongly orthogonal to a
specific $P_{S_{D}}Y_{ref}$ but not to {\it all }$P_{S_{D}}Y$.

\item Restricting attention to the subgroup ${\frak G}_{den}$ of super
multiplication operators contained in $GL\left( S_{D}\right) $, represented
in the form of Eq.$\left( \ref{tden}\right) $,\ leads to the non redundant
parameterization of a path pure states corresponding to different densities
in a 1-1 fashion, (see Eqs.$\left( \ref{tden}\right) $,$\left( \ref{tden2}%
\right) $and $\left( \ref{tden3}\right) )$.
\end{itemize}

\subsection{Example\protect\bigskip}

As an example we consider the pure Independent Particle State $\breve{Y}%
_{ref}\left( \gamma _{0}\right) =\left| \Psi _{0\,IPS}\right\rangle
\left\langle \Psi _{0\,IPS}\right| $, where $\Xi \left( \breve{Y}%
_{ref}\left( \gamma _{0}\right) \right) =\gamma _{0}$, which has the
decomposition 
\begin{equation}
\breve{Y}_{ref}=\breve{Y}_{ref\,K_{OD}}+\breve{Y}_{ref\,K_{D}}+\breve{\Gamma}%
\left( \gamma _{0}\right)
\end{equation}%
{\bf All} pure states with the same density as this reference state are
generated by 
\begin{equation}
\left| \breve{Y}\left( \gamma _{0}\right) \right) =\hat{T}_{fib}\left( \phi
_{K_{D}}{}_{ref}\right) \left| \breve{Y}_{ref}\left( \gamma _{0}\right)
\right) ;\quad \hat{T}_{fib}\left( \phi _{K_{D}}{}_{ref}\right) \in {\frak G}%
_{fib}
\end{equation}%
A {\it full }path of path pure states corresponding to {\bf all} different
densities in a 1-1 fashion is generated by 
\begin{equation}
\left| \breve{Y}_{ref}\left( \gamma \right) \right) =\hat{T}_{den}\left(
\phi _{\Gamma \left( \gamma _{0}\right) }\right) \left| \breve{Y}%
_{ref}\left( \gamma _{0}\right) \right) ;\,\gamma \neq \gamma ^{\prime
},\quad \hat{T}_{den}\left( \phi _{\Gamma \left( \gamma _{0}\right) }\right)
\in {\frak G}_{den}
\end{equation}%
{\bf All }pure states with the density $\gamma $ are generated from $\left| 
\breve{Y}_{ref}\left( \gamma \right) \right) $ by%
\begin{equation}
\left| \breve{Y}\left( \gamma \right) \right) =\hat{T}_{fib}\left( \phi
_{K_{D}}{}_{ref\left( \gamma \right) }\right) \left| \breve{Y}_{ref}\left(
\gamma \right) \right) ;\quad \hat{T}_{fib}\left( \phi _{K_{D}}{}_{ref\left(
\gamma \right) }\right) \in {\frak G}_{fib}
\end{equation}%
A full path can be produced by finding intra fiber transformations $\hat{T}%
_{fib\,IPS}\left( \phi _{S_{D}}\left( \gamma \right) \right) $, fixed for
each fiber, that produce an IPS\ from $\left| \breve{Y}\left( \gamma \right)
\right) $ 
\begin{equation}
\left| \breve{Y}_{IPS}\left( \gamma \right) \right) =\hat{T}%
_{fib\,IPS}\left( \phi _{S_{D}}\left( \gamma \right) \right) \left| \breve{Y}%
\left( \gamma \right) \right) ;\quad \hat{T}_{fib\,IPS}\left( \phi
_{S_{D}}\left( \gamma \right) \right) \in {\frak G}_{fib}
\end{equation}%
The transformations $\hat{T}_{fib\,IPS}\left( \left( \phi _{S_{D}}\left(
\gamma \right) \right) \right) $ are not unique as in general there many
IPS\ in any given fiber/equivalence class. Thus one can produce many full
paths of IPS.

\begin{references}
\bibitem{ideal} A two sided ideal of ${\cal B}\left( {\cal H}^{N}\right) $
is a set ${\cal Y\subseteq B}\left( {\cal H}^{N}\right) $ such that $YX\in 
{\cal B}\left( {\cal H}^{N}\right) $ and $XY\in {\cal B}\left( {\cal H}%
^{N}\right) \forall X\in {\cal B}\left( {\cal H}^{N}\right) $ and $\forall
Y\in {\cal Y}$.

\bibitem{homogen} A groups of transformations acts homogeneously on a space
if fro every $Y_{1}$ and $Y_{2}$ that belong to the space there exists a
group element $g$ such that $Y_{1}=gY_{2}$.
\end{references}
}

\QSubDoc{Include DualPair}{%TCIDATA{Version=5.50.0.2953}
%TCIDATA{LaTeXparent=0,0,EOMN.tex}
                      
%TCIDATA{ChildDefaults=chapter:1,page:1}


\section{Mappings between pairs of Dual Spaces}

Consider a {\em compact} map 
\begin{equation}
A:V\longrightarrow W
\end{equation}%
and its adjoint  
\begin{equation}
A^{\dagger }:W^{d}\longrightarrow V^{d}
\end{equation}%
where $\left\langle V^{d},V\right\rangle $ and $\left\langle
W^{d},W\right\rangle $ are dual space pairs and 
\begin{equation}
\left\langle A^{\dagger }{\bf w}^{d}|{\bf v}\right\rangle =\left\langle {\bf %
w}^{d}|A{\bf v}\right\rangle \quad \forall {\bf v\in }V
\end{equation}%
We then have the following diagram 
\begin{equation}
\frame{$%
\begin{array}{ccc}
V & \stackrel{A}{%
%TCIMACRO{\underset{A^{-1}}{{\LARGE \rightleftarrows }}}%
%BeginExpansion
\mathrel{\mathop{{\LARGE \rightleftarrows }}\limits_{A^{-1}}}%
%EndExpansion
} & W \\ 
&  &  \\ 
V^{d} & \stackrel{A^{\dagger }}{%
%TCIMACRO{\underset{A^{-1\dagger }}{{\LARGE \rightleftarrows }}}%
%BeginExpansion
\mathrel{\mathop{{\LARGE \rightleftarrows }}\limits_{A^{-1\dagger }}}%
%EndExpansion
} & W^{d}%
\end{array}%
$}
\end{equation}%
As the maps are compact their domain and range are closed subspaces and we
can {\it wlog }take 
\begin{eqnarray}
V &=&%
%TCIMACRO{\func{domain}}%
%BeginExpansion
\mathop{\rm domain}%
%EndExpansion
\left\{ A\right\}   \nonumber \\
W^{d} &=&%
%TCIMACRO{\func{domain}}%
%BeginExpansion
\mathop{\rm domain}%
%EndExpansion
\left\{ A^{\dagger }\right\} 
\end{eqnarray}%
Further we can embed $V\hookrightarrow $ $V^{dd}$ by ${\bf v\rightarrow \hat{%
v}}$ 
\begin{equation}
\left\langle {\bf \hat{v}}|{\bf v}^{d}\right\rangle =\left\langle {\bf v}%
^{d}|{\bf v}\right\rangle \quad \forall {\bf v}^{d}\in V^{d}\text{ }
\end{equation}%
so that $V^{dd}\supseteq V$ with equality always occurring for finite
dimensional spaces. The same holds true for the $W$ spaces.

Let

\begin{itemize}
\item $\left\{ \left| k_{j}\right\rangle ;1\leq j\leq m_{1}\right\} $ be a
basis for $\ker A$ and $\left\{ \left\langle k_{j}^{d}\right| ;1\leq j\leq
m_{1}\right\} $ be a biothorgonal dual basis of $\left\{ \ker A\right\} ^{d}$

\item $\left\{ \left| s_{j}\right\rangle ;1\leq j\leq m_{2}\right\} $ be a
basis for $\left\{ \ker A\right\} ^{c}$ and $\left\{ \left\langle
s_{j}^{d}\right| ;1\leq j\leq m_{2}\right\} $ be a biothorgonal dual basis
of $\left\{ \ker A\right\} ^{cd}$

\item $\left\{ \left| \tilde{s}_{j}\right\rangle ;1\leq j\leq m_{2}\right\} $
be a basis for the $%
%TCIMACRO{\func{range}}%
%BeginExpansion
\mathop{\rm range}%
%EndExpansion
A$ and $\left\{ \left\langle \tilde{s}_{j}^{d}\right| ;1\leq j\leq
m_{2}\right\} $ be a biothorgonal dual basis of $\left\{ 
%TCIMACRO{\func{range}}%
%BeginExpansion
\mathop{\rm range}%
%EndExpansion
A\right\} ^{d}$

\item $\left\{ \left| l_{j}\right\rangle ;1\leq j\leq m_{3}\right\} $ be a
basis for $\left\{ 
%TCIMACRO{\func{range}}%
%BeginExpansion
\mathop{\rm range}%
%EndExpansion
A\right\} ^{c}$ and $\left\{ \left\langle l_{j}^{d}\right| ;1\leq j\leq
m_{3}\right\} $ be a biothorgonal dual basis of $\left\{ 
%TCIMACRO{\func{range}}%
%BeginExpansion
\mathop{\rm range}%
%EndExpansion
A\right\} ^{cd}$.
\end{itemize}

Note that 
\begin{eqnarray}
V &=&\ker A\oplus \left\{ \ker A\right\} ^{c}  \nonumber \\
V^{d} &=&\left\{ \ker A\right\} ^{d}\oplus \left\{ \ker A\right\} ^{cd}
\end{eqnarray}

and 
\begin{eqnarray}
W &=&%
%TCIMACRO{\func{range}}%
%BeginExpansion
\mathop{\rm range}%
%EndExpansion
A\oplus \left\{ 
%TCIMACRO{\func{range}}%
%BeginExpansion
\mathop{\rm range}%
%EndExpansion
A\right\} ^{c}  \nonumber \\
W^{d} &=&\left\{ 
%TCIMACRO{\func{range}}%
%BeginExpansion
\mathop{\rm range}%
%EndExpansion
A\right\} ^{d}\oplus \left\{ 
%TCIMACRO{\func{range}}%
%BeginExpansion
\mathop{\rm range}%
%EndExpansion
A\right\} ^{cd}
\end{eqnarray}%
Thus 
\begin{eqnarray}
\dim V &=&m_{1}+m_{2}  \nonumber \\
\dim W &=&m_{2}+m_{3}  \nonumber \\
\dim V-\dim W &=&m_{1}-m_{3}
\end{eqnarray}%
One can observe that $\dim \left\{ 
%TCIMACRO{\func{range}}%
%BeginExpansion
\mathop{\rm range}%
%EndExpansion
A\right\} =\dim \left\{ \ker A\right\} ^{c}=m_{2}$. Associated with these
spaces are the projectors 
\begin{eqnarray}
P_{\ker A} &=&\sum_{1\leq j\leq m_{1}}\left| k_{j}\right\rangle \left\langle
k_{j}^{d}\right|  \nonumber \\
P_{\left\{ \ker A\right\} ^{c}} &=&\sum_{1\leq j\leq m_{2}}\left|
s_{j}\right\rangle \left\langle s_{j}^{d}\right|
\end{eqnarray}%
which have the properties%
\begin{eqnarray}
P_{\ker A}^{2} &=&P_{\ker A}  \nonumber \\
P_{\left\{ \ker A\right\} ^{c}}^{2} &=&P_{\left\{ \ker A\right\} ^{c}} 
\nonumber \\
P_{\left\{ \ker A\right\} ^{c}}P_{\ker A} &=&P_{\ker A}P_{\left\{ \ker
A\right\} ^{c}}=0  \nonumber \\
I_{V} &=&P_{\ker A}+P_{\left\{ \ker A\right\} ^{c}}
\end{eqnarray}%
and 
\begin{eqnarray}
P_{%
%TCIMACRO{\func{range}}%
%BeginExpansion
\mathop{\rm range}%
%EndExpansion
A} &=&\sum_{1\leq j\leq m_{2}}\left| \tilde{s}_{j}\right\rangle \left\langle 
\tilde{s}_{j}^{d}\right|  \nonumber \\
P_{\left\{ 
%TCIMACRO{\func{range}}%
%BeginExpansion
\mathop{\rm range}%
%EndExpansion
A\right\} ^{c}} &=&\sum_{1\leq j\leq m_{3}}\left| l_{j}\right\rangle
\left\langle l_{j}^{d}\right|
\end{eqnarray}%
which have the properties%
\begin{eqnarray}
P_{%
%TCIMACRO{\func{range}}%
%BeginExpansion
\mathop{\rm range}%
%EndExpansion
A}^{2} &=&P_{%
%TCIMACRO{\func{range}}%
%BeginExpansion
\mathop{\rm range}%
%EndExpansion
A}  \nonumber \\
P_{\left\{ 
%TCIMACRO{\func{range}}%
%BeginExpansion
\mathop{\rm range}%
%EndExpansion
A\right\} ^{c}}^{2} &=&P_{\left\{ 
%TCIMACRO{\func{range}}%
%BeginExpansion
\mathop{\rm range}%
%EndExpansion
A\right\} ^{c}}  \nonumber \\
P_{\left\{ 
%TCIMACRO{\func{range}}%
%BeginExpansion
\mathop{\rm range}%
%EndExpansion
A\right\} ^{c}}P_{%
%TCIMACRO{\func{range}}%
%BeginExpansion
\mathop{\rm range}%
%EndExpansion
A} &=&P_{%
%TCIMACRO{\func{range}}%
%BeginExpansion
\mathop{\rm range}%
%EndExpansion
A}P_{\left\{ 
%TCIMACRO{\func{range}}%
%BeginExpansion
\mathop{\rm range}%
%EndExpansion
A\right\} ^{c}}=0  \nonumber \\
I_{W} &=&P_{%
%TCIMACRO{\func{range}}%
%BeginExpansion
\mathop{\rm range}%
%EndExpansion
A}+P_{\left\{ 
%TCIMACRO{\func{range}}%
%BeginExpansion
\mathop{\rm range}%
%EndExpansion
A\right\} ^{c}}
\end{eqnarray}

If $A$ is compact and these bases are canonical for $A$ we obtain the
following expansions for $A$ 
\begin{equation}
A=\sum_{1\leq j\leq m_{2}}\alpha _{j}\left\vert \tilde{s}_{j}\right\rangle
\left\langle s_{j}^{d}\right\vert ;\quad \alpha _{j}>0
\end{equation}%
It is easy to see that $A$ has the properties 
\begin{eqnarray}
AP_{\left\{ \ker A\right\} ^{c}} &=&A  \nonumber \\
P_{%
%TCIMACRO{\func{range}}%
%BeginExpansion
\mathop{\rm range}%
%EndExpansion
A}A &=&A  \nonumber \\
P_{%
%TCIMACRO{\func{range}}%
%BeginExpansion
\mathop{\rm range}%
%EndExpansion
A}AP_{\left\{ \ker A\right\} ^{c}} &=&A
\end{eqnarray}%
By utilizing the canonical expansions of $A$ we can define a unique two
sided generalized inverse 
\begin{equation}
A^{-1}=\sum_{1\leq j\leq m_{2}}\alpha _{j}^{-1}\left\vert s_{j}\right\rangle
\left\langle \tilde{s}_{j}^{d}\right\vert  \label{Ainv}
\end{equation}%
which has the property 
\begin{equation}
AA^{-1}=\sum_{1\leq j\leq m_{2}}\left\vert \tilde{s}_{j}\right\rangle
\left\langle \tilde{s}_{j}^{d}\right\vert =P_{%
%TCIMACRO{\func{range}}%
%BeginExpansion
\mathop{\rm range}%
%EndExpansion
A}
\end{equation}%
and 
\begin{equation}
A^{-1}A=\sum_{1\leq j\leq m_{2}}\left\vert s_{j}\right\rangle \left\langle
s_{j}^{d}\right\vert =P_{\left\{ \ker A\right\} ^{c}}
\end{equation}%
and\ 
\begin{eqnarray}
A^{-1}P_{%
%TCIMACRO{\func{range}}%
%BeginExpansion
\mathop{\rm range}%
%EndExpansion
A} &=&A^{-1}  \nonumber \\
P_{\left\{ \ker A\right\} ^{c}}A^{-1} &=&A^{-1}  \nonumber \\
P_{\left\{ \ker A\right\} ^{c}}A^{-1}P_{%
%TCIMACRO{\func{range}}%
%BeginExpansion
\mathop{\rm range}%
%EndExpansion
A} &=&A^{-1}
\end{eqnarray}%
The definition in eq. $\left( \ref{Ainv}\right) $ can be used as a {\em core}
for left and right inverses of $A$ as is shown in the following. Let 
\begin{eqnarray}
R &:&\left\{ 
%TCIMACRO{\func{range}}%
%BeginExpansion
\mathop{\rm range}%
%EndExpansion
A\right\} ^{c}\longrightarrow \ker A  \nonumber \\
R &=&\sum R_{ij}\left\vert k_{i}\right\rangle \left\langle
l_{j}^{d}\right\vert
\end{eqnarray}%
and define 
\begin{equation}
A_{R}^{-1}=A^{-1}+R
\end{equation}%
then 
\begin{equation}
AA_{R}^{-1}=\sum_{1\leq j\leq m_{2}}\alpha _{j}\left\vert \tilde{s}%
_{j}\right\rangle \left\langle s_{j}^{d}\right\vert \left( \sum_{1\leq j\leq
m_{2}}\alpha _{j}^{-1}\left\vert s_{j}\right\rangle \left\langle \tilde{s}%
_{j}^{d}\right\vert +\sum_{\Sb 1\leq i\leq m_{1}  \\ 1\leq j\leq m_{3} 
\endSb }R_{ij}\left\vert k_{i}\right\rangle \left\langle
l_{j}^{d}\right\vert \right) =P_{%
%TCIMACRO{\func{range}}%
%BeginExpansion
\mathop{\rm range}%
%EndExpansion
A}
\end{equation}%
and 
\begin{eqnarray}
L &:&\ker A\longrightarrow \left\{ 
%TCIMACRO{\func{range}}%
%BeginExpansion
\mathop{\rm range}%
%EndExpansion
A\right\} ^{c}  \nonumber \\
L &=&\sum\Sb 1\leq i\leq m_{3}  \\ 1\leq j\leq m_{1}  \endSb %
L_{ij}\left\vert k_{i}\right\rangle \left\langle l_{j}^{d}\right\vert
\end{eqnarray}%
and define 
\begin{equation}
A_{L}^{-1}=A^{-1}+L
\end{equation}%
then 
\begin{equation}
A_{L}^{-1}A=\left( \sum_{1\leq j\leq m_{2}}\alpha _{j}^{-1}\left\vert
s_{j}\right\rangle \left\langle \tilde{s}_{j}^{d}\right\vert +\sum\Sb 1\leq
i\leq m_{3}  \\ 1\leq j\leq m_{1}  \endSb L_{ij}\left\vert
k_{i}\right\rangle \left\langle l_{j}^{d}\right\vert \right) \sum_{1\leq
j\leq m_{2}}\alpha _{j}\left\vert \tilde{s}_{j}\right\rangle \left\langle
s_{j}^{d}\right\vert =P_{\left\{ \ker A\right\} ^{c}}
\end{equation}%
The preceding relationships are an abstract way of defining left and right
inverses of $A$, based on a {\em unique} core generalized inverse.

The adjoint $A^{\dagger }$ of $A$ can be expanded as 
\begin{equation}
A^{\dagger }=\sum_{1\leq j\leq m_{2}}\alpha _{j}\left\vert \tilde{s}%
_{j}^{d}\right\rangle \left\langle \tilde{s}_{j}\right\vert ;\quad \alpha
_{j}>0
\end{equation}%
We define the orthogonal compliment of a subspace $U\subseteq V$ as a
subspace $U^{\bot }$ contained in the dual space $V^{d}$ as 
\begin{equation}
U^{\bot }=\left\{ {\bf u}^{d};\backepsilon \left\langle {\bf u}^{d}|{\bf u}%
\right\rangle =0\quad \forall {\bf u\in }U\right\}
\end{equation}%
and note that 
\begin{equation}
U^{\bot \bot }\supseteq U
\end{equation}%
Further one has 
\begin{equation}
U^{cd}=U^{\bot }
\end{equation}%
Finally we note a relationship between $\ker A$ and $%
%TCIMACRO{\func{range}}%
%BeginExpansion
\mathop{\rm range}%
%EndExpansion
A^{\dagger }$. If ${\bf v}\in \ker A$ then%
\begin{equation}
\left\langle {\bf w}^{d}|A{\bf v}\right\rangle =0\,\forall {\bf w}%
^{d}\Rightarrow \left\langle A^{\dagger }{\bf w}^{d}|{\bf v}\right\rangle
=0\,\forall {\bf w}^{d}\Rightarrow {\bf v}\in \left\{ 
%TCIMACRO{\func{range}}%
%BeginExpansion
\mathop{\rm range}%
%EndExpansion
A^{\dagger }\right\} ^{\bot }\supseteq \ker A
\end{equation}%
Similarly 
\begin{equation}
\left\langle A^{\dagger }{\bf w}^{d}|{\bf v}\right\rangle =0\,\forall {\bf v}%
\Rightarrow \left\langle {\bf w}^{d}|A{\bf v}\right\rangle =0\,\forall {\bf v%
}\Rightarrow {\bf w}^{d}\in \left\{ 
%TCIMACRO{\func{range}}%
%BeginExpansion
\mathop{\rm range}%
%EndExpansion
A\right\} ^{\bot }\supseteq \ker A^{\dagger }
\end{equation}%
$\lceil $Recall $\left\langle A^{\dagger }{\bf w}^{d}|{\bf v}\right\rangle
=\left\langle {\bf w}^{d}|A{\bf v}\right\rangle \quad \forall {\bf v\rfloor }%
\ \ $

In many applications we have the situation that $%
%TCIMACRO{\func{domain}}%
%BeginExpansion
\mathop{\rm domain}%
%EndExpansion
\left\{ A\right\} \supseteq 
%TCIMACRO{\func{range}}%
%BeginExpansion
\mathop{\rm range}%
%EndExpansion
\left\{ A\right\} $ (equality occurs if $A$ is $1-1$) so that 
\begin{eqnarray}
%TCIMACRO{\func{domain}}%
%BeginExpansion
\mathop{\rm domain}%
%EndExpansion
\left\{ A\right\} &=&\ker \left\{ A\right\} \oplus \ker \left\{ A\right\}
^{c}  \nonumber \\
%TCIMACRO{\func{range}}%
%BeginExpansion
\mathop{\rm range}%
%EndExpansion
\left\{ A\right\} &\cong &\ker \left\{ A\right\} ^{c}
\end{eqnarray}%
thus 
\begin{equation}
%TCIMACRO{\func{range}}%
%BeginExpansion
\mathop{\rm range}%
%EndExpansion
\left\{ A\right\} ^{c}\cong \ker \left\{ A\right\}
\end{equation}%
and 
\begin{equation}
%TCIMACRO{\func{domain}}%
%BeginExpansion
\mathop{\rm domain}%
%EndExpansion
\left\{ A\right\} =%
%TCIMACRO{\func{range}}%
%BeginExpansion
\mathop{\rm range}%
%EndExpansion
\left\{ A\right\} ^{c}\oplus 
%TCIMACRO{\func{range}}%
%BeginExpansion
\mathop{\rm range}%
%EndExpansion
\left\{ A\right\}
\end{equation}%
but note that $\ker \left\{ A\right\} \neq 
%TCIMACRO{\func{range}}%
%BeginExpansion
\mathop{\rm range}%
%EndExpansion
\left\{ A\right\} ^{c}$ unless $\ker \left\{ A\right\} ^{c}$ is an invariant
subspace of $A$, which it is if $A$ is self adjoint.

Another common simplification occurs when $V$ \ and $W$ are self dual i.e. $%
V=V^{d}$ and $W=W^{d}$ and are thus Hilbert spaces. When both of the
preceding simplifications occur the projectors become orthogonal projectors
and all the vectors belong to the same space. Though note that the dual
basis and the normal basis are still in general distinct, unless $A$ is self
adjoint.

\subsection{Trace}

\begin{equation}
A=\sum_{1\leq j\leq m_{2}}\alpha _{j}\left\vert \tilde{s}_{j}\right\rangle
\left\langle s_{j}^{d}\right\vert ;\quad \alpha _{j}>0
\end{equation}%
\begin{equation}
A^{\dagger }=\sum_{1\leq j\leq m_{2}}\alpha _{j}\left\vert \tilde{s}%
_{j}^{d}\right\rangle \left\langle \tilde{s}_{j}\right\vert ;\quad \alpha
_{j}>0
\end{equation}%
\begin{eqnarray*}
AA^{\dagger } &=&\sum_{1\leq j\leq m_{2}}\alpha _{j}\left\vert \tilde{s}%
_{j}\right\rangle \left\langle s_{j}^{d}\right\vert \sum_{1\leq j\leq
m_{2}}\alpha _{j}\left\vert \tilde{s}_{j}^{d}\right\rangle \left\langle 
\tilde{s}_{j}\right\vert = \\
A^{\dagger }A &=&\sum_{1\leq j\leq m_{2}}\alpha _{j}\left\vert \tilde{s}%
_{j}^{d}\right\rangle \left\langle \tilde{s}_{j}\right\vert \sum_{1\leq
j\leq m_{2}}\alpha _{j}\left\vert \tilde{s}_{j}\right\rangle \left\langle
s_{j}^{d}\right\vert =
\end{eqnarray*}

\subsection{Linear Equations}

We can now reexamine the linear equation 
\begin{equation}
A\left| {\bf v}\right\rangle =\left| {\bf w}\right\rangle {\bf ;\quad \left|
v\right\rangle \in V}^{s},\,\left| {\bf w}\right\rangle {\bf \in W}%
^{r};\quad r=m_{2}+m_{3};\quad s=m_{2}+m_{1}  \label{lineq}
\end{equation}%
This equation can be written as%
\begin{eqnarray}
&&\sum_{1\leq j\leq m_{2}}\alpha _{j}\left| \tilde{s}_{j}\right\rangle
\left\langle s_{j}^{d}\right| {\bf v}\rangle =\left| {\bf w}\right\rangle
\,\,{\bf \Rightarrow \,}P_{\left\{ \ker A\right\} ^{c}}{\bf \,}\left| {\bf v}%
\right\rangle =A_{L}^{-1}\left| {\bf w}\right\rangle  \nonumber \\
&\equiv &\left( \sum_{1\leq j\leq m_{2}}\alpha _{j}^{-1}\left|
s_{j}\right\rangle \left\langle \tilde{s}_{j}^{d}\right| +\sum\Sb 1\leq
i\leq m_{3}  \\ 1\leq j\leq m_{1}  \endSb L_{ij}\left| k_{i}\right\rangle
\left\langle l_{j}^{d}\right| \right) \sum_{1\leq j\leq m_{2}}\alpha
_{j}\left| \tilde{s}_{j}\right\rangle \left\langle s_{j}^{d}\right| {\bf v}%
\rangle =P_{\left\{ \ker A\right\} ^{c}}\left| {\bf v}\right\rangle 
\nonumber \\
&{\bf =}&\left( \sum_{1\leq j\leq m_{2}}\alpha _{j}^{-1}\left|
s_{j}\right\rangle \left\langle \tilde{s}_{j}^{d}\right| +\sum\Sb 1\leq
i\leq m_{3}  \\ 1\leq j\leq m_{1}  \endSb L_{ij}\left| k_{i}\right\rangle
\left\langle l_{j}^{d}\right| \right) \left| {\bf w}\right\rangle
\end{eqnarray}%
i.e. 
\begin{equation}
P_{\left\{ \ker A\right\} ^{c}}\left| {\bf v}\right\rangle =\sum_{1\leq
j\leq m_{2}}\alpha _{j}^{-1}\left| s_{j}\right\rangle \left\langle \tilde{s}%
_{j}^{d}\right| {\bf w\rangle +}\sum\Sb 1\leq i\leq m_{3}  \\ 1\leq j\leq
m_{1}  \endSb L_{ij}\left| k_{i}\right\rangle \left\langle l_{j}^{d}\right| 
{\bf w\rangle }
\end{equation}%
which can only be true if the consistency constraint 
\begin{equation}
\left\langle l_{j}^{d}\right| {\bf w}\rangle =0;\quad 1\leq i\leq
m_{3}\Longleftrightarrow P_{\left\{ 
%TCIMACRO{\func{range}}%
%BeginExpansion
\mathop{\rm range}%
%EndExpansion
A\right\} ^{c}}\left| {\bf w}\right\rangle =0
\end{equation}%
holds. The remaining part of the vector $\left| {\bf v}\right\rangle $, i.e. 
$P_{\left\{ \ker A\right\} }\left| {\bf v}\right\rangle $ is arbitrary
unless some further conditions are invoked. Thus the solution to eq. $\left( %
\ref{lineq}\right) $ is 
\begin{equation}
\left| {\bf v}\right\rangle =A^{-1}\left| {\bf w}\right\rangle +\left| {\bf k%
}\right\rangle ;\,\,\left| {\bf k}\right\rangle \in \ker A
\end{equation}%
as long as the condition 
\begin{equation}
P_{\left\{ 
%TCIMACRO{\func{range}}%
%BeginExpansion
\mathop{\rm range}%
%EndExpansion
A\right\} ^{c}}\left| {\bf w}\right\rangle =0
\end{equation}%
is satisfied.

\subsection{Matrix Representations}

Using the preceding biorthogonal basis we obtain 
\begin{eqnarray}
&&A\left| \left( {\bf s,k}\right) \right\rangle =\left| \left( {\bf \tilde{s}%
,l}\right) \right\rangle \left( 
\begin{array}{cc}
\begin{array}{ccc}
\alpha _{1} &  &  \\ 
& \ddots &  \\ 
&  & \alpha _{s}%
\end{array}
& 0 \\ 
0 & 0%
\end{array}%
\right) ;\quad \text{ }\left( m_{1}+m_{3}\right) \times \left(
m_{1}+m_{2}\right) \text{ matrix}  \nonumber \\
&{\Bbb \Rightarrow }&\left\langle \left( {\bf \tilde{s}}^{d},{\bf l}%
^{d}\right) |A\left( {\bf s,k}\right) \right\rangle =\left\langle \left( 
{\bf \tilde{s}}^{d},{\bf l}^{d}\right) |\left( {\bf s,k}\right)
\right\rangle \left( 
\begin{array}{cc}
\begin{array}{ccc}
\alpha _{1} &  &  \\ 
& \ddots &  \\ 
&  & \alpha _{s}%
\end{array}
& 0 \\ 
0 & 0%
\end{array}%
\right) =  \nonumber \\
&&\qquad \qquad \qquad \qquad \quad {\Bbb I}_{\left( m_{1}+m_{3}\right)
}\left( 
\begin{array}{cc}
\begin{array}{ccc}
\alpha _{1} &  &  \\ 
& \ddots &  \\ 
&  & \alpha _{s}%
\end{array}
& 0 \\ 
0 & 0%
\end{array}%
\right) =\left( 
\begin{array}{cc}
\begin{array}{ccc}
\alpha _{1} &  &  \\ 
& \ddots &  \\ 
&  & \alpha _{s}%
\end{array}
& 0 \\ 
0 & 0%
\end{array}%
\right)
\end{eqnarray}%
and 
\begin{equation}
A^{-1}\left( \left| {\bf \tilde{s}}\right\rangle {\bf ,}\left| {\bf l}%
\right\rangle \right) =\left( \left| {\bf s}\right\rangle {\bf ,}\left| {\bf %
k}\right\rangle \right) \left( 
\begin{array}{cc}
\begin{array}{ccc}
\alpha _{1}^{-1} &  &  \\ 
& \ddots &  \\ 
&  & \alpha _{s}^{-1}%
\end{array}
& 0 \\ 
0 & 0%
\end{array}%
\right) ;\quad \text{ }\left( m_{1}+m_{2}\right) \times \left(
m_{1}+m_{3}\right) \text{ matrix}
\end{equation}%
Hence 
\begin{equation}
\left( 
\begin{array}{cc}
\begin{array}{ccc}
\alpha _{1}^{-1} &  &  \\ 
& \ddots &  \\ 
&  & \alpha _{s}^{-1}%
\end{array}
& 0 \\ 
0 & 0%
\end{array}%
\right) \left( 
\begin{array}{cc}
\begin{array}{ccc}
\alpha _{1} &  &  \\ 
& \ddots &  \\ 
&  & \alpha _{s}%
\end{array}
& 0 \\ 
0 & 0%
\end{array}%
\right) =\left( 
\begin{array}{cc}
\begin{array}{ccc}
1 &  &  \\ 
& \ddots &  \\ 
&  & 1%
\end{array}
& 0 \\ 
0 & 0%
\end{array}%
\right) ={\Bbb P}_{\left\{ \ker A\right\} ^{c}}
\end{equation}%
which is $\left( m_{1}+m_{2}\right) \times \left( m_{1}+m_{2}\right) $
matrix, and 
\begin{equation}
\left( 
\begin{array}{cc}
\begin{array}{ccc}
\alpha _{1} &  &  \\ 
& \ddots &  \\ 
&  & \alpha _{s}%
\end{array}
& 0 \\ 
0 & 0%
\end{array}%
\right) \left( 
\begin{array}{cc}
\begin{array}{ccc}
\alpha _{1}^{-1} &  &  \\ 
& \ddots &  \\ 
&  & \alpha _{s}^{-1}%
\end{array}
& 0 \\ 
0 & 0%
\end{array}%
\right) =\left( 
\begin{array}{cc}
\begin{array}{ccc}
1 &  &  \\ 
& \ddots &  \\ 
&  & 1%
\end{array}
& 0 \\ 
0 & 0%
\end{array}%
\right) ={\Bbb P}_{%
%TCIMACRO{\func{range}}%
%BeginExpansion
\mathop{\rm range}%
%EndExpansion
A}
\end{equation}%
which is $\left( m_{1}+m_{3}\right) \times \left( m_{1}+m_{3}\right) $
matrix.

If we choose other bases related to the canonical basis of $A$ by 
\begin{eqnarray}
&&\qquad \quad \left| \left( {\bf s}^{\prime }{\bf ,k}^{\prime }\right)
\right\rangle =S\left| \left( {\bf s,k}\right) \right\rangle =\left| \left( 
{\bf s,k}\right) \right\rangle {\Bbb S}  \nonumber \\
{\Bbb S\,} &{\Bbb =}&\left( 
\begin{array}{cc}
{\Bbb S}_{11} & {\Bbb S}_{12} \\ 
{\Bbb S}_{21} & {\Bbb S}_{22}%
\end{array}%
\right) =\left( 
\begin{array}{cc}
\left\langle {\bf s}^{d}|S{\bf s}\right\rangle & \left\langle {\bf s}^{d}|S%
{\bf k}\right\rangle \\ 
\left\langle {\bf k}^{d}|S{\bf s}\right\rangle & \left\langle {\bf k}^{d}|S%
{\bf k}\right\rangle%
\end{array}%
\right)  \nonumber \\
&&\qquad \quad \,\left| \left( {\bf \tilde{s}}^{\prime }{\bf ,l}^{\prime
}\right) \right\rangle =T\left| \left( {\bf \tilde{s},l}\right)
\right\rangle =\left| \left( {\bf \tilde{s},l}\right) \right\rangle {\Bbb T}
\nonumber \\
{\Bbb T}{\Bbb \,} &{\Bbb =}&\left( 
\begin{array}{cc}
{\Bbb T}_{11} & {\Bbb T}_{12} \\ 
{\Bbb T}_{21} & {\Bbb T}_{22}%
\end{array}%
\right) =\left( 
\begin{array}{cc}
\left\langle {\bf \tilde{s}}^{d}|T{\bf \tilde{s}}\right\rangle & 
\left\langle {\bf \tilde{s}}^{d}|T{\bf l}\right\rangle \\ 
\left\langle {\bf l}^{d}|T{\bf \tilde{s}}\right\rangle & \left\langle {\bf l}%
^{d}|T{\bf l}\right\rangle%
\end{array}%
\right)  \nonumber \\
&&\qquad \quad \Rightarrow \left| \left( {\bf \tilde{s},l}\right)
\right\rangle =\left| \left( {\bf \tilde{s}}^{\prime }{\bf ,l}^{\prime
}\right) \right\rangle {\Bbb T}^{-1}
\end{eqnarray}%
we obtain%
\begin{equation}
A\left| \left( {\bf s}^{\prime }{\bf ,k}^{\prime }\right) \right\rangle
=\left| \left( {\bf \tilde{s}}^{\prime }{\bf ,l}^{\prime }\right)
\right\rangle {\Bbb A}
\end{equation}%
where ${\Bbb A}$ is given by 
\begin{equation}
A\left| \left( {\bf s}^{\prime }{\bf ,k}^{\prime }\right) \right\rangle
=A\left| \left( {\bf s,k}\right) \right\rangle {\Bbb S=\,}\left| \left( {\bf 
\tilde{s},l}\right) \right\rangle {\Bbb \alpha S=\,}\left| \left( {\bf 
\tilde{s}}^{\prime }{\bf ,l}^{\prime }\right) \right\rangle {\Bbb T}^{-1}%
{\Bbb \alpha S}
\end{equation}
\begin{eqnarray}
A\left| \left( {\bf s}^{\prime }{\bf ,k}^{\prime }\right) \right\rangle
&=&\left| \left( {\bf \tilde{s}}^{\prime }{\bf ,l}^{\prime }\right)
\right\rangle \left( 
\begin{array}{cc}
\left( {\Bbb T}^{-1}\right) _{11} & \left( {\Bbb T}^{-1}\right) _{12} \\ 
\left( {\Bbb T}^{-1}\right) _{21} & \left( {\Bbb T}^{-1}\right) _{22}%
\end{array}%
\right) \left( 
\begin{array}{cc}
\begin{array}{ccc}
\alpha _{1} &  &  \\ 
& \ddots &  \\ 
&  & \alpha _{s}%
\end{array}
& 0 \\ 
0 & 0%
\end{array}%
\right) \left( 
\begin{array}{cc}
{\Bbb S}_{11} & {\Bbb S}_{12} \\ 
{\Bbb S}_{21} & {\Bbb S}_{22}%
\end{array}%
\right)  \nonumber \\
&=&\left| \left( {\bf \tilde{s}}^{\prime }{\bf ,l}^{\prime }\right)
\right\rangle \left( 
\begin{array}{cc}
\left( {\Bbb T}^{-1}\right) _{11} & \left( {\Bbb T}^{-1}\right) _{12} \\ 
\left( {\Bbb T}^{-1}\right) _{21} & \left( {\Bbb T}^{-1}\right) _{22}%
\end{array}%
\right) \left( 
\begin{array}{cc}
\alpha {\Bbb S}_{11} & \alpha {\Bbb S}_{12} \\ 
0 & 0%
\end{array}%
\right)  \nonumber \\
&=&\left| \left( {\bf \tilde{s}}^{\prime }{\bf ,l}^{\prime }\right)
\right\rangle \left( 
\begin{array}{cc}
\left( {\Bbb T}^{-1}\right) _{11}\alpha {\Bbb S}_{11} & \left( {\Bbb T}%
^{-1}\right) _{11}\alpha {\Bbb S}_{12} \\ 
\left( {\Bbb T}^{-1}\right) _{21}\alpha {\Bbb S}_{11} & \left( {\Bbb T}%
^{-1}\right) _{21}\alpha {\Bbb S}_{12}%
\end{array}%
\right) =\left| \left( {\bf \tilde{s}}^{\prime }{\bf ,l}^{\prime }\right)
\right\rangle {\Bbb A}
\end{eqnarray}%
matrix

Let ${\Bbb A}$ be a $r\times s$ matrix then its left and right inverses have
the property

\begin{eqnarray}
{\Bbb A}_{L}^{-1}{\Bbb A} &=&{\Bbb P}_{\left\{ \ker A\right\} ^{c}}\leq 
{\Bbb I}_{s}  \nonumber \\
{\Bbb AA}_{R}^{-1} &=&{\Bbb P}_{%
%TCIMACRO{\func{range}}%
%BeginExpansion
\mathop{\rm range}%
%EndExpansion
A}\leq {\Bbb I}_{r}
\end{eqnarray}%
and can be defined as matrix representations of the abstract operators
defined previously using the core generalized inverse.

\section{Solution of Nonlinear Singular Equations.}

Consider the following equation for $x$ 
\begin{equation}
{\cal A}\left( x\right) \left| b\left( x\right) \right\rangle =\left|
c\left( x\right) \right\rangle ;\quad \left| b\left( x\right) \right\rangle
\in V;\,\left| c\left( x\right) \right\rangle \in W  \label{nl}
\end{equation}%
where ${\cal A}\left( x\right) $ is a compact linear operator and $\left|
b\left( x\right) \right\rangle $ and $\left| c\left( x\right) \right\rangle $
are vectors that depends on $x.$ Let ${\cal A}\left( x\right) $ be singular
with a kernel, $\ker {\cal A}\left( x\right) $, and have the expansion 
\begin{equation}
{\cal A}\left( x\right) =\sum_{1\leq j\leq m_{2}}\alpha _{j}\left| \tilde{s}%
_{j}\right\rangle \left\langle s_{j}^{d}\right| ;\quad \alpha _{j}>0
\end{equation}%
The equation $\left( \ref{nl}\right) $ only has solution if $\left| c\left(
x\right) \right\rangle $ is in the range of ${\cal A}\left( x\right) $,
which gives the condition 
\begin{equation}
P_{\left\{ 
%TCIMACRO{\func{range}}%
%BeginExpansion
\mathop{\rm range}%
%EndExpansion
{\cal A}\left( x\right) \right\} ^{c}}\left| c\left( x\right) \right\rangle
=0
\end{equation}%
If this condition is enforced one can form the special solution 
\begin{equation}
{\frak s}\left( x\right) =\sum_{1\leq j\leq m_{2}}\alpha _{j}^{-1}\left|
s_{j}\right\rangle \left\langle \tilde{s}_{j}^{d}\right| \left. c\left(
x\right) \right\rangle
\end{equation}%
by using the generalized inverse ${\cal A}\left( x\right) ^{-1}$ of ${\cal A}%
\left( x\right) $ defined as 
\begin{equation}
{\cal A}\left( x\right) ^{-1}=\sum_{1\leq j\leq m_{2}}\alpha _{j}^{-1}\left|
s_{j}\right\rangle \left\langle \tilde{s}_{j}^{d}\right|
\end{equation}%
If ${\frak s}\left( x\right) $ is such a solution then the general solution
is given by 
\begin{equation}
{\frak g}\left( x\right) ={\frak s}\left( x\right) +\sum_{1\leq i\leq
m_{1}}\lambda _{i}\left| k_{i}\right\rangle
\end{equation}%
as 
\begin{equation}
\sum_{1\leq j\leq m_{2}}\alpha _{j^{\prime }}\left| \tilde{s}_{j^{\prime
}}\right\rangle \left\langle s_{j^{\prime }}^{d}\right| \left\{ \sum_{1\leq
j\leq m_{2}}\alpha _{j}^{-1}\left| s_{j}\right\rangle \left\langle \tilde{s}%
_{j}^{d}\right| \left. c\left( x\right) \right\rangle +\sum_{1\leq i\leq
m_{1}}\lambda _{i}\left| k_{i}\right\rangle \right\} =\left| c\left(
x\right) \right\rangle
\end{equation}%
Thus to summarize the solutions of eq. $\left( \ref{nl}\right) $ are $x\ni $ 
\begin{eqnarray}
&&b\left( x\right) ={\frak s}\left( x\right) +\sum_{1\leq i\leq
m_{1}}\lambda _{i}\left( x\right) \left| k_{i}\left( x\right) \right\rangle 
\nonumber \\
&&{\frak s}\left( x\right) =\sum_{1\leq j\leq m_{2}}\alpha _{j}^{-1}\left(
x\right) \left| s_{j}\left( x\right) \right\rangle \left\langle \tilde{s}%
_{j}^{d}\left( x\right) \right| \left. c\left( x\right) \right\rangle ={\cal %
A}\left( x\right) ^{-1}\left| c\left( x\right) \right\rangle  \nonumber \\
&&P_{\left\{ 
%TCIMACRO{\func{range}}%
%BeginExpansion
\mathop{\rm range}%
%EndExpansion
{\cal A}\left( x\right) \right\} ^{c}}\left| c\left( x\right) \right\rangle
=0
\end{eqnarray}

\section{Antisymmetric Operators}

If the operator Let ${\cal A}$ be a compact antisymmetric operator $\in 
{\cal B}_{0}\left( V\right) $, where $V$ is a Hilbert space then it has the
canonical expansion 
\begin{equation}
{\cal A}=\sum_{1\leq i\leq s}\alpha _{i}\left\{ \left| s_{i}\right\rangle
\left\langle s_{i+s}^{d}\right| -\left| s_{i+s}\right\rangle \left\langle
s_{i}^{d}\right| \right\}  \label{antiexp}
\end{equation}%
If such an operator is also real, ${\cal A}={\cal A}^{\ast }$, it can be
converted into a compact hermitian one, ${\cal H}$, with the properties 
\begin{eqnarray}
{\cal H} &=&i{\cal A}  \nonumber \\
{\cal H}^{\dagger } &=&-i{\cal A}^{t}=i{\cal A}={\cal H}  \nonumber \\
{\cal H}^{\ast } &=&-i{\cal A}=i{\cal A}^{t}={\cal H}^{t}
\end{eqnarray}%
and 
\begin{equation}
{\cal H}=i{\cal A}=\sum_{1\leq i\leq s}\alpha _{i}\left\{ \left|
v_{i}\right\rangle \left\langle v_{i}\right| -\left| v_{i}^{\ast
}\right\rangle \left\langle v_{i}^{\ast }\right| \right\} ;\quad \alpha
_{i}\in {\Bbb R}
\end{equation}%
which is seen by observing that 
\begin{equation}
{\cal H}v=\alpha v\Rightarrow {\cal H}^{\ast }v^{\ast }=\alpha v^{\ast
}\Rightarrow {\cal H}^{t}v^{\ast }=\alpha v^{\ast }\Rightarrow {\cal H}%
v^{\ast }=-\alpha v^{\ast }
\end{equation}%
Thus 
\begin{equation}
{\cal A}=-\sum_{1\leq j\leq s}i\alpha _{j}\left\{ \left| v_{j}\right\rangle
\left\langle v_{j}\right| -\left| v_{j}^{\ast }\right\rangle \left\langle
v_{j}^{\ast }\right| \right\} ;\quad \alpha _{j}\in {\Bbb R}
\end{equation}%
We can define 
\begin{eqnarray}
\left| s_{i}\right\rangle &=&\left| v_{j}+v_{j}^{\ast }\right\rangle 
\nonumber \\
\left| s_{i+s}\right\rangle &=&i\left| v_{j}-v_{j}^{\ast }\right\rangle
\end{eqnarray}%
then 
\begin{eqnarray}
{\cal A}\left| v_{j}+v_{j}^{\ast }\right\rangle &=&-i\alpha _{j}\left|
v_{j}\right\rangle +i\alpha \left| v_{j}^{\ast }\right\rangle =-\alpha
i\left| v_{j}-v_{j}^{\ast }\right\rangle  \nonumber \\
{\cal A}i\left| v_{j}-v_{j}^{\ast }\right\rangle &=&\alpha _{j}\left|
v_{j}\right\rangle +\alpha _{j}\left| v_{j}^{\ast }\right\rangle =\alpha
_{j}\left| v_{j}+v_{j}^{\ast }\right\rangle
\end{eqnarray}%
thus 
\begin{equation}
{\cal A=}\sum_{1\leq i\leq s}\alpha _{i}\left\{ \left| s_{i}\right\rangle
\left\langle s_{i+s}^{d}\right| -\left| s_{i+s}\right\rangle \left\langle
s_{i}^{d}\right| \right\}
\end{equation}%
which justifies the expansion in eq. $\left( \ref{antiexp}\right) $.
}

\QSubDoc{Include SolODE}{%TCIDATA{LaTeXparent=0,0,EOMN.tex}
                      

\section{Solution of Singular ODE}

\subsection{Square Matrices}

If ${\bf A}$ is non-singular then one has the following situation for the
inhomogeneous ODE 
\begin{eqnarray}
{\bf A\dot{x}}\left( t\right) +{\bf B}\left( t\right) &=&{\bf f}\left(
t\right)  \nonumber \\
{\bf x}\left( t_{0}\right) &=&{\bf c}
\end{eqnarray}%
the homogeneous equation 
\begin{equation}
{\bf A\dot{x}}\left( t\right) +{\bf B}\left( t\right) ={\bf 0}
\end{equation}%
has general solution (non unique) 
\begin{equation}
{\bf x}\left( t\right) =e^{-{\bf A}^{-1}{\bf B}t}{\bf q}
\end{equation}%
and the initial value problem 
\begin{eqnarray}
{\bf A\dot{x}}\left( t\right) +{\bf B}\left( t\right) &=&{\bf 0}  \nonumber
\\
{\bf x}\left( t_{0}\right) &=&{\bf c}
\end{eqnarray}%
has the unique solution 
\begin{equation}
{\bf x}\left( t\right) =e^{-{\bf A}^{-1}{\bf B}\left( t-t_{0}\right) }{\bf c}
\end{equation}%
and if ${\bf A}$ is non-singular the unique solution of the inhomogeneous
initial value problem is 
\begin{equation}
{\bf x}\left( t\right) =e^{-{\bf A}^{-1}{\bf B}\left( t-t_{0}\right) }{\bf %
c+A}^{-1}\int_{t_{0}}^{t}e^{-{\bf A}^{-1}{\bf B}\left( t-s\right) }{\bf f}%
\left( s\right) ds
\end{equation}

\begin{definition}
For ${\bf A,B\in }{\Bbb C}^{n\times n}$ and $t_{0}\in {\Bbb R}$, the vector $%
{\bf c\in }{\Bbb C}^{n}$ is said to be a consistent initial vector
associated with $t_{0}$ for the equation ${\bf A\dot{x}}\left( t\right) +%
{\bf B}\left( t\right) ={\bf f}\left( t\right) $ when the initial value
problem ${\bf A\dot{x}}\left( t\right) +{\bf B}\left( t\right) ={\bf f}%
\left( t\right) $, ${\bf x}\left( t_{0}\right) ={\bf c}$ possesses at least
one solution.
\end{definition}

\begin{definition}
The equation ${\bf A\dot{x}}\left( t\right) +{\bf B}\left( t\right) ={\bf f}%
\left( t\right) $ is said to be tractable at the point $t_{0}$ if the
initial value problem ${\bf A\dot{x}}\left( t\right) +{\bf B}\left( t\right)
={\bf f}\left( t\right) $, ${\bf x}\left( t_{0}\right) ={\bf c}$ has a
unique solution for each consistent vector ${\bf c}$ associated with $t_{0}$.
\end{definition}

If the homogenous equation ${\bf A\dot{x}}\left( t\right) +{\bf B}\left(
t\right) ={\bf 0}$ is tractable at some point $t_{0}$ then it is tractable
at every $t$, so we say the equation is tractable.

\begin{theorem}
For ${\bf A,B\in }{\Bbb C}^{n\times n}$ the homogeneous differential
equation 
\begin{equation}
{\bf A\dot{x}}\left( t\right) +{\bf B}\left( t\right) ={\bf 0}
\end{equation}%
is tractable iff there exists a scalar $\lambda $ such that $\left( \lambda 
{\bf A+B}\right) ^{-1}$ exists.
\end{theorem}

\begin{proof}
see p173 ''{\it Generalized Inverses of Linear Transformations'' by S.L.
Campbell and C.D. Meyer Jr (Dover Publications Inc, 1979) }
\end{proof}

\begin{theorem}
Suppose that ${\bf A,B\in }{\Bbb C}^{n\times n}$ are that such $\exists
\lambda \quad \ni \left( \lambda {\bf A+B}\right) ^{-1}.$ Let ${\bf A}%
_{\lambda }=\left( \lambda {\bf A+B}\right) ^{-1}{\bf A,}$ ${\bf B}_{\lambda
}=\left( \lambda {\bf A+B}\right) ^{-1}{\bf B}$ and ${\bf f}_{\lambda
}=\left( \lambda {\bf A+B}\right) ^{-1}{\bf f.}$ $\forall \alpha ,\mu \in 
{\Bbb C}$ \ for which $\left( \alpha {\bf A+B}\right) ^{-1}$ and $\left( \mu 
{\bf A+B}\right) ^{-1}$ $\exists $, the following statements are true

\begin{itemize}
\item ${\bf A}_{\alpha }{\bf A}_{\alpha }^{D}={\bf A}_{\mu }{\bf A}_{\mu
}^{D}$

\item ${\bf A}_{\alpha }^{D}{\bf B}_{\alpha }={\bf A}_{\mu }^{D}{\bf B}_{\mu
}$ and ${\bf A}_{\alpha }{\bf B}_{\alpha }^{D}={\bf A}_{\mu }{\bf B}_{\mu
}^{D}$

\item $%
%TCIMACRO{\func{Ind}}%
%BeginExpansion
\mathop{\rm Ind}%
%EndExpansion
\left( {\bf A}_{\alpha }\right) =%
%TCIMACRO{\func{Ind}}%
%BeginExpansion
\mathop{\rm Ind}%
%EndExpansion
\left( {\bf A}_{\mu }\right) $ and $%
%TCIMACRO{\func{Range}}%
%BeginExpansion
\mathop{\rm Range}%
%EndExpansion
\left( {\bf A}_{\alpha }\right) =%
%TCIMACRO{\func{Range}}%
%BeginExpansion
\mathop{\rm Range}%
%EndExpansion
\left( {\bf A}_{\mu }\right) $

\item ${\bf A}_{\alpha }^{D}{\bf f}_{\alpha }={\bf A}_{\mu }^{D}{\bf f}_{\mu
}$

\item ${\bf B}_{\alpha }^{D}{\bf f}_{\alpha }={\bf B}_{\mu }^{D}{\bf f}_{\mu
}$
\end{itemize}
\end{theorem}

\begin{proof}
see p174 ''{\it Generalized Inverses of Linear Transformations'' by S.L.
Campbell and C.D. Meyer Jr (Dover Publications Inc, 1979) }
\end{proof}

\begin{theorem}
Suppose ${\bf A\dot{x}}\left( t\right) +{\bf B}\left( t\right) ={\bf 0}$ is
tractable. Then the general solution is given by 
\begin{equation}
{\bf x}\left( t\right) =e^{-{\bf A}_{\lambda }^{D}{\bf B}_{\lambda }t}{\bf A}%
_{\lambda }{\bf A}_{\lambda }^{D}{\bf q}
\end{equation}%
A vector ${\bf c\in }{\Bbb C}^{n}$ is a consistent initial vector for the
homogeneous equation iff ${\bf c\in }%
%TCIMACRO{\func{Range}}%
%BeginExpansion
\mathop{\rm Range}%
%EndExpansion
\left( {\bf A}_{\lambda }^{k}\right) =%
%TCIMACRO{\func{Range}}%
%BeginExpansion
\mathop{\rm Range}%
%EndExpansion
\left( {\bf A}_{\lambda }^{D}{\bf A}_{\lambda }\right) .$ Suppose that ${\bf %
f}\left( t\right) $ is k-times continuously differentiable around $t_{0}$
then the non-homogeneous equation ${\bf A\dot{x}}\left( t\right) +{\bf B}%
\left( t\right) ={\bf f}\left( t\right) $ always posses solutions and a
particular solution is given by 
\begin{equation}
{\bf x}\left( t\right) =e^{-{\bf A}_{\lambda }^{D}{\bf B}_{\lambda
}t}\int_{t_{0}}^{t}e^{{\bf A}_{\lambda }^{D}{\bf B}_{\lambda }s}{\bf A}%
_{\lambda }^{D}{\bf f}_{\lambda }\left( s\right) ds+\left( I-{\bf A}%
_{\lambda }{\bf A}_{\lambda }^{D}\right) \sum_{i=0}^{k-1}\left( -1\right)
^{i}\left[ {\bf A}_{\lambda }{\bf B}_{\lambda }^{D}\right] ^{i}{\bf B}%
_{\lambda }^{D}{\bf f}_{\lambda }^{\left( i\right) }\left( t\right)
\end{equation}%
moreover this expression is independent of $\lambda $. The general solution
is given by 
\begin{eqnarray}
{\bf x}\left( t\right) &=&e^{-{\bf A}_{\lambda }^{D}{\bf B}_{\lambda }t}{\bf %
A}_{\lambda }{\bf A_{\lambda }^{D}q+}e^{-{\bf A}_{\lambda }^{D}{\bf B}%
_{\lambda }t}\int_{t_{0}}^{t}e^{{\bf A}_{\lambda }^{D}{\bf B}_{\lambda }s}%
{\bf A}_{\lambda }^{D}{\bf f}_{\lambda }\left( s\right) ds  \nonumber \\
&&\qquad \qquad \qquad \qquad \qquad \qquad \qquad +\left( I-{\bf A}%
_{\lambda }{\bf A}_{\lambda }^{D}\right) \sum_{i=0}^{k-1}\left( -1\right)
^{i}\left[ {\bf A}_{\lambda }{\bf B}_{\lambda }^{D}\right] ^{i}{\bf B}%
_{\lambda }^{D}{\bf f}_{\lambda }^{\left( i\right) }\left( t\right)
\label{sol}
\end{eqnarray}%
\ \ \ \ \ \ \ \ \ \ \ \ \ \ \ \ \ \ \ \ \ \ \ \ \ \ \ \ \ \ \ \ \ \ \ \ \ \
\ \ \ \ \ \ \ \ \ \ \ \ A vector ${\bf c}$ is a consistent initial vector
associated with $t_{0}$ for the inhomogeneous equation iff ${\bf c}$ $\in
\left\{ {\bf w+}%
%TCIMACRO{\func{Range}}%
%BeginExpansion
\mathop{\rm Range}%
%EndExpansion
\left( {\bf A}_{\lambda }^{k}\right) \right\} $, where 
\begin{equation}
{\bf w=}\left( I-{\bf A}_{\lambda }{\bf A}_{\lambda }^{D}\right)
\sum_{i=0}^{k-1}\left( -1\right) ^{i}\left[ {\bf A}_{\lambda }{\bf B}%
_{\lambda }^{D}\right] ^{i}{\bf B}_{\lambda }^{D}{\bf f}_{\lambda }^{\left(
i\right) }\left( t\right)
\end{equation}%
which is independent of $\lambda .$ Furthermore the inhomogeneous equation
is tractable at $t_{0}$ and the unique solution of the initial value problem
with ${\bf x}\left( t\right) ={\bf c}$ is given by eq. $\left( \ref{sol}%
\right) $. [$k=\dim \left\{ \ker {\bf A}\right\} $]
\end{theorem}

\begin{proof}
see p175-176 ''{\it Generalized Inverses of Linear Transformations'' by S.L.
Campbell and C.D. Meyer Jr (Dover Publications Inc, 1979).}
\end{proof}

An important special case is when ${\bf B}$ is invertible. Then we can take $%
\lambda =0$ and 
\begin{eqnarray}
{\bf A}_{\lambda } &=&{\bf B}^{-1}{\bf A}  \nonumber \\
{\bf B}_{\lambda } &=&I  \nonumber \\
{\bf f}_{\lambda } &=&{\bf B}^{-1}{\bf f}
\end{eqnarray}

\subsection{Rectangular Matrices}

\begin{theorem}
Suppose $\left( \lambda {\bf A+B}\right) {\bf \ }${\bf is }1-1 and let 
\begin{equation}
{\bf \hat{A}=}\left( \lambda {\bf A+B}\right) ^{MP}A\text{ and }{\bf \hat{B}=%
}\left( \lambda {\bf A+B}\right) ^{MP}B.
\end{equation}%
Then all the solutions of ${\bf A\dot{x}}\left( t\right) +{\bf B}\left(
t\right) ={\bf 0}$ are of the form 
\begin{equation}
{\bf x}\left( t\right) =e^{-{\bf \hat{A}}^{D}{\bf \hat{B}}t}{\bf q}
\end{equation}%
where 
\begin{equation}
{\bf q\in }%
%TCIMACRO{\func{Range}}%
%BeginExpansion
\mathop{\rm Range}%
%EndExpansion
\left( {\bf \hat{A}}^{D}{\bf \hat{A}}\right)
\end{equation}
and 
\begin{equation}
\left[ {\bf I-}\left( \lambda {\bf A+B}\right) \left( \lambda {\bf A+B}%
\right) ^{MP}\right] {\bf A\hat{A}}^{D}\left[ {\bf \hat{A}}^{Dm}{\bf \hat{B}}%
^{m}\right] {\bf q=0}\text{ for }m=0,1,\ldots ,n
\end{equation}

\begin{corollary}
If $\left( \lambda {\bf A+B}\right) {\bf \ }${\bf is }1-1, and $\ker \left(
\lambda ^{\ast }{\bf A}^{\dagger }{\bf +B}^{\dagger }\right) $=$\ker \left( 
{\bf A}^{\dagger }\right) \cap \ker \left( {\bf B}^{\dagger }\right) $ Then
all the solutions of ${\bf A\dot{x}}\left( t\right) +{\bf B}\left( t\right) =%
{\bf 0}$ are of the form 
\begin{equation}
{\bf x}\left( t\right) =e^{-{\bf \hat{A}}^{D}{\bf \hat{B}}t}{\bf \hat{A}}^{D}%
{\bf \hat{A}q}
\end{equation}%
where ${\bf q}$ is arbitrary.

\begin{theorem}
Suppose $\left( \lambda {\bf A+B}\right) {\bf \ }${\bf is }1-1 and ${\bf A%
\dot{x}}\left( t\right) +{\bf B}\left( t\right) ={\bf f}$ is consistent.
Then all the solutions of ${\bf A\dot{x}}\left( t\right) +{\bf B}\left(
t\right) ={\bf f}$ are of the form 
\begin{eqnarray}
{\bf x}\left( t\right) &=&e^{-{\bf \hat{A}}^{D}{\bf \hat{B}}t}{\bf \hat{A}}%
^{D}{\bf \hat{A}q+\hat{A}}^{D}e^{-{\bf \hat{A}}^{D}{\bf \hat{B}}%
t}\int_{t_{0}}^{t}e^{{\bf \hat{A}}^{D}{\bf \hat{B}}s}{\bf A}^{D}{\bf \hat{f}}%
\left( s\right) ds  \nonumber \\
&&\qquad \qquad \qquad \qquad \qquad \qquad \qquad +\left( {\bf I}-{\bf 
\hat{A}\hat{A}}^{D}\right) \sum_{i=0}^{k-1}\left( -1\right) ^{i}\left[ {\bf 
\hat{A}\hat{B}}^{D}\right] ^{i}{\bf \hat{B}}^{D}{\bf \hat{f}}^{\left(
i\right) }\left( t\right)
\end{eqnarray}%
where ${\bf \hat{f}=}\left( \lambda {\bf A+B}\right) ^{MP}{\bf f}$
\end{theorem}

\begin{theorem}
Suppose $\left( \lambda {\bf A+B}\right) {\bf \ }${\bf is }1-1 and $\ker
\left( \lambda ^{\ast }{\bf A}^{\dagger }{\bf +B}^{\dagger }\right) $=$\ker
\left( {\bf A}^{\dagger }\right) \cap \ker \left( {\bf B}^{\dagger }\right) $%
. Then ${\bf A\dot{x}}\left( t\right) +{\bf B}\left( t\right) ={\bf f}$ is
consistent $iff$ 
\begin{equation}
\left[ {\bf I-}\left( \lambda {\bf A+B}\right) \left( \lambda {\bf A+B}%
\right) ^{MP}\right] {\bf f=0}
\end{equation}
\end{theorem}
\end{corollary}
\end{theorem}
}

\QSubDoc{Include TaySer}{%TCIDATA{LaTeXparent=0,0,EOMN.tex}
                      
%TCIDATA{ChildDefaults=chapter:1,page:1}


\section{Taylor Series Expansion for Functionals}

\subsection{Functionals, Function Spaces and Change of Basis}

We are considering nonlinear (in general) functionals ${\cal F}%
:F\longrightarrow {\Bbb C}$, where $F$ is a seperable linear space of
functions $\left\{ \varkappa \right\} $ that has a basis (thus has a
countable basis $\left\{ \left| {\bf \varkappa }_{j}\right\rangle ;1\leq
j\leq n\right\} $). Each function can then be expanded as 
\begin{equation}
\left| \varkappa \right\rangle =\sum_{1\leq j\leq n}\varkappa _{j}\left| 
{\bf \varkappa }_{j}\right\rangle \equiv \sum_{1\leq j\leq n}\left| {\bf %
\varkappa }_{j}\right\rangle \left\langle {\bf \varkappa }_{j}^{d}|\varkappa
\right\rangle \equiv \sum_{1\leq j\leq n}\left\langle {\bf \varkappa }%
_{j}^{d}|\varkappa \right\rangle \left| {\bf \varkappa }_{j}\right\rangle
\equiv \sum_{1\leq j\leq n}\varkappa \left( {\bf \varkappa }_{j}^{d}\right)
\left| {\bf \varkappa }_{j}\right\rangle
\end{equation}%
Note that the function $\varkappa \left( {\bf \varkappa }_{j}^{d}\right) $
is {\em not linear in ''j''}, i.e. $\varkappa \left( {\bf \varkappa }%
_{j_{1}}^{d}+{\bf \varkappa }_{j_{2}}^{d}\right) \neq ${\em \ }$\varkappa
\left( {\bf \varkappa }_{j_{1}}^{d}\right) +\varkappa \left( {\bf \varkappa }%
_{j_{2}}^{d}\right) $, though it is linear wrt to dual vectors belonging to $%
F^{d}$. We also assume we have a continuous decomposition in terms of an
overcomplete coordinate basis $\left\{ \left| {\sf z}\right\rangle \right\} $
\begin{equation}
\left| \varkappa \right\rangle =\int \varkappa \left( {\sf z}\right) \left| 
{\sf z}\right\rangle d{\sf z}\equiv \int \left| {\sf z}\right\rangle
\left\langle {\sf z}|\varkappa \right\rangle d{\sf z}\equiv \int
\left\langle {\sf z}|\varkappa \right\rangle \left| {\sf z}\right\rangle d%
{\sf z}\equiv \int \varkappa \left( {\sf z}\right) \left| {\sf z}%
\right\rangle d{\sf z.}
\end{equation}%
again we note that $\varkappa \left( {\sf z}\right) $ is not linear in the
continuum of index labels ''${\sf z}$''. $\lceil $The functions $\left\{
\left| {\sf z}\right\rangle \right\} $ and the
functions/distributions/hyperfunctions $\left\{ \left\langle {\sf z}\right|
\right\} $ can be quite different it is just their labelling sets that are
the same$\rfloor $. From the preceding two expansions we have the change of
basis formulae 
\begin{eqnarray}
\left| {\bf \varkappa }_{j}\right\rangle &=&\int \left| {\sf z}\right\rangle
\left\langle {\sf z}|{\bf \varkappa }_{j}\right\rangle d{\sf z}\equiv \int
\left| {\sf z}\right\rangle \varkappa _{j}\left( {\sf z}\right) d{\sf z} 
\nonumber \\
&\Rightarrow &\frac{\partial }{\partial \varkappa _{j}}=\int \varkappa
_{j}\left( {\sf z}\right) \frac{\delta }{\delta \varkappa \left( {\sf z}%
\right) }d{\sf z\equiv }\int \frac{\partial \varkappa \left( {\sf z}\right) 
}{\partial \varkappa _{j}}\frac{\delta }{\delta \varkappa \left( {\sf z}%
\right) }d{\sf z}  \label{change1}
\end{eqnarray}%
and%
\begin{eqnarray}
\left| {\sf z}\right\rangle &=&\sum_{1\leq j\leq n}\left| {\bf \varkappa }%
_{j}\right\rangle \left\langle {\bf \varkappa }_{j}^{d}|{\sf z}\right\rangle
\equiv \sum_{1\leq j\leq n}\left| {\bf \varkappa }_{j}\right\rangle
\varkappa _{j}^{d}\left( {\sf z}\right)  \nonumber \\
&\Rightarrow &\frac{\delta }{\delta \varkappa \left( {\sf z}\right) }%
=\sum_{1\leq j\leq n}\varkappa _{j}^{d}\left( {\sf z}\right) \frac{\partial 
}{\partial \varkappa _{j}}\equiv \sum_{1\leq j\leq n}\frac{\delta \varkappa
_{j}}{\delta \varkappa \left( {\sf z}\right) }\frac{\partial }{\partial
\varkappa _{j}}  \label{change2}
\end{eqnarray}%
Note in eq. $\left( \ref{change1}\right) \frac{\partial \varkappa \left( 
{\sf z}\right) }{\partial \varkappa _{j}}$ {\bf is defined }to be equal to $%
\varkappa _{j}\left( {\sf z}\right) =\left\langle {\sf z}|{\bf \varkappa }%
_{j}\right\rangle $ and are the ''matrix elements'' of the Jacobian matrix
of the transformation from the $\left\{ \left| {\sf z}\right\rangle \right\} 
$ to the $\left\{ \left| {\bf \varkappa }_{j}\right\rangle \right\} $
coordinate system. The Jacobian matrix elements of the transformation from
the $\left\{ \left| {\bf \varkappa }_{j}\right\rangle \right\} $ to the $%
\left\{ \left| {\sf z}\right\rangle \right\} $ coordinate system are $%
\varkappa _{j}^{d}\left( {\sf z}\right) =\left\langle {\bf \varkappa }%
_{j}^{d}|{\sf z}\right\rangle $.

Functionals can be expanded in either basis around the point $\varkappa _{0}$
as 
\begin{eqnarray}
{\cal F}\left( \varkappa \right) &=&{\cal F}\left( \varkappa _{0}\right)
+\sum_{1\leq j\leq n}{\cal F}_{1j}\left( \varkappa _{j}-\varkappa
_{0j}\right) +\sum_{1\leq j_{1},j_{2}\leq n}{\cal F}_{2j_{1}j_{2}}\left(
\varkappa _{j_{1}}-\varkappa _{0j_{1}}\right) \left( \varkappa
_{j_{2}}-\varkappa _{0j_{2}}\right) +\cdots  \nonumber \\
{\cal F}\left( \varkappa \right) &=&{\cal F}\left( \varkappa _{0}\right)
+\int {\cal F}_{1}\left( {\sf z}\right) \left( \varkappa \left( {\sf z}%
\right) -\varkappa _{0}\left( {\sf z}\right) \right) d{\sf z}  \nonumber \\
&&\qquad \quad +\int {\cal F}_{2}\left( {\sf z}_{1},{\sf z}_{2}\right)
\left( \varkappa \left( {\sf z}_{1}\right) -\varkappa _{0}\left( {\sf z}%
_{1}\right) \right) \left( \varkappa \left( {\sf z}_{2}\right) -\varkappa
_{0}\left( {\sf z}_{2}\right) \right) d{\sf z}_{1}d{\sf z}_{2}+\cdots
\end{eqnarray}%
where 
\begin{eqnarray}
{\cal F}_{1j} &=&\frac{\partial {\cal F}\left( \varkappa _{0}\right) }{%
\partial \varkappa _{j}},{\cal F}_{2j_{1}j_{2}}=\frac{1}{2}\frac{\partial
^{2}{\cal F}\left( \varkappa _{0}\right) }{\partial \varkappa
_{j_{1}}\partial \varkappa _{j_{2}}},\cdots ,{\cal F}_{rj_{1}\ldots j_{r}}=%
\frac{1}{r!}\frac{\partial ^{r}{\cal F}\left( \varkappa _{0}\right) }{%
\partial \varkappa _{j_{1}}\cdots \partial \varkappa _{j_{r}}}  \nonumber \\
{\cal F}_{1}\left( {\sf z}\right) &=&\frac{\delta {\cal F}\left( \varkappa
_{0}\right) }{\delta \varkappa \left( {\sf z}\right) },{\cal F}_{2}\left( 
{\sf z}_{1},{\sf z}_{2}\right) =\frac{1}{2}\frac{\delta ^{2}{\cal F}\left(
\varkappa _{0}\right) }{\delta \varkappa \left( {\sf z}_{1}\right) \delta
\varkappa \left( {\sf z}_{2}\right) },\cdots ,{\cal F}_{r}\left( {\sf z}%
_{1},\ldots ,{\sf z}_{r}\right) =\frac{1}{r!}\frac{\delta ^{r}{\cal F}\left(
\varkappa _{0}\right) }{\delta \varkappa \left( {\sf z}_{1}\right) \cdots
\delta \varkappa \left( {\sf z}_{r}\right) }
\end{eqnarray}%
From eq. $\left( \ref{change1}\right) $%
\begin{equation}
\frac{\partial {\cal F}\left( \varkappa _{0}\right) }{\partial \varkappa _{j}%
}=\int \varkappa _{j}\left( {\sf z}\right) \frac{\delta {\cal F}\left(
\varkappa _{0}\right) }{\delta \varkappa \left( {\sf z}\right) }d{\sf z}%
\equiv \int \frac{\partial \varkappa \left( {\sf z}\right) }{\partial
\varkappa _{j}}\frac{\delta {\cal F}\left( \varkappa _{0}\right) }{\delta
\varkappa \left( {\sf z}\right) }d{\sf z}
\end{equation}%
and from eq. $\left( \ref{change2}\right) $%
\begin{equation}
\frac{\delta {\cal F}\left( \varkappa _{0}\right) }{\delta \varkappa \left( 
{\sf z}\right) }=\sum_{1\leq j\leq n}\varkappa _{j}^{d}\left( {\sf z}\right) 
\frac{\partial {\cal F}\left( \varkappa _{0}\right) }{\partial \varkappa _{j}%
}\equiv \sum_{1\leq j\leq n}\frac{\delta \varkappa _{j}}{\delta \varkappa
\left( {\sf z}\right) }\frac{\partial {\cal F}\left( \varkappa _{0}\right) }{%
\partial \varkappa _{j}}
\end{equation}%
$\lceil $Note that in these expressions $\frac{\delta }{\delta \varkappa
\left( {\sf z}\right) }$signifies partial derrivative operation in the $%
\left| {\sf z}\right\rangle $ direction$\rfloor $.

Time dependence can be introduced by 
\begin{eqnarray}
\left| {\bf \varkappa }\left( t\right) \right\rangle &=&\sum_{1\leq j\leq
n}\varkappa _{j}\left( t\right) \left| {\bf \varkappa }_{j}\right\rangle
\Rightarrow \\
\varkappa \left( {\sf z},t\right) &\equiv &\left\langle {\sf z}|{\bf %
\varkappa }\left( t\right) \right\rangle =\sum_{1\leq j\leq n}\varkappa
_{j}\left( t\right) \left\langle {\sf z}|{\bf \varkappa }_{j}\right\rangle \\
&=&\sum_{1\leq j\leq n}\varkappa _{j}\left\langle {\sf z}|{\bf \varkappa }%
_{j}\left( t\right) \right\rangle =\sum_{1\leq j\leq n}\varkappa
_{j}\left\langle {\sf z}\left( t\right) |{\bf \varkappa }_{j}\right\rangle
\end{eqnarray}%
If we introduce a time evolution operator ${\cal T}\left( t_{2},t_{1}\right) 
$ then%
\begin{equation}
\left| {\bf \varkappa }\left( t_{2}\right) \right\rangle ={\cal T}\left(
t_{2},t_{1}\right) \left| {\bf \varkappa }\left( t_{1}\right) \right\rangle
=Te^{\int_{t_{1}}^{t_{2}}\Gamma \left( t\right) dt}\left| {\bf \varkappa }%
\left( t_{1}\right) \right\rangle
\end{equation}%
and 
\begin{equation}
\left| {\bf \varkappa }_{\lambda }\left( t\right) \right\rangle
=Te^{\int_{0}^{t}\Gamma \left( t^{\prime }\right) dt^{\prime }}\left| {\bf %
\varkappa }_{\lambda }\right\rangle \equiv {\cal T}\left( t\right) \left| 
{\bf \varkappa }_{\lambda }\right\rangle ={\cal T}\left( t\right) {\cal V}%
_{\lambda }\left| {\bf \varkappa }_{0}\right\rangle \equiv {\cal T}_{\lambda
}\left( t\right) \left| {\bf \varkappa }_{0}\right\rangle
\end{equation}%
where ${\cal V}_{\lambda }$ is a transformation from the reference function $%
{\bf \varkappa }_{0}$ to the function ${\bf \varkappa }_{\lambda }$ giving 
\begin{equation}
\varkappa _{\lambda }\left( {\sf z},t\right) =\left\langle {\sf z}|{\bf %
\varkappa }_{\lambda }\left( t\right) \right\rangle =\left\langle {\sf z}|%
{\cal T}_{\lambda }\left( t\right) {\bf \varkappa }_{0}\right\rangle
=\left\langle {\cal T}_{\lambda }\left( t\right) ^{\dagger }{\sf z}|{\bf %
\varkappa }_{0}\right\rangle \equiv \left\langle {\sf z}_{\lambda }\left(
t\right) |{\bf \varkappa }_{0}\right\rangle =\varkappa \left( {\sf z}%
_{\lambda }\left( t\right) \right)
\end{equation}%
Thus as one varies over all {\sf ''}${\sf z"}$ all possible coordinate paths 
$\left\{ {\sf z}_{\lambda }\left( t\right) \right\} $ are produced. These
paths are in general discontinuous wrt time variation. The time dependence
of the functions $\left\{ \varkappa \right\} $ can thus be described in
terms of ${\bf \varkappa }\left( t\right) \left( {\sf z}\right) \equiv
\varkappa \left( {\sf z},t\right) $ for any $\left\{ {\sf z},t\right\} $ or $%
\varkappa \left( {\sf z}_{\lambda }\left( t\right) \right) $ for any
coordinate path ${\sf z}_{\lambda }\left( t\right) $. In the following we
usually use the form $\varkappa \left( {\sf z},t\right) $.

\subsection{The Action Functional}

In particular we consider the action functional 
\begin{equation}
{\cal A}={\cal A}\left( \varkappa \right) ={\cal A}\left( \varkappa ,\tau
\right) ={\cal A}\left( {\cal V},\omega ,\tau \right) ={\cal A}\left( {\cal V%
}_{fib},\omega _{fib},{\cal V}_{den},\omega _{den},\tau \right)
\end{equation}%
given in terms of the Lagrangian by 
\[
{\cal A}\left( \varkappa \right) =\int_{t_{i}}^{t_{f}}{\cal L}\left(
\varkappa ,\dot{\varkappa},t\right) dt=\int_{t_{i}}^{t_{f}}\frac{d{\cal A}%
\left( \varkappa \left( t\right) \right) }{dt}dt=\int_{t_{i}}^{t_{f}}{\cal 
\dot{A}}\left( \varkappa \left( t\right) \right) dt 
\]%
The action is defined over the closed time interval $\tau \equiv \left[
t_{i},t_{f}\right] $. Clearly the variables $\left\{ \dot{\varkappa}\right\} 
$ are dependent on $\left\{ \varkappa \right\} $ over the open time interval 
$\left( t_{i},t_{f}\right) $ and if the fixed endpoint conditions on $%
\left\{ \varkappa \right\} $ are invoked $\lceil \varkappa \left( {\sf z}%
,t_{i}\right) =\varkappa _{i}\left( {\sf z}\right) $ and $\varkappa \left( 
{\sf z},t_{f}\right) =\varkappa _{f}\left( {\sf z}\right) \rfloor $ the
action will only depend on the {\em independent }variables $\left\{
\varkappa \right\} $. On the other hand the Lagrangian is local in time and
at each instant $\varkappa =$ $\varkappa \left( {\sf z},t\right) $ and $\dot{%
\varkappa}=$ $\dot{\varkappa}\left( {\sf z},t\right) $ can be thought of as
independent functions of ${\sf z}$ and thus as independent variables $\lceil 
$this relates to the distinction between $\varkappa \left( {\sf z},t\right) $
and $\varkappa \left( {\sf z}_{\lambda }\left( t\right) \right) \rfloor $.
Small variations of the trajectories $\varkappa $ produce changes in the
action and Lagrangian expressed as 
\begin{eqnarray}
{\cal A}\left( \varkappa +\delta \varkappa \right) &=&\int_{t_{i}}^{t_{f}}%
{\cal \dot{A}}\left( \varkappa \left( t\right) +\delta \varkappa \left(
t\right) \right) dt=\int_{t_{i}}^{t_{f}}{\cal L}\left( \varkappa +\delta
\varkappa ,\dot{\varkappa}+\delta \dot{\varkappa},t\right) dt \\
{\cal A}\left( \varkappa +\delta \varkappa \right) &=&\sum_{0\leq m\leq
\infty }\delta ^{\left( m\right) }{\cal A}\left( \varkappa \right) \\
\delta ^{\left( m\right) }{\cal A}\left( \varkappa \right)
&=&\int_{t_{i}}^{t_{f}}\delta ^{\left( m\right) }{\cal L}\left( \varkappa ,%
\dot{\varkappa},t\right) dt=\int_{t_{i}}^{t_{f}}\delta ^{\left( m\right) }%
{\cal \dot{A}}\left( \varkappa \left( t\right) \right) dt \\
{\cal L}\left( \varkappa +\delta \varkappa ,\dot{\varkappa}+\delta \dot{%
\varkappa},t\right) &=&\sum_{0\leq i\leq \infty }\delta ^{\left( m\right) }%
{\cal L}\left( \varkappa ,\dot{\varkappa},t\right)
\end{eqnarray}%
where $\delta ^{\left( 0\right) }{\cal A}\left( \varkappa \right) \equiv 
{\cal A}\left( \varkappa \right) $ and $\delta ^{\left( 1\right) }{\cal A}%
\left( \varkappa \right) \equiv \delta {\cal A}\left( \varkappa \right) $
and similarly for ${\cal L}$. Each order of variation can be computed in
terms of functional derivatives as 
\begin{eqnarray}
&&\delta ^{\left( m\right) }{\cal L}\left( \varkappa +\delta \varkappa ,\dot{%
\varkappa}+\delta \dot{\varkappa},t\right) =\sum\Sb m_{1}+m_{2}=m  \\ 0\leq
m_{1}\leq m;0\leq m_{2}\leq m  \endSb \frac{1}{m_{1}!m_{2}!}D^{m_{1}+m_{2}}%
{\cal L}\left( \varkappa ,\dot{\varkappa},t\right) \left( \otimes
^{m_{1}}\delta \varkappa \otimes ^{m_{2}}\delta \dot{\varkappa}\right) 
\nonumber \\
&&D^{m_{1}+m_{2}}{\cal L}\left( \varkappa ,\dot{\varkappa},t\right) \left(
\otimes ^{m_{1}}\delta \varkappa \otimes ^{m_{2}}\delta \dot{\varkappa}%
\right) =  \nonumber \\
&&\int \left\{ \frac{\delta ^{m_{1}+m_{2}}{\cal L}\left( \varkappa ,\dot{%
\varkappa},t\right) }{\delta \varkappa \left( {\sf z}_{1}\right) \cdots
\delta \varkappa \left( {\sf z}_{m_{1}}\right) \cdots \delta \dot{\varkappa}%
\left( {\sf z}_{m_{1}+1}\right) \cdots \delta \dot{\varkappa}\left( {\sf z}%
_{m_{2}}\right) }\delta \varkappa \left( {\sf z}_{1},t\right) \cdots \delta
\varkappa \left( {\sf z}_{m_{1}},t\right) \delta \dot{\varkappa}\left( {\sf z%
}_{m_{1}+1},t\right) \cdots \delta \dot{\varkappa}\left( {\sf z}%
_{m_{2}},t\right) \right\}  \nonumber \\
&&\qquad \qquad \qquad \qquad \qquad \qquad \qquad \qquad \qquad \qquad
\qquad \qquad \qquad \qquad \qquad \qquad \qquad \qquad \times d{\sf z}%
_{1}\cdots d{\sf z}_{m}
\end{eqnarray}%
and 
\begin{eqnarray}
&&{\cal A}\left( \varkappa +\delta \varkappa \right) =\sum_{0\leq m\leq
\infty }\delta ^{\left( m\right) }{\cal A}\left( \varkappa \right)  \nonumber
\\
&=&\sum_{0\leq m\leq \infty }\int_{t_{i}}^{t_{f}}\delta ^{\left( m\right) }%
{\cal \dot{A}}\left( \varkappa \left( t\right) \right) dt=\sum_{0\leq m\leq
\infty }\int_{t_{i}}^{t_{f}}\frac{\delta ^{m}{\cal \dot{A}}\left( \varkappa
\left( t\right) \right) }{\delta \varkappa \left( {\sf z}_{1}\right) \cdots
\delta \varkappa \left( {\sf z}_{m}\right) }\delta \varkappa \left( {\sf z}%
_{1},t\right) \cdots \delta \varkappa \left( {\sf z}_{m},t\right) d{\sf z}%
_{1}\cdots d{\sf z}_{m}dt
\end{eqnarray}%
In particular the first order variation is given as 
\begin{eqnarray}
\delta ^{\left( 1\right) }{\cal A}\left( \varkappa \right)
&=&\int_{t_{i}}^{t_{f}}\delta ^{\left( 1\right) }{\cal \dot{A}}\left(
\varkappa \left( t\right) \right) dt=\int_{t_{i}}^{t_{f}}\frac{\delta {\cal 
\dot{A}}\left( \varkappa \left( t\right) \right) }{\delta \varkappa \left( 
{\sf z}\right) }\delta \varkappa \left( {\sf z}\right) d{\sf z}dt  \label{A1}
\\
&=&\int_{t_{i}}^{t_{f}}\delta ^{\left( 1\right) }{\cal L}\left( \varkappa ,%
\dot{\varkappa},t\right) dt=\int_{t_{i}}^{t_{f}}\left\{ \frac{\delta {\cal L}%
\left( \varkappa ,\dot{\varkappa},t\right) }{\delta \varkappa \left( {\sf z}%
\right) }\delta \varkappa \left( {\sf z}\right) +\frac{\delta {\cal L}\left(
\varkappa ,\dot{\varkappa},t\right) }{\delta \dot{\varkappa}\left( {\sf z}%
\right) }\delta \dot{\varkappa}\left( {\sf z}\right) \right\} d{\sf z}dt
\label{L1}
\end{eqnarray}%
The dependence of $\dot{\varkappa}$ on $\varkappa $ over the time interval $%
\left( t_{i},t_{f}\right) $ plus the fixed end point conditions can be
displayed by using integration by parts wrt time on the $\dot{\varkappa}$
term in eq. $\left( \ref{L1}\right) $ to give 
\begin{eqnarray}
&&\int_{t_{i}}^{t_{f}}\frac{\delta {\cal L}\left( \varkappa ,\dot{\varkappa}%
,t\right) }{\delta \dot{\varkappa}\left( {\sf z}\right) }\delta \dot{%
\varkappa}\left( {\sf z},t\right) d{\sf z}dt  \nonumber \\
&&\left. {}\right. \qquad =\int \left\{ \left. \delta \varkappa \left( {\sf z%
},t\right) \frac{\delta {\cal L}\left( \varkappa ,\dot{\varkappa},t\right) }{%
\delta \dot{\varkappa}\left( {\sf z}\right) }\right|
_{t_{i}}^{t_{f}}\right\} d{\sf z}-\int_{t_{i}}^{t_{f}}\frac{d}{dt}\left\{ 
\frac{\delta {\cal L}\left( {\cal \varkappa },{\cal \dot{\varkappa}}%
,t\right) }{\delta {\cal \dot{\varkappa}}\left( {\sf z}\right) }\right\}
\delta \varkappa \left( {\sf z},t\right) d{\sf z}dt
\end{eqnarray}%
which when we constrain the end points of the trajectories to be fixed
becomes 
\begin{equation}
=-\int_{t_{i}}^{t_{f}}\frac{d}{dt}\left\{ \frac{\delta {\cal L}\left( {\cal %
\varkappa },{\cal \dot{\varkappa}},t\right) }{\delta {\cal \dot{\varkappa}}%
\left( {\sf z}\right) }\right\} \delta \varkappa \left( {\sf z},t\right) d%
{\sf z}dt;\qquad \text{ }\delta \varkappa \left( t_{f}\right) =\delta
\varkappa \left( t_{i}\right) =0
\end{equation}%
producing 
\begin{eqnarray}
\delta ^{\left( 1\right) }{\cal A}\left( \varkappa \right)
&=&\int_{t_{i}}^{t_{f}}\delta ^{\left( 1\right) }{\cal L}\left( \varkappa ,%
\dot{\varkappa},t\right) dt=\int_{t_{i}}^{t_{f}}\left\{ \frac{\delta {\cal L}%
\left( {\cal \varkappa },{\cal \dot{\varkappa}},t\right) }{\delta {\cal %
\varkappa }\left( {\sf z}\right) }-\frac{d}{dt}\left\{ \frac{\delta {\cal L}%
\left( {\cal \varkappa },{\cal \dot{\varkappa}},t\right) }{\delta {\cal \dot{%
\varkappa}}\left( {\sf z}\right) }\right\} \right\} \delta \varkappa \left( 
{\sf z},t\right) d{\sf z}dt  \nonumber \\
&=&\int_{t_{i}}^{t_{f}}\delta ^{\left( 1\right) }{\cal \dot{A}}\left(
\varkappa \left( t\right) \right) dt=\int_{t_{i}}^{t_{f}}\frac{\delta {\cal 
\dot{A}}\left( \varkappa \left( t\right) \right) }{\delta \varkappa \left( 
{\sf z}\right) }\delta \varkappa \left( {\sf z}\right) d{\sf z}dt
\end{eqnarray}%
and we can make the identifications 
\begin{eqnarray}
\delta ^{\left( 1\right) }{\cal \dot{A}}\left( \varkappa \left( t\right)
\right) &\equiv &\int \left\{ \frac{\delta {\cal L}\left( {\cal \varkappa },%
{\cal \dot{\varkappa}},t\right) }{\delta {\cal \varkappa }\left( {\sf z}%
\right) }-\frac{d}{dt}\left\{ \frac{\delta {\cal L}\left( {\cal \varkappa },%
{\cal \dot{\varkappa}},t\right) }{\delta {\cal \dot{\varkappa}}\left( {\sf z}%
\right) }\right\} \right\} \delta \varkappa \left( {\sf z},t\right) d{\sf z}
\nonumber \\
\frac{\delta {\cal \dot{A}}\left( \varkappa \left( t\right) \right) }{\delta
\varkappa \left( {\sf z}\right) } &\equiv &\frac{\delta {\cal L}\left( {\cal %
\varkappa },{\cal \dot{\varkappa}},t\right) }{\delta {\cal \varkappa }\left( 
{\sf z}\right) }-\frac{d}{dt}\left\{ \frac{\delta {\cal L}\left( {\cal %
\varkappa },{\cal \dot{\varkappa}},t\right) }{\delta {\cal \dot{\varkappa}}%
\left( {\sf z}\right) }\right\}  \label{D1A}
\end{eqnarray}%
The second order variation of the action is 
\begin{eqnarray}
&&\delta ^{\left( 2\right) }{\cal A}\left( \varkappa \right) =\frac{1}{2}%
\int_{t_{i}}^{t_{f}}\frac{\delta ^{2}{\cal \dot{A}}\left( \varkappa \left(
t\right) \right) }{\delta \varkappa \left( {\sf z}_{1}\right) \delta
\varkappa \left( {\sf z}_{2}\right) }\delta \varkappa \left( {\sf z}%
_{1},t\right) \delta \varkappa \left( {\sf z}_{2},t\right) d{\sf z}_{1}d{\sf %
z}_{2}dt  \nonumber \\
&&\left. {}\right. =\frac{1}{2}\int_{t_{i}}^{t_{f}}\left\{ \frac{\delta ^{2}%
{\cal L}\left( {\cal \varkappa },{\cal \dot{\varkappa}},t\right) }{\delta
\varkappa \left( {\sf z}_{1}\right) \delta \varkappa \left( {\sf z}%
_{2}\right) }\delta \varkappa \left( {\sf z}_{1},t\right) \delta \varkappa
\left( {\sf z}_{2},t\right) +2\frac{\delta ^{2}{\cal L}\left( {\cal %
\varkappa },{\cal \dot{\varkappa}},t\right) }{\delta \varkappa \left( {\sf z}%
_{1}\right) \delta \dot{\varkappa}\left( {\sf z}_{2}\right) }\delta
\varkappa \left( {\sf z}_{1},t\right) \delta \dot{\varkappa}\left( {\sf z}%
_{2},t\right) \right.  \nonumber \\
&&\qquad \qquad \qquad \qquad \qquad \qquad \qquad +\left. \frac{\delta ^{2}%
{\cal L}\left( {\cal \varkappa },{\cal \dot{\varkappa}},t\right) }{\delta 
\dot{\varkappa}\left( {\sf z}_{1}\right) \delta \dot{\varkappa}\left( {\sf z}%
_{2}\right) }\delta \dot{\varkappa}\left( {\sf z}_{1},t\right) \delta \dot{%
\varkappa}\left( {\sf z}_{2},t\right) \right\} d{\sf z}_{1}d{\sf z}_{2}dt
\label{Avar2}
\end{eqnarray}%
Using integration by parts on the second term of the Lagrangian we obtain 
\begin{eqnarray}
&&\int_{t_{i}}^{t_{f}}\frac{\delta ^{2}{\cal L}\left( {\cal \varkappa },%
{\cal \dot{\varkappa}},t\right) }{\delta \varkappa \left( {\sf z}_{1}\right)
\delta \dot{\varkappa}\left( {\sf z}_{2}\right) }\delta \varkappa \left( 
{\sf z}_{1},t\right) \delta \dot{\varkappa}\left( {\sf z}_{2},t\right) d{\sf %
z}_{1}d{\sf z}_{2}dt=  \nonumber \\
&&\int \left\{ \left. \left\{ \delta \varkappa \left( {\sf z}_{1},t\right)
\delta \varkappa \left( {\sf z}_{2},t\right) \frac{\delta ^{2}{\cal L}\left(
\varkappa ,\dot{\varkappa},t\right) }{\delta \varkappa \left( {\sf z}%
_{1}\right) \delta \dot{\varkappa}\left( {\sf z}_{2}\right) }\right\}
\right| _{t_{i}}^{t_{f}}\right\} d{\sf z}_{1}d{\sf z}_{2}  \nonumber \\
&&\qquad \qquad \qquad \qquad \qquad -\int_{t_{i}}^{t_{f}}\frac{d}{dt}%
\left\{ \frac{\delta ^{2}{\cal L}\left( \varkappa ,\dot{\varkappa},t\right) 
}{\delta \varkappa \left( {\sf z}_{1}\right) \delta \dot{\varkappa}\left( 
{\sf z}_{2}\right) }\right\} \delta \varkappa \left( {\sf z}_{1},t\right)
\delta \varkappa \left( {\sf z}_{2},t\right) d{\sf z}_{1}d{\sf z}_{2}dt 
\nonumber \\
&&\left. {}\right. \qquad \qquad \qquad \qquad \qquad \quad
=-\int_{t_{i}}^{t_{f}}\frac{d}{dt}\left\{ \frac{\delta ^{2}{\cal L}\left(
\varkappa ,\dot{\varkappa},t\right) }{\delta \varkappa \left( {\sf z}%
_{1}\right) \delta \dot{\varkappa}\left( {\sf z}_{2}\right) }\right\} \delta
\varkappa \left( {\sf z}_{1},t\right) \delta \varkappa \left( {\sf z}%
_{2},t\right) d{\sf z}_{1}d{\sf z}_{2}dt
\end{eqnarray}%
so that eq. $\left( \ref{Avar2}\right) $ becomes 
\begin{eqnarray}
\delta ^{\left( 2\right) }{\cal A}\left( \varkappa \right) &=&\frac{1}{2}%
\int_{t_{i}}^{t_{f}}\left\{ \left\{ \frac{\delta ^{2}{\cal L}\left( {\cal %
\varkappa },{\cal \dot{\varkappa}},t\right) }{\delta \varkappa \left( {\sf z}%
_{1}\right) \delta \varkappa \left( {\sf z}_{2}\right) }-\frac{d}{dt}\left\{ 
\frac{\delta ^{2}{\cal L}\left( \varkappa ,\dot{\varkappa},t\right) }{\delta
\varkappa \left( {\sf z}_{1}\right) \delta \dot{\varkappa}\left( {\sf z}%
_{2}\right) }\right\} \right\} \delta \varkappa \left( {\sf z}_{1},t\right)
\delta \varkappa \left( {\sf z}_{2},t\right) \right.  \nonumber \\
&&\qquad \qquad \qquad \qquad \qquad +\left. \frac{\delta ^{2}{\cal L}\left( 
{\cal \varkappa },{\cal \dot{\varkappa}},t\right) }{\delta \dot{\varkappa}%
\left( {\sf z}_{1}\right) \delta \dot{\varkappa}\left( {\sf z}_{2}\right) }%
\delta \dot{\varkappa}\left( {\sf z}_{1},t\right) \delta \dot{\varkappa}%
\left( {\sf z}_{2},t\right) \right\} d{\sf z}_{1}d{\sf z}_{2}dt
\label{Avar2a}
\end{eqnarray}%
It should be possible to remove the explicit dependence on $\delta \dot{%
\varkappa}$ in the second order variation of the action as we argued that
the value of $\dot{\varkappa}$ is determined by $\varkappa $ over the
interval $\left( t_{i},t_{f}\right) $ To this end one can use integration by
parts on the third term of the Lagrangian to obtain 
\begin{eqnarray}
&&\int_{t_{i}}^{t_{f}}\frac{\delta ^{2}{\cal L}\left( {\cal \varkappa },%
{\cal \dot{\varkappa}},t\right) }{\delta \dot{\varkappa}\left( {\sf z}%
_{1}\right) \delta \dot{\varkappa}\left( {\sf z}_{2}\right) }\delta \dot{%
\varkappa}\left( {\sf z}_{1},t\right) \delta \dot{\varkappa}\left( {\sf z}%
_{2},t\right) d{\sf z}_{1}d{\sf z}_{2}dt=  \nonumber \\
&&\int \left\{ \left. \delta \varkappa \left( {\sf z}_{1},t\right) \delta 
\dot{\varkappa}\left( {\sf z}_{2},t\right) \frac{\delta ^{2}{\cal L}\left( 
{\cal \varkappa },{\cal \dot{\varkappa}},t\right) }{\delta \dot{\varkappa}%
\left( {\sf z}_{1}\right) \delta \dot{\varkappa}\left( {\sf z}_{2}\right) }%
\right| _{t_{i}}^{t_{f}}\right\} d{\sf z}_{1}d{\sf z}_{2}  \nonumber \\
&&\qquad \qquad \qquad \qquad \qquad -\int_{t_{i}}^{t_{f}}\frac{d}{dt}%
\left\{ \frac{\delta ^{2}{\cal L}\left( \varkappa ,\dot{\varkappa},t\right) 
}{\delta \dot{\varkappa}\left( {\sf z}_{1}\right) \delta \dot{\varkappa}%
\left( {\sf z}_{2}\right) }\right\} \delta \varkappa \left( {\sf z}%
_{1},t\right) \delta \dot{\varkappa}\left( {\sf z}_{2},t\right) d{\sf z}_{1}d%
{\sf z}_{2}dt  \nonumber \\
&&\left. {}\right. \qquad \qquad \qquad \qquad \qquad \qquad \qquad
=-\int_{t_{i}}^{t_{f}}\frac{d}{dt}\left\{ \frac{\delta ^{2}{\cal L}\left(
\varkappa ,\dot{\varkappa},t\right) }{\delta \dot{\varkappa}\left( {\sf z}%
_{1}\right) \delta \dot{\varkappa}\left( {\sf z}_{2}\right) }\right\} \delta
\varkappa \left( {\sf z}_{1},t\right) \delta \dot{\varkappa}\left( {\sf z}%
_{2},t\right) d{\sf z}_{1}d{\sf z}_{2}dt
\end{eqnarray}%
Then we can use integration by parts again to obtain 
\begin{eqnarray}
&&-\int_{t_{i}}^{t_{f}}\frac{d}{dt}\left\{ \frac{\delta ^{2}{\cal L}\left(
\varkappa ,\dot{\varkappa},t\right) }{\delta \dot{\varkappa}\left( {\sf z}%
_{1}\right) \delta \dot{\varkappa}\left( {\sf z}_{2}\right) }\right\} \delta
\varkappa \left( {\sf z}_{1},t\right) \delta \dot{\varkappa}\left( {\sf z}%
_{2},t\right) d{\sf z}_{1}d{\sf z}_{2}dt  \nonumber \\
&&\left. {}\right. =-\int \left\{ \left. \delta \varkappa \left( {\sf z}%
_{1},t\right) \delta \varkappa \left( {\sf z}_{2},t\right) \frac{\delta ^{2}%
{\cal L}\left( {\cal \varkappa },{\cal \dot{\varkappa}},t\right) }{\delta 
\dot{\varkappa}\left( {\sf z}_{1}\right) \delta \dot{\varkappa}\left( {\sf z}%
_{2}\right) }\right| _{t_{i}}^{t_{f}}\right\} d{\sf z}_{1}d{\sf z}_{2} 
\nonumber \\
&&\qquad \qquad \qquad \qquad \qquad +\int_{t_{i}}^{t_{f}}\frac{d^{2}}{dt^{2}%
}\left\{ \frac{\delta ^{2}{\cal L}\left( \varkappa ,\dot{\varkappa},t\right) 
}{\delta \dot{\varkappa}\left( {\sf z}_{1}\right) \delta \dot{\varkappa}%
\left( {\sf z}_{2}\right) }\right\} \delta \varkappa \left( {\sf z}%
_{1},t\right) \delta \varkappa \left( {\sf z}_{2},t\right) d{\sf z}_{1}d{\sf %
z}_{2}dt  \nonumber \\
&&\left. {}\right. \qquad \qquad \qquad \qquad \qquad \qquad \qquad
=\int_{t_{i}}^{t_{f}}\frac{d^{2}}{dt^{2}}\left\{ \frac{\delta ^{2}{\cal L}%
\left( \varkappa ,\dot{\varkappa},t\right) }{\delta \dot{\varkappa}\left( 
{\sf z}_{1}\right) \delta \dot{\varkappa}\left( {\sf z}_{2}\right) }\right\}
\delta \varkappa \left( {\sf z}_{1},t\right) \delta \varkappa \left( {\sf z}%
_{2},t\right) d{\sf z}_{1}d{\sf z}_{2}dt
\end{eqnarray}%
Hence 
\begin{eqnarray}
\delta ^{\left( 2\right) }{\cal A}\left( \varkappa \right) &=&\frac{1}{2}%
\int_{t_{i}}^{t_{f}}\left\{ \frac{\delta ^{2}{\cal L}\left( {\cal \varkappa }%
,{\cal \dot{\varkappa}},t\right) }{\delta \varkappa \left( {\sf z}%
_{1}\right) \delta \varkappa \left( {\sf z}_{2}\right) }-\frac{d}{dt}\left\{ 
\frac{\delta ^{2}{\cal L}\left( \varkappa ,\dot{\varkappa},t\right) }{\delta
\varkappa \left( {\sf z}_{1}\right) \delta \dot{\varkappa}\left( {\sf z}%
_{2}\right) }\right\} \right.  \nonumber \\
&&\qquad \qquad \qquad \qquad \qquad \left. +\frac{d^{2}}{dt^{2}}\left\{ 
\frac{\delta ^{2}{\cal L}\left( \varkappa ,\dot{\varkappa},t\right) }{\delta 
\dot{\varkappa}\left( {\sf z}_{1}\right) \delta \dot{\varkappa}\left( {\sf z}%
_{2}\right) }\right\} \right\} \delta \varkappa \left( {\sf z}_{1},t\right)
\delta \varkappa \left( {\sf z}_{2},t\right) d{\sf z}_{1}d{\sf z}_{2}dt
\label{Avar2b}
\end{eqnarray}
In the case we are dealing with viz Lagrangians linear in the velocity eq. $%
\left( \ref{Avar2a}\right) $ does not contain any $\delta \dot{\varkappa}$
dependence as $\left\{ \frac{\delta ^{2}{\cal L}\left( \varkappa ,\dot{%
\varkappa},t\right) }{\delta \dot{\varkappa}\left( {\sf z}_{1}\right) \delta 
\dot{\varkappa}\left( {\sf z}_{2}\right) }\right\} =0\forall {\sf z}_{1},%
{\sf z}_{2}$ and thus 
\begin{equation}
\delta ^{\left( 2\right) }{\cal A}\left( \varkappa \right) =\frac{1}{2}%
\int_{t_{i}}^{t_{f}}\left\{ \left\{ \frac{\delta ^{2}{\cal L}\left( {\cal %
\varkappa },{\cal \dot{\varkappa}},t\right) }{\delta \varkappa \left( {\sf z}%
_{1}\right) \delta \varkappa \left( {\sf z}_{2}\right) }-\frac{d}{dt}\left\{ 
\frac{\delta ^{2}{\cal L}\left( \varkappa ,\dot{\varkappa},t\right) }{\delta
\varkappa \left( {\sf z}_{1}\right) \delta \dot{\varkappa}\left( {\sf z}%
_{2}\right) }\right\} \right\} \delta \varkappa \left( {\sf z}_{1},t\right)
\delta \varkappa \left( {\sf z}_{2},t\right) \right\} d{\sf z}_{1}d{\sf z}%
_{2}dt
\end{equation}%
and 
\begin{equation}
\frac{\delta ^{2}{\cal \dot{A}}\left( \varkappa \left( t\right) \right) }{%
\delta \varkappa \left( {\sf z}_{1}\right) \delta \varkappa \left( {\sf z}%
_{2}\right) }=\frac{\delta ^{2}{\cal L}\left( {\cal \varkappa },{\cal \dot{%
\varkappa}},t\right) }{\delta \varkappa \left( {\sf z}_{1}\right) \delta
\varkappa \left( {\sf z}_{2}\right) }-\frac{d}{dt}\left\{ \frac{\delta ^{2}%
{\cal L}\left( \varkappa ,\dot{\varkappa},t\right) }{\delta \varkappa \left( 
{\sf z}_{1}\right) \delta \dot{\varkappa}\left( {\sf z}_{2}\right) }\right\}
\end{equation}

\subsection{Stationary Action}

The Action is stationary if 
\begin{equation}
\delta {\cal A}\left( \varkappa \right) =\int_{t_{i}}^{t_{f}}\delta ^{\left(
1\right) }{\cal L}\left( \varkappa ,\dot{\varkappa},t\right)
dt=\int_{t_{i}}^{t_{f}}\delta ^{\left( 1\right) }{\cal \dot{A}}\left(
\varkappa \left( t\right) \right) dt=0\,\text{ }
\end{equation}%
which implies that 
\begin{eqnarray}
\delta {\cal A}\left( \varkappa \right) &=&\int_{t_{i}}^{t_{f}}\left\{ \frac{%
\delta {\cal L}\left( {\cal \varkappa },{\cal \dot{\varkappa}},t\right) }{%
\delta {\cal \varkappa }\left( {\sf z}\right) }-\frac{d}{dt}\left\{ \frac{%
\delta {\cal L}\left( {\cal \varkappa },{\cal \dot{\varkappa}},t\right) }{%
\delta {\cal \dot{\varkappa}}\left( {\sf z}\right) }\right\} \right\} \delta
\varkappa \left( {\sf z},t\right) dt=0\quad \forall \delta \varkappa
\label{vara1} \\
&\equiv &\int_{t_{i}}^{t_{f}}\frac{\delta {\cal \dot{A}}\left( \varkappa
\left( t\right) \right) }{\delta {\cal \varkappa }\left( {\sf z}\right) }%
\delta \varkappa \left( {\sf z},t\right) dt=0\quad \forall \delta \varkappa
\label{vara2}
\end{eqnarray}%
By the fundamental lemma of the calculus of variations the eqs. $\left( \ref%
{vara1}\right) $ and $\left( \ref{vara2}\right) $ gives the Euler-Lagrange
equation%
\begin{equation}
\frac{\delta {\cal \dot{A}}\left( \varkappa \left( t\right) \right) }{\delta 
{\cal \varkappa }\left( {\sf z}\right) }=\frac{\delta {\cal L}\left( {\cal %
\varkappa },{\cal \dot{\varkappa}},t\right) }{\delta {\cal \varkappa }\left( 
{\sf z}\right) }-\frac{d}{dt}\left\{ \frac{\delta {\cal L}\left( {\cal %
\varkappa },{\cal \dot{\varkappa}},t\right) }{\delta {\cal \dot{\varkappa}}%
\left( {\sf z}\right) }\right\} =0
\end{equation}%
which gives necessary and sufficient conditions for stationary action. {\bf %
Note it is }$\frac{\delta {\cal \dot{A}}\left( \varkappa \left( t\right)
\right) }{\delta {\cal \varkappa }\left( {\sf z}\right) }=0$ {\bf that is
the stationary condition NOT }$\frac{\delta {\cal L}\left( {\cal \varkappa },%
{\cal \dot{\varkappa}},t\right) }{\delta {\cal \varkappa }\left( {\sf z}%
\right) }=0$.

$\lceil ${\bf A note on notation: }The explicit display of $``t"$ in ${\cal L%
}\left( {\cal \varkappa },{\cal \dot{\varkappa}},t\right) $ indicates that
not all of the time dependence of ${\cal L}$ is due to the time dependent
functions $\left\{ {\cal \varkappa },{\cal \dot{\varkappa}}\right\} $. This
will occur if the Hamiltonian is time dependent e.g. depends on time
dependent external potentials i.e. describes an open system. In most cases
however we consider only closed systems as we describe all sources of fields
to be part of the system. The parameter functions, ${\cal \varkappa }$, are
functions of the spatial and spin coordinates ${\sf z}$ and time coordinate $%
t$ as described above$\rfloor $

Our Lagrangian is 
\begin{equation}
{\cal L}\left( {\cal \varkappa }\left( t\right) ,{\cal \dot{\varkappa}}%
\left( t\right) \right) ={\cal N}\left( {\cal \varkappa }\left( t\right)
\right) ^{-1}\left\{ {\cal T}\left( {\cal \varkappa }\left( t\right) ,{\cal 
\dot{\varkappa}}\left( t\right) \right) -{\cal H}\left( {\cal \varkappa }%
\left( t\right) \right) \right\} =\int {\cal \dot{\varkappa}}\left( {\sf z}%
,t\right) {\cal Z}\left( {\sf z},t\right) d{\sf z}-{\cal E}\left( {\cal %
\varkappa }\left( t\right) \right)
\end{equation}%
where 
\begin{eqnarray}
{\cal N}\left( {\cal \varkappa }\left( t\right) \right) &=&\left\langle \Psi
\left( {\cal \varkappa }\left( t\right) \right) |\Psi \left( {\cal \varkappa 
}\left( t\right) \right) \right\rangle \\
{\cal T}\left( {\cal \varkappa }\left( t\right) ,{\cal \dot{\varkappa}}%
\left( t\right) \right) &=&\left\{ \left\langle \Psi \left( {\cal \varkappa }%
\left( t\right) \right) |\dot{\Psi}\left( {\cal \varkappa }\left( t\right)
\right) \right\rangle -\left\langle \dot{\Psi}\left( {\cal \varkappa }\left(
t\right) \right) |\Psi \left( {\cal \varkappa }\left( t\right) \right)
\right\rangle \right\} =\int {\cal T}\left( {\sf z},t\right) {\cal \dot{%
\varkappa}}\left( {\sf z},t\right) d{\sf z} \\
{\cal T}\left( {\sf z},t\right) &=&\frac{i}{2}\left\{ \left\langle \Psi
\left( {\cal \varkappa }\left( t\right) \right) \left| \frac{\delta \Psi
\left( {\cal \varkappa }\left( t\right) \right) }{\delta {\cal \varkappa }%
\left( {\sf z}\right) }\right. \right\rangle -\left\langle \left. \frac{%
\delta \Psi \left( {\cal \varkappa }\left( t\right) \right) }{\delta {\cal %
\varkappa }\left( {\sf z}\right) }\right| \Psi \left( {\cal \varkappa }%
\left( t\right) \right) \right\rangle \right\} \\
{\cal Z}\left( {\sf z},t\right) &=&\frac{i}{2}\left\langle \Psi \left( {\cal %
\varkappa }\left( t\right) \right) |\Psi \left( {\cal \varkappa }\left(
t\right) \right) \right\rangle ^{-1}\left\{ \left\langle \Psi \left( {\cal %
\varkappa }\left( t\right) \right) \left| \frac{\delta \Psi \left( {\cal %
\varkappa }\left( t\right) \right) }{\delta {\cal \varkappa }\left( {\sf z}%
\right) }\right. \right\rangle -\left\langle \left. \frac{\delta \Psi \left( 
{\cal \varkappa }\left( t\right) \right) }{\delta {\cal \varkappa }\left( 
{\sf z}\right) }\right| \Psi \left( {\cal \varkappa }\left( t\right) \right)
\right\rangle \right\} \\
{\cal Z}\left( {\sf z},t\right) &=&\frac{{\cal T}\left( {\sf z},t\right) }{%
{\cal N}\left( {\cal \varkappa }\left( t\right) \right) } \\
{\cal E}\left( {\cal \varkappa }\left( t\right) \right) &=&\frac{%
\left\langle \Psi \left( {\cal \varkappa }\left( t\right) \right) |H\Psi
\left( {\cal \varkappa }\left( t\right) \right) \right\rangle }{\left\langle
\Psi \left( {\cal \varkappa }\left( t\right) \right) |\Psi \left( {\cal %
\varkappa }\left( t\right) \right) \right\rangle }=\frac{{\cal H}\left( 
{\cal \varkappa }\left( t\right) \right) }{{\cal N}\left( {\cal \varkappa }%
\left( t\right) \right) } \\
{\cal H}\left( {\cal \varkappa }\left( t\right) \right) &=&\left\langle \Psi
\left( {\cal \varkappa }\left( t\right) \right) |H\Psi \left( {\cal %
\varkappa }\left( t\right) \right) \right\rangle
\end{eqnarray}%
In order to calculate the first derrivative of the action we need $\frac{%
\delta {\cal L}\left( {\cal \varkappa },{\cal \dot{\varkappa}},t\right) }{%
\delta {\cal \varkappa }\left( {\sf z}\right) }$ and $\frac{d}{dt}\left\{ 
\frac{\delta {\cal L}\left( {\cal \varkappa },{\cal \dot{\varkappa}}%
,t\right) }{\delta {\cal \dot{\varkappa}}\left( {\sf z}\right) }\right\} $,
we first note that 
\begin{equation}
{\cal L}\left( {\cal \varkappa },{\cal \dot{\varkappa}},t\right) =\frac{%
{\cal T}\left( {\cal \varkappa }\left( t\right) ,{\cal \dot{\varkappa}}%
\left( t\right) \right) -{\cal H}\left( {\cal \varkappa }\left( t\right)
\right) }{{\cal N}\left( {\cal \varkappa }\left( t\right) \right) }
\end{equation}%
which leads to 
\begin{equation}
\frac{\delta {\cal L}\left( {\cal \varkappa },{\cal \dot{\varkappa}}%
,t\right) }{\delta {\cal \varkappa }\left( {\sf z}_{1}\right) }=\frac{1}{%
{\cal N}\left( {\cal \varkappa }\left( t\right) \right) }\left\{ \frac{%
\delta \left\{ {\cal T}\left( {\cal \varkappa }\left( t\right) ,{\cal \dot{%
\varkappa}}\left( t\right) \right) -{\cal H}\left( {\cal \varkappa }\left(
t\right) \right) \right\} }{\delta {\cal \varkappa }\left( {\sf z}%
_{1}\right) }-{\cal L}\left( {\cal \varkappa },{\cal \dot{\varkappa}}%
,t\right) \frac{\delta {\cal N}\left( {\cal \varkappa }\left( t\right)
\right) }{\delta {\cal \varkappa }\left( {\sf z}_{1}\right) }\right\}
\end{equation}%
and 
\begin{eqnarray}
\frac{d}{dt}\left\{ \frac{\delta {\cal L}\left( {\cal \varkappa },{\cal \dot{%
\varkappa}},t\right) }{\delta {\cal \dot{\varkappa}}\left( {\sf z}%
_{1}\right) }\right\} &=&\frac{d{\cal Z}\left( {\sf z}_{1},t\right) }{dt}=%
\frac{d}{dt}\left\{ \frac{{\cal T}\left( {\sf z}_{1},t\right) }{{\cal N}%
\left( {\cal \varkappa }\left( t\right) \right) }\right\}  \nonumber \\
&=&\frac{1}{{\cal N}\left( {\cal \varkappa }\left( t\right) \right) }\left\{ 
\frac{d{\cal T}\left( {\sf z}_{1},t\right) }{dt}-{\cal Z}\left( {\sf z}%
_{1},t\right) \frac{d{\cal N}\left( {\cal \varkappa }\left( t\right) \right) 
}{dt}\right\}
\end{eqnarray}%
thus 
\begin{eqnarray}
\frac{\delta {\cal \dot{A}}\left( \varkappa \left( t\right) \right) }{\delta
\varkappa \left( {\sf z}_{1}\right) } &=&\frac{\delta {\cal L}\left( {\cal %
\varkappa },{\cal \dot{\varkappa}},t\right) }{\delta {\cal \varkappa }\left( 
{\sf z}_{1}\right) }-\frac{d}{dt}\left\{ \frac{\delta {\cal L}\left( {\cal %
\varkappa },{\cal \dot{\varkappa}},t\right) }{\delta {\cal \dot{\varkappa}}%
\left( {\sf z}_{1}\right) }\right\} =\frac{\delta {\cal L}\left( {\cal %
\varkappa },{\cal \dot{\varkappa}},t\right) }{\delta {\cal \varkappa }\left( 
{\sf z}_{1}\right) }-\frac{d{\cal Z}\left( {\sf z}_{1},t\right) }{dt} 
\nonumber \\
&=&\frac{1}{{\cal N}\left( {\cal \varkappa }\left( t\right) \right) }\left\{ 
\frac{\delta \left\{ {\cal T}\left( {\cal \varkappa }\left( t\right) ,{\cal 
\dot{\varkappa}}\left( t\right) \right) -{\cal H}\left( {\cal \varkappa }%
\left( t\right) \right) \right\} }{\delta {\cal \varkappa }\left( {\sf z}%
_{1}\right) }-{\cal L}\left( {\cal \varkappa },{\cal \dot{\varkappa}}%
,t\right) \frac{\delta {\cal N}\left( {\cal \varkappa }\left( t\right)
\right) }{\delta {\cal \varkappa }\left( {\sf z}_{1}\right) }\right\} 
\nonumber \\
&&-\frac{1}{{\cal N}\left( {\cal \varkappa }\left( t\right) \right) }\left\{ 
\frac{d{\cal T}\left( {\sf z}_{1},t\right) }{dt}-{\cal Z}\left( {\sf z}%
_{1},t\right) \frac{d{\cal N}\left( {\cal \varkappa }\left( t\right) \right) 
}{dt}\right\}
\end{eqnarray}%
We can also express the first derrivative of ${\cal \dot{A}}\left( \varkappa
\left( t\right) \right) $ in a more compact form as 
\begin{eqnarray}
\frac{\delta {\cal \dot{A}}\left( \varkappa \left( t\right) \right) }{\delta
\varkappa \left( {\sf z}_{1}\right) } &=&\int {\cal \dot{\varkappa}}\left( 
{\sf z},t\right) \left\{ \frac{\delta {\cal Z}\left( {\sf z},t\right) }{%
\delta \varkappa \left( {\sf z}_{1}\right) }-\frac{\delta {\cal Z}\left( 
{\sf z}_{1},t\right) }{\delta {\cal \varkappa }\left( {\sf z}\right) }%
\right\} d{\sf z}-\frac{\delta {\cal E}\left( {\cal \varkappa }\left(
t\right) \right) }{\delta \varkappa \left( {\sf z}_{1}\right) } \\
&\equiv &\int \eta \left( {\sf z}_{1},{\sf z},t\right) {\cal \dot{\varkappa}}%
\left( {\sf z},t\right) d{\sf z}-\frac{\delta {\cal E}\left( {\cal \varkappa 
}\left( t\right) \right) }{\delta \varkappa \left( {\sf z}_{1}\right) }
\end{eqnarray}%
after making the substitution 
\begin{equation}
\frac{d{\cal Z}\left( {\sf z}_{1},t\right) }{dt}=\int {\cal \dot{\varkappa}}%
\left( {\sf z},t\right) \frac{\delta {\cal Z}\left( {\sf z}_{1},t\right) }{%
\delta {\cal \varkappa }\left( {\sf z}\right) }d{\sf z}
\end{equation}%
where the metric tensor field $\eta $ is defined as 
\begin{equation}
\eta \left( {\sf z}_{1},{\sf z},t\right) =\frac{\delta {\cal Z}\left( {\sf z}%
,t\right) }{\delta \varkappa \left( {\sf z}_{1}\right) }-\frac{\delta {\cal Z%
}\left( {\sf z}_{1},t\right) }{\delta {\cal \varkappa }\left( {\sf z}\right) 
}
\end{equation}%
The expression for the metric tensor can be further developed by observing
that 
\begin{equation}
\frac{\delta {\cal Z}\left( {\sf z},t\right) }{\delta \varkappa \left( {\sf z%
}_{1}\right) }=\frac{1}{{\cal N}\left( {\cal \varkappa }\left( t\right)
\right) }\left\{ \frac{\delta {\cal T}\left( {\sf z},t\right) }{\delta
\varkappa \left( {\sf z}_{1}\right) }-{\cal Z}\left( {\sf z},t\right) \frac{%
\delta {\cal N}\left( \varkappa \left( t\right) \right) }{\delta \varkappa
\left( {\sf z}_{1}\right) }\right\}
\end{equation}%
and 
\begin{eqnarray}
\frac{\delta {\cal T}\left( {\sf z},t\right) }{\delta \varkappa \left( {\sf z%
}_{1}\right) } &=&\frac{i}{2}\left\{ \left\langle \frac{\delta \Psi \left( 
{\cal \varkappa }\left( t\right) \right) }{\delta {\cal \varkappa }\left( 
{\sf z}_{1}\right) }\left| \frac{\delta \Psi \left( {\cal \varkappa }\left(
t\right) \right) }{\delta {\cal \varkappa }\left( {\sf z}\right) }\right.
\right\rangle -\left\langle \left. \frac{\delta \Psi \left( {\cal \varkappa }%
\left( t\right) \right) }{\delta {\cal \varkappa }\left( {\sf z}\right) }%
\right| \frac{\delta \Psi \left( {\cal \varkappa }\left( t\right) \right) }{%
\delta {\cal \varkappa }\left( {\sf z}_{1}\right) }\right\rangle \right. 
\nonumber \\
&&+\left. \left\langle \Psi \left( {\cal \varkappa }\left( t\right) \right)
\left| \frac{\delta ^{2}\Psi \left( {\cal \varkappa }\left( t\right) \right) 
}{\delta {\cal \varkappa }\left( {\sf z}_{1}\right) \delta {\cal \varkappa }%
\left( {\sf z}\right) }\right. \right\rangle -\left\langle \left. \frac{%
\delta ^{2}\Psi \left( {\cal \varkappa }\left( t\right) \right) }{\delta 
{\cal \varkappa }\left( {\sf z}_{1}\right) \delta {\cal \varkappa }\left( 
{\sf z}\right) }\right| \Psi \left( {\cal \varkappa }\left( t\right) \right)
\right\rangle \right\} \\
\frac{\delta {\cal N}\left( \varkappa \left( t\right) \right) }{\delta
\varkappa \left( {\sf z}_{1}\right) } &=&\left\langle \Psi \left( {\cal %
\varkappa }\left( t\right) \right) \left| \frac{\delta \Psi \left( {\cal %
\varkappa }\left( t\right) \right) }{\delta {\cal \varkappa }\left( {\sf z}%
_{1}\right) }\right. \right\rangle +\left\langle \left. \frac{\delta \Psi
\left( {\cal \varkappa }\left( t\right) \right) }{\delta {\cal \varkappa }%
\left( {\sf z}_{1}\right) }\right| \Psi \left( {\cal \varkappa }\left(
t\right) \right) \right\rangle
\end{eqnarray}%
and thus%
\begin{eqnarray}
&&\frac{\delta {\cal T}\left( \left( {\sf z},t\right) \right) }{\delta
\varkappa \left( {\sf z}_{1}\right) }-\frac{\delta {\cal T}\left( \left( 
{\sf z},t\right) \right) }{\delta \varkappa \left( {\sf z}\right) }= 
\nonumber \\
&&\left. {}\right. \qquad \qquad \qquad i\left\{ \left\langle \frac{\delta
\Psi \left( {\cal \varkappa }\left( t\right) \right) }{\delta {\cal %
\varkappa }\left( {\sf z}_{1}\right) }\left| \frac{\delta \Psi \left( {\cal %
\varkappa }\left( t\right) \right) }{\delta {\cal \varkappa }\left( {\sf z}%
\right) }\right. \right\rangle -\left\langle \left. \frac{\delta \Psi \left( 
{\cal \varkappa }\left( t\right) \right) }{\delta {\cal \varkappa }\left( 
{\sf z}\right) }\right| \frac{\delta \Psi \left( {\cal \varkappa }\left(
t\right) \right) }{\delta {\cal \varkappa }\left( {\sf z}_{1}\right) }%
\right\rangle \right\}  \nonumber \\
&&\left. {}\right. \qquad \qquad \qquad =-2%
%TCIMACRO{\func{Im}}%
%BeginExpansion
\mathop{\rm Im}%
%EndExpansion
\left\langle \frac{\delta \Psi \left( {\cal \varkappa }\left( t\right)
\right) }{\delta {\cal \varkappa }\left( {\sf z}_{1}\right) }\left| \frac{%
\delta \Psi \left( {\cal \varkappa }\left( t\right) \right) }{\delta {\cal %
\varkappa }\left( {\sf z}\right) }\right. \right\rangle
\end{eqnarray}%
and 
\begin{eqnarray}
&&{\cal Z}\left( {\sf z},t\right) \frac{\delta {\cal N}\left( \varkappa
\left( t\right) \right) }{\delta \varkappa \left( {\sf z}_{1}\right) }-{\cal %
Z}\left( {\sf z}_{1},t\right) \frac{\delta {\cal N}\left( \varkappa \left(
t\right) \right) }{\delta \varkappa \left( {\sf z}\right) }=  \nonumber \\
&&\frac{i}{2}\left\langle \Psi \left( {\cal \varkappa }\left( t\right)
\right) |\Psi \left( {\cal \varkappa }\left( t\right) \right) \right\rangle
^{-1}\left\{ \left\{ \left\langle \Psi \left( {\cal \varkappa }\left(
t\right) \right) \left| \frac{\delta \Psi \left( {\cal \varkappa }\left(
t\right) \right) }{\delta {\cal \varkappa }\left( {\sf z}\right) }\right.
\right\rangle -\left\langle \left. \frac{\delta \Psi \left( {\cal \varkappa }%
\left( t\right) \right) }{\delta {\cal \varkappa }\left( {\sf z}\right) }%
\right| \Psi \left( {\cal \varkappa }\left( t\right) \right) \right\rangle
\right\} \times \right.  \nonumber \\
&&\qquad \qquad \qquad \qquad \qquad \qquad \qquad \left\{ \left\langle \Psi
\left( {\cal \varkappa }\left( t\right) \right) \left| \frac{\delta \Psi
\left( {\cal \varkappa }\left( t\right) \right) }{\delta {\cal \varkappa }%
\left( {\sf z}_{1}\right) }\right. \right\rangle +\left\langle \left. \frac{%
\delta \Psi \left( {\cal \varkappa }\left( t\right) \right) }{\delta {\cal %
\varkappa }\left( {\sf z}_{1}\right) }\right| \Psi \left( {\cal \varkappa }%
\left( t\right) \right) \right\rangle \right\}  \nonumber \\
&&-\left\{ \left\langle \Psi \left( {\cal \varkappa }\left( t\right) \right)
\left| \frac{\delta \Psi \left( {\cal \varkappa }\left( t\right) \right) }{%
\delta {\cal \varkappa }\left( {\sf z}_{1}\right) }\right. \right\rangle
-\left\langle \left. \frac{\delta \Psi \left( {\cal \varkappa }\left(
t\right) \right) }{\delta {\cal \varkappa }\left( {\sf z}_{1}\right) }%
\right| \Psi \left( {\cal \varkappa }\left( t\right) \right) \right\rangle
\right\} \times  \nonumber \\
&&\qquad \qquad \qquad \qquad \qquad \qquad \qquad \left. \left\{
\left\langle \Psi \left( {\cal \varkappa }\left( t\right) \right) \left| 
\frac{\delta \Psi \left( {\cal \varkappa }\left( t\right) \right) }{\delta 
{\cal \varkappa }\left( {\sf z}\right) }\right. \right\rangle +\left\langle
\left. \frac{\delta \Psi \left( {\cal \varkappa }\left( t\right) \right) }{%
\delta {\cal \varkappa }\left( {\sf z}\right) }\right| \Psi \left( {\cal %
\varkappa }\left( t\right) \right) \right\rangle \right\} \right\}  \nonumber
\\
&=&i\left\langle \Psi \left( {\cal \varkappa }\left( t\right) \right) |\Psi
\left( {\cal \varkappa }\left( t\right) \right) \right\rangle ^{-1}\left\{
\left\langle \Psi \left( {\cal \varkappa }\left( t\right) \right) \left| 
\frac{\delta \Psi \left( {\cal \varkappa }\left( t\right) \right) }{\delta 
{\cal \varkappa }\left( {\sf z}\right) }\right. \right\rangle \left\langle
\left. \frac{\delta \Psi \left( {\cal \varkappa }\left( t\right) \right) }{%
\delta {\cal \varkappa }\left( {\sf z}_{1}\right) }\right| \Psi \left( {\cal %
\varkappa }\left( t\right) \right) \right\rangle \right.  \nonumber \\
&&\qquad \qquad \qquad \qquad \qquad \qquad \qquad -\left. \left\langle
\left. \frac{\delta \Psi \left( {\cal \varkappa }\left( t\right) \right) }{%
\delta {\cal \varkappa }\left( {\sf z}\right) }\right| \Psi \left( {\cal %
\varkappa }\left( t\right) \right) \right\rangle \left\langle \Psi \left( 
{\cal \varkappa }\left( t\right) \right) \left| \frac{\delta \Psi \left( 
{\cal \varkappa }\left( t\right) \right) }{\delta {\cal \varkappa }\left( 
{\sf z}_{1}\right) }\right. \right\rangle \right\}  \nonumber \\
&&\left. {}\right. \qquad \qquad \qquad \qquad \qquad =-\frac{2}{{\cal N}%
\left( {\cal \varkappa }\left( t\right) \right) }%
%TCIMACRO{\func{Im}}%
%BeginExpansion
\mathop{\rm Im}%
%EndExpansion
\left\langle \Psi \left( {\cal \varkappa }\left( t\right) \right) \left| 
\frac{\delta \Psi \left( {\cal \varkappa }\left( t\right) \right) }{\delta 
{\cal \varkappa }\left( {\sf z}\right) }\right. \right\rangle \left\langle
\left. \frac{\delta \Psi \left( {\cal \varkappa }\left( t\right) \right) }{%
\delta {\cal \varkappa }\left( {\sf z}_{1}\right) }\right| \Psi \left( {\cal %
\varkappa }\left( t\right) \right) \right\rangle
\end{eqnarray}%
giving the following form for the metric tensor 
\begin{eqnarray}
\eta \left( {\sf z}_{1},{\sf z},t\right) &=&-\frac{2}{{\cal N}\left( {\cal %
\varkappa }\left( t\right) \right) }%
%TCIMACRO{\func{Im}}%
%BeginExpansion
\mathop{\rm Im}%
%EndExpansion
\left\{ \left\langle \frac{\delta \Psi \left( {\cal \varkappa }\left(
t\right) \right) }{\delta {\cal \varkappa }\left( {\sf z}_{1}\right) }\left| 
\frac{\delta \Psi \left( {\cal \varkappa }\left( t\right) \right) }{\delta 
{\cal \varkappa }\left( {\sf z}\right) }\right. \right\rangle \right. 
\nonumber \\
&&\qquad \qquad \left. -\frac{1}{{\cal N}\left( {\cal \varkappa }\left(
t\right) \right) }\left\langle \Psi \left( {\cal \varkappa }\left( t\right)
\right) \left| \frac{\delta \Psi \left( {\cal \varkappa }\left( t\right)
\right) }{\delta {\cal \varkappa }\left( {\sf z}\right) }\right.
\right\rangle \left\langle \left. \frac{\delta \Psi \left( {\cal \varkappa }%
\left( t\right) \right) }{\delta {\cal \varkappa }\left( {\sf z}_{1}\right) }%
\right| \Psi \left( {\cal \varkappa }\left( t\right) \right) \right\rangle
\right\}  \label{metric}
\end{eqnarray}%
The energy derrivative can be further developed to give 
\begin{equation}
\frac{\delta {\cal E}\left( {\cal \varkappa }\left( t\right) \right) }{%
\delta \varkappa \left( {\sf z}_{1}\right) }=\frac{1}{{\cal N}\left( {\cal %
\varkappa }\left( t\right) \right) }\left\{ \frac{\delta {\cal H}\left(
\varkappa \left( t\right) \right) }{\delta \varkappa \left( {\sf z}%
_{1}\right) }-{\cal E}\left( {\cal \varkappa }\left( t\right) \right) \frac{%
\delta {\cal N}\left( \varkappa \left( t\right) \right) }{\delta \varkappa
\left( {\sf z}_{1}\right) }\right\}
\end{equation}%
where 
\begin{equation}
\frac{\delta {\cal H}\left( \varkappa \left( t\right) \right) }{\delta
\varkappa \left( {\sf z}_{1}\right) }=\left\langle \Psi \left( {\cal %
\varkappa }\left( t\right) \right) \left| H\frac{\delta \Psi \left( {\cal %
\varkappa }\left( t\right) \right) }{\delta \varkappa \left( {\sf z}%
_{1}\right) }\right. \right\rangle +\left\langle \left. \frac{\delta \Psi
\left( {\cal \varkappa }\left( t\right) \right) }{\delta \varkappa \left( 
{\sf z}_{1}\right) }\right| H\Psi \left( {\cal \varkappa }\left( t\right)
\right) \right\rangle
\end{equation}

\subsection{Minimum Condition}

The condition for a local minimum is 
\begin{eqnarray}
&&\delta ^{\left( 2\right) }{\cal A}\left( \varkappa \right)
=\int_{t_{i}}^{t_{f}}\delta ^{\left( 2\right) }{\cal \dot{A}}\left(
\varkappa \left( t\right) \right) dt=\int_{t_{i}}^{t_{f}}\delta ^{\left(
2\right) }{\cal L}\left( \varkappa ,\dot{\varkappa},t\right) dt=  \nonumber
\\
&&\frac{1}{2}\int_{t_{i}}^{t_{f}}\left\{ \left( 
\begin{array}{cc}
\delta \varkappa \left( {\sf z}_{1}\right) & \delta \dot{\varkappa}\left( 
{\sf z}_{2}\right)%
\end{array}%
\right) \left( 
\begin{array}{cc}
\frac{\delta ^{2}{\cal L}\left( \varkappa ,\dot{\varkappa},t\right) }{\delta
\varkappa \left( {\sf z}_{1}\right) \delta \varkappa \left( {\sf z}%
_{2}\right) } & \frac{\delta ^{2}{\cal L}\left( \varkappa ,\dot{\varkappa}%
,t\right) }{\delta \varkappa \left( {\sf z}_{1}\right) \delta \dot{\varkappa}%
\left( {\sf z}_{2}\right) } \\ 
\frac{\delta ^{2}{\cal L}\left( \varkappa ,\dot{\varkappa},t\right) }{\delta 
\dot{\varkappa}\left( {\sf z}_{1}\right) \delta \varkappa \left( {\sf z}%
_{2}\right) } & \frac{\delta ^{2}{\cal L}\left( \varkappa ,\dot{\varkappa}%
,t\right) }{\delta \dot{\varkappa}\left( {\sf z}_{1}\right) \delta \dot{%
\varkappa}\left( {\sf z}_{2}\right) }%
\end{array}%
\right) \left( 
\begin{array}{c}
\delta \varkappa \left( {\sf z}_{1}\right) \\ 
\delta \dot{\varkappa}\left( {\sf z}_{2}\right)%
\end{array}%
\right) \right\} d{\sf z}_{1}d{\sf z}_{2}dt\geq 0\forall \delta \varkappa
\end{eqnarray}%
which is the condition for a positive semi definite integral kernel. In our
case this condition becomes 
\begin{equation}
\left\{ \frac{\delta ^{2}{\cal \dot{A}}\left( \varkappa \left( t\right)
\right) }{\delta \varkappa \left( {\sf z}_{1}\right) \delta \varkappa \left( 
{\sf z}_{2}\right) }\right\} =\left\{ \frac{\delta ^{2}{\cal L}\left( {\cal %
\varkappa },{\cal \dot{\varkappa}},t\right) }{\delta \varkappa \left( {\sf z}%
_{1}\right) \delta \varkappa \left( {\sf z}_{2}\right) }-\frac{d}{dt}\left\{ 
\frac{\delta ^{2}{\cal L}\left( \varkappa ,\dot{\varkappa},t\right) }{\delta
\varkappa \left( {\sf z}_{1}\right) \delta \dot{\varkappa}\left( {\sf z}%
_{2}\right) }\right\} \right\} \geq 0  \label{ActHess}
\end{equation}%
as an operator on the space of functions $\left\{ \delta \varkappa \left( 
{\sf z},t\right) \right\} $. The second derrivative is given by: 
\begin{eqnarray}
\frac{\delta ^{2}{\cal \dot{A}}\left( \varkappa \left( t\right) \right) }{%
\delta \varkappa \left( {\sf z}_{2}\right) \delta \varkappa \left( {\sf z}%
_{1}\right) } &=&\int \frac{\delta \eta \left( {\sf z}_{1},{\sf z},t\right) 
}{\delta \varkappa \left( {\sf z}_{2}\right) }{\cal \dot{\varkappa}}\left( 
{\sf z},t\right) d{\sf z}-\frac{\delta ^{2}{\cal E}\left( {\cal \varkappa }%
\left( t\right) \right) }{\delta \varkappa \left( {\sf z}_{2}\right) \delta
\varkappa \left( {\sf z}_{1}\right) }  \label{ha} \\
&=&\frac{\delta ^{2}{\cal L}\left( {\cal \varkappa },{\cal \dot{\varkappa}}%
,t\right) }{\delta \varkappa \left( {\sf z}_{1}\right) \delta \varkappa
\left( {\sf z}_{2}\right) }-\frac{d}{dt}\left\{ \frac{\delta ^{2}{\cal L}%
\left( \varkappa ,\dot{\varkappa},t\right) }{\delta \varkappa \left( {\sf z}%
_{1}\right) \delta \dot{\varkappa}\left( {\sf z}_{2}\right) }\right\}
\label{hb}
\end{eqnarray}%
Eq. $\left( \ref{hb}\right) $ can be expresed as 
\begin{equation}
\frac{\delta ^{2}{\cal L}\left( {\cal \varkappa },{\cal \dot{\varkappa}}%
,t\right) }{\delta \varkappa \left( {\sf z}_{1}\right) \delta \varkappa
\left( {\sf z}_{2}\right) }-\frac{d}{dt}\left\{ \frac{\delta ^{2}{\cal L}%
\left( \varkappa ,\dot{\varkappa},t\right) }{\delta \varkappa \left( {\sf z}%
_{1}\right) \delta \dot{\varkappa}\left( {\sf z}_{2}\right) }\right\} =\frac{%
\delta ^{2}{\cal L}\left( {\cal \varkappa },{\cal \dot{\varkappa}},t\right) 
}{\delta \varkappa \left( {\sf z}_{1}\right) \delta \varkappa \left( {\sf z}%
_{2}\right) }-\frac{d}{dt}\left\{ \frac{\delta {\cal Z}\left( {\sf z}%
_{2},t\right) }{\delta \varkappa \left( {\sf z}_{1}\right) }\right\}
\end{equation}%
By using the chain rule on $\frac{\delta {\cal \dot{A}}\left( \varkappa
\left( t\right) \right) }{\delta \varkappa \left( {\sf z}_{1}\right) }$ in
the form 
\begin{equation}
\frac{\delta {\cal L}\left( {\cal \varkappa },{\cal \dot{\varkappa}}%
,t\right) }{\delta {\cal \varkappa }\left( {\sf z}_{2}\right) }-\frac{d{\cal %
Z}\left( {\sf z}_{2},t\right) }{dt}
\end{equation}%
we can also produce 
\begin{eqnarray}
&&\frac{\delta }{\delta \varkappa \left( {\sf z}_{2}\right) }\left\{ \frac{%
\delta {\cal L}\left( {\cal \varkappa },{\cal \dot{\varkappa}},t\right) }{%
\delta {\cal \varkappa }\left( {\sf z}_{1}\right) }-\frac{d{\cal Z}\left( 
{\sf z}_{1},t\right) }{dt}\right\} =\frac{\delta }{\delta \varkappa \left( 
{\sf z}_{2}\right) }\left\{ \frac{\delta {\cal L}\left( {\cal \varkappa },%
{\cal \dot{\varkappa}},t\right) }{\delta {\cal \varkappa }\left( {\sf z}%
_{1}\right) }-\int \frac{\delta {\cal Z}\left( {\sf z}_{1},t\right) }{\delta
\varkappa \left( {\sf z}\right) }{\cal \dot{\varkappa}}\left( {\sf z}%
,t\right) d{\sf z}\right\}  \nonumber \\
&=&\frac{\delta ^{2}{\cal L}\left( {\cal \varkappa },{\cal \dot{\varkappa}}%
,t\right) }{\delta \varkappa \left( {\sf z}_{2}\right) \delta {\cal %
\varkappa }\left( {\sf z}_{1}\right) }-\int \frac{\delta ^{2}{\cal Z}\left( 
{\sf z}_{1},t\right) }{\delta \varkappa \left( {\sf z}_{2}\right) \delta
\varkappa \left( {\sf z}\right) }{\cal \dot{\varkappa}}\left( {\sf z}%
,t\right) d{\sf z}  \label{hc} \\
&=&\int \frac{\delta ^{2}{\cal Z}\left( {\sf z},t\right) }{\delta \varkappa
\left( {\sf z}_{2}\right) \delta {\cal \varkappa }\left( {\sf z}_{1}\right) }%
{\cal \dot{\varkappa}}\left( {\sf z},t\right) d{\sf z-}\frac{\delta ^{2}%
{\cal E}\left( {\cal \varkappa }\left( t\right) \right) }{\delta \varkappa
\left( {\sf z}_{2}\right) \delta {\cal \varkappa }\left( {\sf z}_{1}\right) }%
-\int \frac{\delta ^{2}{\cal Z}\left( {\sf z}_{1},t\right) }{\delta
\varkappa \left( {\sf z}_{2}\right) \delta \varkappa \left( {\sf z}\right) }%
{\cal \dot{\varkappa}}\left( {\sf z},t\right) d{\sf z}  \nonumber \\
&=&\int \left\{ \frac{\delta ^{2}{\cal Z}\left( {\sf z},t\right) }{\delta
\varkappa \left( {\sf z}_{2}\right) \delta {\cal \varkappa }\left( {\sf z}%
_{1}\right) }-\frac{\delta ^{2}{\cal Z}\left( {\sf z}_{1},t\right) }{\delta
\varkappa \left( {\sf z}_{2}\right) \delta \varkappa \left( {\sf z}\right) }%
\right\} {\cal \dot{\varkappa}}\left( {\sf z},t\right) d{\sf z-}\frac{\delta
^{2}{\cal E}\left( {\cal \varkappa }\left( t\right) \right) }{\delta
\varkappa \left( {\sf z}_{2}\right) \delta {\cal \varkappa }\left( {\sf z}%
_{1}\right) }
\end{eqnarray}%
which is eq. $\left( \ref{ha}\right) $ again.

We will now consider eq.$\left( \ref{hb}\right) $%
\begin{eqnarray}
&&\frac{\delta ^{2}{\cal L}\left( {\cal \varkappa },{\cal \dot{\varkappa}}%
,t\right) }{\delta \varkappa \left( {\sf z}_{1}\right) \delta \varkappa
\left( {\sf z}_{2}\right) }  \nonumber \\
&&\left. {}\right. =\frac{\delta }{\delta \varkappa \left( {\sf z}%
_{2}\right) }\frac{1}{{\cal N}\left( {\cal \varkappa }\left( t\right)
\right) }\left\{ \frac{\delta \left\{ {\cal T}\left( {\cal \varkappa }\left(
t\right) ,{\cal \dot{\varkappa}}\left( t\right) \right) -{\cal H}\left( 
{\cal \varkappa }\left( t\right) \right) \right\} }{\delta {\cal \varkappa }%
\left( {\sf z}_{1}\right) }-{\cal L}\left( {\cal \varkappa },{\cal \dot{%
\varkappa}},t\right) \frac{\delta {\cal N}\left( {\cal \varkappa }\left(
t\right) \right) }{\delta {\cal \varkappa }\left( {\sf z}_{1}\right) }%
\right\} \\
&&\left. {}\right. =\frac{1}{{\cal N}\left( {\cal \varkappa }\left( t\right)
\right) }\left\{ \frac{\delta ^{2}\left\{ {\cal T}\left( {\cal \varkappa }%
\left( t\right) ,{\cal \dot{\varkappa}}\left( t\right) \right) -{\cal H}%
\left( {\cal \varkappa }\left( t\right) \right) \right\} }{\delta \varkappa
\left( {\sf z}_{2}\right) \delta {\cal \varkappa }\left( {\sf z}_{1}\right) }%
\right. -\frac{\delta {\cal L}\left( {\cal \varkappa },{\cal \dot{\varkappa}}%
,t\right) }{\delta \varkappa \left( {\sf z}_{2}\right) }\frac{\delta {\cal N}%
\left( {\cal \varkappa }\left( t\right) \right) }{\delta {\cal \varkappa }%
\left( {\sf z}_{1}\right) }  \nonumber \\
&&\qquad \qquad \qquad \qquad \qquad \qquad \left. -{\cal L}\left( {\cal %
\varkappa },{\cal \dot{\varkappa}},t\right) \frac{\delta ^{2}{\cal N}\left( 
{\cal \varkappa }\left( t\right) \right) }{\delta \varkappa \left( {\sf z}%
_{2}\right) \delta {\cal \varkappa }\left( {\sf z}_{1}\right) }-\frac{\delta 
{\cal L}\left( {\cal \varkappa },{\cal \dot{\varkappa}},t\right) }{\delta
\varkappa \left( {\sf z}_{1}\right) }\frac{\delta {\cal N}\left( {\cal %
\varkappa }\left( t\right) \right) }{\delta {\cal \varkappa }\left( {\sf z}%
_{2}\right) }\right\}  \label{hb1}
\end{eqnarray}%
and 
\begin{eqnarray}
&&\frac{d}{dt}\left\{ \frac{\delta ^{2}{\cal L}\left( \varkappa ,\dot{%
\varkappa},t\right) }{\delta \varkappa \left( {\sf z}_{1}\right) \delta \dot{%
\varkappa}\left( {\sf z}_{2}\right) }\right\} =\frac{d}{dt}\left\{ \frac{%
\delta {\cal Z}\left( {\sf z}_{2},t\right) }{\delta \varkappa \left( {\sf z}%
_{1}\right) }\right\} =\frac{d}{dt}\left\{ \frac{\delta }{\delta \varkappa
\left( {\sf z}_{1}\right) }\left\{ \frac{{\cal T}\left( {\sf z}_{2},t\right) 
}{{\cal N}\left( {\cal \varkappa }\left( t\right) \right) }\right\} \right\}
\nonumber \\
&&\left. {}\right. =\frac{d}{dt}\left\{ \frac{1}{{\cal N}\left( {\cal %
\varkappa }\left( t\right) \right) }\left\{ \frac{\delta {\cal T}\left( {\sf %
z}_{2},t\right) }{\delta \varkappa \left( {\sf z}_{1}\right) }-{\cal Z}%
\left( {\sf z}_{2},t\right) \frac{\delta {\cal N}\left( {\cal \varkappa }%
\left( t\right) \right) }{\delta \varkappa \left( {\sf z}_{1}\right) }%
\right\} \right\}  \nonumber \\
&&\left. {}\right. =\frac{1}{{\cal N}\left( {\cal \varkappa }\left( t\right)
\right) }\left\{ \frac{d}{dt}\frac{\delta {\cal T}\left( {\sf z}%
_{2},t\right) }{\delta \varkappa \left( {\sf z}_{1}\right) }-\frac{d{\cal Z}%
\left( {\sf z}_{2},t\right) }{dt}\frac{\delta {\cal N}\left( {\cal \varkappa 
}\left( t\right) \right) }{\delta \varkappa \left( {\sf z}_{1}\right) }%
\right.  \nonumber \\
&&\left. {}\right. \qquad \qquad \qquad \qquad \qquad \qquad \left. -{\cal Z}%
\left( {\sf z}_{2},t\right) \frac{d}{dt}\frac{\delta {\cal N}\left( {\cal %
\varkappa }\left( t\right) \right) }{\delta \varkappa \left( {\sf z}%
_{1}\right) }-\frac{\delta {\cal Z}\left( {\sf z}_{2},t\right) }{\delta
\varkappa \left( {\sf z}_{1}\right) }\frac{d{\cal N}\left( {\cal \varkappa }%
\left( t\right) \right) }{dt}\right\}  \label{hb2}
\end{eqnarray}%
rearranging and combining eqs. $\left( \ref{hb1}\right) $ and $\left( \ref%
{hb2}\right) $ one obtains 
\begin{eqnarray}
&&\frac{\delta ^{2}{\cal L}\left( {\cal \varkappa },{\cal \dot{\varkappa}}%
,t\right) }{\delta \varkappa \left( {\sf z}_{1}\right) \delta \varkappa
\left( {\sf z}_{2}\right) }-\frac{d}{dt}\left\{ \frac{\delta ^{2}{\cal L}%
\left( \varkappa ,\dot{\varkappa},t\right) }{\delta \varkappa \left( {\sf z}%
_{1}\right) \delta \dot{\varkappa}\left( {\sf z}_{2}\right) }\right\} = 
\nonumber \\
&&\frac{1}{{\cal N}\left( {\cal \varkappa }\left( t\right) \right) }\left\{ 
\frac{\delta ^{2}\left\{ {\cal T}\left( {\sf z},t\right) -{\cal H}\left( 
{\cal \varkappa }\left( t\right) \right) \right\} }{\delta \varkappa \left( 
{\sf z}_{2}\right) \delta {\cal \varkappa }\left( {\sf z}_{1}\right) }-{\cal %
L}\left( {\cal \varkappa },{\cal \dot{\varkappa}},t\right) \frac{\delta ^{2}%
{\cal N}\left( {\cal \varkappa }\left( t\right) \right) }{\delta \varkappa
\left( {\sf z}_{2}\right) \delta {\cal \varkappa }\left( {\sf z}_{1}\right) }%
\right.  \nonumber \\
&&\qquad \qquad \qquad \qquad \qquad \qquad \qquad \qquad -\frac{d}{dt}\frac{%
\delta {\cal T}\left( {\sf z}_{2},t\right) }{\delta \varkappa \left( {\sf z}%
_{1}\right) }+{\cal Z}\left( {\sf z}_{2},t\right) \frac{d}{dt}\frac{\delta 
{\cal N}\left( {\cal \varkappa }\left( t\right) \right) }{\delta \varkappa
\left( {\sf z}_{1}\right) }  \nonumber \\
&&+\frac{d{\cal Z}\left( {\sf z}_{2},t\right) }{dt}\frac{\delta {\cal N}%
\left( {\cal \varkappa }\left( t\right) \right) }{\delta \varkappa \left( 
{\sf z}_{1}\right) }-\frac{\delta {\cal L}\left( {\cal \varkappa },{\cal 
\dot{\varkappa}},t\right) }{\delta \varkappa \left( {\sf z}_{2}\right) }%
\frac{\delta {\cal N}\left( {\cal \varkappa }\left( t\right) \right) }{%
\delta {\cal \varkappa }\left( {\sf z}_{1}\right) }  \nonumber \\
&&\qquad \qquad \qquad \qquad \qquad \left. -\frac{\delta {\cal L}\left( 
{\cal \varkappa },{\cal \dot{\varkappa}},t\right) }{\delta \varkappa \left( 
{\sf z}_{1}\right) }\frac{\delta {\cal N}\left( {\cal \varkappa }\left(
t\right) \right) }{\delta {\cal \varkappa }\left( {\sf z}_{2}\right) }+\frac{%
\delta {\cal Z}\left( {\sf z}_{2},t\right) }{\delta \varkappa \left( {\sf z}%
_{1}\right) }\frac{d{\cal N}\left( {\cal \varkappa }\left( t\right) \right) 
}{dt}\right\}
\end{eqnarray}%
we will next consider eq. $\left( \ref{ha}\right) $ using the expression eq. 
$\left( \ref{metric}\right) $ for the metric we obtain 
\begin{eqnarray}
&&\frac{\delta \eta \left( {\sf z}_{1},{\sf z},t\right) }{\delta \varkappa
\left( {\sf z}_{2}\right) }=\frac{\delta }{\delta \varkappa \left( {\sf z}%
_{2}\right) }\left\{ -\frac{2}{{\cal N}\left( {\cal \varkappa }\left(
t\right) \right) }%
%TCIMACRO{\func{Im}}%
%BeginExpansion
\mathop{\rm Im}%
%EndExpansion
\left\{ \left\langle \frac{\delta \Psi \left( {\cal \varkappa }\left(
t\right) \right) }{\delta {\cal \varkappa }\left( {\sf z}_{1}\right) }\left| 
\frac{\delta \Psi \left( {\cal \varkappa }\left( t\right) \right) }{\delta 
{\cal \varkappa }\left( {\sf z}\right) }\right. \right\rangle \right. \right.
\nonumber \\
&&\qquad \qquad \qquad \qquad \left. \left. -\frac{1}{{\cal N}\left( {\cal %
\varkappa }\left( t\right) \right) }\left\langle \Psi \left( {\cal \varkappa 
}\left( t\right) \right) \left| \frac{\delta \Psi \left( {\cal \varkappa }%
\left( t\right) \right) }{\delta {\cal \varkappa }\left( {\sf z}\right) }%
\right. \right\rangle \left\langle \left. \frac{\delta \Psi \left( {\cal %
\varkappa }\left( t\right) \right) }{\delta {\cal \varkappa }\left( {\sf z}%
_{1}\right) }\right| \Psi \left( {\cal \varkappa }\left( t\right) \right)
\right\rangle \right\} \right\}
\end{eqnarray}%
which gives 
\begin{eqnarray}
&&-\frac{2}{{\cal N}\left( {\cal \varkappa }\left( t\right) \right) }%
%TCIMACRO{\func{Im}}%
%BeginExpansion
\mathop{\rm Im}%
%EndExpansion
\left\{ \left\langle \frac{\delta ^{2}\Psi \left( {\cal \varkappa }\left(
t\right) \right) }{\delta \varkappa \left( {\sf z}_{2}\right) \delta {\cal %
\varkappa }\left( {\sf z}_{1}\right) }\left| \frac{\delta \Psi \left( {\cal %
\varkappa }\left( t\right) \right) }{\delta {\cal \varkappa }\left( {\sf z}%
\right) }\right. \right\rangle +\left\langle \frac{\delta \Psi \left( {\cal %
\varkappa }\left( t\right) \right) }{\delta {\cal \varkappa }\left( {\sf z}%
_{1}\right) }\left| \frac{\delta ^{2}\Psi \left( {\cal \varkappa }\left(
t\right) \right) }{\delta \varkappa \left( {\sf z}_{2}\right) \delta {\cal %
\varkappa }\left( {\sf z}\right) }\right. \right\rangle \right.  \nonumber \\
&&-\frac{1}{{\cal N}\left( {\cal \varkappa }\left( t\right) \right) }\left\{
\left\langle \frac{\delta \Psi \left( {\cal \varkappa }\left( t\right)
\right) }{\delta \varkappa \left( {\sf z}_{2}\right) }\left| \frac{\delta
\Psi \left( {\cal \varkappa }\left( t\right) \right) }{\delta {\cal %
\varkappa }\left( {\sf z}\right) }\right. \right\rangle \left\langle \left. 
\frac{\delta \Psi \left( {\cal \varkappa }\left( t\right) \right) }{\delta 
{\cal \varkappa }\left( {\sf z}_{1}\right) }\right| \Psi \left( {\cal %
\varkappa }\left( t\right) \right) \right\rangle \right.  \nonumber \\
&&+\left\langle \Psi \left( {\cal \varkappa }\left( t\right) \right) \left| 
\frac{\delta ^{2}\Psi \left( {\cal \varkappa }\left( t\right) \right) }{%
\delta \varkappa \left( {\sf z}_{2}\right) \delta {\cal \varkappa }\left( 
{\sf z}\right) }\right. \right\rangle \left\langle \left. \frac{\delta \Psi
\left( {\cal \varkappa }\left( t\right) \right) }{\delta {\cal \varkappa }%
\left( {\sf z}_{1}\right) }\right| \Psi \left( {\cal \varkappa }\left(
t\right) \right) \right\rangle  \nonumber \\
&&+\left\langle \Psi \left( {\cal \varkappa }\left( t\right) \right) \left| 
\frac{\delta \Psi \left( {\cal \varkappa }\left( t\right) \right) }{\delta 
{\cal \varkappa }\left( {\sf z}\right) }\right. \right\rangle \left\langle
\left. \frac{\delta ^{2}\Psi \left( {\cal \varkappa }\left( t\right) \right) 
}{\delta \varkappa \left( {\sf z}_{2}\right) \delta {\cal \varkappa }\left( 
{\sf z}_{1}\right) }\right| \Psi \left( {\cal \varkappa }\left( t\right)
\right) \right\rangle  \nonumber \\
&&+\left\langle \Psi \left( {\cal \varkappa }\left( t\right) \right) \left| 
\frac{\delta \Psi \left( {\cal \varkappa }\left( t\right) \right) }{\delta 
{\cal \varkappa }\left( {\sf z}\right) }\right. \right\rangle \left\langle
\left. \frac{\delta \Psi \left( {\cal \varkappa }\left( t\right) \right) }{%
\delta {\cal \varkappa }\left( {\sf z}_{1}\right) }\right| \frac{\delta \Psi
\left( {\cal \varkappa }\left( t\right) \right) }{\delta \varkappa \left( 
{\sf z}_{2}\right) }\right\rangle  \nonumber \\
&&+\left\{ \frac{2}{{\cal N}\left( {\cal \varkappa }\left( t\right) \right) }%
\left\{ \left\langle \frac{\delta \Psi \left( {\cal \varkappa }\left(
t\right) \right) }{\delta {\cal \varkappa }\left( {\sf z}_{1}\right) }\left| 
\frac{\delta \Psi \left( {\cal \varkappa }\left( t\right) \right) }{\delta 
{\cal \varkappa }\left( {\sf z}\right) }\right. \right\rangle \right. \right.
\nonumber \\
&&-\left. \left. \frac{1}{{\cal N}\left( {\cal \varkappa }\left( t\right)
\right) }\left\langle \Psi \left( {\cal \varkappa }\left( t\right) \right)
\left| \frac{\delta \Psi \left( {\cal \varkappa }\left( t\right) \right) }{%
\delta {\cal \varkappa }\left( {\sf z}\right) }\right. \right\rangle
\left\langle \left. \frac{\delta \Psi \left( {\cal \varkappa }\left(
t\right) \right) }{\delta {\cal \varkappa }\left( {\sf z}_{1}\right) }%
\right| \Psi \left( {\cal \varkappa }\left( t\right) \right) \right\rangle
\right\} \frac{\delta {\cal N}\left( {\cal \varkappa }\left( t\right)
\right) }{\delta \varkappa \left( {\sf z}_{2}\right) }\right\}
\end{eqnarray}%
and 
\begin{eqnarray}
&&\frac{\delta ^{2}{\cal E}\left( {\cal \varkappa }\left( t\right) \right) }{%
\delta \varkappa \left( {\sf z}_{2}\right) \delta \varkappa \left( {\sf z}%
_{1}\right) }=\frac{\delta }{\delta \varkappa \left( {\sf z}_{2}\right) }%
\frac{1}{{\cal N}\left( {\cal \varkappa }\left( t\right) \right) }\left\{ 
\frac{\delta {\cal H}\left( \varkappa \left( t\right) \right) }{\delta
\varkappa \left( {\sf z}_{1}\right) }-{\cal E}\left( {\cal \varkappa }\left(
t\right) \right) \frac{\delta {\cal N}\left( \varkappa \left( t\right)
\right) }{\delta \varkappa \left( {\sf z}_{1}\right) }\right\}  \nonumber \\
&=&\frac{1}{{\cal N}\left( {\cal \varkappa }\left( t\right) \right) }\left\{ 
\frac{\delta ^{2}{\cal H}\left( {\cal \varkappa }\left( t\right) \right) }{%
\delta \varkappa \left( {\sf z}_{2}\right) \delta \varkappa \left( {\sf z}%
_{1}\right) }-{\cal E}\left( {\cal \varkappa }\left( t\right) \right) \frac{%
\delta ^{2}{\cal N}\left( {\cal \varkappa }\left( t\right) \right) }{\delta
\varkappa \left( {\sf z}_{2}\right) \delta \varkappa \left( {\sf z}%
_{1}\right) }\right. \\
&&\qquad \qquad \qquad \qquad \qquad \qquad \qquad \left. -\frac{\delta 
{\cal E}\left( \varkappa \left( t\right) \right) }{\delta \varkappa \left( 
{\sf z}_{2}\right) }\frac{\delta {\cal N}\left( \varkappa \left( t\right)
\right) }{\delta \varkappa \left( {\sf z}_{1}\right) }-\frac{\delta {\cal E}%
\left( \varkappa \left( t\right) \right) }{\delta \varkappa \left( {\sf z}%
_{1}\right) }\frac{\delta {\cal N}\left( \varkappa \left( t\right) \right) }{%
\delta \varkappa \left( {\sf z}_{2}\right) }\right\}
\end{eqnarray}%
hence 
\begin{eqnarray}
&&\int \frac{\delta \eta \left( {\sf z}_{1},{\sf z},t\right) }{\delta
\varkappa \left( {\sf z}_{2}\right) }{\cal \dot{\varkappa}}\left( {\sf z}%
,t\right) d{\sf z}-\frac{\delta ^{2}{\cal E}\left( {\cal \varkappa }\left(
t\right) \right) }{\delta \varkappa \left( {\sf z}_{2}\right) \delta
\varkappa \left( {\sf z}_{1}\right) } \\
&=&\int \frac{1}{{\cal N}\left( {\cal \varkappa }\left( t\right) \right) }%
\left\{ {\frak K}\left( {\sf z}_{1},{\sf z}_{2}\right) -{\frak K}\left( {\sf %
z},{\sf z}_{2}\right) -\frac{1}{{\cal N}\left( {\cal \varkappa }\left(
t\right) \right) }\frac{\delta {\cal Z}\left( {\sf z},t\right) }{\delta
\varkappa \left( {\sf z}_{1}\right) }\frac{\delta {\cal N}\left( {\cal %
\varkappa }\left( t\right) \right) }{\delta \varkappa \left( {\sf z}%
_{2}\right) }\right\} {\cal \dot{\varkappa}}\left( {\sf z},t\right) d{\sf z}
\\
&&-\frac{1}{{\cal N}\left( {\cal \varkappa }\left( t\right) \right) }\left\{ 
{\frak E}\left( {\sf z}_{1},{\sf z}_{2}\right) -\frac{1}{{\cal N}\left( 
{\cal \varkappa }\left( t\right) \right) }\frac{\delta {\cal E}\left(
\varkappa \left( t\right) \right) }{\delta \varkappa \left( {\sf z}%
_{1}\right) }\frac{\delta {\cal N}\left( \varkappa \left( t\right) \right) }{%
\delta \varkappa \left( {\sf z}_{2}\right) }\right\} \\
&=&\frac{1}{{\cal N}\left( {\cal \varkappa }\left( t\right) \right) }\left\{
\int \left\{ {\frak K}\left( {\sf z}_{1},{\sf z}_{2}\right) -{\frak K}\left( 
{\sf z},{\sf z}_{2}\right) \right\} {\cal \dot{\varkappa}}\left( {\sf z}%
,t\right) d{\sf z-}{\frak E}\left( {\sf z}_{1},{\sf z}_{2}\right) -\frac{1}{%
{\cal N}\left( {\cal \varkappa }\left( t\right) \right) }\frac{\delta {\cal 
\dot{A}}\left( {\sf z},t\right) }{\delta \varkappa \left( {\sf z}_{1}\right) 
}\frac{\delta {\cal N}\left( {\cal \varkappa }\left( t\right) \right) }{%
\delta \varkappa \left( {\sf z}_{2}\right) }\right\}
\end{eqnarray}%
\begin{equation}
\frac{\delta ^{2}{\cal L}\left( \varkappa ,\dot{\varkappa},t\right) }{\delta
\varkappa \left( {\sf z}_{2}\right) \delta \dot{\varkappa}\left( {\sf z}%
_{1}\right) }=\frac{\delta {\cal Z}\left( {\sf z}_{1},t\right) }{\delta
\varkappa \left( {\sf z}_{2}\right) }
\end{equation}

\bigskip

\begin{eqnarray}
&&\frac{\delta {\cal \dot{A}}\left( \varkappa \left( t\right) \right) }{%
\delta \varkappa \left( {\sf z}_{1}\right) }=\int {\cal \dot{\varkappa}}%
\left( {\sf z},t\right) \left\{ \frac{\delta {\cal Z}\left( {\sf z},t\right) 
}{\delta \varkappa \left( {\sf z}_{1}\right) }-\frac{\delta {\cal Z}\left( 
{\sf z}_{1},t\right) }{\delta {\cal \varkappa }\left( {\sf z}\right) }%
\right\} d{\sf z}-\frac{\delta {\cal E}\left( {\cal \varkappa }\left(
t\right) \right) }{\delta \varkappa \left( {\sf z}_{1}\right) }= \\
&&\int {\cal \dot{\varkappa}}\left( {\sf z},t\right) \frac{1}{{\cal N}\left( 
{\cal \varkappa }\left( t\right) \right) }\left\{ \frac{\delta {\cal T}%
\left( {\sf z},t\right) }{\delta \varkappa \left( {\sf z}_{1}\right) }-{\cal %
Z}\left( {\sf z},t\right) \frac{\delta {\cal N}\left( {\sf z},t\right) }{%
\delta \varkappa \left( {\sf z}_{1}\right) }-\frac{\delta {\cal T}\left( 
{\sf z}_{1},t\right) }{\delta \varkappa \left( {\sf z}\right) }-{\cal Z}%
\left( {\sf z}_{1},t\right) \frac{\delta {\cal N}\left( {\sf z}_{1},t\right) 
}{\delta \varkappa \left( {\sf z}\right) }\right\} d{\sf z} \\
&&\frac{-1}{{\cal N}\left( {\cal \varkappa }\left( t\right) \right) }\left\{ 
\frac{\delta {\cal H}\left( {\sf z},t\right) }{\delta \varkappa \left( {\sf z%
}_{1}\right) }-{\cal E}\left( {\sf z},t\right) \frac{\delta {\cal N}\left( 
{\sf z},t\right) }{\delta \varkappa \left( {\sf z}_{1}\right) }\right\} = \\
&&\frac{1}{{\cal N}\left( {\cal \varkappa }\left( t\right) \right) }\left\{
\int {\cal \dot{\varkappa}}\left( {\sf z},t\right) \left\{ \frac{\delta 
{\cal T}\left( {\sf z},t\right) }{\delta \varkappa \left( {\sf z}_{1}\right) 
}-{\cal Z}\left( {\sf z},t\right) \frac{\delta {\cal N}\left( {\sf z}%
,t\right) }{\delta \varkappa \left( {\sf z}_{1}\right) }-\frac{\delta {\cal T%
}\left( {\sf z}_{1},t\right) }{\delta \varkappa \left( {\sf z}\right) }-%
{\cal Z}\left( {\sf z}_{1},t\right) \frac{\delta {\cal N}\left( {\sf z}%
_{1},t\right) }{\delta \varkappa \left( {\sf z}\right) }\right\} d{\sf z}%
\right. \\
&&\left. -\left\{ \frac{\delta {\cal H}\left( {\sf z},t\right) }{\delta
\varkappa \left( {\sf z}_{1}\right) }-{\cal E}\left( {\sf z},t\right) \frac{%
\delta {\cal N}\left( {\sf z},t\right) }{\delta \varkappa \left( {\sf z}%
_{1}\right) }\right\} \right\}
\end{eqnarray}%
\begin{eqnarray}
&&\frac{\delta ^{2}{\cal \dot{A}}\left( \varkappa \left( t\right) \right) }{%
\delta \varkappa \left( {\sf z}_{2}\right) \delta \varkappa \left( {\sf z}%
_{1}\right) }=-\frac{1}{{\cal N}\left( {\cal \varkappa }\left( t\right)
\right) }\frac{\delta {\cal \dot{A}}\left( \varkappa \left( t\right) \right) 
}{\delta \varkappa \left( {\sf z}_{1}\right) }\frac{\delta {\cal N}\left( 
{\sf z},t\right) }{\delta \varkappa \left( {\sf z}_{2}\right) }+ \\
&&\frac{1}{{\cal N}\left( {\cal \varkappa }\left( t\right) \right) }\left\{
\int {\cal \dot{\varkappa}}\left( {\sf z},t\right) \left\{ \frac{\delta ^{2}%
{\cal T}\left( {\sf z},t\right) }{\delta \varkappa \left( {\sf z}_{2}\right)
\delta \varkappa \left( {\sf z}_{1}\right) }-\frac{\delta {\cal Z}\left( 
{\sf z},t\right) }{\delta \varkappa \left( {\sf z}_{2}\right) }\frac{\delta 
{\cal N}\left( {\sf z},t\right) }{\delta \varkappa \left( {\sf z}_{1}\right) 
}-{\cal Z}\left( {\sf z},t\right) \frac{\delta ^{2}{\cal N}\left( {\sf z}%
,t\right) }{\delta \varkappa \left( {\sf z}_{2}\right) \delta \varkappa
\left( {\sf z}_{1}\right) }\right. \right. \\
&&\left. -\frac{\delta ^{2}{\cal T}\left( {\sf z}_{1},t\right) }{\delta
\varkappa \left( {\sf z}_{2}\right) \delta \varkappa \left( {\sf z}\right) }+%
\frac{\delta {\cal Z}\left( {\sf z}_{1},t\right) }{\delta \varkappa \left( 
{\sf z}_{2}\right) }\frac{\delta {\cal N}\left( {\sf z}_{1},t\right) }{%
\delta \varkappa \left( {\sf z}\right) }+{\cal Z}\left( {\sf z}_{1},t\right) 
\frac{\delta ^{2}{\cal N}\left( {\sf z}_{1},t\right) }{\delta \varkappa
\left( {\sf z}_{2}\right) \delta \varkappa \left( {\sf z}\right) }\right\} d%
{\sf z} \\
&&-\left\{ \frac{\delta ^{2}{\cal H}\left( {\cal \varkappa }\left( t\right)
\right) }{\delta \varkappa \left( {\sf z}_{2}\right) \delta \varkappa \left( 
{\sf z}_{1}\right) }-\frac{\delta {\cal E}\left( {\cal \varkappa }\left(
t\right) \right) }{\delta \varkappa \left( {\sf z}_{2}\right) }\frac{\delta 
{\cal N}\left( {\cal \varkappa }\left( t\right) \right) }{\delta \varkappa
\left( {\sf z}_{1}\right) }-{\cal E}\left( {\sf z},t\right) \frac{\delta ^{2}%
{\cal N}\left( {\cal \varkappa }\left( t\right) \right) }{\delta \varkappa
\left( {\sf z}_{2}\right) \delta \varkappa \left( {\sf z}_{1}\right) }%
\right\} \\
&=&\frac{1}{{\cal N}\left( {\cal \varkappa }\left( t\right) \right) }\left\{ 
\frac{\delta {\cal \dot{A}}\left( \varkappa \left( t\right) \right) }{\delta
\varkappa \left( {\sf z}_{2}\right) }\frac{\delta {\cal N}\left( {\cal %
\varkappa }\left( t\right) \right) }{\delta \varkappa \left( {\sf z}%
_{1}\right) }-\frac{\delta {\cal \dot{A}}\left( \varkappa \left( t\right)
\right) }{\delta \varkappa \left( {\sf z}_{1}\right) }\frac{\delta {\cal N}%
\left( {\cal \varkappa }\left( t\right) \right) }{\delta \varkappa \left( 
{\sf z}_{2}\right) }\right. + \\
&&\int {\cal \dot{\varkappa}}\left( {\sf z},t\right) \left\{ \frac{\delta
^{2}{\cal T}\left( {\sf z},t\right) }{\delta \varkappa \left( {\sf z}%
_{2}\right) \delta \varkappa \left( {\sf z}_{1}\right) }-\frac{\delta ^{2}%
{\cal T}\left( {\sf z}_{1},t\right) }{\delta \varkappa \left( {\sf z}%
_{2}\right) \delta \varkappa \left( {\sf z}\right) }\right. \\
&&+\left. {\cal Z}\left( {\sf z}_{1},t\right) \frac{\delta ^{2}{\cal N}%
\left( {\sf z}_{1},t\right) }{\delta \varkappa \left( {\sf z}_{2}\right)
\delta \varkappa \left( {\sf z}\right) }-{\cal Z}\left( {\sf z},t\right) 
\frac{\delta ^{2}{\cal N}\left( {\sf z},t\right) }{\delta \varkappa \left( 
{\sf z}_{2}\right) \delta \varkappa \left( {\sf z}_{1}\right) }\right\} d%
{\sf z} \\
&&-\left. \left\{ \frac{\delta ^{2}{\cal H}\left( {\cal \varkappa }\left(
t\right) \right) }{\delta \varkappa \left( {\sf z}_{2}\right) \delta
\varkappa \left( {\sf z}_{1}\right) }-{\cal E}\left( {\sf z},t\right) \frac{%
\delta ^{2}{\cal N}\left( {\cal \varkappa }\left( t\right) \right) }{\delta
\varkappa \left( {\sf z}_{2}\right) \delta \varkappa \left( {\sf z}%
_{1}\right) }\right\} \right\}
\end{eqnarray}

\subsection{Many Variables}

In terms of the variables ${\cal V}_{fib},\omega _{fib},\omega _{den},{\cal V%
}_{den}$ the first order variation of the Lagrangian about the functions $%
{\cal V}_{fib},\omega _{fib},\omega _{den},{\cal V}_{den}$ is 
\begin{eqnarray}
\delta ^{\left( 1\right) }{\cal L}\left( {\cal V}_{fib},\omega _{fib},\omega
_{den},{\cal V}_{den}\right) &{\cal =}&  \nonumber \\
&&\int \left\{ \frac{\delta {\cal L}\left( \varkappa ,\dot{\varkappa}%
,t\right) }{\delta {\cal V}_{fib}\left( {\sf z}\right) }\delta {\cal V}%
_{fib}\left( {\sf z}\right) +\frac{\delta {\cal L}\left( \varkappa ,\dot{%
\varkappa},t\right) }{\delta {\cal \dot{V}}_{fib}\left( {\sf z}\right) }%
\delta {\cal \dot{V}}_{fib}\left( {\sf z}\right) \right\} d{\sf z}  \nonumber
\\
&&+\int \left\{ \frac{\delta {\cal L}\left( \varkappa ,\dot{\varkappa}%
,t\right) }{\delta {\cal \omega }_{fib}\left( {\sf z}\right) }\delta {\cal %
\omega }_{fib}\left( {\sf z}\right) +\frac{\delta {\cal L}\left( \varkappa ,%
\dot{\varkappa},t\right) }{\delta {\cal \dot{\omega}}_{fib}\left( {\sf z}%
\right) }\delta {\cal \dot{\omega}}_{fib}\left( {\sf z}\right) \right\} d%
{\sf z}  \nonumber \\
&&+\int \left\{ \frac{\delta {\cal L}\left( \varkappa ,\dot{\varkappa}%
,t\right) }{\delta {\cal \omega }_{den}\left( {\sf z}\right) }\delta {\cal %
\omega }_{den}\left( {\sf z}\right) +\frac{\delta {\cal L}\left( \varkappa ,%
\dot{\varkappa},t\right) }{\delta {\cal \dot{\omega}}_{den}\left( {\sf z}%
\right) }\delta {\cal \dot{\omega}}_{den}\left( {\sf z}\right) \right\} d%
{\sf z}  \nonumber \\
&&+\int \left\{ \frac{\delta {\cal L}\left( \varkappa ,\dot{\varkappa}%
,t\right) }{\delta {\cal V}_{den}\left( {\sf z}\right) }\delta {\cal V}%
_{den}\left( {\sf z}\right) +\frac{\delta {\cal L}\left( \varkappa ,\dot{%
\varkappa},t\right) }{\delta {\cal \dot{V}}_{den}\left( {\sf z}\right) }%
\delta {\cal \dot{V}}_{den}\left( {\sf z}\right) \right\} d{\sf z}
\end{eqnarray}%
and the stationarity condition becomes

\begin{eqnarray}
&&\delta _{{\cal V}_{fib}}{\cal A}=\int \left\{ \frac{\delta {\cal L}\left(
\varkappa ,\dot{\varkappa},t\right) }{\delta {\cal V}_{fib}\left( {\sf z}%
\right) }\delta {\cal V}_{fib}\left( {\sf z}\right) +\frac{\delta {\cal L}%
\left( \varkappa ,\dot{\varkappa},t\right) }{\delta {\cal \dot{V}}%
_{fib}\left( {\sf z}\right) }\delta {\cal \dot{V}}_{fib}\left( {\sf z}%
\right) \right\} =0  \nonumber \\
&&\delta _{{\cal \omega }_{fib}}{\cal A}=\int \left\{ \frac{\delta {\cal L}%
\left( \varkappa ,\dot{\varkappa},t\right) }{\delta {\cal \omega }%
_{fib}\left( {\sf z}\right) }\delta {\cal \omega }_{fib}\left( {\sf z}%
\right) +\frac{\delta {\cal L}\left( \varkappa ,\dot{\varkappa},t\right) }{%
\delta {\cal \dot{\omega}}_{fib}\left( {\sf z}\right) }\delta {\cal \dot{%
\omega}}_{fib}\left( {\sf z}\right) \right\} =0  \nonumber \\
&&\delta _{{\cal \omega }_{den}}{\cal A}=\int \left\{ \frac{\delta {\cal L}%
\left( \varkappa ,\dot{\varkappa},t\right) }{\delta {\cal \omega }%
_{den}\left( {\sf z}\right) }\delta {\cal \omega }_{den}\left( {\sf z}%
\right) +\frac{\delta {\cal L}\left( \varkappa ,\dot{\varkappa},t\right) }{%
\delta {\cal \dot{\omega}}_{den}\left( {\sf z}\right) }\delta {\cal \dot{%
\omega}}_{den}\left( {\sf z}\right) \right\} =0  \nonumber \\
&&\delta _{{\cal V}_{den}}{\cal A}=\int \left\{ \frac{\delta {\cal L}\left(
\varkappa ,\dot{\varkappa},t\right) }{\delta {\cal V}_{den}\left( {\sf z}%
\right) }\delta {\cal V}_{den}\left( {\sf z}\right) +\frac{\delta {\cal L}%
\left( \varkappa ,\dot{\varkappa},t\right) }{\delta {\cal \dot{V}}%
_{den}\left( {\sf z}\right) }\delta {\cal \dot{V}}_{den}\left( {\sf z}%
\right) \right\} =0
\end{eqnarray}%
which are equivalent to the Euler-Lagrange equations 
\begin{eqnarray}
\frac{\delta {\cal L}\left( {\cal \varkappa },{\cal \dot{\varkappa}}\right) 
}{\delta {\cal V}_{fib}\left( {\sf z}\right) }-\frac{d}{dt}\left\{ \frac{%
\delta {\cal L}\left( {\cal \varkappa },{\cal \dot{\varkappa}}\right) }{%
\delta {\cal \dot{V}}_{fib}\left( {\sf z}\right) }\right\} &=&0  \nonumber \\
\frac{\delta {\cal L}\left( {\cal \varkappa },{\cal \dot{\varkappa}}\right) 
}{\delta \omega _{fib}\left( {\sf z}\right) }-\frac{d}{dt}\left\{ \frac{%
\delta {\cal L}\left( {\cal \varkappa },{\cal \dot{\varkappa}}\right) }{%
\delta \dot{\omega}_{fib}\left( {\sf z}\right) }\right\} &=&0  \nonumber \\
\frac{\delta {\cal L}\left( {\cal \varkappa },{\cal \dot{\varkappa}}\right) 
}{\delta \omega _{den}\left( {\sf z}\right) }-\frac{d}{dt}\left\{ \frac{%
\delta {\cal L}\left( {\cal \varkappa },{\cal \dot{\varkappa}}\right) }{%
\delta \dot{\omega}_{den}\left( {\sf z}\right) }\right\} &=&0  \nonumber \\
\frac{\delta {\cal L}\left( {\cal \varkappa },{\cal \dot{\varkappa}}\right) 
}{\delta {\cal V}_{den}\left( {\sf z}\right) }-\frac{d}{dt}\left\{ \frac{%
\delta {\cal L}\left( {\cal \varkappa },{\cal \dot{\varkappa}}\right) }{%
\delta {\cal \dot{V}}_{den}\left( {\sf z}\right) }\right\} &=&0
\end{eqnarray}%
Note that each ${\cal V}_{fib}\left( {\sf z}\right) ,\omega _{fib}\left( 
{\sf z}\right) ,\omega _{den}\left( {\sf z}\right) ,{\cal V}_{den}\left( 
{\sf z}\right) $ acts like an independent variable. The condition for a
minimum is then 
\begin{equation}
\left( 
\begin{array}{cccccccc}
\frac{\delta ^{2}{\cal L}\left( \varkappa ,\dot{\varkappa},t\right) }{\delta 
{\cal V}_{fib}\delta {\cal V}_{fib}} & \cdots & \cdots & \frac{\delta ^{2}%
{\cal L}\left( \varkappa ,\dot{\varkappa},t\right) }{\delta {\cal V}%
_{fib}\delta {\cal V}_{den}} & \cdots & \cdots & \cdots & \cdots \\ 
\frac{\delta ^{2}{\cal L}\left( \varkappa ,\dot{\varkappa},t\right) }{\delta 
{\cal \omega }_{fib}\delta {\cal V}_{fib}} & \cdots & \cdots & \frac{\delta
^{2}{\cal L}\left( \varkappa ,\dot{\varkappa},t\right) }{\delta {\cal \omega 
}_{fib}\delta {\cal V}_{den}} &  &  &  &  \\ 
\frac{\delta ^{2}{\cal L}\left( \varkappa ,\dot{\varkappa},t\right) }{\delta 
{\cal \omega }_{den}\delta {\cal V}_{fib}} & \cdots & \cdots & \frac{\delta
^{2}{\cal L}\left( \varkappa ,\dot{\varkappa},t\right) }{\delta {\cal \omega 
}_{den}\delta {\cal V}_{den}} &  &  &  &  \\ 
\frac{\delta ^{2}{\cal L}\left( \varkappa ,\dot{\varkappa},t\right) }{\delta 
{\cal V}_{den}\delta {\cal V}_{fib}} & \vdots & \vdots & \frac{\delta ^{2}%
{\cal L}\left( \varkappa ,\dot{\varkappa},t\right) }{\delta {\cal V}%
_{den}\delta {\cal V}_{den}} & \vdots & \vdots & \vdots & \vdots \\ 
&  &  &  &  &  &  &  \\ 
\vdots &  &  & \vdots &  &  &  & \vdots \\ 
&  &  &  &  &  &  &  \\ 
\cdots & \cdots & \cdots & \cdots & \cdots & \cdots & \cdots & \cdots%
\end{array}%
\right) \geq 0
\end{equation}%
in a pointwise fashion. If we consider the variations $\delta \varkappa
=\left( \delta {\cal V}_{fib},\delta {\cal \omega }_{fib},\delta {\cal %
\omega }_{den},0\right) $ and find trajectories where 
\begin{eqnarray}
&&\int \left\{ \frac{\delta {\cal L}\left( \varkappa ,\dot{\varkappa}%
,t\right) }{\delta {\cal V}_{fib}\left( {\sf z}\right) }\delta {\cal V}%
_{fib}\left( {\sf z}\right) +\frac{\delta {\cal L}\left( \varkappa ,\dot{%
\varkappa},t\right) }{\delta {\cal \dot{V}}_{fib}\left( {\sf z}\right) }%
\delta {\cal \dot{V}}_{fib}\left( {\sf z}\right) \right\} =0\quad \forall
\delta {\cal V}_{fib}  \nonumber \\
&&\int \left\{ \frac{\delta {\cal L}\left( \varkappa ,\dot{\varkappa}%
,t\right) }{\delta {\cal \omega }_{fib}\left( {\sf z}\right) }\delta {\cal %
\omega }_{fib}\left( {\sf z}\right) +\frac{\delta {\cal L}\left( \varkappa ,%
\dot{\varkappa},t\right) }{\delta {\cal \dot{\omega}}_{fib}\left( {\sf z}%
\right) }\delta {\cal \dot{\omega}}_{fib}\left( {\sf z}\right) \right\}
=0\quad \forall \delta {\cal \omega }_{fib}  \nonumber \\
&&\int \left\{ \frac{\delta {\cal L}\left( \varkappa ,\dot{\varkappa}%
,t\right) }{\delta {\cal \omega }_{den}\left( {\sf z}\right) }\delta {\cal %
\omega }_{den}\left( {\sf z}\right) +\frac{\delta {\cal L}\left( \varkappa ,%
\dot{\varkappa},t\right) }{\delta {\cal \dot{\omega}}_{den}\left( {\sf z}%
\right) }\delta {\cal \dot{\omega}}_{den}\left( {\sf z}\right) \right\}
=0\quad \forall \delta {\cal \omega }_{den}
\end{eqnarray}%
and the $6\times 6$ Hessian is positive definite almost everywhere then the
action is a local minimum along those directions for that fixed value of $%
{\cal V}_{den}$.
}

\QSubDoc{Include ComplexForm}{%TCIDATA{LaTeXparent=0,0,EOMN.tex}
                      

\section{Complex forms of the EOM}

In order to write the EOM in a complex analytic form that explicitly refers
to vector bundle structure one must rearrange the order of the variables in
eq. $\left( \ref{EOMVW}\right) $and introduce an explicit pairing of the $%
\omega $ variables into $\omega _{fib}({\cal V}_{den})$ and $\omega _{den}(%
{\cal V}_{den})$ to obtain 
\begin{eqnarray}
&&{\Bbb J}\left( 
\begin{array}{cccc}
0 & \eta _{{\cal V}_{fib}\omega _{fib}} & 0 & \eta _{{\cal V}_{fib}\omega
_{den}} \\ 
\eta _{\omega _{fib}{\cal V}_{fib}} & 0 & \eta _{\omega _{fib}{\cal V}_{den}}
& 0 \\ 
0 & \eta _{{\cal V}_{den}\omega _{fib}} & 0 & \eta _{{\cal V}_{den}\omega
_{den}} \\ 
\eta _{\omega _{den}{\cal V}_{fib}} & 0 & \eta _{\omega _{den}{\cal V}_{den}}
& 0%
\end{array}%
\right) {\Bbb J}^{\dagger }{\Bbb J}\left( 
\begin{array}{c}
{\cal \dot{V}}_{fib} \\ 
\dot{\omega}_{fib} \\ 
{\cal \dot{V}}_{den} \\ 
\dot{\omega}_{den}%
\end{array}%
\right) ={\Bbb J}\left( 
\begin{array}{c}
\nabla _{{\cal V}_{fib}}{\cal E} \\ 
\nabla _{\omega _{fib}}{\cal E} \\ 
\nabla _{{\cal V}_{den}}{\cal E} \\ 
\nabla _{\omega _{den}}{\cal E}%
\end{array}%
\right) \\
&\equiv &\left( 
\begin{array}{cccc}
0 & \eta _{{\cal T}_{fib}{\cal T}_{fib}^{\ast }} & 0 & \eta _{{\cal T}_{fib}%
{\cal T}_{den}^{\ast }} \\ 
\eta _{{\cal T}_{fib}^{\ast }{\cal T}_{fib}} & 0 & \eta _{{\cal T}%
_{fib}^{\ast }{\cal T}_{den}} & 0 \\ 
0 & \eta _{{\cal T}_{den}{\cal T}_{fib}^{\ast }} & 0 & \eta _{{\cal T}_{den}%
{\cal T}_{den}^{\ast }} \\ 
\eta _{{\cal T}_{den}^{\ast }{\cal T}_{fib}} & 0 & \eta _{{\cal T}%
_{den}^{\ast }{\cal T}_{den}} & 0%
\end{array}%
\right) \left( 
\begin{array}{c}
{\cal \dot{T}}_{fib} \\ 
{\cal \dot{T}}_{fib}^{\ast } \\ 
{\cal \dot{T}}_{den} \\ 
{\cal \dot{T}}_{den}^{\ast }%
\end{array}%
\right) =\left( 
\begin{array}{c}
\nabla {\cal E}_{{\cal T}_{fib}} \\ 
\nabla {\cal E}_{{\cal T}_{fib}^{\ast }} \\ 
\nabla {\cal E}_{{\cal T}_{den}} \\ 
\nabla {\cal E}_{{\cal T}_{den}^{\ast }}%
\end{array}%
\right)
\end{eqnarray}%
where 
\begin{equation}
{\Bbb J}=\frac{1}{\sqrt{2}}\left( 
\begin{array}{cc}
\left( 
\begin{array}{cc}
{\Bbb I}_{fib} & i{\Bbb I}_{fib} \\ 
{\Bbb I}_{fib} & -i{\Bbb I}_{fib}%
\end{array}%
\right) & 0 \\ 
0 & \left( 
\begin{array}{cc}
{\Bbb I}_{den} & i{\Bbb I}_{den} \\ 
{\Bbb I}_{den} & -i{\Bbb I}_{den}%
\end{array}%
\right)%
\end{array}%
\right)
\end{equation}%
and 
\begin{eqnarray}
\left( 
\begin{array}{cc}
\left( 
\begin{array}{cc}
{\Bbb I}_{fib} & i{\Bbb I}_{fib} \\ 
{\Bbb I}_{fib} & -i{\Bbb I}_{fib}%
\end{array}%
\right) & 0 \\ 
0 & \left( 
\begin{array}{cc}
{\Bbb I}_{den} & i{\Bbb I}_{den} \\ 
{\Bbb I}_{den} & -i{\Bbb I}_{den}%
\end{array}%
\right)%
\end{array}%
\right) \left( 
\begin{array}{c}
{\cal V}_{fib} \\ 
\omega _{fib} \\ 
{\cal V}_{den} \\ 
\omega _{den}%
\end{array}%
\right) &=&\left( 
\begin{array}{c}
{\cal T}_{fib} \\ 
{\cal T}_{fib}^{\ast } \\ 
{\cal T}_{den} \\ 
{\cal T}_{den}^{\ast }%
\end{array}%
\right) \\
\frac{1}{2}\left( 
\begin{array}{cc}
\left( 
\begin{array}{cc}
{\Bbb I}_{fib} & i{\Bbb I}_{fib} \\ 
{\Bbb I}_{fib} & -i{\Bbb I}_{fib}%
\end{array}%
\right) & 0 \\ 
0 & \left( 
\begin{array}{cc}
{\Bbb I}_{den} & i{\Bbb I}_{den} \\ 
{\Bbb I}_{den} & -i{\Bbb I}_{den}%
\end{array}%
\right)%
\end{array}%
\right) ^{\dagger }\left( 
\begin{array}{c}
{\cal T}_{fib} \\ 
{\cal T}_{fib}^{\ast } \\ 
{\cal T}_{den} \\ 
{\cal T}_{den}^{\ast }%
\end{array}%
\right) &=&\left( 
\begin{array}{c}
{\cal V}_{fib} \\ 
\omega _{fib} \\ 
{\cal V}_{den} \\ 
\omega _{den}%
\end{array}%
\right)
\end{eqnarray}%
One can also express the full EOM eq. $\left( \ref{FullEOM}\right) $ in
complex analytic form by 
\begin{equation}
{\Bbb J}\left( 
\begin{array}{ccc}
0 & 0 & \eta _{{\cal V}_{den}\omega } \\ 
0 & 0 & \eta _{{\cal V}_{fib}\omega } \\ 
\eta _{\omega {\cal V}_{den}} & \eta _{{\cal V}_{den}\omega } & 0%
\end{array}%
\right) {\Bbb J}^{\dagger }{\Bbb J}\left( 
\begin{array}{c}
{\cal \dot{V}}_{den} \\ 
{\cal \dot{V}}_{fib} \\ 
\dot{\omega}%
\end{array}%
\right) ={\Bbb J}\left( 
\begin{array}{c}
\nabla _{{\cal V}_{den}}{\cal E} \\ 
\nabla _{{\cal V}_{fib}}{\cal E} \\ 
\nabla _{\omega }{\cal E}%
\end{array}%
\right)
\end{equation}%
where 
\begin{equation}
{\Bbb J}=\frac{1}{\sqrt{2}}\left( 
\begin{array}{cccc}
{\Bbb I}_{den} & 0 & i{\Bbb I}_{den} & 0 \\ 
0 & {\Bbb I}_{fib} & 0 & i{\Bbb I}_{fib} \\ 
{\Bbb I}_{den} & 0 & -i{\Bbb I}_{den} & 0 \\ 
0 & {\Bbb I}_{fib} & 0 & -i{\Bbb I}_{fib}%
\end{array}%
\right)
\end{equation}%
and 
\begin{eqnarray}
\left( 
\begin{array}{cccc}
{\Bbb I}_{den} & 0 & i{\Bbb I}_{den} & 0 \\ 
0 & {\Bbb I}_{fib} & 0 & i{\Bbb I}_{fib} \\ 
{\Bbb I}_{den} & 0 & -i{\Bbb I}_{den} & 0 \\ 
0 & {\Bbb I}_{fib} & 0 & -i{\Bbb I}_{fib}%
\end{array}%
\right) \left( 
\begin{array}{c}
{\cal V}_{den} \\ 
{\cal V}_{fib} \\ 
{\cal \omega }_{den} \\ 
{\cal \omega }_{fib}%
\end{array}%
\right) &=&\left( 
\begin{array}{c}
{\cal T} \\ 
{\cal T}^{\ast }%
\end{array}%
\right) \\
\frac{1}{2}\left( 
\begin{array}{cccc}
{\Bbb I}_{den} & 0 & i{\Bbb I}_{den} & 0 \\ 
0 & {\Bbb I}_{fib} & 0 & i{\Bbb I}_{fib} \\ 
{\Bbb I}_{den} & 0 & -i{\Bbb I}_{den} & 0 \\ 
0 & {\Bbb I}_{fib} & 0 & -i{\Bbb I}_{fib}%
\end{array}%
\right) ^{\dagger }\left( 
\begin{array}{c}
{\cal T} \\ 
{\cal T}^{\ast }%
\end{array}%
\right) &=&\left( 
\begin{array}{c}
{\cal V}_{den} \\ 
{\cal V}_{fib} \\ 
{\cal \omega }_{den} \\ 
{\cal \omega }_{fib}%
\end{array}%
\right)
\end{eqnarray}

\subsection{Partitioning of the EOM wrt ${\cal V}_{fib}({\cal V}_{den})$, $%
\protect\omega _{fib}({\cal V}_{den})$ and ${\cal V}_{den},$ $\protect\omega %
_{den}({\cal V}_{den})$}

We can partition these EOM 
\begin{equation}
\left( 
\begin{array}{cccc}
0 & \eta _{{\cal V}_{den}\omega _{den}} & 0 & \eta _{{\cal V}_{den}\omega
_{fib}} \\ 
\eta _{\omega _{den}{\cal V}_{den}} & 0 & \eta _{\omega _{den}{\cal V}_{fib}}
& 0 \\ 
0 & \eta _{{\cal V}_{fib}\omega _{den}} & 0 & \eta _{{\cal V}_{fib}\omega
_{fib}} \\ 
\eta _{\omega _{fib}{\cal V}_{den}} & 0 & \eta _{\omega _{fib}{\cal V}_{fib}}
& 0%
\end{array}%
\right) \left( 
\begin{array}{c}
{\cal \dot{V}}_{den} \\ 
\dot{\omega}_{den} \\ 
{\cal \dot{V}}_{fib} \\ 
\dot{\omega}_{fib}%
\end{array}%
\right) =\left( 
\begin{array}{c}
\nabla _{{\cal V}_{den}}{\cal E} \\ 
\nabla _{\omega _{den}}{\cal E} \\ 
\nabla _{{\cal V}_{fib}}{\cal E} \\ 
\nabla _{\omega _{fib}}{\cal E}%
\end{array}%
\right)
\end{equation}%
between the ''$fib$'' variables ${\cal V}_{fib}({\cal V}_{den})$, $\omega
_{fib}({\cal V}_{den})$ and the ''$den$'' variables to obtain equations for
fiber variable trajectories in the fixed time dependent effective fields
produced by the fixed density trajectory as 
\begin{equation}
\left( 
\begin{array}{c}
{\cal \dot{V}}_{fib} \\ 
\dot{\omega}_{fib}%
\end{array}%
\right) =\left( 
\begin{array}{cc}
0 & \eta _{{\cal V}_{fib}\omega _{fib}} \\ 
\eta _{\omega _{fib}{\cal V}_{fib}} & 0%
\end{array}%
\right) ^{-1}\left[ \left( 
\begin{array}{c}
\nabla _{{\cal V}_{fib}}{\cal E} \\ 
\nabla _{\omega _{fib}}{\cal E}%
\end{array}%
\right) -\left( 
\begin{array}{cc}
0 & \eta _{{\cal V}_{fib}{\cal \omega }_{den}} \\ 
\eta _{{\cal \omega }_{fib}{\cal V}_{den}} & 0%
\end{array}%
\right) \left( 
\begin{array}{c}
{\cal \dot{V}}_{den} \\ 
\dot{\omega}_{den}%
\end{array}%
\right) \right]
\end{equation}%
These equations can be further simplified to give 
\begin{eqnarray}
{\cal \dot{V}}_{fib} &=&\eta _{{\cal V}_{fib}\omega _{fib}}^{-1}\left[
\nabla _{\omega _{fib}}{\cal E}-\eta _{{\cal \omega }_{fib}{\cal V}_{den}}%
{\cal \dot{V}}_{den}\right]  \nonumber \\
\dot{\omega}_{fib} &=&\eta _{\omega _{fib}{\cal V}_{fib}}^{-1}\left[ \nabla
_{{\cal V}_{fib}}{\cal E}-\eta _{{\cal V}_{fib}{\cal \omega }_{den}}\dot{%
\omega}_{den}\right]  \label{FEOM}
\end{eqnarray}%
Using the real variable trajectories $\left\{ {\cal V}_{fib}\left( t\right)
,\omega _{fib}\left( t\right) ,{\cal V}_{den}\left( t\right) ,\omega
_{den}\left( t\right) \right\} $ we can obtain the following equations for $%
``den"$ variable trajectories in the fixed time dependent effective fields
produced by the $``fib"$ variable trajectories 
\begin{eqnarray}
&&\left( 
\begin{array}{c}
{\cal \dot{V}}_{den} \\ 
\dot{\omega}_{den}%
\end{array}%
\right) =\left[ \left( 
\begin{array}{cc}
0 & \eta _{{\cal V}_{den}\omega _{den}} \\ 
\eta _{\omega _{den}{\cal V}_{den}} & 0%
\end{array}%
\right) -\right.  \nonumber \\
&&\left. \left( 
\begin{array}{cc}
0 & \eta _{{\cal V}_{den}\omega _{fib}} \\ 
\eta _{\omega _{den}{\cal V}_{fib}} & 0%
\end{array}%
\right) \left( 
\begin{array}{cc}
0 & \eta _{{\cal V}_{fib}\omega _{fib}} \\ 
\eta _{\omega _{fib}{\cal V}_{fib}} & 0%
\end{array}%
\right) ^{-1}\left( 
\begin{array}{cc}
0 & \eta _{{\cal V}_{fib}\omega _{den}} \\ 
\eta _{\omega _{fib}{\cal V}_{den}} & 0%
\end{array}%
\right) \right] ^{-1}\times  \nonumber \\
&&\left[ \left( 
\begin{array}{c}
\nabla _{{\cal V}_{den}}{\cal E} \\ 
\nabla _{\omega _{den}}{\cal E}%
\end{array}%
\right) -\left( 
\begin{array}{cc}
0 & \eta _{{\cal V}_{den}\omega _{fib}} \\ 
\eta _{\omega _{den}{\cal V}_{fib}} & 0%
\end{array}%
\right) \left( 
\begin{array}{cc}
0 & \eta _{{\cal V}_{fib}\omega _{fib}} \\ 
\eta _{\omega _{fib}{\cal V}_{fib}} & 0%
\end{array}%
\right) ^{-1}\left( 
\begin{array}{c}
\nabla _{{\cal V}_{fib}}{\cal E} \\ 
\nabla _{\omega _{fib}}{\cal E}%
\end{array}%
\right) \right]
\end{eqnarray}%
which can be further simplified to give 
\begin{eqnarray}
&&\left( 
\begin{array}{c}
{\cal \dot{V}}_{den} \\ 
\dot{\omega}_{den}%
\end{array}%
\right) =\left[ \left( 
\begin{array}{cc}
0 & \eta _{{\cal V}_{den}\omega _{den}} \\ 
\eta _{\omega _{den}{\cal V}_{den}} & 0%
\end{array}%
\right) \right.  \nonumber \\
&&\left. -\left( 
\begin{array}{cc}
0 & \eta _{{\cal V}_{den}\omega _{fib}}\eta _{\omega _{fib}{\cal V}%
_{fib}}^{-1}\eta _{{\cal V}_{fib}\omega _{den}} \\ 
\eta _{\omega _{den}{\cal V}_{fib}}\eta _{{\cal V}_{fib}\omega
_{fib}}^{-1}\eta _{\omega _{fib}{\cal V}_{den}} & 0%
\end{array}%
\right) \right] ^{-1}\times  \nonumber \\
&&\left[ \left( 
\begin{array}{c}
\nabla _{{\cal V}_{den}}{\cal E} \\ 
\nabla _{\omega _{den}}{\cal E}%
\end{array}%
\right) -\left( 
\begin{array}{cc}
\eta _{{\cal V}_{den}\omega _{fib}}\eta _{\omega _{fib}{\cal V}_{fib}}^{-1}
& 0 \\ 
0 & \eta _{\omega _{den}{\cal V}_{fib}}\eta _{{\cal V}_{fib}\omega
_{fib}}^{-1}%
\end{array}%
\right) \left( 
\begin{array}{c}
\nabla _{{\cal V}_{fib}}{\cal E} \\ 
\nabla _{\omega _{fib}}{\cal E}%
\end{array}%
\right) \right]
\end{eqnarray}%
and finally 
\begin{eqnarray}
\dot{\omega}_{den} &=&\left[ \eta _{\omega _{den}{\cal V}_{den}}-\eta _{%
{\cal V}_{den}\omega _{fib}}\eta _{\omega _{fib}{\cal V}_{fib}}^{-1}\eta _{%
{\cal V}_{fib}\omega _{den}}\right] ^{-1}\left[ \nabla _{{\cal V}_{den}}%
{\cal E}-\eta _{{\cal V}_{den}\omega _{fib}}\eta _{\omega _{fib}{\cal V}%
_{fib}}^{-1}\nabla _{{\cal V}_{fib}}{\cal E}\right]  \label{OEOM} \\
{\cal \dot{V}}_{den} &=&\left[ \eta _{{\cal V}_{den}\omega _{den}}-\eta
_{\omega _{den}{\cal V}_{fib}}\eta _{{\cal V}_{fib}\omega _{fib}}^{-1}\eta
_{\omega _{fib}{\cal V}_{den}}\right] ^{-1}\left[ \nabla _{\omega _{den}}%
{\cal E}-\eta _{\omega _{den}{\cal V}_{fib}}\eta _{{\cal V}_{fib}\omega
_{fib}}^{-1}\nabla _{\omega _{fib}}{\cal E}\right]  \label{VEOM1}
\end{eqnarray}

\begin{enumerate}
\item We start with $\left( {\cal V}_{den},\omega _{den}^{0},{\cal V}%
_{fib}^{0},\omega _{fib}^{0}\right) $ then from eq. $\left( \ref{FEOM}%
\right) $

\item we obtain $\left( {\cal V}_{fib}^{\omega ^{0}},\omega _{fib}^{\omega
^{0}}\right) $ which is a self consistent solution for $\left( {\cal V}%
_{den},\omega _{den}^{0}\right) $ then from eq. $\left( \ref{OEOM}\right) $

\item we obtain $\omega _{den}^{1}$ which is a self consistent solution for $%
\left( {\cal V}_{den},{\cal V}_{fib}^{\omega ^{0}},\omega _{fib}^{\omega
^{0}}\right) $ then from eq. $\left( \ref{FEOM}\right) $

\item we obtain $\left( {\cal V}_{fib}^{\omega ^{1}},\omega _{fib}^{\omega
^{1}}\right) $ which is a self consistent solution for $\left( {\cal V}%
_{den},\omega _{den}^{1}\right) $ then

\item we continue finding solutions first to eq. $\left( \ref{FEOM}\right) $
and then from eq. $\left( \ref{OEOM}\right) $ to produce 
\begin{eqnarray}
\left( {\cal V}_{fib}^{\omega ^{0}},\omega _{fib}^{\omega ^{0}}\right) ,
&\ldots &,\left( {\cal V}_{fib}^{\omega ^{\#}},\omega _{fib}^{\omega
^{\#}}\right) \\
\omega _{den}^{1} &,\ldots &,\omega _{den}^{\#}
\end{eqnarray}%
where $\omega _{den}^{\#}$ gives $\left( {\cal V}_{fib}^{\omega
^{\#}},\omega _{fib}^{\omega ^{\#}}\right) $ and $\left( {\cal V}%
_{fib}^{\omega ^{\#}},\omega _{fib}^{\omega ^{\#}}\right) $ gives $\omega
_{den}^{\#}$ self consistently. We then
\end{enumerate}

change ${\cal V}_{den}$ to an arbitrary ${\cal V}_{den}^{\prime }$ and go
through steps 1-5 again (after choosing arbitrary/convenient new $\left(
\omega _{den}^{0},{\cal V}_{fib}^{0},\omega _{fib}^{0}\right) $ or determine
a ${\cal V}_{den}^{1}$ by solving eq. $\left( \ref{VEOM1}\right) $ and then
through steps 1-5 again, but starting with $\left( \omega _{den}^{\#},{\cal V%
}_{fib}^{\omega ^{\#}},\omega _{fib}^{\omega ^{\#}}\right) $. However this
partitioning {\bf does }not lead to constrained solutions to the full
equation of motion. Instead we must use the partitioning described in the
following section.

{\bf However this form of partitioning and solution process does not give a
solution of the constrained action problem. }

\subsection{Partitioning of the EOM wrt ${\cal V}_{den},{\cal V}_{fib}({\cal %
V}_{den})$ and $\protect\omega _{den}({\cal V}_{den}),$ $\protect\omega %
_{fib}({\cal V}_{den})$}

We expand the EOM 
\begin{equation}
\left( 
\begin{array}{cc}
\begin{array}{cc}
0 & 0 \\ 
0 & 0%
\end{array}
& 
\begin{array}{cc}
\eta _{{\cal V}_{den}\omega _{den}} & \eta _{{\cal V}_{den}\omega _{fib}} \\ 
\eta _{{\cal V}_{fib}\omega _{den}} & \eta _{{\cal V}_{fib}\omega _{fib}}%
\end{array}
\\ 
\begin{array}{cc}
-\eta _{{\cal V}_{den}\omega _{den}}^{t} & -\eta _{{\cal V}_{fib}\omega
_{den}}^{t} \\ 
-\eta _{{\cal V}_{den}\omega _{fib}}^{t} & -\eta _{{\cal V}_{fib}\omega
_{fib}}^{t}%
\end{array}
& 
\begin{array}{cc}
0 & 0 \\ 
0 & 0%
\end{array}%
\end{array}%
\right) \left( 
\begin{array}{c}
{\cal \dot{V}}_{den} \\ 
{\cal \dot{V}}_{fib} \\ 
\dot{\omega}_{den} \\ 
\dot{\omega}_{fib}%
\end{array}%
\right) =\left( 
\begin{array}{c}
\nabla _{{\cal V}_{den}}{\cal E} \\ 
\nabla _{{\cal V}_{fib}}{\cal E} \\ 
\nabla _{\omega _{den}}{\cal E} \\ 
\nabla _{\omega _{fib}}{\cal E}%
\end{array}%
\right)
\end{equation}%
to obtain 
\begin{equation}
\left( 
\begin{array}{c}
\dot{\omega}_{den} \\ 
\dot{\omega}_{fib}%
\end{array}%
\right) =\left( 
\begin{array}{cc}
\eta _{{\cal V}_{den}\omega _{den}} & \eta _{{\cal V}_{den}\omega _{fib}} \\ 
\eta _{{\cal V}_{fib}\omega _{den}} & \eta _{{\cal V}_{fib}\omega _{fib}}%
\end{array}%
\right) ^{-1}\left( 
\begin{array}{c}
\nabla _{{\cal V}_{den}}{\cal E} \\ 
\nabla _{{\cal V}_{fib}}{\cal E}%
\end{array}%
\right)  \label{w}
\end{equation}%
and 
\begin{equation}
\left( 
\begin{array}{cc}
-\eta _{{\cal V}_{den}\omega _{den}}^{t} & -\eta _{{\cal V}_{fib}\omega
_{den}}^{t} \\ 
-\eta _{{\cal V}_{den}\omega _{fib}}^{t} & -\eta _{{\cal V}_{fib}\omega
_{fib}}^{t}%
\end{array}%
\right) \left( 
\begin{array}{c}
{\cal \dot{V}}_{den} \\ 
{\cal \dot{V}}_{fib}%
\end{array}%
\right) =\left( 
\begin{array}{c}
\nabla _{\omega _{den}}{\cal E} \\ 
\nabla _{\omega _{fib}}{\cal E}%
\end{array}%
\right)  \label{VFD}
\end{equation}%
We then partition eq. $\left( \ref{VFD}\right) $ to obtain 
\begin{equation}
{\cal \dot{V}}_{fib}=-\eta _{{\cal V}_{fib}\omega _{fib}}^{-t}\left\{ \nabla
_{\omega _{fib}}{\cal E+}\eta _{{\cal V}_{den}\omega _{fib}}^{t}{\cal \dot{V}%
}_{den}\right\}  \label{vfib}
\end{equation}%
and 
\begin{equation}
{\cal \dot{V}}_{den}=-\eta _{{\cal V}_{den}\omega _{den}}^{-t}\left\{ \nabla
_{\omega _{den}}{\cal E+}\eta _{{\cal V}_{fib}\omega _{den}}^{t}{\cal \dot{V}%
}_{fib}\right\}  \label{vden}
\end{equation}%
We then solve eqs. $\left( \ref{w}\right) ,\left( \ref{VFD}\right) $, and $%
\left( \ref{vden}\right) $ for a fixed trajectory ${\cal V}_{den}$ by

\begin{enumerate}
\item Starting with $\left( {\cal V}_{den},{\cal V}_{fib}^{0},\omega
_{den}^{0},\omega _{fib}^{0}\right) $ we use eq. $\left( \ref{w}\right) $

\item to obtain $\left( {\cal \omega }_{fib}^{{\cal V}_{fib}^{0}},\omega
_{den}^{{\cal V}_{fib}^{0}}\right) $ which is a self consistent solution for 
$\left( {\cal V}_{den},{\cal V}_{fib}^{0}\right) $ then from eq. $\left( \ref%
{vfib}\right) $

\item we obtain ${\cal V}_{fib}^{1}$ which is a self consistent solution for 
$\left( {\cal V}_{den},{\cal \omega }_{fib}^{{\cal V}_{fib}^{0}},\omega
_{den}^{{\cal V}_{fib}^{0}}\right) $ then from eq. $\left( \ref{w}\right) $

\item we obtain $\left( {\cal \omega }_{fib}^{{\cal V}_{fib}^{1}},\omega
_{den}^{{\cal V}_{fib}^{1}}\right) $ which is a self consistent solution for 
$\left( {\cal V}_{den},{\cal V}_{fib}^{1}\right) $ then

\item we continue finding solutions first to eq. $\left( \ref{w}\right) $
and then from eq. $\left( \ref{vfib}\right) $ to produce 
\begin{eqnarray}
\left( {\cal \omega }_{fib}^{{\cal V}_{fib}^{0}},\omega _{den}^{{\cal V}%
_{fib}^{0}}\right) , &\ldots &,\left( {\cal \omega }_{fib}^{{\cal V}%
_{fib}^{\#}},\omega _{den}^{{\cal V}_{fib}^{\#}}\right) \\
{\cal V}_{fib}^{0} &,\ldots &,{\cal V}_{fib}^{\#}
\end{eqnarray}%
where ${\cal V}_{fib}^{\#}$ gives $\left( {\cal \omega }_{fib}^{{\cal V}%
_{fib}^{\#}},\omega _{den}^{{\cal V}_{fib}^{\#}}\right) $ and $\left( {\cal %
\omega }_{fib}^{{\cal V}_{fib}^{\#}},\omega _{den}^{{\cal V}%
_{fib}^{\#}}\right) $ gives ${\cal V}_{fib}^{\#}$ self consistently. We then
\end{enumerate}

change ${\cal V}_{den}$ to an arbitrary ${\cal V}_{den}^{\prime }$ and go
through steps 1-5 again (after choosing arbitrary/convenient new $\left( 
{\cal V}_{fib}^{0},\omega _{den}^{0},\omega _{fib}^{0}\right) $ or determine
a ${\cal V}_{den}^{1}$ by solving eq. $\left( \ref{vden}\right) $ and then
through steps 1-5 again, but starting with $\left( {\cal V}_{fib}^{\#},{\cal %
\omega }_{fib}^{{\cal V}_{fib}^{\#}},\omega _{den}^{{\cal V}%
_{fib}^{\#}}\right) $.

Thus eqs. $\left( \ref{w}\right) $ and $\left( \ref{vfib}\right) $ determine
time dependent functionals of the density, while eq. $\left( \ref{vden}%
\right) $ restricted to functionals that are solutions to eqs. $\left( \ref%
{vfib}\right) $ and $\left( \ref{w}\right) $ are {\em the exact }EOM for the
density trajectory. The sequence described above can be displayed as 
\begin{eqnarray}
&&\left( {\cal V}_{den},{\cal V}_{fib}^{0},\omega _{den}^{0},\omega
_{fib}^{0}\right) \longrightarrow \left( {\cal V}_{den},{\cal V}_{fib}^{0},%
{\cal \omega }_{fib}^{{\cal V}_{fib}^{0}},\omega _{den}^{{\cal V}%
_{fib}^{0}}\right) \longrightarrow \left( {\cal V}_{den},{\cal V}_{fib}^{1},%
{\cal \omega }_{fib}^{{\cal V}_{fib}^{0}},\omega _{den}^{{\cal V}%
_{fib}^{0}}\right) \longrightarrow  \nonumber \\
&&\left( {\cal V}_{den},{\cal V}_{fib}^{1},{\cal \omega }_{fib}^{{\cal V}%
_{fib}^{1}},\omega _{den}^{{\cal V}_{fib}^{1}}\right) \longrightarrow \cdots
\nonumber \\
&&\left( {\cal V}_{den},{\cal V}_{fib}^{k},{\cal \omega }_{fib}^{{\cal V}%
_{fib}^{k-1}},\omega _{den}^{{\cal V}_{fib}^{k-1}}\right) \longrightarrow
\left( {\cal V}_{den},{\cal V}_{fib}^{k},{\cal \omega }_{fib}^{{\cal V}%
_{fib}^{k}},\omega _{den}^{{\cal V}_{fib}^{k}}\right) \longrightarrow \left( 
{\cal V}_{den},{\cal V}_{fib}^{k+1},{\cal \omega }_{fib}^{{\cal V}%
_{fib}^{k}},\omega _{den}^{{\cal V}_{fib}^{k}}\right) \longrightarrow 
\nonumber \\
&&\left. {}\right. \qquad \qquad \qquad \qquad \qquad \qquad \qquad \qquad
\qquad \qquad \qquad \cdots \longrightarrow \left( {\cal V}_{den},{\cal V}%
_{fib}^{\#},{\cal \omega }_{fib}^{{\cal V}_{fib}^{\#}},\omega _{den}^{{\cal V%
}_{fib}^{\#}}\right)
\end{eqnarray}%
{\bf However this form of partitioning and solution process does not give a
solution of the constrained action problem. }
}

\QSubDoc{Include Antisym}{\input{Antisym.tex}}

\QSubDoc{Include Group}{%TCIDATA{LaTeXparent=0,0,EOMN.tex}
                      

\section{Commutative Group ${\frak G}_{\ker }$\protect\bigskip }

We can look for a unitary subgroup of ${\cal U}({\cal H}^{N})$ such that it
maps equivalence classes $\left[ f\right] _{N}$ to themselves and sends
positive operators to positive operators. If it exists it has the properties%
\begin{eqnarray}
&&Tr\{\Phi ^{\dagger }\Phi UYU^{\dagger }\}=Tr\{\Phi ^{\dagger }\Phi
Y\}\quad \forall \,Y\in {\cal B}_{1}({\cal H}^{N})  \nonumber \\
&\Rightarrow &Tr\{\Phi ^{\dagger }\Phi UD^{N}U^{\dagger }\}=Tr\{\Phi
^{\dagger }\Phi D^{N}\}\quad \forall \,D^{N}\in {\cal B}_{1}({\cal H}^{N}) 
\nonumber \\
&\Rightarrow &\left\langle U\Psi |\Phi ^{\dagger }\Phi U\Psi \right\rangle
=\left\langle \Psi |\Phi ^{\dagger }\Phi \Psi \right\rangle \quad \forall
\Psi \in {\cal H}^{N}  \nonumber \\
&\Leftrightarrow &U^{\dagger }\Phi ^{\dagger }\Phi U=\Phi ^{\dagger }\Phi
\Longleftrightarrow \left[ \Phi ^{\dagger }\Phi ,U\right] =0  \nonumber \\
&\Longleftrightarrow &U=f\left( \Phi ^{\dagger }\Phi \right) \Rightarrow
U=\exp i\int \omega \left( \underline{\kappa }\right) \prod_{1\leq \iota
\leq N}\Phi \left( \kappa _{i}\right) ^{\dagger }\Phi \left( \kappa
_{i}\right) d\underline{\kappa }
\end{eqnarray}%
i.e. it is a unitary commutative subgroup, ${\frak G}_{\Xi }$, of ${\cal U}(%
{\cal H}^{N})$ and it is diagonal in the coordinate representation, {\bf %
which is its defining representation}. Hence, trivially, it has the
representation, ${\cal R}$, in the space of bounded operators acting in $%
{\cal H}^{N}$%
\begin{eqnarray}
{\cal R} &:&{\cal G}_{\Xi }\longrightarrow {\cal B}\left( {\cal H}^{N}\right)
\nonumber \\
{\cal R}\left( g\right) &=&\exp i\int \omega _{g}\left( \underline{\kappa }%
\right) \prod \Phi \left( \kappa _{i}\right) ^{\dagger }\Phi \left( \kappa
_{i}\right) d\underline{\kappa };\quad g\in {\cal G}_{\Xi }  \label{unitary}
\end{eqnarray}%
This representation has the property that ${\cal R}\left( {\cal G}_{\Xi
}\right) :$ $\left[ f\right] _{N}\longrightarrow \left[ f\right] _{N}\quad
\forall f$ \ i.e. every ${\cal R}\left( g\right) $ maps all the members of
an equivalence to the same equivalence class and the same representation
works for all equivalence classes, explicitly $UYU^{\dagger }\in \left[ f%
\right] _{N}\Longleftrightarrow Y\in \left[ f\right] _{N}$ for $U\in {\cal R}%
\left( {\cal G}_{\Xi }\right) $. (Note a pure state $\left| \Psi
\right\rangle $ defines an element of an equivalence class by $\left| \Psi
\right\rangle \left\langle \Psi \right| $).

${\cal R}\left( {\cal G}_{\Xi }\right) $ does not produce all pure states
with the same density as a reference pure state, when acting on that
reference state. This is proved by noting that $\langle $\underline{$\kappa $%
}${\bf |}{\cal U}\left( \omega \right) \Psi \rangle =e^{i\omega \left( 
\underline{\kappa }\right) }\Psi \left( \underline{\kappa }\right) $ \ does
not range over all pure state wavefunctions that have the same density for a
fixed $\Psi $ and varying $\omega .$ This can be seen from the observations
that $e^{{\cal V}\left( \underline{\kappa }\right) }\Psi \left( \underline{%
\kappa }\right) $ has the same density as $\Psi \left( \underline{\kappa }%
\right) $ when $e^{{\cal V}\left( \underline{\kappa }\right) }=\left( 1+\phi
\left( \underline{\kappa }\right) \right) ^{\frac{1}{2}}$ and $\phi \left( 
\underline{\kappa }\right) $ is strongly orthogonal to $\left| \Psi \left( 
\underline{\kappa }\right) \right| ^{2}$ and noting that $e^{{\cal V}\left( 
\underline{\kappa }\right) }\Psi \left( \underline{\kappa }\right) $ cannot
be obtained from $\Psi \left( \underline{\kappa }\right) $ by multiplying by 
$e^{i\omega \left( \underline{\kappa }\right) }.$ Thus the action of the
unitary commutative subgroup does not produce all pure states with the same
density.

If one drops the unitary requirement {\it it is} possible find
transformations, {\it that are diagonal in the coordinate representation, }%
which produce all pure states with a given density, $\gamma $, from a pure
reference state $\left| \Psi \right\rangle $ of that density{\it . }These
transformations are of the form%
\begin{equation}
\breve{T}=\exp \left\{ \int \left\{ {\cal V}\left( \underline{\kappa }%
\right) +i\omega \left( \underline{\kappa }\right) \right\} \prod_{1\leq
\iota \leq N}\Phi \left( \kappa _{i}\right) ^{\dagger }\Phi \left( \kappa
_{i}\right) d\underline{\kappa }\right\} {\bf \quad \Rightarrow }\breve{T}%
\in GL\left( {\cal H}^{N}\right)
\end{equation}%
where $e^{{\cal V}\left( \underline{\kappa }\right) }=\left( 1+\phi \left( 
\underline{\kappa }\right) \right) ^{\frac{1}{2}}$ and $\phi \left( 
\underline{\kappa }\right) $ is strongly orthogonal to $\left| \Psi \left( 
\underline{\kappa }\right) \right| ^{2}$. The proof of this assertion is by
construction and by noting that these transformations produce all pure
states from any reference pure state. This can be seen by noting that

\begin{equation}
\left[ \breve{T}\Psi \right] \left( \underline{\kappa }\right) =e^{{\cal V}%
\left( \underline{\kappa }\right) +i\omega \left( \underline{\kappa }\right)
}\Psi \left( \underline{\kappa }\right)
\end{equation}%
and by considering the local norm/phase form of the initial and final
wavefunctions%
\begin{eqnarray}
\langle \underline{\kappa }\left| \Psi \right\rangle &=&N\left( \underline{%
\kappa }\right) e^{i\theta \left( \underline{\kappa }\right) } \\
\langle \underline{\kappa }\left| \breve{T}\Psi \right\rangle &=&N^{\prime
}\left( \underline{\kappa }\right) e^{i\theta ^{\prime }\left( \underline{%
\kappa }\right) }
\end{eqnarray}%
These transformations form{\it \ }{\bf a defining representation of a
commutative subgroup}{\it \ }${\frak G}_{fib}={\cal G}_{\ker }\times {\cal G}%
_{\Xi }$. This is the{\bf \ }abstract group whose representation ${\cal R}%
\left( {\frak G}_{fib}\right) $ maps each equivalence class 
\begin{equation}
\lbrack f]_{N}\equiv \left\{ \left( \breve{k},\breve{\Gamma}\left( f\right)
\right) ;\breve{k}\in \ker \Xi ,\text{ }\breve{\Gamma}\left( f\right) \in
\ker \Xi ^{c},\text{fixed }f\right\}
\end{equation}%
to itself for {\bf any }$f$ i.e. 
\begin{equation}
{\cal R}\left( g\right) \left( \breve{k},\breve{\Gamma}\left( f\right)
\right) =\left( \breve{k}^{\prime },\breve{\Gamma}\left( f\right) \right)
;\quad g\in {\frak G}_{fib};\,{\cal R}\left( g\right) \breve{k}=\breve{k}%
^{\prime };{\cal R}\left( g\right) \breve{\Gamma}\left( f\right) =\breve{%
\Gamma}\left( f\right) \,
\end{equation}%
The representation is 
\begin{eqnarray}
{\cal R} &:&{\frak G}_{fib}\longrightarrow {\cal B}\left( {\cal H}^{N}\right)
\nonumber \\
{\cal R}\left( g\right) &=&\exp \int \left\{ {\cal V}\left( \underline{%
\kappa }\right) +i\omega \left( \underline{\kappa }\right) \right\} \prod
\Phi \left( \kappa _{i}\right) ^{\dagger }\Phi \left( \kappa _{i}\right) d%
\underline{\kappa ;}\qquad g\in {\frak G}_{fib};
\end{eqnarray}%
where $\left\{ {\cal V}\left( \underline{\kappa }\right) \right\} $ have the
property 
\begin{equation}
{\cal V}\left( \underline{\kappa }\right) \Gamma \left( f\right) \left( 
\underline{\kappa }\right) =0\quad a.e.
\end{equation}%
wrt $\Gamma \left( f\right) \left( \underline{\kappa }\right) \quad \forall
f\in L_{1}\left( Z\right) .$ Explicitly $TYT^{\dagger }\in \left[ f\right]
_{N}\Longleftrightarrow Y\in \left[ f\right] _{N}$ for $T\in {\cal R}\left( 
{\frak G}_{fib}\right) $. The group ${\frak G}_{fib}$ does {\it not }act
homogeneously in $\left[ f\right] _{N}$, i.e. it does not generate all of $%
\left[ f\right] _{N}$ from a reference element, but it does generate all
pure states in $\left[ f\right] _{N}$ from a reference pure state in $\left[
f\right] _{N}$ i.e. all pure states lie on a single orbit of ${\frak G}%
_{fib} $ in $\left[ f\right] _{N}$ in contrast to orbits of ${\cal G}_{\Xi }$%
. ({\bf Note }that{\bf \ }the groups ${\cal G}_{\Xi }$ and ${\frak G}_{fib}$
act in ${\frak B}_{vect}$ through their representations as operators in $%
{\cal B}_{1}\left( {\cal H}^{N}\right) $.) Agian {\bf one} representation
suffices for {\bf all }equivalence classes.

The set of pure state wavefunctions $\left\{ \Phi _{\gamma }\left( 
\underline{\kappa }\right) \right\} $ associated with pure states in an
equivalence class $\left[ \gamma \right] _{N}$ can all be expressed in terms
of a reference wavefunction $\Psi \left( \underline{\kappa }\right) $, and
the non redundant parameterizing functions ${\cal V}$ and $\omega $ as $%
\left\{ \Phi \left( \underline{\kappa }\right) =e^{{\cal V}\left( \underline{%
\kappa }\right) +i\omega \left( \underline{\kappa }\right) }\Psi \left( 
\underline{\kappa }\right) =\left\langle \underline{\kappa }\left| e^{{\cal V%
}+i\breve{\omega}}\right. \Psi \right\rangle \right\} .$ If the reference
state is normalized then the transformed states $\left| e^{{\cal V}+i\breve{%
\omega}}\Psi \right\rangle $ $\equiv \left| \Psi \left( {\cal V},\omega
\right) \right\rangle $ are normalized. The question of {\bf whether there
exists} a set of unitary operators $\subset U\left( {\cal H}^{N}\right)
\subset {\cal B}\left( {\cal H}^{N}\right) $ that form a group and map $%
\left\{ \left| \Psi \left( {\cal V},\omega \right) \right\rangle
\left\langle \Psi \left( {\cal V},\omega \right) \right| \right\} $ $\subset %
\left[ \gamma \right] _{N}$ to themselves is an open question. Clearly such
a set of operators would have be noncommutative and the set would depend on $%
\gamma $ and correspond to different (but probably equivalent) unitary
representations of some abstract group.

Luckily we do not need a unitary solution so the commutative non unitary one
obtained above suffices and has the added nice property of mapping diagonal
operators in the coordinate representation to diagonal operators and
nondiagonal ones to nondiagonal ones, as well as being commutative.
}

\begin{references}
\bibitem{Geroch} Robert Geroch, {\it Mathematical Physics, p57} (University
of Chicago Press 1985).
\end{references}

\end{document}
